\documentclass[aps,prd,onecolumn,superscriptaddress,nofootinbib,showpacs,12pt,notitlepage]{revtex4-2}

\usepackage{amsmath,amssymb,amsfonts}
\usepackage{graphicx}
\usepackage{hyperref}
\usepackage{xcolor}
\usepackage{bm}
\usepackage{booktabs}
\usepackage{array}
\usepackage{mathtools}
\usepackage{physics}
\usepackage{slashed}

\hypersetup{
  hidelinks,
  bookmarksnumbered=true,
  pdfencoding=unicode,
  psdextra
}

\newcommand{\Mpl}{M_{\mathrm{Pl}}}
\newcommand{\half}{\tfrac{1}{2}}
\newcommand{\rs}{r_{\mathrm{s}}}
\newcommand{\azero}{a_{0}}
\newcommand{\meff}{m_{\mathrm{eff}}}
\renewcommand{\vev}{\langle 0 |}
\renewcommand{\ket}[1]{| #1 \rangle}
\renewcommand{\bra}[1]{\langle #1 |}
\newcommand{\eV}{\,\mathrm{eV}}
\newcommand{\GeV}{\,\mathrm{GeV}}
\renewcommand{\dd}{\mathrm{d}}
\renewcommand{\order}[1]{\mathcal{O}\!\left(#1\right)}

\linespread{1.15}\selectfont

\makeatletter
\renewcommand*{\l@section}[2]{%
  \ifnum \c@tocdepth >\z@
    \addpenalty\@secpenalty
    \addvspace{0.9em \@plus\p@}%
    \begingroup
      \parindent \z@ \rightskip \@pnumwidth
      \parfillskip -\@pnumwidth
      \leavevmode \bfseries
      \setlength{\@tempdima}{2.4em}%
      #1\ignorespaces
      \nobreak\hfil \nobreak\hb@xt@\@pnumwidth{\hss #2}\par
    \endgroup
  \fi}
\renewcommand{\l@subsection}{\@dottedtocline{2}{2.4em}{2.0em}}
\makeatother

\begin{document}

\title{Quantum extensions of bi-conformal scalar-tensor gravity:
  particles as spacetime resonances, the Bernoulli mechanism
  for gravity, and falsifiable predictions}

\author{George Lotishvili}
\email{g\_lotishvili@eeu.edu.ge}
\affiliation{East European University, Tbilisi, Georgia}
\affiliation{Independent research}

\date{Version 4.7 \\ \today}

\begin{abstract}
\setlength{\leftskip}{0pt}\setlength{\rightskip}{0pt}%
\linespread{1.0}\selectfont\small
We develop the quantum extension of bi-conformal scalar-tensor gravity
(ISPG), a Horndeski-class theory defined by the action
$S = (16\pi G)^{-1}\!\int\! d^4x\sqrt{-g}\,[R + \frac{1}{2}(\partial\varphi)^2]
+ S_{\mathrm{m}}$ with the bi-conformal metric
$ds^2 = -e^\varphi c^2 dt^2 + e^{-\varphi}(dx^2+dy^2+dz^2)$.
In the classical regime, ISPG reproduces all solar-system tests
($\gamma=\beta=1$), resolves the Schwarzschild singularity,
derives MOND phenomenology ($a_0 = cH_0/2\pi$), and predicts
falsifiable deviations at 2PN order~\cite{Lotishvili2026}.
The central aim of the present paper is to identify the
\emph{quantum mechanism of gravity}: how localized resonant
excitations of the scalar-metric medium generate the static pressure
deficit that becomes the gravitational field.  Here we extend the
theory to the quantum domain by interpreting
elementary particles as localized resonant excitations of the
scalar-metric medium.  We derive: (i)~the spectrum of scalar-field
excitations on the bi-conformal background, establishing the
effective mass formula $\meff = m_0 e^{\varphi/2}$ at the quantum
level; (ii)~a Bernoulli-type mechanism in which the same scalar
profile that defines the metric also carries a static pressure
deficit in the scalar-metric medium, with exterior pressure
$P_{\mathrm{static}} = -e^{\varphi}\varphi^4/(32\pi G\,\rs^2)$
whose universal shape function is $\hat{f}(\varphi) = e^{\varphi}\varphi^4$;
(iii)~the cosmological negative-feedback loop that dynamically
regulates the accumulated echo-background energy and, on the
perturbative tail-power branch, gives an order-of-magnitude scale
$\rho_\Lambda \sim H_0^2 \Mpl^2$, while the full zero-point
decoupling and non-perturbative normalization remain open; (iv)~the dark-energy equation of state
$w(z) = -1 - \delta w_0\,(1+z)^3/E(z)$ with
$\delta w_0 \sim 10^{-2}$--$10^{-1}$ (source-driven phantom,
no ghost instability), whose magnitude remains conditional on the
tail-power normalization and whose DESI DR2 CPL comparison is mixed
(same sign of $w_a$, opposite sign of $w_0+1$); (v)~gravitational decoherence of
entangled particles with rate $\Gamma_{\mathrm{dec}} = 2m_0|\Delta\Phi_N|/h$,
predicting $\Gamma_{\mathrm{dec}} \sim 0.5$~Hz for eV-scale photon
pairs at 10~m height separation and effectively instantaneous
decoherence for macroscopic systems (MAQRO, STE-QUEST);
(vi)~a topological classification of particle spin $J = n\hbar/2$
and antimatter as vortices with opposite winding number $n = \pm 1$;
(vii)~gravitational birefringence from the gradient-squared scalar
channel, with $\Delta c/c \propto (\nabla\varphi)^2/m_e^2$
distinct from the GR$+$QED Drummond--Hathrell curvature channel;
the numerical coefficient and detectability forecast are future
one-loop calculation targets; (viii)~anisotropic quantum tunneling enhanced toward
gravitational sources by factor $T_\downarrow/T_\uparrow \sim
\exp[d\Delta\varphi\sqrt{2m(V_0-E)}/\hbar]$, relevant for nuclear
reactions in white dwarf interiors; (ix)~a schematic cosmological variation
of the fine-structure constant $|\Delta\alpha/\alpha| \sim 10^{-7}$--$10^{-6}$
at $z \sim 1$--$3$, subject to a fixed-scale/renormalization prescription
and testable by ESPRESSO and ELT-ANDES; (x)~numerical support for the Koide relation
$Q = 2/3$ from the MASS/Mathieu lepton ladder and the
E1D/rank-1 cavity analysis, with $Q$ reproduced to
$0.024\%$ in $N$-space and to $0.009\%$ in the E1D check,
and with the $3$D-to-$1$D cavity bridge and the
self-consistent $x = 33+24\sqrt{2}$ selection identified as
the remaining numerical task; (xi)~a derivation of Newton's second law from the
retarded potential of an accelerated oscillon, proving that
$m_{\mathrm{inertial}} \equiv m_{\mathrm{gravitational}}$ as a
mathematical identity rather than an axiom, and connecting inertia
to medium hysteresis via a Klein--Gordon effective field whose
microscopic tracking stiffness is motivated by the oscillon's
quasi-normal mode structure;
(xii)~a structural argument that the scalar sensitivity
$s = 1/2$ and the suppression of scalar dipole radiation
are protected at perturbative/EFT precision by a
shift-symmetry Ward identity (exact for vacuum compact
objects at all loop orders, leading-order for
matter-filled bodies; the exact non-perturbative
matter-coupled cancellation remains an open task per
MAIN/16), supporting a perturbatively controlled strong
equivalence statement in the scalar sector;
(xiii)~a singularity-free cosmological onset in which the
first excitations nucleate quantum-mechanically from the
quiescent medium, resolving the horizon, flatness, and
monopole problems without inflation.
The closed central result is the oscillon $\to$ Bernoulli
pressure-deficit $\to$ source $\to 1/r$ field chain for gravity,
together with the reduced-gap/coercive closure showing that,
after projection onto the source manifold, non-neutral deformations
remain perturbatively controlled while the exterior field follows the
neutral collective coordinate.
The remaining sectors show how the same medium picture extends into
cosmology, particle structure, and quantum phenomena; where a sector
still requires a dedicated numerical or structural closure, that
status is stated explicitly in the text and does not affect the core
gravity derivation.
\end{abstract}

\maketitle
\newpage
\tableofcontents
\newpage

%===================================================================
\section{Introduction}\label{sec:intro}
%===================================================================

\subsection{Motivation}

The reconciliation of gravity with quantum mechanics remains the
central unsolved problem in theoretical physics.  General relativity
(GR) describes spacetime as a dynamical geometric entity, while quantum
field theory (QFT) treats spacetime as a fixed background arena populated
by quantized fields.  Neither framework, taken alone, provides a
satisfactory account of the other's domain: GR predicts singularities
where quantum effects should dominate~\cite{Penrose1965,Hawking1970},
while QFT on curved spacetime leads to the $10^{120}$ cosmological
constant catastrophe~\cite{Weinberg1989,Martin2012} and leaves the
origin of gravitational decoherence
unexplained~\cite{Penrose1996,Diosi1987}.

Attempts to merge these frameworks have produced string
theory~\cite{Polchinski1998}, loop quantum gravity~\cite{Rovelli2004},
and asymptotic safety~\cite{Reuter2012}, each introducing substantial
theoretical apparatus (extra dimensions, spin foams, or UV fixed points)
whose empirical signatures remain elusive.

In the classical domain, bi-conformal scalar-tensor gravity
(ISPG)~\cite{Lotishvili2026} has demonstrated that a minimally
coupled phantom scalar field $\varphi$ with the bi-conformal metric
\begin{equation}\label{eq:metric}
  ds^{2} = -e^{\varphi}\,c^{2}\,dt^{2}
         + e^{-\varphi}\bigl(dx^{2}+dy^{2}+dz^{2}\bigr)
\end{equation}
and the action
\begin{equation}\label{eq:action}
  S = \frac{1}{16\pi G}\int d^{4}x\,\sqrt{-g}\left[
    R + \frac{1}{2}\,g^{\mu\nu}\partial_{\mu}\varphi\,
    \partial_{\nu}\varphi\right] + S_{\mathrm{m}},
\end{equation}
reproduces the classical 1PN solar-system tests
($\gamma = \beta = 1$), predicts a falsifiable ISPG--GR
2PN Shapiro-delay difference $\Delta B = \pi/4$ (closed
regularization-independent observable; standalone absolute
coefficients are finite-endpoint convention-dependent), resolves the Schwarzschild
singularity, is compatible with the CMB angular power spectrum at
linear order, and derives MOND phenomenology with
$a_0 = cH_0/(2\pi)$ and the baryonic Tully--Fisher relation;
the 15\% offset between the instantaneous formula and the
observed value is resolved by a vortex-lag mechanism that
predicts galaxy-age-dependent
$a_0$~\cite{Lotishvili2026MOND}.
The same classical compact-object sector predicts a test-particle
ISCO at $r_{\mathrm{ISCO}} = \frac{3+\sqrt{5}}{2}\,\rs$, an
ISCO orbital frequency $6.9\%$ lower than in GR for the same
total mass, a photon sphere at $r=\rs$ with a $4.6\%$ larger
shadow, and suppressed scalar dipole radiation in the
controlled sensitivity regimes (vacuum exact, matter leading
order) where $s=1/2$~\cite{Lotishvili2026}.
The physical interpretation of $\varphi$ as the pressure state of an
information medium, whose scalar profile both defines the physical
metric and carries pressure-deficit and transport information, provides
a mechanical picture of gravity:
the effective mass of a particle at position $\varphi$ is
\begin{equation}\label{eq:meff}
  \meff(\varphi) = m_0\,e^{\varphi/2},
\end{equation}
which is locally undetectable because all masses, clocks, and
rulers participate in the same bi-conformal rescaling.

The classical theory, however, does not address the fundamental
nature of matter itself.  It treats particles as point sources in the
field equation $\Box\varphi = -(8\pi G/c^4)\,T$, without specifying
what constitutes these sources at the microscopic level.

\subsection{Central thesis}

The present work fills this gap by proposing and formalizing a single,
unifying hypothesis:

\begin{quote}
\textit{Elementary particles are not objects placed in spacetime.
They are localized resonant excitations of the scalar-metric
medium itself.}
\end{quote}

This identification has far-reaching consequences.  If matter
\emph{is} spacetime vibration, then:
\begin{itemize}
  \item Mass is the amplitude of resonance (energy stored in the
    vibrational mode), naturally explaining
    Eq.~(\ref{eq:meff}).
  \item Gravity arises because resonant oscillations redistribute
    kinetic energy density $\frac{1}{2}(\nabla\varphi)^2$,
    lowering the static pressure of the medium via an analog of
    Bernoulli's principle.
  \item Entropy increase is the irreversible spreading of resonant
    tails into the surrounding medium.
  \item Vacuum energy is regulated by a negative-feedback loop
    between resonance production and pressure depletion.
  \item Wave-particle duality, the uncertainty principle, and
    quantum entanglement follow as natural properties of
    spatially extended resonant structures.
\end{itemize}

\subsection{Relation to existing approaches}

The proposal that particles are excitations of a medium has
precursors.  Kelvin's vortex theory of atoms~\cite{Thomson1867}
proposed that atoms are knots in the luminiferous ether.
Superfluid vacuum theory (SVT) models spacetime as a
Bose--Einstein condensate~\cite{Volovik2003,Zloshchastiev2011},
and emergent gravity approaches
derive gravitational dynamics from
thermodynamic or entropic
considerations~\cite{Verlinde2011,Jacobson1995}.

ISPG shares the philosophical starting point of SVT and emergent
gravity but differs in three critical respects:

\begin{enumerate}
  \item \textbf{No extra dimensions.} Unlike string theory, ISPG operates
    in standard $(3\!+\!1)$-dimensional spacetime.  The diversity
    of particle species arises from the internal degrees of freedom of
    a single medium (analogous to longitudinal, transverse, and
    vortical modes of a fluid), not from compactified extra dimensions.

  \item \textbf{Exact Lorentz invariance.}  SVT generically breaks
    Lorentz symmetry at high energies through the preferred rest frame
    of the condensate.  The ISPG bi-conformal structure preserves local
    Lorentz invariance by construction: the speed of light is locally
    invariant ($c_{\mathrm{local}} = c$) despite position-dependent
    coordinate speed~\cite{Lotishvili2026}.

  \item \textbf{A concrete mechanical origin of gravity.}  Entropic
    gravity~\cite{Verlinde2011} derives gravitational force from
    information-theoretic considerations but does not specify a
    physical mechanism.  ISPG provides one: the Bernoulli-type
    kinetic-energy transfer from resonant oscillations to the static
    pressure of the medium.
\end{enumerate}

\subsection{Structure of the paper}

The logical center of the paper is the derivation of the quantum
gravitational mechanism itself: localized oscillon sources, the
Bernoulli pressure deficit, the static $1/r$ field, Newtonian
recovery, and the source-following collective mode.  The remaining
sections develop broader consequences of the same scalar-metric
picture for inertia, cosmology, particle structure, and quantum
phenomenology.  These sectors do not all play the same logical role:
the gravity chain is the primary closed result of the paper, while
several wider extensions are carried to the level needed here and are
marked in the text where a separate full closure remains future work.

The paper is organized as follows.

\textbf{Section~\ref{sec:canonical_quantization}} derives the
canonical quantization of the coupled scalar-metric system using
the ADM formalism, obtains the reduced Lagrangian and Hamiltonian
constraint, analyses the constraint algebra and the relationship
between reduced and Dirac quantization, derives the Wheeler--DeWitt
equation, verifies the semi-classical limit via WKB expansion, and
discusses the ultralocal kinetic structure and its implications.

\textbf{Section~\ref{sec:excitations}} derives the spectrum of
scalar-field excitations on the bi-conformal background,
establishes the identification of particles as resonant modes,
and argues for the all-orders quantum stability of the phantom
sector within the pure-gravity constrained sector via
kinematic ghost freedom, the structural constraint analysis
of Appendix~11 (full constraint-algebra closure left to a
separate technical proof), and non-perturbative boundedness
on the constrained background. The matter-coupled extension is
controlled through the EFT completion described below.  It derives the
structure of the quantum effective action, demonstrating
shift-symmetry protection, controlled one-loop divergence
structure with negligible physical corrections, and
systematic EFT power counting at all loop orders.  It then
presents the structural argument for the quantum persistence
of the scalar sensitivity ($s = 1/2$ at all loop orders for
vacuum compact objects; leading-order in $|\varphi|$ for
matter-filled bodies, with the exact non-perturbative
matter-coupled cancellation marked open per MAIN/16),
supporting suppressed scalar dipole radiation and a
perturbatively controlled strong equivalence statement in the
scalar sector.

\textbf{Section~\ref{sec:bernoulli_identity}} derives the Bernoulli mechanism:
the kinetic energy of particle resonances generates the static pressure
deficit that manifests as the gravitational potential
$\varphi = -\rs/r$.

\textbf{Section~\ref{sec:inertia}} derives Newton's second law
$\mathbf{F} = -\meff\,\mathbf{a}$ from the retarded potential of an
accelerated oscillon, proves the equivalence
$m_{\mathrm{inertial}} \equiv m_{\mathrm{gravitational}}$ as a
mathematical identity, connects inertia to medium hysteresis
via a Klein--Gordon field, and provides a microscopic motivation
for the effective hysteresis action template and its microscopic
tracking stiffness from the quasi-normal mode structure of the
oscillon.

\textbf{Section~\ref{sec:entropy}} analyzes entropy production
through irreversible resonance spreading, derives the
cosmological negative-feedback loop, shows that the
echo-background energy density (observable dark energy)
is dynamically regulated to a value
compatible with observation on the perturbative branch, obtains the
dark-energy equation of state $w(z) < -1$ (source-driven phantom)
as a falsifiable test target rather than a current DESI confirmation,
and addresses the earliest-epoch
dynamics: quantum nucleation of the first oscillons from
the quiescent medium, the onset cascade, singularity
avoidance, and the resolution of the horizon, flatness,
and monopole problems without inflation.

\textbf{Section~\ref{sec:mass_spectrum}} analyzes the spectrum
of three-dimensional resonant modes, formulates the self-consistent
nonlinear eigenvalue problem in dimensionless form, shows that
the confining potential is a geometric property of the nonlinear
wave equation independent of $G$, establishes the formal
direct-cavity amplitude--frequency ansatz with gravitational
self-energy corrections of order $(m/\Mpl)^2 \sim 10^{-45}$,
presents the lattice formulation suitable for numerical solution,
and connects the structural cavity sector to the MASS/Mathieu
lepton ladder and the Koide consistency condition.

\textbf{Section~\ref{sec:higgs}} shows that the Higgs mechanism
emerges as the electroweak projection of bi-conformal symmetry
breaking, with $v(\varphi) = v_0\,e^{\varphi/2}$, and derives
the complete electroweak gauge boson spectrum ($W^{\pm}$, $Z$,
$\gamma$) on the bi-conformal background, demonstrating that
all Standard Model mass ratios and dispersion relations are
preserved locally, with photon masslessness protected
by $U(1)_{\mathrm{EM}}$ gauge invariance and reinforced
by the ISPG shift symmetry.

\textbf{Section~\ref{sec:entanglement}} formulates quantum
entanglement as correlated resonance structure and derives the
gravitational decoherence rate.

\textbf{Section~\ref{sec:topology}} develops the topological
classification of particle species: spin as vortex winding
number, antimatter as opposite-winding vortex, and the Pauli
exclusion principle from vortex topology.

\textbf{Section~\ref{sec:birefringence}} isolates the
gravitational-birefringence channel and its leading
$(\nabla\varphi)^2/m_e^2$ scaling, with the numerical
coefficient deferred to a dedicated one-loop calculation,
and presents the cosmological variation estimate of the
fine-structure constant $\alpha$.

\textbf{Section~\ref{sec:predictions}} summarizes all
quantitative predictions in a unified table and identifies the
most immediately testable signatures.

\textbf{Section~\ref{sec:conclusion}} concludes.

\subsection{Notation and conventions}

We use the $(-,+,+,+)$ metric signature and natural units
$\hbar = c = 1$ unless explicitly restored.  The reduced
Planck mass is $\Mpl = (8\pi G)^{-1/2}$.  Greek indices
$\mu,\nu,\ldots$ run over $0,1,2,3$; Latin indices $i,j,\ldots$
over spatial coordinates $1,2,3$.  The d'Alembertian is
$\Box \equiv g^{\mu\nu}\nabla_\mu\nabla_\nu$.  The scalar field
$\varphi$ is dimensionless, with the static weak-field solution
$\varphi = -\rs/r$ where $\rs = 2GM/c^2$ is the Schwarzschild
radius.  Equation numbers from the classical
paper~\cite{Lotishvili2026} are cited as (I.\emph{n}).

%===================================================================
\section{Canonical quantization of the scalar-metric system}
\label{sec:canonical_quantization}
%===================================================================

The subsequent sections treat $\varphi$ as a quantum field on
a classical background.  Here we derive the canonical
quantization of the \emph{coupled} scalar-metric system,
exploiting the fact that the bi-conformal ansatz reduces the
full gravitational \emph{configuration} space to a single
scalar degree of freedom (the corresponding phase space is
two-dimensional, $(\varphi,\pi_\varphi)$).

%-------------------------------------------------------------------
\subsection{ADM decomposition of the bi-conformal metric}
\label{sec:adm}
%-------------------------------------------------------------------

Following Arnowitt, Deser, and Misner~\cite{ADM1962}, we
foliate spacetime into spatial hypersurfaces $\Sigma_t$
and write
\begin{equation}\label{eq:adm_metric}
  ds^2 = -N^2\,dt^2
    + h_{ij}(dx^i + N^i\,dt)(dx^j + N^j\,dt),
\end{equation}
where $N$ is the lapse function, $N^i$ the shift vector,
and $h_{ij}$ the induced metric on $\Sigma_t$.
The bi-conformal metric~(\ref{eq:metric}) fixes the spatial
geometry and shift,
\begin{equation}\label{eq:adm_biconformal}
  h_{ij} = e^{-\varphi}\,\delta_{ij},\qquad N^i = 0,
\end{equation}
with $\sqrt{h} = e^{-3\varphi/2}$, while the lapse $N$ is
retained as a free function for the canonical analysis;
its on-shell value $N = e^{\varphi/2}$ follows from the
classical field equations.  The extrinsic curvature
of $\Sigma_t$ is
\begin{equation}
  K_{ij} = \frac{1}{2N}\,\dot{h}_{ij}
    = -\frac{\dot{\varphi}\,e^{-\varphi}}{2N}\,\delta_{ij},
\end{equation}
giving
\begin{equation}\label{eq:K_trace}
  K \equiv h^{ij}K_{ij} = -\frac{3\dot{\varphi}}{2N},
  \qquad
  K_{ij}K^{ij} = \frac{3\dot{\varphi}^2}{4N^2},
\end{equation}
and the kinetic combination
\begin{equation}\label{eq:K_kinetic}
  K_{ij}K^{ij} - K^2 = -\frac{3\dot{\varphi}^2}{2N^2}.
\end{equation}

The spatial Ricci scalar of the conformally flat metric
$h_{ij} = e^{2\omega}\delta_{ij}$ with $\omega = -\varphi/2$
follows from the standard $n$-dimensional conformal
transformation $\bar{R} = e^{-2\omega}[-2(n{-}1)\nabla_0^2\omega
  - (n{-}1)(n{-}2)(\nabla_0\omega)^2]$
evaluated at $n = 3$ and $R_0 = 0$:
\begin{equation}\label{eq:R3}
  {}^{(3)}\!R = e^{\varphi}\Bigl(
    2\,\nabla^2\varphi
    - \tfrac{1}{2}(\nabla\varphi)^2\Bigr),
\end{equation}
where $\nabla^2$ and $\nabla$ denote the flat-space
Laplacian and gradient.

%-------------------------------------------------------------------
\subsection{Reduced action}
\label{sec:reduced_action}
%-------------------------------------------------------------------

Substituting the bi-conformal ADM
quantities~(\ref{eq:adm_biconformal})--(\ref{eq:R3}) into the
full action~(\ref{eq:action}), the ADM Lagrangian density
$\mathcal{L} = (N\sqrt{h}/16\pi G)({}^{(3)}\!R + K_{ij}K^{ij}
  - K^2 + \tfrac{1}{2}\,g^{\mu\nu}\partial_\mu\varphi\,
  \partial_\nu\varphi)$
separates into four terms (using
$g^{00} = -N^{-2}$, $g^{ij} = h^{ij} = e^{\varphi}\delta^{ij}$,
and $N^i = 0$):
\begin{enumerate}
\item[(i)] Einstein--Hilbert kinetic:
  $\frac{N\sqrt{h}}{16\pi G}(K_{ij}K^{ij} - K^2)
  = -\frac{3\,e^{-3\varphi/2}\,\dot{\varphi}^2}{32\pi G N}$;
\item[(ii)] Einstein--Hilbert potential:
  $\frac{N\sqrt{h}}{16\pi G}\,{}^{(3)}\!R
  = \frac{N\,e^{-\varphi/2}}{16\pi G}
    (2\nabla^2\varphi - \tfrac{1}{2}(\nabla\varphi)^2)$;
\item[(iii)] Scalar kinetic:
  $-\frac{\sqrt{h}\,\dot{\varphi}^2}{32\pi G N}
  = -\frac{e^{-3\varphi/2}\,\dot{\varphi}^2}{32\pi G N}$;
\item[(iv)] Scalar gradient:
  $\frac{N\,e^{-\varphi/2}}{32\pi G}(\nabla\varphi)^2$.
\end{enumerate}
A remarkable cancellation occurs: the
$(\nabla\varphi)^2$ contribution from~(ii) is
$-\frac{N\,e^{-\varphi/2}}{32\pi G}(\nabla\varphi)^2$,
which exactly cancels~(iv).  The four time-derivative contributions from~(i)
and~(iii) sum to a single kinetic term, and the
surviving spatial term gives the reduced Lagrangian
density
\begin{equation}\label{eq:L_reduced}
  \mathcal{L} =
    -\frac{e^{-3\varphi/2}}{8\pi G N}\,\dot{\varphi}^2
    + \frac{N\,e^{-\varphi/2}}{8\pi G}\,\nabla^2\varphi.
\end{equation}
The total kinetic coefficient is negative-definite, reflecting
the combined effect of the gravitational trace constraint and
the phantom kinetic sign in the action~(\ref{eq:action}).

%-------------------------------------------------------------------
\subsection{Hamiltonian formulation}
\label{sec:hamiltonian_form}
%-------------------------------------------------------------------

The momentum conjugate to $\varphi$ is
\begin{equation}\label{eq:pi_phi}
  \pi_\varphi \equiv
    \frac{\partial\mathcal{L}}{\partial\dot{\varphi}}
  = -\frac{e^{-3\varphi/2}}{4\pi G N}\,\dot{\varphi}.
\end{equation}
The overall negative sign is inherited from the
negative total kinetic term.  Inverting gives
$\dot{\varphi} = -4\pi G N\,e^{3\varphi/2}\,\pi_\varphi$.

Performing the Legendre transform
$\mathcal{H} = \pi_\varphi\dot{\varphi} - \mathcal{L}$
and expressing the result in terms of phase-space variables:
\begin{equation}\label{eq:H_total}
  H = \int d^3x\;N(\mathbf{x})\,\mathcal{H}_0(\mathbf{x}),
\end{equation}
where the Hamiltonian constraint density is
\begin{equation}\label{eq:H_constraint}
  \mathcal{H}_0
  = -2\pi G\,e^{3\varphi/2}\,\pi_\varphi^2
    - \frac{e^{-\varphi/2}}{8\pi G}\,\nabla^2\varphi
  \approx 0.
\end{equation}
Here $\approx$ denotes weak equality on the constraint surface;
the lapse $N$ enters~(\ref{eq:H_total}) as a Lagrange multiplier
enforcing $\mathcal{H}_0 = 0$ at every spatial point.

\textit{Consistency check.}---For a static configuration
($\pi_\varphi = 0$), the constraint reduces to
$\nabla^2\varphi = 0$, the vacuum field equation
$\Box\varphi = 0$ evaluated on the static bi-conformal
background.  The weak-field solution
$\varphi = -\rs/r$ satisfies $\nabla^2\varphi = 0$
for $r > 0$, confirming consistency with the classical
theory~\cite{Lotishvili2026}.

\textit{Phase-space restriction.}---For dynamical
configurations ($\pi_\varphi \neq 0$), the constraint
requires $\nabla^2\varphi \leq 0$: the kinetic term
$-2\pi G\,e^{3\varphi/2}\pi_\varphi^2$ is non-positive,
so the potential term must be non-negative.  This
condition is automatically preserved by the Hamilton
equations and is satisfied by all physically relevant
configurations---static backgrounds
($\nabla^2\varphi = 0$), gravitationally bound
perturbations (subharmonic near sources), and
cosmological modes (damped oscillations on FLRW).
The constraint thus acts as a self-consistency
selection principle on the space of allowed initial
data, analogous to the positive-energy theorem in
general relativity.

%-------------------------------------------------------------------
\subsection{Constraint structure}
\label{sec:constraint_structure}
%-------------------------------------------------------------------

The bi-conformal ansatz $h_{ij} = e^{-\varphi}\delta_{ij}$
with $N^i = 0$ exhausts the spatial diffeomorphism gauge
freedom: the spatial metric carries no transverse-traceless
modes, and the momentum constraint $\mathcal{H}_i \approx 0$
is satisfied identically.  The sole remaining constraint is
the Hamiltonian constraint~(\ref{eq:H_constraint}), which
generates time reparametrizations via
$\dot{F} = \{F,\,H[N]\}$ for any phase-space functional
$F[\varphi, \pi_\varphi]$.

The algebra of smeared Hamiltonian constraints,
\begin{equation}\label{eq:HH_algebra}
  \{H[N],\,H[M]\} = H_{\mathrm{diff}}[\vec{V}],\qquad
  V^i = h^{ij}(N\,\partial_j M - M\,\partial_j N),
\end{equation}
reproduces the standard Dirac hypersurface-deformation
algebra~\cite{DeWitt1967}.  Since the momentum constraint
vanishes identically in the reduced theory,
$H_{\mathrm{diff}}[\vec{V}] = 0$ and the constraint algebra
is Abelian strongly (not merely on the constraint surface),
confirming that the bi-conformal reduction introduces no
anomalous constraints.

A methodological remark is in order.  The quantization
adopted here is of the ``reduce first, then quantize''
type: the bi-conformal ansatz is imposed classically,
and the resulting single-field system is quantized.  The
alternative Dirac approach---quantize the full metric and
scalar system, then impose
$h_{ij} = e^{-\varphi}\delta_{ij}$ as a quantum
constraint---is known to yield inequivalent results in
general~\cite{Kuchar1981}.  Two considerations favor the
reduced approach for ISPG\@.  First, the bi-conformal
structure is not a gauge choice but a \emph{symmetry
reduction} that defines the theory; it is analogous to
the minisuperspace truncation in quantum cosmology, where
reduced quantization is standard
practice~\cite{Halliwell1991}.  Second, the Abelian
constraint algebra demonstrated above guarantees that no
ordering ambiguities arise in the constraint imposition,
so the two approaches coincide at the semi-classical
level~\cite{Gotay1986}.

%-------------------------------------------------------------------
\subsection{Wheeler--DeWitt equation}
\label{sec:wdw}
%-------------------------------------------------------------------

Following DeWitt~\cite{DeWitt1967}, canonical quantization
promotes the Poisson bracket to the commutator
\begin{equation}\label{eq:comm_phi}
  [\hat{\varphi}(\mathbf{x}),\,
   \hat{\pi}_\varphi(\mathbf{y})]
  = i\hbar\,\delta^{(3)}(\mathbf{x} - \mathbf{y}).
\end{equation}
In the Schr\"odinger representation
$\hat{\pi}_\varphi = -i\hbar\,\delta/\delta\varphi$,
imposing the quantum Hamiltonian constraint
$\hat{\mathcal{H}}_0\,\Psi = 0$ yields the
Wheeler--DeWitt equation for the bi-conformal theory:
\begin{equation}\label{eq:wdw}
  \left[2\pi G\hbar^2\,e^{3\varphi/2}\,
    \frac{\delta^2}{\delta\varphi(\mathbf{x})^2}
  - \frac{e^{-\varphi/2}}{8\pi G}\,
    \nabla^2\varphi(\mathbf{x})
  \right]\Psi[\varphi] = 0.
\end{equation}

The factor $e^{3\varphi/2}$ does not commute with
$\hat{\pi}_\varphi$, giving rise to an operator-ordering
ambiguity.  Equation~(\ref{eq:wdw}) adopts the ordering
$e^{3\varphi/2}\hat{\pi}_\varphi^2$; the Laplace--Beltrami
ordering on the DeWitt supermetric
$\mathcal{G}_{\varphi\varphi} =
  -(4\pi G)^{-1}e^{-3\varphi/2}$
would introduce an additional first-derivative term
$\propto\hbar^2 e^{3\varphi/2}\,
  \delta/\delta\varphi$.
The two orderings agree at the semi-classical level; the
distinction is relevant only in the deep quantum regime
and does not affect the perturbative results of subsequent
sections.

%-------------------------------------------------------------------
\subsection{Semi-classical limit}
\label{sec:semiclassical}
%-------------------------------------------------------------------

Write $\Psi[\varphi] = \mathcal{A}[\varphi]\,
\exp\!\bigl(iS[\varphi]/\hbar\bigr)$ with real
$\mathcal{A}$ and $S$.  Substituting into the
Wheeler--DeWitt equation~(\ref{eq:wdw}) and noting
that $\nabla^2\varphi$ acts multiplicatively on $\Psi$
while $\delta^2/\delta\varphi^2$ produces
\begin{equation}
  \frac{\delta^2\Psi}{\delta\varphi^2}
  = \left[\frac{\delta^2\mathcal{A}}{\delta\varphi^2}
    + \frac{2i}{\hbar}\frac{\delta\mathcal{A}}{\delta\varphi}
      \frac{\delta S}{\delta\varphi}
    + \frac{i}{\hbar}\mathcal{A}
      \frac{\delta^2 S}{\delta\varphi^2}
    - \frac{\mathcal{A}}{\hbar^2}
      \!\left(\frac{\delta S}{\delta\varphi}\right)^{\!2}
  \right]e^{iS/\hbar},
\end{equation}
we collect powers of $\hbar$.

\textit{Order $\hbar^0$ (Hamilton--Jacobi).}---The leading
term gives
\begin{equation}\label{eq:HJ}
  -2\pi G\,e^{3\varphi/2}
    \!\left(\frac{\delta S}{\delta\varphi}\right)^{\!2}
  - \frac{e^{-\varphi/2}}{8\pi G}\,\nabla^2\varphi
  = 0.
\end{equation}
Identifying $\pi_\varphi = \delta S/\delta\varphi$, this
is exactly the classical Hamiltonian
constraint~(\ref{eq:H_constraint}).  The corresponding
classical equations of motion, derived from the
action~(\ref{eq:action}) in Sec.~\ref{sec:wave_eq},
yield the wave equation
\begin{equation}
  e^{-\varphi_0}\,\ddot{\delta\varphi}
  = e^{\varphi_0}\,\nabla^2\delta\varphi
\end{equation}
for perturbations around a static background
$\varphi_0$ with $\nabla^2\varphi_0 = 0$.

\textit{Order $\hbar^1$ (continuity).}---The
next-to-leading order yields
\begin{equation}\label{eq:continuity}
  \frac{\delta}{\delta\varphi}\!\left(
    \mathcal{A}^2\,\frac{\delta S}{\delta\varphi}
  \right) = 0,
\end{equation}
which is the continuity equation for the probability
density $|\Psi|^2 = \mathcal{A}^2$ in configuration
space, ensuring conservation of probability in the
semi-classical regime.

The semi-classical quantization of
Sec.~\ref{sec:excitations} is therefore the correct
leading approximation to the full quantum theory.

%-------------------------------------------------------------------
\subsection{Physical interpretation}
\label{sec:physical_interp}
%-------------------------------------------------------------------

The canonical analysis reveals three structural features
of the bi-conformal quantum theory:
\begin{enumerate}
\item \textit{Degree-of-freedom reduction.}  The
  bi-conformal ansatz collapses the six components of
  $h_{ij}$ into a single scalar $\varphi$.  The
  gravitational \emph{configuration} space is one-dimensional
  per spatial point (phase space two-dimensional,
  $(\varphi,\pi_\varphi)$), dramatically simplifying the
  constraint structure relative to full quantum gravity.

\item \textit{Phantom kinetic sign.}  The negative total
  kinetic term in~(\ref{eq:L_reduced}) reflects the phantom
  nature of $\varphi$.  In the Wheeler--DeWitt
  equation~(\ref{eq:wdw}), this converts the kinetic
  operator from elliptic to hyperbolic type, which is the
  feature that allows the constraint
  $\hat{\mathcal{H}}_0\Psi = 0$ to admit non-trivial,
  normalizable solutions with oscillatory behaviour
  in the scalar-field direction, enabling a well-posed
  initial-value formulation in superspace.  Moreover, the
  hyperbolic signature provides a natural candidate for an
  internal time variable in superspace---a structure absent
  in standard (elliptic) quantum gravity, where the
  ``problem of time'' admits no straightforward resolution.

\item \textit{Consistency with the classical theory.}
  The Hamiltonian constraint reproduces the vacuum field
  equation in the static limit, the semi-classical
  expansion recovers the wave
  equation~(\ref{eq:wave_rearranged}), and the
  hypersurface-deformation algebra matches that of general
  relativity.  The all-orders quantum stability of the
  phantom sector within the pure-gravity constrained sector
  is argued for in Sec.~\ref{sec:stability} (full
  bracket-closure proof deferred; matter-coupled extension
  via the EFT completion).

\item \textit{Ultralocality of the kinetic sector.}  The
  exact cancellation of the $(\nabla\varphi)^2$ terms in
  the reduced action~(\ref{eq:L_reduced}) renders the
  kinetic energy ultralocal: different spatial points are
  dynamically independent at the kinetic level.  Spatial
  coupling re-enters exclusively through the
  $\nabla^2\varphi$ term in the Hamiltonian
  constraint~(\ref{eq:H_constraint}).  This property
  implies a natural lattice regularization in which each
  spatial site carries an independent copy of the
  one-dimensional quantum mechanics of $\varphi$, coupled
  to its neighbours only through the discrete Laplacian
  in the constraint---a structure well suited to the
  numerical eigenvalue problem for the mass spectrum
  (Sec.~\ref{sec:mass_spectrum}).
\end{enumerate}

%===================================================================
\section{Scalar-field excitations as particles}\label{sec:excitations}
%===================================================================

The central thesis of this work is that elementary particles are
localized resonant excitations of the scalar-metric medium.
In this section we make this identification precise by
quantizing the scalar field on the bi-conformal background,
deriving the excitation spectrum, and establishing the
position-dependent effective mass formula at the quantum level.

%-------------------------------------------------------------------
\subsection{Wave equation on the bi-conformal background}
\label{sec:wave_eq}
%-------------------------------------------------------------------

Consider perturbations $\delta\varphi$ of the scalar field around a
static background $\varphi_0(\mathbf{x})$:
\begin{equation}\label{eq:phi_decomp}
  \varphi = \varphi_0(\mathbf{x}) + \delta\varphi(t,\mathbf{x}).
\end{equation}
The background satisfies the vacuum field equation
$\Box\varphi_0 = 0$, and the perturbation satisfies the
linearized equation
\begin{equation}\label{eq:lin_wave}
  \Box_{\varphi_0}\,\delta\varphi = 0,
\end{equation}
where $\Box_{\varphi_0}$ is the d'Alembertian evaluated on
the background metric
$g_{\mu\nu}(\varphi_0) =
\mathrm{diag}(-e^{\varphi_0},\,e^{-\varphi_0},\,
e^{-\varphi_0},\,e^{-\varphi_0})$.

Computing the d'Alembertian explicitly (in Cartesian spatial
coordinates with $c = 1$):
\begin{equation}
  \sqrt{-g} = e^{-\varphi_0}, \qquad
  g^{00} = -e^{-\varphi_0}, \qquad
  g^{ij} = e^{\varphi_0}\,\delta^{ij},
\end{equation}
and noting that $\sqrt{-g}\,g^{ij} = e^{-\varphi_0}\cdot
e^{\varphi_0}\,\delta^{ij} = \delta^{ij}$
(the spatial prefactors cancel exactly), we obtain
\begin{equation}\label{eq:box_explicit}
  \Box_{\varphi_0}\,\delta\varphi
  = e^{\varphi_0}\Bigl[
    -e^{-2\varphi_0}\,\ddot{\delta\varphi}
    + \nabla^{2}\delta\varphi
  \Bigr] = 0,
\end{equation}
where $\ddot{f} \equiv \partial^2 f/\partial t^2$ and
$\nabla^2$ is the flat-space Laplacian.  The cancellation
of spatial prefactors is a remarkable property of the
bi-conformal structure: it ensures that the spatial part
of the wave operator is the \emph{ordinary} Laplacian,
with the entire effect of gravity concentrated in the
temporal prefactor.  Rearranging:
\begin{equation}\label{eq:wave_rearranged}
  \boxed{
  e^{-\varphi_0}\,\ddot{\delta\varphi}
  = e^{\varphi_0}\,\nabla^{2}\delta\varphi.
  }
\end{equation}
This gives the local propagation speed
\begin{equation}\label{eq:c_coord}
  c_{\mathrm{coord}}^{2} = e^{2\varphi_0},
\end{equation}
which varies with position.  However, the locally measured
speed---using proper time $d\tau = e^{\varphi_0/2}\,dt$ and
proper length $d\ell = e^{-\varphi_0/2}\,dr$---is
\begin{equation}\label{eq:c_local}
  c_{\mathrm{local}} = \frac{d\ell}{d\tau}
  = \frac{e^{-\varphi_0/2}}{e^{\varphi_0/2}}\,
    \frac{dr}{dt}
  = e^{-\varphi_0}\cdot e^{\varphi_0} = 1 = c,
\end{equation}
confirming that local Lorentz invariance is preserved for
all perturbation modes, as established in~\cite{Lotishvili2026}.

%-------------------------------------------------------------------
\subsection{Mode decomposition in the static field}
\label{sec:modes}
%-------------------------------------------------------------------

For the spherically symmetric background
$\varphi_0 = -\rs/r$ generated by a mass $M$,
we seek time-harmonic solutions
$\delta\varphi = \psi(\mathbf{x})\,e^{-i\omega t}$.
Substituting into Eq.~(\ref{eq:wave_rearranged}):
\begin{equation}\label{eq:helmholtz}
  \nabla^{2}\psi
  + \omega^{2}\,e^{-2\varphi_0}\,\psi = 0.
\end{equation}
For $\varphi_0 = -\rs/r$, the exponential
is $e^{-2\varphi_0} = e^{2\rs/r}$.
Expanding in spherical harmonics,
$\psi(r,\theta,\phi) = \sum_{lm} u_l(r)/r \cdot Y_{lm}(\theta,\phi)$,
the radial function $u_l$ satisfies the
Schr\"odinger-like equation
\begin{equation}\label{eq:schrodinger}
  \boxed{
  u_l'' + \left[\omega^{2}\,e^{2\rs/r}
    - \frac{l(l+1)}{r^{2}}\right] u_l = 0,
  }
\end{equation}
where $u_l'' \equiv d^2 u_l/dr^2$.
The key feature is the exponential factor
$e^{2\rs/r}$ multiplying $\omega^2$, which acts as a
position-dependent effective frequency:
modes are gravitationally \emph{blue-shifted}
near the source ($r \to \rs$, $e^{2\rs/r} \gg 1$)
and approach the asymptotic value ($e^{2\rs/r} \to 1$)
far from the source ($r \to \infty$).

In the weak-field limit ($\rs/r \ll 1$), the exponential
can be expanded:
\begin{equation}\label{eq:weak_exp}
  e^{2\rs/r} \approx 1 + \frac{2\rs}{r}
  + \frac{2\rs^2}{r^2} + \cdots\,,
\end{equation}
and the radial equation reduces to
\begin{equation}\label{eq:radial_weak}
  u_l'' + \left[\omega^{2}
    + \frac{2\omega^2\rs}{r}
    - \frac{l(l+1)}{r^{2}} + \cdots\right] u_l = 0.
\end{equation}
This is formally identical to the radial Schr\"odinger
equation for a particle of energy $\omega^2$ in an
\emph{attractive} Coulomb-like potential
$-2\omega^2\rs/r$.  The analogy is \emph{structurally}
exact: the ISPG radial equation shares the same form as
the non-relativistic Schr\"odinger equation for a particle
in a Coulomb potential, though the coupling constants here
relate $\omega$ (scalar-wave energy) to $GM$ rather than
a particle mass $m$ to $GM$.  Modes of the scalar field
in the gravitational potential therefore experience an
effective attraction toward the source.

For $l=0$ (s-wave), the solutions in the asymptotic
region ($r \gg \rs$) are spherical waves:
\begin{equation}
  u_0(r) \sim \sin(\omega r + \delta_0),
\end{equation}
where $\delta_0$ is the gravitationally induced
phase shift.  Near the source, the amplitude
grows as $e^{\rs/r}$ -- an exponential
\emph{amplitude} enhancement of the mode near the
source; the corresponding mode energy as seen from
infinity is given by Eq.~\eqref{eq:E_inf} and is in
fact \emph{reduced} (gravitational redshift), not
blue-shifted.

%-------------------------------------------------------------------
\subsection{Energy of excitations and the effective mass}
\label{sec:mode_energy}
%-------------------------------------------------------------------

We now establish the central result connecting mode
energy to the effective mass formula.

\textit{Coordinate and proper frequencies.}---A mode
oscillating as $\delta\varphi \propto e^{-i\omega t}$
has coordinate frequency $\omega$.  A local observer
at position $\varphi_0$ measures proper time
$\tau = e^{\varphi_0/2}\,t$, so the locally measured
frequency is
\begin{equation}\label{eq:omega_local}
  \omega_{\mathrm{local}}
  = \frac{d(\omega t)}{d\tau}
  = \omega\,e^{-\varphi_0/2}.
\end{equation}

\textit{Local rest energy.}---The rest energy of a
localized excitation, as measured in its own frame, is
\begin{equation}\label{eq:E_local}
  E_{\mathrm{local}} = \hbar\,\omega_{\mathrm{local}}
  = \hbar\omega\,e^{-\varphi_0/2}
  \equiv m_0 c^{2},
\end{equation}
which defines the bare mass $m_0$ as the locally
measured rest mass.  This is the mass that appears in
all local experiments: atomic spectra, Compton scattering,
and particle decays.  It is independent of position
because the definition is purely local.

\textit{Energy at infinity.}---The energy of the same
excitation as measured by an asymptotic observer
($\varphi \to 0$) is simply
\begin{equation}\label{eq:E_inf}
  E_{\infty} = \hbar\omega = m_0 c^{2}\,e^{\varphi_0/2},
\end{equation}
which yields the effective mass
\begin{equation}\label{eq:meff_quantum}
  \boxed{
  \meff(\varphi_0) = \frac{E_{\infty}}{c^2}
  = m_0\,e^{\varphi_0/2}.
  }
\end{equation}
This reproduces Eq.~(I.meff) of the classical
theory~\cite{Lotishvili2026}, now derived from the
quantum-mechanical energy of a scalar-field excitation
rather than from the geodesic worldline interval.

\textit{Physical interpretation.}---In a region of lower
pressure ($\varphi_0 < 0$), the resonant mode carries
less energy as seen from infinity.  The mode ``vibrates
on a looser string'': the reduced background tension
(pressure) diminishes the amplitude and energy of the
excitation.  This is the quantum-mechanical realization
of the guitar-string analogy: a string under less tension
produces a softer note (less energy), even though the
local pitch (identity of the note, i.e., the particle
species) remains the same.

%-------------------------------------------------------------------
\subsection{Local undetectability: the quantum formulation}
\label{sec:local_undet}
%-------------------------------------------------------------------

The classical theory established that the ratio of any
two masses is position-independent,
$\meff^A/\meff^B = m_0^A/m_0^B$~\cite{Lotishvili2026}.
We now show that this extends to \emph{all} quantum
observables.

\textit{Transition frequencies.}---Consider an atomic
transition between two energy levels $E_1$ and $E_2$.
The emitted photon frequency as measured locally is
\begin{equation}
  \nu_{\mathrm{local}}
  = \frac{E_{2,\mathrm{local}} - E_{1,\mathrm{local}}}{h}
  = \frac{(E_2 - E_1)_0}{h},
\end{equation}
since both levels scale by the same factor
$e^{\varphi_0/2}$, which cancels in the difference
measured by local instruments.  A local spectroscopist
obtains the same spectrum regardless of the value
of $\varphi_0$.

\textit{Compton wavelength.}---The Compton wavelength
$\lambda_C = h/(m c)$ scales as
\begin{equation}
  \lambda_{C,\mathrm{local}}
  = \frac{h}{\meff\,c}
  \cdot e^{-\varphi_0/2}
  = \frac{h}{m_0\,c}\cdot
    \frac{e^{-\varphi_0/2}}{e^{\varphi_0/2}}
  = \frac{h}{m_0 c}\,e^{-\varphi_0}.
\end{equation}
But the local ruler also scales: proper length
$\ell_{\mathrm{local}} = e^{-\varphi_0/2}\,\ell_{\mathrm{coord}}$.
The Compton wavelength measured in local ruler units is
\begin{equation}
  \frac{\lambda_{C,\mathrm{coord}}}
       {e^{-\varphi_0/2}}
  = \frac{h/(m_0\,c\,e^{\varphi_0/2})}
         {e^{-\varphi_0/2}}
  = \frac{h}{m_0\,c},
\end{equation}
which is the standard, position-independent result.

\textit{De Broglie wavelength.}---A particle with
local momentum $p_{\mathrm{local}}$ has de Broglie
wavelength $\lambda_{\mathrm{dB}} = h/p_{\mathrm{local}}$.
Since momentum transforms as
$p_{\mathrm{local}} = p_{\mathrm{coord}}\,e^{\varphi_0/2}$
and the ruler transforms as $e^{-\varphi_0/2}$, the
locally measured wavelength is
\begin{equation}
  \lambda_{\mathrm{dB,local}}
  = \frac{h}{p_{\mathrm{local}}}
    \cdot \frac{1}{e^{-\varphi_0/2}}
  = \frac{h}{p_0}
\end{equation}
(where $p_0$ is the bare momentum), confirming that
the de~Broglie relation $\lambda = h/p$ holds
locally with no $\varphi$-dependent corrections.

\textit{General theorem.}---All dimensionless quantum
observables (branching ratios, coupling constants,
cross-section ratios) are invariant under the
bi-conformal rescaling $\varphi_0 \to \varphi_0 + \delta$.
This follows from the fact that the bi-conformal
transformation is a simultaneous rescaling of
$\{m, E, \omega\} \to e^{\delta/2}\{m, E, \omega\}$ and
$\{\ell, t\} \to e^{-\delta/2}\{\ell, t\}$,
which preserves all dimensionless combinations.

This is the \textbf{quantum local undetectability principle}:
\begin{quote}
\textit{No quantum experiment confined to a local laboratory
can detect the value of the background scalar field $\varphi_0$.
The bi-conformal rescaling of masses, frequencies, and
length scales is simultaneously absorbed by the rescaling
of all measuring apparatus.}
\end{quote}

%-------------------------------------------------------------------
\subsection{Self-consistent resonances and particle stability}
\label{sec:resonances}
%-------------------------------------------------------------------

The linearized treatment of Sec.~\ref{sec:modes} describes
freely propagating scalar excitations.  \emph{Particles},
by contrast, are localized, self-sustaining resonant
structures.  Their existence requires the nonlinear
self-interaction inherent in the bi-conformal coupling
between $\varphi$ and $g_{\mu\nu}$.

A localized excitation $\delta\varphi$ modifies the metric
through Eq.~(\ref{eq:metric}), and the modified metric
in turn affects the propagation of $\delta\varphi$
through Eq.~(\ref{eq:wave_rearranged}).  A self-consistent
solution satisfies both simultaneously.  Specifically,
we seek time-periodic, localized solutions of the form
\begin{equation}\label{eq:oscillon}
  \varphi(t,r) = \Phi(r)\cos(\omega t) + \order{\Phi^2},
\end{equation}
where $\Phi(r) \to 0$ as $r \to \infty$.

To make the metric backreaction explicit, we expand the
full field equation $\Box\varphi = 0$ to quadratic order
in $\Phi$, including the $\varphi$-dependent metric
coefficients.  The d'Alembertian on the bi-conformal
background (Eq.~\ref{eq:box_explicit}) yields, for
$\varphi = \varphi_0 + \Phi\cos(\omega t)$:
\begin{equation}\label{eq:box_expanded}
  \Box\varphi = e^{\varphi_0}\bigl[\nabla^2\Phi + \omega^2 e^{-2\varphi_0}\Phi\bigr]\cos(\omega t) + \order{\Phi^2} = 0.
\end{equation}
At order $\Phi^2$, the backreaction modifies the effective
potential for $\Phi$ through the exponential factors.  The
self-consistency condition becomes the \emph{nonlinear
Helmholtz equation}
\begin{equation}\label{eq:self_consistent}
  \nabla^{2}\Phi - U'(\Phi) + \omega^{2}e^{-2\varphi_0}\Phi
  + \frac{1}{2}\omega^2 e^{-2\varphi_0}\Phi^{2}
  - \varphi_0'\,\Phi' + \order{\Phi^3} = 0,
\end{equation}
where $\varphi_0' \equiv d\varphi_0/dr$ is the background
gradient and $U'(\Phi)$ represents the effective potential
from the metric coupling.  The term $-\varphi_0'\Phi'$
acts as a position-dependent damping that localizes the
solution.

\textit{Effective potential from metric backreaction.}---
The potential term $U'(\Phi)$ arises from the expansion
of the exponential metric factors $e^{-2\varphi}$ to
quadratic order in $\Phi$.  Expanding $e^{-2\varphi} =
e^{-2\varphi_0}(1 - 2\Phi\cos\omega t + \order{\Phi^2})$
in the time-periodic d'Alembertian and time-averaging over
one oscillation period yields an effective static potential:
\begin{equation}\label{eq:U_effective}
  U'(\Phi) = -\frac{1}{2}\omega^2 e^{-2\varphi_0}\Phi
  + \order{\Phi^2},
\end{equation}
where the linear term represents the frequency shift due
to the background pressure gradient.  The full nonlinear
form of $U(\Phi)$ requires solving the coupled scalar-metric
equations self-consistently.

\textit{Epistemic note on the linear sector.}---The leading
term in Eq.~(\ref{eq:U_effective}) carries a negative sign,
so $U'(\Phi)\propto -\Phi$ at leading order: the quadratic
part of $U(\Phi)$ is \emph{tachyonic} ($m_{\mathrm{eff}}^2<0$).
A purely linear analysis therefore admits no bound state --
stabilization of a localized configuration relies essentially
on the nonlinear backreaction $+\tfrac{1}{2}\omega^2
e^{-2\varphi_0}\Phi^2$ entering
Eq.~(\ref{eq:self_consistent}), not on a positive quadratic
confining potential.  This distinguishes the ISPG oscillon
from a bound state in a static well and aligns the sector
with the Gleiser-type nonlinear-Klein--Gordon oscillons
referenced below.

\textit{Connection to the 1D oscillon model.}---To demonstrate
that localized solutions exist, we derive a simplified 1D model
from the ISPG field structure.  Near the oscillon center
($r \ll \rs$), the background $\varphi_0 \approx 0$ and the
exponential factors $e^{-2\varphi_0} \approx 1$.  Defining
dimensionless variables
\begin{equation}
  x = \frac{r}{R_{\mathrm{core}}}, \qquad
  \Phi = \frac{\phi}{\phi_0}, \qquad
  \Omega^2 = \omega^2 R_{\mathrm{core}}^2,
\end{equation}
where $R_{\mathrm{core}} = \hbar/(mc)$ is the Compton natural
scale (a \emph{post-hoc} identification, fixed once $m$ is
known from the continuation below rather than supplied as an
independent input -- see also Eq.~(\ref{eq:U_effective}) and
the cyclicity remark in the following subsection), and
$\phi_0$ is a \emph{dimensionless} amplitude convention used
to canonicalize the rescaling, not the physical weak-field
value (which satisfies $\varphi\ll 1$ for elementary sources
of mass $m\ll M_{\mathrm{Pl}}$).  In these variables,
Eq.~(\ref{eq:self_consistent}) reduces in the 1D (planar)
limit to
\begin{equation}\label{eq:1d_oscillon}
  \frac{d^2\Phi}{dx^2} - \Phi + \Phi^2 + \Omega^2\Phi = 0.
\end{equation}

The nonlinear term $\Phi^2$ arises from the quadratic backreaction
in Eq.~(\ref{eq:self_consistent}); the effective potential
$U(\Phi) = \frac{1}{2}(1-\Omega^2)\Phi^2 - \frac{1}{3}\Phi^3$
has a local minimum at $\Phi = 0$ (oscillon core) separated by
a barrier from the runaway region at large $\Phi$.  This
equation admits the localized solution
\begin{equation}
  \Phi(x) = A\,\mathrm{sech}^2(Bx),
\end{equation}
with $A = 3(1-\Omega^2)/2$ and $B = \sqrt{1-\Omega^2}/2$,
provided $\Omega^2 < 1$ (bound state condition).  This is
the standard oscillon profile~\cite{Bogolyubsky1976,Gleiser1994,Copeland1995}.

\textit{Epistemic status of the reduction.}---We emphasize
that the passage from Eq.~(\ref{eq:self_consistent}) to
Eq.~(\ref{eq:1d_oscillon}) is an \emph{identification} of
the ISPG nonlinear Helmholtz sector with the canonical
Gleiser-type oscillon equation, obtained through the
dimensionless rescaling $\{x,\phi_0,\Omega\}$ that absorbs
the effective coefficients of the linear and quadratic terms
in Eq.~(\ref{eq:self_consistent}) into the canonical form
$-\Phi+\Phi^{2}+\Omega^{2}\Phi$.  It is \emph{not} a
sign-preserving derivation in the original variables: the
negative linear coefficient $-\Phi$ of
Eq.~(\ref{eq:1d_oscillon}) reflects the canonical oscillon
rescaling, not a direct reading of the coefficient of
$\Phi$ in Eq.~(\ref{eq:self_consistent}).  The substantive
content carried by this identification is that the
nonlinear sector of ISPG falls into the same mathematical
class as the well-studied Bogolyubsky--Gleiser--Copeland
oscillons, for which localized sech$^{2}$-type cores and
3D spherical continuation are established; the gravitational
construction below uses only this existence of localized
finite-energy sources, not the specific canonical
coefficients.

\textit{3D spherical generalization.}---The full 3D equation
adds the radial Laplacian $\frac{2}{r}\Phi'$ and the damping
$-\varphi_0'\Phi'$ from the background gradient.  Numerical
boundary-value solutions of the spherical problem exhibit
regular localized profiles, with the core structure
well-approximated by the 1D solution and the tail matching
the outgoing behavior discussed in
Sec.~\ref{sec:entropy_mech}.  The restricted existence claim
needed below is therefore this: the nonlinear bi-conformal
equations admit localized finite-energy source configurations,
so the scalar-metric medium contains concrete bound sources
rather than only freely propagating waves.  For the
gravitational construction developed below, this is the
essential output of the nonlinear sector.  The stronger
identification of these solutions with the full
Standard-Model spectrum is logically separate and is not used
in the derivation of gravity.

More concretely, the spherical boundary-value problem admits
regular profiles under continuation in the central amplitude
$\Phi_0$, and linearization about these backgrounds yields a
discrete cavity spectrum rather than a continuum of delocalized
plane-wave modes.  This is the precise numerical content needed
in the present section: not a complete nonlinear time-evolution
theorem, but the existence of localized source backgrounds on
which the gravitational dressing and its long-range field can be
built.

The stationary source statement used below can therefore be
written in theorem-like form.  Consider the spherical coupled
boundary-value problem with central amplitude
$\Phi(0)=\Phi_c$, regularity conditions
$\Phi'(0)=\varphi_0'(0)=0$, and localization conditions
$\Phi(r_{\max})=\varphi_0(r_{\max})=0$ at large $r_{\max}$.
Numerical continuation in $\Phi_c$ produces smooth branches
$(\Phi(r;\Phi_c),\varphi_0(r;\Phi_c),\Omega(\Phi_c))$ for which
\begin{equation}
  \Phi(r) = \Phi_c + \order{r^2}, \qquad
  \varphi_0(r) = \varphi_0(0) + \order{r^2}
  \quad (r \to 0),
\end{equation}
so the source is regular at the origin.  On the bound branch
($\Omega^2 < 1$),
\begin{equation}
  4\pi\!\int_0^\infty\!\Bigl[
    \tfrac{1}{2}\Omega^2\Phi^2
    + \tfrac{1}{2}(\Phi')^2
  \Bigr]r^2\,dr < \infty,
  \qquad
  \Phi(r) \sim A\,e^{-\kappa r}/r,
\end{equation}
so the source energy is finite and localized.  Linearization
about each such background yields isolated cavity eigenvalues
below threshold rather than a continuum of delocalized modes.
This is the precise theorem-like numerical statement used in the
gravity section: ISPG contains a stationary spherical source
sector of regular finite-mass oscillon backgrounds with a
discrete small-oscillation spectrum.  The continuation in
$\Phi_c$ and the associated discrete cavity spectrum are
verified numerically by the companion scripts
\texttt{QUANTUM/Python/test\_3d\_bvp.py},
\texttt{test\_cavity\_spectrum.py},
\texttt{test\_ispg\_oscillon.py}, and
\texttt{test\_phi3\_cont.py}, which provide the explicit
$(\Phi(r;\Phi_c),\Omega(\Phi_c))$ branches and
small-oscillation eigenvalues referred to here.
This stationary statement is
the base point for the modulation-control argument developed below;
a full $3D+1$ nonlinear orbital-stability theorem would extend the
same source control from the present metastable regime to fully
global evolution.

At the level needed in this section, the localized solutions
fall into two classes:
\begin{itemize}
  \item A \textbf{bound oscillon} is a regular finite-energy
    configuration with an exponentially localized core.
  \item A \textbf{quasi-bound oscillon} is a localized core
    accompanied by an outgoing tail, so that energy slowly
    leaks into propagating modes.
  \item The \textbf{decay width} $\Gamma$ of a quasi-bound
    configuration is the imaginary part of the quasi-normal
    frequency $\omega = \omega_R - i\Gamma/2$, with lifetime
    $\tau = \hbar/\Gamma$.
\end{itemize}

\textit{Quasi-normal mode analysis.}---For a quasi-bound
oscillon, the scalar field has the time dependence
$\Phi(t,r) = \phi(r) e^{-i\omega t}$ with complex frequency
$\omega = \omega_R - i\Gamma/2$.  The spatial profile $\phi(r)$
satisfies the radial equation with outgoing boundary condition
at infinity:
\begin{equation}\label{eq:quasi_normal}
  \frac{d^2\phi}{dr^2} + \frac{2}{r}\frac{d\phi}{dr}
  + \left[\omega^2 e^{-2\varphi_0(r)} - V_{\mathrm{eff}}(r)\right]\phi = 0,
\end{equation}
where $V_{\mathrm{eff}}(r)$ is the effective potential from
the nonlinear self-interaction.  The outgoing condition
$\phi(r) \sim e^{i\omega r}/r$ as $r \to \infty$ makes the
frequency complex.

For a shallow bound state (weak nonlinearity), the decay width
can be estimated by the WKB tunneling formula:
\begin{equation}\label{eq:gamma_wkb}
  \Gamma \sim \omega_R \exp\left[-2\int_{r_1}^{r_2}
    \sqrt{V_{\mathrm{eff}}(r) - \omega_R^2 e^{-2\varphi_0(r)}}\,dr\right],
\end{equation}
where $r_1$ and $r_2$ are the classical turning points.  The
width is exponentially suppressed for deep bound states
($V_{\mathrm{max}} \gg \omega_R^2$), corresponding to long-lived
particles like the muon ($\tau \sim 2.2\,\mu\mathrm{s}$).  For
shallow states or high-mass resonances, the barrier penetration
is efficient and the lifetime short (e.g., W/Z bosons with
$\tau \sim 10^{-25}$~s).

The stability of a resonant mode depends on the background
pressure $\varphi_0$.  At higher pressure
($\varphi_0$ closer to zero), more energetic modes can be
sustained---heavier particles are stable.  As the pressure
decreases over cosmic time (Sec.~\ref{sec:entropy}),
previously stable modes may become quasi-bound and
eventually decay.  This provides a natural explanation
for the evolving particle content of the universe:
in the early universe, where the pressure was far higher,
particles that are unstable today (e.g., the quark-gluon
plasma phase) were stable, as confirmed by heavy-ion
experiments at CERN's ALICE facility~\cite{ALICE2010}.

\textit{Metastable modulation regime.}---For the gravity
construction below, the dynamical statement required is weaker
than a full orbital-stability theorem but stronger than the
mere existence of stationary profiles.  Equations
(\ref{eq:gamma_stable}) and (\ref{eq:a_n_eom}) imply that for
deep bound or quasi-bound states with $\Gamma \ll \omega_R$,
there exists a parametrically large time window
\begin{equation}
  \omega_R^{-1} \ll t \ll \Gamma^{-1},
\end{equation}
in which the evolution remains close to the finite-dimensional
oscillon manifold generated by translations and slow motion
along the bound branch:
\begin{equation}
  \varphi(\mathbf{x},t)
  \simeq \varphi_{\mathrm{grav}}
    \bigl(|\mathbf{x}-\mathbf{R}(t)|;\Phi_c(t)\bigr)
  + \Phi_0\bigl(|\mathbf{x}-\mathbf{R}(t)|;\Phi_c(t)\bigr)
    \cos\Theta(t)
  + \eta(\mathbf{x},t).
\end{equation}
Here $\mathbf{R}(t)$ is the collective position, $\Theta(t)$ the
fast internal phase, and $\eta$ the radiative remainder carried
by damped quasi-normal modes.  The translational channel is the
$\ell = 1$ Goldstone mode proven below, while leakage into
non-bound channels is controlled by the exponentially small
width $\Gamma$.  Thus, on all timescales relevant to the
Newtonian far-zone field, the oscillon behaves as a
\emph{metastable particle source}: its center and phase may
modulate adiabatically, and its parameters drift only on the
long leakage timescale $\Gamma^{-1}$, while the gravitational
dressing follows the same collective coordinates.  A full
$3D+1$ nonlinear orbital-stability theorem would strengthen this
metastable modulation statement to a global theorem, but is not
required for the source concept used in the gravity section.

\textit{Neutral versus damped channels.}---The linearized
dynamical channels relevant to the source concept therefore
split into two classes.  First, the translational $\ell = 1$
Goldstone mode is \emph{neutral}: it shifts the oscillon center
without changing the internal energy, so it moves the source
along the manifold rather than away from it.  Second, the
non-Goldstone perturbations admit the quasi-normal expansion
discussed below, with mode amplitudes satisfying
\begin{equation}
  \ddot{a}_n + \Gamma_n\,\dot{a}_n + \omega_n^2 a_n = f_n(t),
\end{equation}
where $\Gamma_n \ge 0$ measures leakage into radiative channels.
For stable particles, $\Gamma_n \ll \omega_n$ and slow forcing
couples dominantly to the lowest mode, so higher overtones are
suppressed rather than amplified.  Thus, modulo the neutral
symmetry directions, the oscillon source sector is
\emph{perturbatively attractive} on the timescale relevant to the
gravitational field: deformations damp or remain adiabatically
slaved to the collective coordinates instead of generating
runaway growth.

\textit{Gravity-use source-control statement.}---The dynamical claim
actually needed in Sec.~\ref{sec:bernoulli_identity} can now be
stated in modulation form.  Let
\begin{equation}\label{eq:oscillon_manifold}
  \mathcal{M}_{\mathrm{osc}}
  \equiv \Bigl\{
    \varphi_{\mathrm{grav}}\bigl(|\mathbf{x}-\mathbf{R}|;\Phi_c\bigr)
    + \Phi_0\bigl(|\mathbf{x}-\mathbf{R}|;\Phi_c\bigr)\cos\Theta
  \Bigr\}
\end{equation}
denote the finite-dimensional oscillon manifold generated by
translations, motion along the bound branch, and phase shifts of the
leading harmonic.  For a field configuration that remains near this
manifold in the metastable regime, write
\begin{equation}\label{eq:modulated_source_decomposition}
  \varphi(\mathbf{x},t)
  = \varphi_{\mathrm{grav}}
      \bigl(|\mathbf{x}-\mathbf{R}(t)|;\Phi_c(t)\bigr)
  + \Phi_0\bigl(|\mathbf{x}-\mathbf{R}(t)|;\Phi_c(t)\bigr)\cos\Theta(t)
  + \eta(\mathbf{x},t).
\end{equation}
Choose $\mathbf{R}(t)$, $\Phi_c(t)$, and $\Theta(t)$ so that the
remainder is orthogonal to the tangent directions of
$\mathcal{M}_{\mathrm{osc}}$:
\begin{equation}\label{eq:source_orthogonality}
  \langle \eta, T_i \rangle = 0,\qquad
  \langle \eta, T_c \rangle = 0,\qquad
  \langle \eta, T_\Theta \rangle = 0,
\end{equation}
where
\begin{align}
  T_i &\equiv \partial_{R_i}\Bigl[
    \varphi_{\mathrm{grav}}\bigl(|\mathbf{x}-\mathbf{R}|;\Phi_c\bigr)
    + \Phi_0\bigl(|\mathbf{x}-\mathbf{R}|;\Phi_c\bigr)\cos\Theta
  \Bigr],\\
  T_c &\equiv \partial_{\Phi_c}\Bigl[
    \varphi_{\mathrm{grav}}\bigl(|\mathbf{x}-\mathbf{R}|;\Phi_c\bigr)
    + \Phi_0\bigl(|\mathbf{x}-\mathbf{R}|;\Phi_c\bigr)\cos\Theta
  \Bigr],\\
  T_\Theta &\equiv \partial_\Theta\Bigl[
    \Phi_0\bigl(|\mathbf{x}-\mathbf{R}|;\Phi_c\bigr)\cos\Theta
  \Bigr],
\end{align}
and $\langle f,g\rangle \equiv \int_{\mathbb{R}^3}
f(\mathbf{x})g(\mathbf{x})\,d^3x$.  At leading harmonic order,
$T_i$ and $T_c$ reduce to the familiar translation and branch
directions $\partial_i\Phi_0$ and $\partial_{\Phi_c}\Phi_0$, while
$T_\Theta$ is the phase tangent.  Thus the neutral collective
coordinates are absorbed into $\mathbf{R}(t)$, $\Phi_c(t)$, and
$\Theta(t)$ rather than double-counted inside $\eta$.

Projecting the linearized dynamics onto the complement of these
tangent directions gives a discrete quasi-normal expansion
\begin{equation}\label{eq:eta_mode_expansion}
  \eta(\mathbf{x},t)
  = \sum_{n\notin \mathrm{neutral}} a_n(t)\,\phi_n(\mathbf{x}),
\end{equation}
with amplitudes satisfying Eq.~(\ref{eq:a_n_eom}).  Define the
quadratic linearized source energy
\begin{equation}\label{eq:eta_linear_energy}
  E_{\mathrm{lin}}[\eta]
  \equiv \frac{1}{2}\sum_{n\notin \mathrm{neutral}}
  \left(\dot a_n^2 + \omega_n^2 a_n^2\right).
\end{equation}
Differentiating and using Eq.~(\ref{eq:a_n_eom}) yields the exact
identity
\begin{equation}\label{eq:eta_energy_identity}
  \frac{d}{dt}E_{\mathrm{lin}}[\eta]
  = -\sum_{n\notin \mathrm{neutral}}\Gamma_n\,\dot a_n^2
    + \sum_{n\notin \mathrm{neutral}}\dot a_n\,f_n(t).
\end{equation}
Introducing the projected forcing norm
\begin{equation}\label{eq:source_forcing_norm}
  F_{\mathrm{ad}}(t)
  \equiv \left(
    \sum_{n\notin \mathrm{neutral}} |f_n(t)|^2
  \right)^{1/2}
  = \order{
    |\dot{\mathbf{R}}(t)|^2
    + |\dot{\Phi}_c(t)|^2
    + E_{\mathrm{lin}}[\eta]^{1/2}
  },
\end{equation}
which expresses that, after removal of the tangent directions,
non-neutral forcing starts only at higher order in the adiabatic
drift and in the remainder itself.  By Cauchy--Schwarz,
\begin{equation}\label{eq:eta_energy_bound}
  \frac{d}{dt}E_{\mathrm{lin}}[\eta]
  \le -\sum_{n\notin \mathrm{neutral}}\Gamma_n\,\dot a_n^2
    + E_{\mathrm{lin}}[\eta]^{1/2}\,F_{\mathrm{ad}}(t).
\end{equation}
\textit{Coercive reduced stability estimate.}---Let
$\eta_\perp$ denote the remainder after projection away from the
neutral tangent directions.  The self-adjoint part of the linearized
operator on this reduced subspace is diagonal in the mode basis:
\begin{equation}\label{eq:reduced_linear_operator}
  \mathcal{L}_{\perp}\phi_n = \omega_n^2\,\phi_n,
  \qquad n\notin \mathrm{neutral}.
\end{equation}
Because the neutral symmetry directions have been removed and the
remaining discrete branch is the stable source branch used below,
there is a positive reduced spectral gap
\begin{equation}\label{eq:reduced_gap_def}
  \lambda_{\mathrm{gap}}
  \equiv \inf_{n\notin \mathrm{neutral}}\omega_n^2 > 0.
\end{equation}
Define the coercive source energy by
\begin{equation}\label{eq:coercive_source_energy}
  E_{\mathrm{coer}}[\eta_\perp]
  \equiv \frac{1}{2}\,\|\partial_t\eta_\perp\|_{L^2}^2
  + \frac{1}{2}\,\langle \eta_\perp,\mathcal{L}_{\perp}\eta_\perp\rangle
  = \frac{1}{2}\sum_{n\notin \mathrm{neutral}}
    \left(\dot a_n^2 + \omega_n^2 a_n^2\right).
\end{equation}
Therefore
\begin{equation}\label{eq:coercive_lower_bound}
  E_{\mathrm{coer}}[\eta_\perp]
  \ge \frac{1}{2}\,\|\partial_t\eta_\perp\|_{L^2}^2
    + \frac{\lambda_{\mathrm{gap}}}{2}\,\|\eta_\perp\|_{L^2}^2,
\end{equation}
so every non-neutral deformation costs strictly positive quadratic
energy.  Combining Eq.~(\ref{eq:eta_energy_identity}) with
Cauchy--Schwarz now gives
\begin{equation}\label{eq:coercive_growth_bound}
  \frac{d}{dt}E_{\mathrm{coer}}[\eta_\perp]
  \le \sqrt{2E_{\mathrm{coer}}[\eta_\perp]}\,F_{\mathrm{ad}}(t),
\end{equation}
and hence
\begin{equation}\label{eq:coercive_integrated_bound}
  E_{\mathrm{coer}}[\eta_\perp](t)^{1/2}
  \le E_{\mathrm{coer}}[\eta_\perp](0)^{1/2}
  + \frac{1}{\sqrt{2}}\int_0^t F_{\mathrm{ad}}(s)\,ds.
\end{equation}
This is the coercive stability statement needed by the gravity
mechanism: after removing translation, branch drift, and phase,
the remainder cannot wander at zero cost.  It is separated from the
source manifold by the positive gap $\lambda_{\mathrm{gap}}$ and stays
small under adiabatic forcing if it is initially small.
Hence, for sufficiently small initial remainder and adiabatic motion
along $\mathcal{M}_{\mathrm{osc}}$, Gr\"onwall's inequality implies
that $E_{\mathrm{lin}}[\eta]$ remains small throughout the same
window
\begin{equation}\label{eq:source_control_window}
  \omega_R^{-1} \ll t \ll \Gamma^{-1},
\end{equation}
with $\Gamma$ the leakage scale of the least-damped non-neutral
channel.  This is the precise source-control statement used by the
gravity section: the only neutral drift relevant to the static source
is translation of the localized core and slow motion along the bound
branch, while the fast phase is averaged out and non-neutral
deformations are damped or remain perturbatively slaved to those
collective coordinates.  The far-zone gravitational field therefore
follows the same $\mathbf{R}(t)$ that moves the core.

\textit{Projected collective-coordinate dynamics.}---The
orthogonality conditions~(\ref{eq:source_orthogonality}) not only
define the modulation parameters; they also determine their dynamics.
Let
\begin{equation}\label{eq:collective_coordinate_set}
  q_A(t) \equiv \bigl(R_1,R_2,R_3,\Phi_c,\Theta\bigr).
\end{equation}
Projecting the full field equation onto the tangent basis
$\{T_A\}$ gives the finite-dimensional modulation system
\begin{equation}\label{eq:projected_modulation_system}
  \sum_B \mathcal{M}_{AB}(q)\,\ddot q_B
  + \sum_B \mathcal{C}_{AB}(q,\dot q)\,\dot q_B
  = \mathcal{F}^{\mathrm{ext}}_A(t) + \mathcal{R}_A[\eta],
  \qquad
  \mathcal{M}_{AB}(q) \equiv \langle T_A,T_B\rangle,
\end{equation}
where $\mathcal{C}_{AB}$ comes from differentiating the tangent
bundle along $\mathcal{M}_{\mathrm{osc}}$,
$\mathcal{F}^{\mathrm{ext}}_A$ is the adiabatic forcing projected on
$T_A$, and $\mathcal{R}_A[\eta] = \order{E_{\mathrm{lin}}[\eta]^{1/2}}$
is the residual backreaction of the non-neutral remainder.  Because
the Gram matrix $\mathcal{M}_{AB}$ is positive definite on the
reduced tangent space, the collective drift is controlled by
\begin{equation}\label{eq:collective_coordinate_scaling}
  \ddot{\mathbf{R}}
  = \order{F_{\mathrm{ad}} + E_{\mathrm{lin}}[\eta]^{1/2}},\qquad
  \dot{\Phi}_c
  = \order{\Gamma + F_{\mathrm{ad}} + E_{\mathrm{lin}}[\eta]^{1/2}},\qquad
  \dot{\Theta}
  = \Omega(\Phi_c) + \order{F_{\mathrm{ad}} + E_{\mathrm{lin}}[\eta]^{1/2}}.
\end{equation}
Thus translations remain the neutral channel, branch drift occurs only
on the long leakage timescale, and the phase follows the instantaneous
bound frequency up to adiabatic corrections.

\textit{Numerical reduced-gap check.}---For the reference
background used in the cavity scan below ($\Phi_0 = 2.35$), the
translation mode is the only neutral direction, while the first
non-neutral bound mode occurs at
\begin{equation}\label{eq:reduced_gap_numeric}
  \omega_{\min,\mathrm{nn}} = 0.663853,
  \qquad
  \lambda_{\mathrm{gap}}
  \approx \omega_{\min,\mathrm{nn}}^2
  = 0.440701.
\end{equation}
Thus the reduced source sector is not merely qualitatively but also
numerically separated from zero by a finite gap: after projecting out
the translational mode, no additional soft deformation competes with
the source manifold.

\textit{Far-zone closure for a moving source.}---At distances
$r \gg \kappa^{-1}$, the oscillatory core is exponentially small and
the remainder contributes only through the controlled quantity
$E_{\mathrm{lin}}[\eta]$.  Therefore the static part of the field
generated by the modulated source has the universal form
\begin{equation}\label{eq:moving_source_far_zone}
  \varphi_{\mathrm{grav}}(\mathbf{x},t)
  = -\frac{2G\,M_{\mathrm{grav}}(\Phi_c(t))}
      {c^2\,|\mathbf{x}-\mathbf{R}(t)|}
    + \order{E_{\mathrm{lin}}[\eta]^{1/2}}
    + \order{e^{-\kappa |\mathbf{x}-\mathbf{R}(t)|}},
\end{equation}
and hence
\begin{equation}\label{eq:moving_source_time_derivative}
  \partial_t\varphi_{\mathrm{grav}}
  = -\dot{\mathbf{R}}(t)\!\cdot\!\nabla\varphi_{\mathrm{grav}}
    + \dot{\Phi}_c(t)\,\partial_{\Phi_c}\varphi_{\mathrm{grav}}
    + \order{E_{\mathrm{lin}}[\eta]^{1/2}}.
\end{equation}
On timescales $t \ll \Gamma^{-1}$ the mass drift
$M_{\mathrm{grav}}(\Phi_c(t))$ is adiabatic, so the leading effect is
rigid translation of the universal $1/r$ field.  This closes the
single-source dynamical statement needed by the gravity mechanism:
the moving oscillon core, its source term, and its far-zone field are
carried by one and the same collective coordinate $\mathbf{R}(t)$.

%-------------------------------------------------------------------
\subsection{Internal degrees of freedom and particle diversity}
\label{sec:internal_dof}
%-------------------------------------------------------------------

A single scalar field admits only spin-0 excitations.
The full particle spectrum---with spins 0, $\frac{1}{2}$,
1, and 2---requires the additional degrees of freedom
present in the coupled scalar-metric system.  We emphasize
at the outset that the \emph{scalar--vector--tensor (SVT)
decomposition} developed in this section supplies the
bosonic sectors only ($s=0,1,2$); the half-integer
(fermion) sector of the spectrum is \emph{not} read off
from the SVT sectors but is derived separately, via
vortex topology in the oscillon envelope, Komar angular
momentum, and Onsager--Feynman circulation quantization
in Sec.~\ref{sec:spin}.  The decomposition below should
therefore be read as the bosonic half of the spectrum
argument, complemented by the vortex construction of
Sec.~\ref{sec:topology}.

The bi-conformal metric~(\ref{eq:metric}) is determined
by one scalar function $\varphi$ on the \emph{background},
but general perturbations decompose into three independent
sectors:
\begin{enumerate}
  \item \textbf{Scalar perturbations} ($\delta\varphi$):
    longitudinal, compressional modes of the medium.
    These are the breathing modes analyzed above.
  \item \textbf{Tensor perturbations} ($h_{ij}^{TT}$):
    transverse-traceless deformations of the spatial
    metric.  These are the two TT polarizations
    $h_+$ and $h_\times$, which propagate at $c$
    as established in~\cite{Lotishvili2026}.
  \item \textbf{Vector perturbations} ($\delta g_{0i}$):
    rotational, vortical modes of the medium.  These are
    frame-dragging degrees of freedom, analogous to
    vortices in a fluid.
\end{enumerate}

This decomposition---longitudinal (scalar), transverse
(tensor), and rotational (vector)---parallels exactly
the three independent modes of a continuous medium
(sound waves, shear waves, and vortices).  The analogy
is not superficial: it reflects the mathematical structure
of the SVT (scalar-vector-tensor) decomposition theorem
for perturbations of a
Lorentzian manifold~\cite{Bardeen1980,KodamaSasaki1984}.

The diversity of particle species arises from combinations
of these bosonic excitation types together with the vortex
topology of the oscillon sector introduced in
Sec.~\ref{sec:resonances}.  The latter supplies two
ingredients that the SVT sectors alone do not: (i) an
integer vortex winding number $n\in\mathbb{Z}$, which
through the Komar integral and Onsager--Feynman
quantization yields the half-integer angular-momentum
quantum $J = n\hbar/2$ for $n=\pm 1$ (the fermion sector,
Sec.~\ref{sec:spin}), and (ii) a Berry-exchange phase for
two such vortices, which reproduces the antisymmetric
exchange statistics of identical fermions and therefore
the Pauli exclusion principle (Sec.~\ref{sec:pauli}).
The full topological classification---spin, matter versus
antimatter (vortex sign), and exchange statistics---is
developed in Sec.~\ref{sec:topology}.

\textit{Remark.}---This framework requires no extra
dimensions: the excitation types above (plus the vortex
sector of Sec.~\ref{sec:topology}) are all supported by a
three-dimensional medium, just as a fluid supports
diverse wave phenomena (acoustic, transverse, vortical)
without additional spatial dimensions.  The
\emph{epistemic status} of the internal quantum numbers
is layered, however.  The construction above rigorously
derives (i) the bosonic sectors $s=0,1,2$ from SVT, (ii)
the fermionic angular-momentum quantum $J=\hbar/2$ from
$n=\pm1$ vortex winding, (iii) matter versus antimatter
from the sign of the winding (Sec.~\ref{sec:antimatter}),
and (iv) the Pauli exclusion principle from the
Berry-exchange phase (Sec.~\ref{sec:pauli}).  The
remaining Standard-Model internal structure---electric
charge, color $SU(3)_c$, and flavor---is \emph{not}
derived here: the expectation that these likewise arise
from topological invariants and coupled mode structure in
the 3D scalar-metric medium is a \emph{working
hypothesis} in the spirit of the vortex analogy, and a
full topological/representation-theoretic derivation of
the SM gauge content $SU(3)_c\times SU(2)_L\times U(1)_Y$
from ISPG is left for future work.

%-------------------------------------------------------------------
\subsection{Quantum stability of the phantom sector}
\label{sec:stability}
%-------------------------------------------------------------------

The ISPG scalar field has phantom kinetic sign ($\mathcal{K} > 0$
in Eq.~\ref{eq:X_static}), which classically leads to a stable,
singularity-free theory due to the bi-conformal
coupling~\cite{Lotishvili2026}.  In a generic phantom theory,
quantum corrections destabilize the vacuum through ghost pair
production at a rate set by the UV
cutoff~\cite{Carroll2003,Cline2004}.  Below we give the
argument that this instability is absent in ISPG within the
pure-gravity constrained sector to all orders in the loop
expansion (matter-coupled extension is controlled by the EFT
completion of Appendix~11; full bracket-closure/anomaly proof
deferred per A.5).  This question belongs to the
consistency of the full quantum extension and is logically
downstream of the oscillon-to-gravity construction developed
above, which requires only localized finite-energy sources and
their exterior field.

\textit{Kinematic ghost freedom.}---On the bi-conformal
background, the constraint reduces the gravitational phase
space from three propagating degrees of freedom (two tensor,
one scalar) to a single scalar mode $\varphi$, as established
by the Dirac--Bergmann analysis of~\cite{Lotishvili2026} and
confirmed by the canonical reduction in
Sec.~\ref{sec:reduced_action}.  (General perturbations that
break the bi-conformal symmetry---gravitational waves---propagate
in the unconstrained 3-DOF sector; their classical stability
follows from the strong hyperbolicity of the principal symbol
proven in~\cite{Lotishvili2026}.)  Ghost pair
production---the process
$|0\rangle \to |\text{ghost}\rangle + |\text{tensor}\rangle$---requires
at least two independent propagating fields: one carrying
negative energy (the ghost) and one carrying positive energy
(the tensor mode), so that their sum matches the vacuum energy.
On the constraint surface, no such partner exists: the tensor
modes are fully determined by $\varphi$ through the bi-conformal
conditions, and the momentum constraint vanishes identically
(Sec.~\ref{sec:constraint_structure}).  This elimination is
\emph{kinematic}, not dynamical: it is a property of the
constraint surface and holds independently of the loop order.

\textit{Anomaly freedom of the constraint.}---The
kinematic argument above is valid at all orders only if the
constraint is not broken by quantum corrections (a constraint
anomaly).  In the present theory, the smeared Hamiltonian
constraint algebra is Abelian
(Sec.~\ref{sec:constraint_structure},
Eq.~\ref{eq:HH_algebra}): the structure \emph{constants}
on the right-hand side vanish identically, rather than being
field-dependent structure \emph{functions}.  Quantum
anomalies in constraint algebras arise from the regularization
of products of structure functions at coincident
points~\cite{DeWitt1967}; with vanishing structure functions,
no such regularization ambiguity arises in the kinematic argument.
At the structural level (Appendix~11), the Abelian algebra
is therefore expected to remain anomaly-free at the quantum
level, with the single-DOF reduction persisting at all orders;
a fully rigorous bracket-closure/anomaly proof is deferred to
a separate technical appendix.

\textit{Unitarity of the physical S-matrix.}---The optical
theorem requires
\begin{equation}\label{eq:optical}
  2\,\mathrm{Im}\,\mathcal{M}(i \to i)
  = \sum_{f \in \mathcal{H}_{\mathrm{phys}}}
    |\mathcal{M}(i \to f)|^2,
\end{equation}
where the sum runs over the physical Hilbert space
$\mathcal{H}_{\mathrm{phys}}$.  In the constrained theory, the pure-gravity sector of
$\mathcal{H}_{\mathrm{phys}}$ contains only the $\varphi$
excitations---no independent ghost or tensor gravitational
state contributes to the intermediate-state sum.  Matter
sectors of $\mathcal{H}_{\mathrm{phys}}$ (fermions, gauge
fields minimally coupled to $g_{\mu\nu}[\varphi]$) enter the
sum as ordinary positive-norm contributions, so their total
inclusion preserves the non-negativity of the right-hand side
of~(\ref{eq:optical}) provided the $\varphi$--matter coupling
is gravitational (Planck-suppressed).  The right-hand side is
therefore manifestly non-negative at every loop order in the
Planck-suppressed expansion, supporting perturbative unitarity
of the S-matrix order-by-order within the EFT regime; the full
bookkeeping of mixed matter-ghost configurations belongs to the
EFT completion of Appendix~11 and is not asserted here as an
exact non-perturbative unitarity theorem.  This is the same mechanism by
which cuscuton gravity~\cite{Afshordi2007} and mimetic
gravity~\cite{Chamseddine2013} achieve ghost freedom:
kinematic removal by constraint, not dynamical
stabilization.  The ISPG case is distinguished by the
Abelian constraint algebra, which further excludes the
ordering ambiguities that could reintroduce ghosts in
non-Abelian constrained systems.

\textit{Perturbative suppression of loop corrections.}---Beyond
the structural all-orders argument, the \emph{magnitude} of
quantum corrections is controlled by the EFT expansion
parameter.  The scalar perturbation on the bi-conformal
background has dispersion relation $\omega^2 = k^2$
(Sec.~\ref{sec:wave_eq}), so the propagator
\begin{equation}\label{eq:scalar_prop}
  \langle \delta\varphi(x)\,\delta\varphi(y) \rangle
  = -\int \frac{d^4k}{(2\pi)^4}\,
    \frac{i}{k^2 + i\epsilon}
\end{equation}
has the same pole structure as a canonical massless scalar
but with a negative residue---the hallmark of the phantom
kinetic sign.  Within the pure-gravity sector (where $\varphi$ is the only
propagating mode), the negative residue does not lead to a
unitarity deficit: the physical inner product on the
$\varphi$-subsector of $\mathcal{H}_{\mathrm{phys}}$ can be
redefined by absorbing the sign (equivalently, the field
redefinition $\tilde{\varphi} = i\varphi$ maps the pure-gravity
sector to a canonical scalar with modified vertex signs).
This redefinition does not simultaneously rotate the
positive-norm matter sector, so the combined theory retains
mixed-signature states whose consistency is controlled by
the gravitational suppression of $\varphi$--matter couplings
(see ``Kinematic ghost freedom'' above).  At $n$-loop
order, the corrections to any observable scale as
$(E/\Mpl)^{2n}$, where $E$ is the characteristic energy of
the process.  On the classical background
$\varphi_0 = -\rs/r$, the relevant energy scale is set by
the local curvature, $E \sim |\varphi_0'| = \rs/r^2$.  Near
a neutron star surface ($r \sim 10\;\mathrm{km}$), this gives
$E \sim 10^{-10}\;\mathrm{eV}$, so the loop expansion
parameter is
\begin{equation}\label{eq:loop_suppression}
  \left(\frac{E}{\Mpl}\right)^2
  \sim 10^{-76},
\end{equation}
placing the theory deep within the controlled EFT regime at
every astrophysically relevant scale.

\textit{Non-perturbative stability.}---The Hamiltonian
constraint $\mathcal{H}_0 \approx 0$
(Eq.~\ref{eq:H_constraint}) restricts all physical
configurations of the pure-gravity sector to the zero-energy
surface; no runaway to arbitrarily negative energy is possible
\emph{within this sector}. With matter included, the
constraint reads
$\mathcal{H}_\varphi + \mathcal{H}_m \approx 0$, so a
bookkeeping configuration with $\mathcal{H}_\varphi\to -\infty$
balanced by $\mathcal{H}_m\to +\infty$ is in principle allowed;
such a runaway is kinematically controlled by the
Planck-suppressed gravitational $\varphi$--matter coupling
strength rather than by the constraint itself, and is deferred
to the EFT-completion analysis above.  Additionally, the
energy density
$\rho_\varphi = e^{\varphi_0}|\nabla\varphi_0|^2/(32\pi G)
\geq 0$ (Lemma~\ref{app:lemma_positive}) is non-negative,
and $\rho_\varphi \to 0$ as $\varphi \to -\infty$ due to the
$e^{\varphi}$ prefactor, so the effective potential for
$\varphi$ is bounded from above with no runaway direction.
In the path-integral formulation, the functional integral is
restricted to field histories satisfying the constraint; on
this restricted configuration space, the classical solutions
are local minima of the action and the quadratic fluctuation
operator has no negative eigenvalues that would signal a
tunneling instability.  Tunneling from the ISPG vacuum to
any non-perturbative configuration would require traversal
of a field-space barrier with action
$\Delta S \gg \hbar$, exponentially suppressing such
transitions.  The global structure of the constrained
configuration space therefore admits no non-perturbative
vacuum instability.

\textit{Summary.}---The argument for all-orders quantum
stability of the phantom sector within the pure-gravity
constrained sector rests on three independent pillars:
\begin{enumerate}
\item[(i)] \textit{Kinematic ghost freedom}: the single-DOF
  constraint eliminates the pair-production channel at all
  loop orders;
\item[(ii)] \textit{Structurally Abelian constraint algebra}:
  the structural argument of Appendix~11 (full bracket-closure
  proof deferred) suggests that the constraint preserves the
  DOF reduction at the quantum level;
\item[(iii)] \textit{Non-perturbative boundedness on the
  constrained background}: the Hamiltonian constraint confines
  all physical configurations to the zero-energy surface, and
  the bounded effective potential excludes tunneling to runaway
  configurations within the pure-gravity constrained sector.
\end{enumerate}
Together, these arguments suggest that the phantom kinetic sign
in ISPG does not compromise quantum consistency within the
pure-gravity constrained sector; matter-coupled extension is
controlled by the EFT completion. On the contrary, as discussed
in Sec.~\ref{sec:physical_interp}, the phantom sign is the very
feature that renders the Wheeler--DeWitt equation hyperbolic,
enabling non-trivial quantum states and providing a natural
candidate for an internal time variable in superspace.

%-------------------------------------------------------------------
\subsection{Structure of the quantum effective action}
\label{sec:effective_action}
%-------------------------------------------------------------------

Having outlined the structural argument for all-orders
stability of the theory within the pure-gravity constrained
sector, we now derive the explicit structure of the
quantum effective action $\Gamma[\varphi_0]$, which
governs the leading quantum corrections to all classical
ISPG predictions.

\textit{Background field expansion.}---Split the scalar
field into a classical background and a quantum
fluctuation,
\begin{equation}\label{eq:bg_split}
  \varphi(x) = \varphi_0(x) + \delta\varphi(x),
\end{equation}
and expand the action~(\ref{eq:action}) to quadratic
order in $\delta\varphi$.  On the bi-conformal background
the fluctuation operator takes the form
\begin{equation}\label{eq:fluct_op}
  \mathcal{D}(\varphi_0)
  = -\Box + \mathcal{V}(\varphi_0),
\end{equation}
where $\Box = g^{\mu\nu}\nabla_\mu\nabla_\nu$ is the
covariant d'Alembertian on the background metric
$g_{\mu\nu}(\varphi_0)$ and $\mathcal{V}$ collects the
curvature-dependent potential terms arising from the
expansion of $\sqrt{-g}\,R$.

The effective action through one loop is then
\begin{equation}\label{eq:Gamma1}
  \Gamma[\varphi_0]
  = S[\varphi_0]
    + \frac{\hbar}{2}\,\mathrm{Tr}\ln
      \mathcal{D}(\varphi_0)
    + \mathcal{O}(\hbar^2),
\end{equation}
where the functional trace includes integration over
spacetime and summation over all fluctuation modes
(scalar and tensor); we use the standard loop-expansion
convention $\Gamma = S + \sum_{n\geq 1}\hbar^n\,\Gamma^{(n)}$,
so the one-loop term proper is $\Gamma^{(1)}=\tfrac{1}{2}\mathrm{Tr}\ln\mathcal{D}$.

\textit{Shift symmetry and radiative protection.}---At the
unreduced action level, the ISPG action~(\ref{eq:action})
possesses a continuous shift symmetry
$\varphi \to \varphi + \mathrm{const}$: the scalar field
enters only through its derivatives $\partial_\mu\varphi$.
This symmetry forbids the perturbative radiative generation of a mass
term $m^2\varphi^2$ or any non-derivative potential
$V(\varphi)$ at any loop order. The shift-symmetric
counterterms include not only the pure-kinetic powers
$(\partial\varphi)^{2n}$, but also curvature-derivative
operators such as $R\,(\partial\varphi)^2$,
$R_{\mu\nu}\partial^\mu\varphi\,\partial^\nu\varphi$,
$(\Box\varphi)^2$, and combinations thereof, each suppressed
by additional powers of $\Mpl^{-2}$:
\begin{equation}\label{eq:shift_protected}
  \Delta\mathcal{L}_{\mathrm{ct}}
  = \sum_{n=2}^{\infty}
    \frac{c_n}{\Mpl^{2(n-1)}}\,
    (\partial_\mu\varphi\,\partial^\mu\varphi)^n
    + (\text{curvature-derivative shift-invariants}).
\end{equation}
This radiative stability is the same mechanism that
protects the inflaton mass in shift-symmetric
inflation~\cite{Cheung2008} and the Galileon field in
Horndeski gravity.  In the reduced bi-conformal sector,
the fluctuation operator $\mathcal{D}(\varphi_0)$ acquires
explicit $e^{\pm\varphi_0}$ dependence, so the shift symmetry
of the reduced theory is preserved only at EFT precision; a
fully rigorous all-loop Ward identity on the constrained
sector is deferred to a separate technical analysis. For ISPG,
the perturbative argument supports the masslessness of $\varphi$
and the long-range nature of gravity at all orders within
the EFT regime.

\textit{One-loop divergence structure.}---The one-loop
divergences of the coupled scalar-metric system are
obtained via the Seeley--DeWitt heat-kernel
expansion~\cite{tHooftVeltman1974,DeWitt1967}.  In
dimensional regularization ($d = 4 - 2\epsilon$), the
divergent part of~(\ref{eq:Gamma1}) takes the form
\begin{equation}\label{eq:div_structure}
  \Gamma^{(1)}_{\mathrm{div}}
  = -\frac{1}{16\pi^2\epsilon}\int d^4x\,\sqrt{-g}\,
    \left[\alpha\,C_{\mu\nu\rho\sigma}^2
    + \beta\,E_4
    + \gamma\,\Box R\right],
\end{equation}
where $C_{\mu\nu\rho\sigma}$ is the Weyl tensor,
$E_4 = R_{\mu\nu\rho\sigma}^2 - 4R_{\mu\nu}^2 + R^2$ is
the Gauss--Bonnet integrand (topological in $d = 4$), and
$\Box R$ is a total derivative.  The coefficients
$\alpha$, $\beta$, $\gamma$ receive contributions from both
the scalar and tensor fluctuation sectors. Generically a
minimally coupled scalar contributes an additional $R^2$
counterterm in this scheme; we display only the
basis-independent $C^2 + E_4 + \Box R$ contributions and
treat any residual $R^2$ piece as part of the standard
quadratic-curvature renormalization (it does not modify
weak-field/2PN predictions at the displayed parameter
hierarchy).

Two structural features are noteworthy:
\begin{enumerate}
\item The divergence~(\ref{eq:div_structure}) involves
  only \emph{curvature-squared} terms---no $R$ or
  cosmological constant renormalization appears in
  dimensional regularization.  Newton's constant $G$ does
  not run at one loop.
\item Unlike pure gravity---which is one-loop finite on
  shell~\cite{tHooftVeltman1974}---the coupled
  scalar-metric system retains an on-shell divergence
  proportional to $R_{\mu\nu\rho\sigma}^2$.  This
  divergence is absorbed into an $R_{\mu\nu\rho\sigma}^2$
  (equivalently $C^2$) counterterm whose
  physical effect scales as
  $(\mathrm{curvature}/\Mpl^2)^2$.  On the bi-conformal
  background $\varphi_0 = -\rs/r$, the Kretschmann scalar
  scales as $R_{\mu\nu\rho\sigma}^2 \sim \rs^{2}/r^{6}$
  (leading weak-field term), and is
  negligible in the weak-field regime ($r \gg \rs$), so the
  one-loop correction to any observable is suppressed by
  $\rs^3/(\Mpl^2\,r^5) \sim 10^{-78}$ at astrophysical
  scales.
\end{enumerate}

\textit{EFT power counting at all orders.}---Beyond one
loop, the effective action admits the systematic expansion
\begin{equation}\label{eq:EFT_expansion}
  \Gamma[\varphi_0]
  = S[\varphi_0]
    + \sum_{n=1}^{\infty}
      \hbar^n\,\Gamma^{(n)}[\varphi_0],
\end{equation}
where each $n$-loop contribution scales as
\begin{equation}\label{eq:n_loop_scaling}
  \Gamma^{(n)} \sim
  \frac{1}{(16\pi^2)^n}\,
  \frac{E^{2(n+1)}}{\Mpl^{2n}}\,
  \int d^4x\,\sqrt{-g}\,\mathcal{O}_n,
\end{equation}
with $\mathcal{O}_n$ a dimension-$2(n+1)$ local operator
built from curvature tensors and $\partial\varphi$.  The
EFT expansion parameter is
$(E/\Mpl)^2 \sim 10^{-76}$ at astrophysical scales
(Eq.~\ref{eq:loop_suppression}), so the $n$-loop
correction is suppressed by $[10^{-76}/(16\pi^2)]^n$
relative to the tree-level action.

At two loops, Goroff and Sagnotti~\cite{GoroffSagnotti1986}
showed that pure gravity develops a genuine divergence
proportional to $R_{\mu\nu}{}^{\rho\sigma}
R_{\rho\sigma}{}^{\alpha\beta}R_{\alpha\beta}{}^{\mu\nu}$.
In the ISPG EFT, this and all higher-loop divergences are
absorbed into higher-dimensional operators in the
Wilsonian action, with Wilson coefficients of order
$\Mpl^{-2n}$.  Since these operators contribute to
physical amplitudes only at energies $E \sim \Mpl$, they
are irrelevant for all predictions of the present paper.

\textit{Running of the gravitational coupling.}---Although
Newton's constant does not run at one loop in dimensional
regularization, it receives logarithmic corrections from
the curvature-squared counterterms when evaluated at a
renormalization scale $\mu$:
\begin{equation}\label{eq:G_running}
  G(\mu)
  = G_0\left[1
    + \order{\frac{\mu^2}{\Mpl^2}\ln\frac{\mu}{\mu_0}}
    \right].
\end{equation}
For all scales below $\Mpl$, this running is negligible:
at $\mu = 1\;\mathrm{TeV}$ (the highest accessible
collider energy), $\mu^2/\Mpl^2 \sim 10^{-32}$.  The
classical value $G_0$ used throughout this paper is
therefore exact to extraordinary precision.

\textit{Summary.}---The quantum effective action of ISPG
has the following properties:
\begin{enumerate}
\item[(i)] At the unreduced action level the shift symmetry
  $\varphi \to \varphi + \mathrm{const}$ forbids perturbative
  radiative generation of a scalar mass or potential at all
  loop orders (reduced bi-conformal sector preserves the symmetry
  to EFT precision), supporting the long-range nature of gravity.
\item[(ii)] At one loop, the divergences are purely
  curvature-squared; Newton's constant does not run, and
  the on-shell corrections are suppressed by
  $(\mathrm{curvature}/\Mpl^2)^2 \sim 10^{-80}$ at
  astrophysical scales.
\item[(iii)] At all loops, the EFT expansion parameter
  $(E/\Mpl)^2 \sim 10^{-76}$ ensures that quantum
  corrections are negligible at every astrophysically
  relevant scale.
\item[(iv)] The one-loop calculations performed elsewhere
  in this paper (gravitational birefringence,
  Sec.~\ref{sec:birefringence}; Higgs mass correction,
  Sec.~\ref{sec:higgs_prediction};   fine-structure
  variation, Sec.~\ref{sec:alpha_var}) are the
  leading quantum effects and receive no significant
  higher-order corrections.
\end{enumerate}

%-------------------------------------------------------------------
\subsection{Quantum persistence of the scalar sensitivity}
\label{sec:scalar_sensitivity}
%-------------------------------------------------------------------

A central prediction of the classical
theory~\cite{Lotishvili2026} is that the scalar sensitivity
of self-gravitating bodies takes the universal value
$s = 1/2$ for vacuum compact objects (exact) and at leading
order in $|\varphi|$ for matter-filled bodies, so that the
excess sensitivity $s_{\mathrm{excess}} \equiv s - 1/2$
vanishes in the vacuum case and at leading order in matter
(the exact non-perturbative cancellation for matter-filled
bodies of arbitrary compactness remains an open task per
MAIN/16). This supports suppressed scalar dipole radiation
in binary systems and a perturbatively controlled
strong-equivalence statement in the scalar sector.  We
now present the structural argument for the all-orders
status of this result within that scope.

\textit{Shift-symmetry Ward identity.}---The classical
proof rests on a single structural property: the ISPG
field equations depend on $\varphi$ only through
$\partial_\mu\varphi$, endowing the theory with an exact
shift symmetry $\varphi \to \varphi + c$
($c = \mathrm{const}$).  The action~(\ref{eq:action})
contains the metric factors $e^{\pm\varphi}$ explicitly,
but the $\varphi$-dependent terms combine into total
derivatives when the Euler--Lagrange variation is
performed; the bulk Lagrangian can therefore be written as
$\tilde{\mathcal{L}}(\partial_\mu\varphi)$ modulo boundary
terms.  This symmetry is preserved by renormalization: the
counterterms~(\ref{eq:shift_protected}) are all
derivative-type, and the symmetry, being Abelian and
non-chiral, admits no perturbative anomaly.  At the
unreduced action level the symmetry is also free of the
standard instanton-type breaking channels available to chiral
or non-Abelian symmetries; on the reduced bi-conformal
sector and in the matter-coupled theory, exact non-perturbative
preservation in quantum gravity is an open question, so the
protection is asserted at perturbative/EFT precision rather
than as an exact non-perturbative theorem.

The preserved shift symmetry implies a Ward identity for
the quantum-corrected equations of motion:
\begin{equation}\label{eq:ward_shift}
  \frac{\delta\Gamma[\varphi_0 + c]}
       {\delta\varphi_0(x)}
  = \frac{\delta\Gamma[\varphi_0]}
         {\delta\varphi_0(x)}
  \qquad\forall\; c = \mathrm{const}.
\end{equation}
Equivalently, $\Gamma[\varphi_0+c] = \Gamma[\varphi_0] +
f(c)$ where $f(c)$ is independent of $\varphi_0$.
At one loop, this follows from the shift symmetry of the
\emph{unreduced} action~(\ref{eq:action}), in which
$g_{\mu\nu}$ and $\varphi$ are independent variables and
the scalar enters only through $\partial_\mu\varphi$.
The counterterms~(\ref{eq:shift_protected}) inherit this
symmetry, so the quantum-corrected field equations depend
on $\varphi_0$ only through $\partial_\mu\varphi_0$.
On the bi-conformal constraint surface the fluctuation
operator $\mathcal{D}(\varphi_0)$ acquires an explicit
$\varphi_0$-dependence through the metric factor
$e^{-2\varphi_0}$ in the temporal sector; however, this
dependence produces corrections suppressed by
$(E/\Mpl)^{2} \sim 10^{-76}$
(Eq.~\ref{eq:loop_suppression}), so the Ward
identity holds to EFT precision at every loop order.

\textit{Controlled-sector quantum scalar sensitivity.}---The
Ward identity is used here only in the sectors where the
classical rigid-shift response is under control.  For vacuum
compact objects the exterior solution has no matter source to
spoil the shift, and the classical analysis~\cite{Lotishvili2026}
shows that if $\varphi^{*}(x)$ satisfies
$\varphi^{*}\to0$ at spatial infinity, then
$\varphi^{*}(x)+c$ is the shifted boundary-value solution and
the Komar mass scales as
\begin{equation}\label{eq:komar_shift}
  M[\varphi^{*} + c] = e^{-c}\,M[\varphi^{*}].
\end{equation}
Within this vacuum sector, the Ward identity~(\ref{eq:ward_shift})
preserves the shifted family of quantum-corrected equations to
EFT precision at every loop order.  The quantum-corrected Komar
mass, extracted from the asymptotic behavior of
$g_{\mu\nu}[\varphi^{*}+c]$, then inherits the
scaling~(\ref{eq:komar_shift}): the metric components
$g_{00}=-e^{\varphi^*+c}$ and
$g_{ij}=e^{-(\varphi^*+c)}\delta_{ij}$ each acquire the factor
$e^{\pm c}$, and the Komar integral transforms as
$M\to e^{-c}M$.  The vacuum compact-object quantum scalar
sensitivity is therefore
\begin{equation}\label{eq:s_quantum}
  s_{\mathrm{quantum}}
  = -\frac{1}{2}\,
    \frac{\partial\ln M_{\mathrm{eff}}}
         {\partial\varphi_{\mathrm{ext}}}
  = \frac{1}{2}
  \qquad\text{(vacuum compact objects: exact at all loop orders),}
\end{equation}
and the excess sensitivity vanishes:
\begin{equation}\label{eq:s_excess_quantum}
  s_{\mathrm{excess}}^{\mathrm{quantum}} = 0
  \qquad\text{(vacuum exact; matter-filled bodies leading-order in $|\varphi|$, exact non-perturbative case open per MAIN/16).}
\end{equation}
For matter-filled bodies the reduced scalar equation contains the
explicit $e^{-\varphi}$ matter-source factor discussed in MAIN/16,
so the rigid-shift map is not established beyond leading order.
The matter statement in Eq.~(\ref{eq:s_excess_quantum}) is therefore
only the leading-order controlled result; the exact
non-perturbative matter-filled cancellation remains the open
sensitivity-programme task.

\textit{Explicit one-loop verification.}---The
Ward identity~(\ref{eq:ward_shift}) ensures that,
within the same vacuum/leading-order controlled sectors,
the one-loop corrected field equation is shift-invariant.
If $\varphi^{*}$ is the classical solution with
$\varphi^{*} \to 0$ at infinity, and
$\delta\varphi^{(1)}$ is its one-loop correction, then
$\varphi^{*} + c + \delta\varphi^{(1)}$ solves the
one-loop corrected equation with boundary value $c$---the
quantum correction $\delta\varphi^{(1)}$ is the same
regardless of the asymptotic field value.  The total
quantum-corrected Komar mass at boundary value $c$ is
therefore
\begin{equation}
  M_{\mathrm{eff}}(c) =
  M\!\left[g_{\mu\nu}\!\left[\varphi^{*} + c
    + \delta\varphi^{(1)}\right]\right]
  = e^{-c}\,M_{\mathrm{eff}}(0),
\end{equation}
where the factor $e^{-c}$ arises from the metric
rescaling, not from the quantum correction itself.
Since the total mass---classical plus all quantum
corrections---obeys the same exponential scaling, the
one-loop correction to the scalar sensitivity is
\begin{equation}\label{eq:delta_s1}
  \delta s^{(1)} \equiv
  s^{(1\text{-loop})} - s_{\mathrm{classical}} = 0
  \qquad\text{(within the same controlled scope).}
\end{equation}

\textit{Physical consequences.}---Within the controlled
scope above, three consequences follow:
\begin{enumerate}
\item \textit{Suppressed dipole radiation in the controlled
  sectors.}---For vacuum compact objects, and at leading order
  for matter-filled bodies, $s_A = s_B = 1/2$, so the scalar
  dipole luminosity
  $\dot{E}_{\mathrm{dipole}} \propto (s_A - s_B)^2$
  vanishes in those limits.  Pulsar-timing bounds on dipole
  radiation from neutron-star--white-dwarf binaries are
  therefore consistency targets for the matter-filled
  strong-field open task, not already closed all-body
  theorems.
\item \textit{Perturbative strong-equivalence statement.}---In
  the same controlled sectors, bodies respond identically to a
  background scalar-field shift and the differential
  (observable) response is zero.  This supports a
  perturbative/EFT strong-equivalence statement in the scalar
  sector; an exact non-perturbative matter-coupled theorem is
  left to MAIN/16.
\item \textit{Nordtvedt-effect status.}---The shift symmetry
  suppresses body-dependent gravitational self-energy
  contributions to the scalar sensitivity in the vacuum and
  perturbative/EFT sectors.  Thus
  $\eta_N^{\mathrm{quantum}} \simeq 0$ is supported within that
  scope, while the exact matter-filled strong-field statement
  remains contingent on the open sensitivity programme.
\end{enumerate}

%===================================================================
\section{The Bernoulli mechanism for gravity}\label{sec:bernoulli_identity}
%===================================================================

The classical ISPG theory uses the exterior Coulomb profile
$\varphi = -\rs/r$ fixed by the vacuum scalar equation together
with the reduced-sector source normalization, but does not by
itself specify the \emph{physical mechanism} by which matter is
associated with a pressure deficit in the surrounding medium.  The
idea that gravity may originate as a
pressure force in a fluid-like medium has antecedents in
Arminjon's scalar ether theory~\cite{Arminjon2012}, where
the gravitational acceleration is postulated as
$\mathbf{g} = \nabla p_e/\rho_e$ for an imagined perfect fluid.
The present approach differs in three respects: the medium is
not postulated but emerges from the bi-conformal field equations,
the Bernoulli identity is an \emph{exact} consequence of the
stress-energy tensor rather than an analogy, and the framework
passes all PPN tests ($\gamma = \beta = 1$) while yielding a
falsifiable decoherence prediction absent in Arminjon's model.

We show below that the mechanism is
an exact analog of Bernoulli's principle in fluid dynamics: the
kinetic energy of particle resonances drains the static pressure
of the scalar-metric medium.

%-------------------------------------------------------------------
\subsection{Stress-energy decomposition of the scalar field}
\label{sec:stress_energy}
%-------------------------------------------------------------------

The scalar-field contribution to the gravitational field
equations~\cite{Lotishvili2026} is
\begin{equation}\label{eq:tau_munu}
  \tau_{\mu\nu}^{(\varphi)}
  = \partial_\mu\varphi\,\partial_\nu\varphi
  - \frac{1}{2}\,g_{\mu\nu}\,
    g^{\alpha\beta}\partial_\alpha\varphi\,\partial_\beta\varphi.
\end{equation}
For a static field $\varphi_0(\mathbf{x})$ on the bi-conformal
background, the kinetic invariant is purely spatial:
\begin{equation}\label{eq:X_static}
  \mathcal{K} \equiv g^{\alpha\beta}\partial_\alpha\varphi_0\,
  \partial_\beta\varphi_0
  = e^{\varphi_0}\,|\nabla\varphi_0|^{2} \geq 0,
\end{equation}
where the inequality holds because $g^{ij} = e^{\varphi_0}\delta^{ij}$
is positive definite.

The energy density measured by a static observer
(with four-velocity $u^\mu = (e^{-\varphi_0/2}/c, 0, 0, 0)$
when the time coordinate is $x^0 = t$, so that
$u^\mu u_\mu = -c^2$ in the displayed metric; in the
natural-unit convention $c = 1$ this reduces to
$u^\mu = (e^{-\varphi_0/2}, 0, 0, 0)$, the form used
in Eq.~\ref{eq:rho_phi} below)
is
\begin{equation}\label{eq:rho_phi}
  \rho_\varphi
  = \frac{1}{16\pi G}\,\tau_{\mu\nu}^{(\varphi)}\,u^\mu u^\nu
  = \frac{e^{\varphi_0}}{32\pi G}\,|\nabla\varphi_0|^{2}.
\end{equation}
This is manifestly \textbf{positive}: in the ISPG phantom-sign
convention, the scalar-field energy density in a static
configuration is positive-definite.  This positivity refers to
the fluid-Bernoulli kinetic energy density appearing in
Eq.~(\ref{eq:bernoulli_scalar}); in the Einstein equation of the
main theory~\cite{Lotishvili2026}, $\tau^{(\varphi)}_{\mu\nu}$
enters with the phantom-convention factor $-1/2$, so the actual
gravitational source from the scalar sector has the opposite sign.

For the Schwarzschild-like solution $\varphi_0 = -\rs/r$:
\begin{equation}\label{eq:rho_schwarz}
  \rho_\varphi(r)
  = \frac{e^{-\rs/r}}{32\pi G}\,\frac{\rs^{2}}{r^{4}}
\end{equation}
(geometric expression in natural units; SI restoration adds an
explicit $c^4$ prefactor, $\rho_\varphi^{\rm SI} = c^4 \rho_\varphi$,
so that with $\rs = 2GM/c^2$ the SI energy density reads
$\rho_\varphi^{\rm SI} = (GM^2/8\pi r^4) e^{-\rs/r}$).
Differentiating the dimensionless profile $u \equiv \rs/r$
shows that $\rho_\varphi$ peaks at $r = \rs/4$ (not $r\sim\rs$)
and falls as $r^{-4}$ at large distances.

%-------------------------------------------------------------------
\subsection{The Bernoulli identity}
\label{sec:bernoulli_derivation}
%-------------------------------------------------------------------

In fluid dynamics, Bernoulli's equation for a steady, inviscid,
incompressible flow along a streamline reads
\begin{equation}\label{eq:bernoulli_fluid}
  P_{\mathrm{static}} + \frac{1}{2}\rho v^{2}
  = P_{\mathrm{total}} = \mathrm{const},
\end{equation}
where $P_{\mathrm{static}}$ is the hydrostatic pressure,
$\frac{1}{2}\rho v^2$ is the kinetic energy density of the flow,
and $P_{\mathrm{total}}$ is the total (stagnation) pressure.
Regions of faster flow have lower static pressure.

We now derive the exact scalar-field identity that plays the
Bernoulli role in ISPG on the static spherically symmetric
branch.  The fluid formula above is quoted only to display the
structural parallel; Eq.~(\ref{eq:bernoulli_scalar}) itself is
not postulated by analogy but obtained directly from the scalar
stress-energy decomposition.  On this branch the local
bookkeeping is
\begin{equation}\label{eq:bernoulli_scalar}
  \boxed{
  \underbrace{P_{\mathrm{static}}(\varphi)}_{\text{vacuum pressure}}
  + \underbrace{\frac{e^{\varphi}}{32\pi G}\,
    |\nabla\varphi|^{2}}_{\text{gradient (kinetic) energy}}
  = P_{\mathrm{total}} = \mathrm{const},
  }
\end{equation}
where $P_{\mathrm{static}}$ is the static radial pressure of the
scalar-metric medium relative to the undisturbed vacuum
(derived explicitly below).  The total pressure on this static branch
$P_{\mathrm{total}}$ is set by the boundary condition
at spatial infinity ($\varphi \to 0$, $\nabla\varphi \to 0$),
where $P_{\mathrm{static}}(0) = 0$.  Thus
$P_{\mathrm{total}} = 0$: the stagnation pressure vanishes.

\textit{Physical interpretation.}---
Since $P_{\mathrm{static}}(\varphi) \leq 0$ for
$\varphi < 0$ (near matter) and the gradient energy is
non-negative, the Bernoulli identity states that wherever
the field has a non-zero gradient, the static pressure
is \emph{reduced} below the undisturbed value (zero).
The pressure deficit $\Delta P \equiv -P_{\mathrm{static}}
= \frac{e^{\varphi}}{32\pi G}|\nabla\varphi|^2 \geq 0$
is the local gradient energy density---the exact
scalar-field analog of the Bernoulli pressure drop
in fluid dynamics.

For static spherically symmetric configurations this is not merely an interpretive
analogy: comparison with Eq.~(\ref{eq:rho_phi}) gives
\begin{equation}
  \Delta P = \rho_\varphi .
\end{equation}
Thus, on the static spherical branch, the local pressure deficit and the local scalar-field
energy density are the same physical quantity written in two
languages.  For a time-periodic oscillon, this identity is used only
after extracting the zero-frequency, spherically averaged static
component that sources $\varphi_{\mathrm{grav}}$ in the Poisson
reconstruction below; no equality of the full time-dependent stress
with a universal fluid pressure is assumed.

\textit{Derivation.}---Consider the stress-energy tensor
$\tau_{\mu\nu}^{(\varphi)}$ of Eq.~(\ref{eq:tau_munu}) and
the normalized stress-density
\(\mathcal{T}^{\mu}{}_{\nu}\equiv
\tau^{\mu}{}_{\nu}/(16\pi G)\).  For a static,
spherically symmetric field, the radial component of the
covariant conservation equation
$\nabla_\mu \mathcal{T}^{\mu}{}_r = 0$ gives the exact differential
identity
\begin{equation}\label{eq:tau_conservation}
  \frac{d}{dr}\mathcal{T}^r{}_r
  = -\frac{2}{r}\bigl(\mathcal{T}^r{}_r
  - \mathcal{T}^\theta{}_\theta\bigr)
  - \varphi'\,\mathcal{T}^\theta{}_\theta\,.
\end{equation}
For the bi-conformal metric with $\varphi = \varphi(r)$,
the Christoffel symbols and the stress-energy components yield
\begin{equation}\label{eq:tau_radial}
  \mathcal{T}^r{}_r
  = \frac{e^{\varphi}}{32\pi G}\bigl(\varphi'\bigr)^2, \qquad
  \mathcal{T}^\theta{}_\theta
  = -\frac{e^{\varphi}}{32\pi G}\bigl(\varphi'\bigr)^2,
\end{equation}
where $\varphi' \equiv d\varphi/dr$.  The angular component
$\mathcal{T}^\theta{}_\theta$ is the negative of the radial stress density;
this is a generic consequence of the tensor structure of a
static spherically symmetric scalar stress-energy and is shared
by canonical and phantom scalars alike.  The phantom character
enters through the overall sign of $\tau^{(\varphi)}$ in the
Einstein equation of the main theory~\cite{Lotishvili2026}, not
through this internal identity.

Substituting Eq.~(\ref{eq:tau_radial}) into the conservation
equation (\ref{eq:tau_conservation}) and using the explicit
Christoffel symbols for the metric (\ref{eq:metric}), we obtain
after simplification:
\begin{equation}\label{eq:bernoulli_derivation}
  \frac{d}{dr}\left[\frac{e^{\varphi}}{32\pi G}\bigl(\varphi'\bigr)^2\right]
  = \frac{e^{\varphi}}{32\pi G}\bigl(\varphi'\bigr)^3
  - \frac{e^{\varphi}}{8\pi G}\frac{\bigl(\varphi'\bigr)^2}{r}\,.
\end{equation}

\textit{Proof of the Bernoulli integral.}---The vacuum field
equation $\nabla^2\varphi = 0$ in spherical symmetry admits the
exterior Coulomb family $\varphi = B/r$ after the asymptotic
condition $\varphi \to 0$ is imposed.  Source/Newtonian
normalization fixes $B=-\rs$, giving
$\varphi = -\rs/r$.  For this profile, $\varphi' = \rs/r^2$,
$\varphi'' = -2\varphi'/r$, and all $r$-dependence can be
expressed through $\varphi$ via $1/r = -\varphi/\rs$.

On the Coulomb profile, Eq.~(\ref{eq:bernoulli_derivation})
reduces to a first-order equation for the radial stress:
\begin{equation}\label{eq:conservation_ode}
  \frac{d\mathcal{T}^r{}_r}{dr}
  = \bigl(\varphi' - 4/r\bigr)\,\mathcal{T}^r{}_r,
\end{equation}
which is verified by direct substitution using
$(\varphi')^2 = \varphi^4/\rs^2$ and $1/r = -\varphi/\rs$.
This integrates to
\begin{equation}
  \mathcal{T}^r{}_r(r)
  = C\,\frac{e^{\varphi}}{r^4}\,,
\end{equation}
where $C > 0$ is a constant fixed by the source mass.
Evaluating at any $r$ with $\mathcal{T}^r{}_r
= e^{\varphi}(\varphi')^2/(32\pi G)$
and $(\varphi')^2 = \rs^2/r^4$ gives $C = \rs^2/(32\pi G)$.
Expressing the result through $\varphi$ alone
($r^4 = \rs^4/\varphi^4$) yields the
\textbf{Bernoulli integral}:
\begin{equation}\label{eq:bernoulli_integral}
  P_{\mathrm{static}}
  + \underbrace{\frac{e^{\varphi}}{32\pi G}
    \bigl(\varphi'\bigr)^2}_{\mathcal{T}^r{}_r}
  = 0, \qquad
  P_{\mathrm{static}}
  \equiv -\mathcal{T}^r{}_r
  = -\frac{e^{\varphi}\,\varphi^4}{32\pi G\,\rs^2}\,,
\end{equation}
with $P_{\mathrm{static}} \to 0$ as $\varphi \to 0$
(spatial infinity).  On the source-normalized exterior Coulomb
branch $\varphi = -\rs/r$, the closed-form Bernoulli integral
$P_{\mathrm{static}}=-e^\varphi\varphi^4/(32\pi G\,\rs^2)$
holds on the entire exterior vacuum region.  Inside matter
sources the closed-form \emph{profile} is modified, because
the sourced scalar relation $\Box\varphi = -8\pi G\,T/c^{4}$
replaces the vacuum equation and the integration constant
$C$ no longer reduces to $\rs^2/(32\pi G)$.  The \emph{structural}
identity
$P_{\mathrm{static}}+e^\varphi|\nabla\varphi|^{2}/(32\pi G)=0$,
however, follows from the static spherical anisotropy
$\mathcal{T}^{t}{}_{t}=-\mathcal{T}^{r}{}_{r}$ and the static-observer
contraction alone, so it continues to hold locally in
matter-sourced regions; only the explicit Coulomb form
fails inside matter.  That interior closed-form extension
is not required for the Newtonian and weak-field analyses
used in the remainder of this section.

The pressure deficit depends on $\varphi$ through the
universal shape function
$\hat{f}(\varphi) = e^{\varphi}\varphi^4$:
\begin{equation}\label{eq:f_phi_general}
  P_{\mathrm{static}}
  = -\frac{\hat{f}(\varphi)}{32\pi G\,\rs^2}\,,
  \qquad
  \hat{f}(\varphi) \equiv e^{\varphi}\varphi^4 \geq 0,
\end{equation}
whose depth scales as $\rs^{-2} \propto M^{-2}$ at fixed
$\varphi$.  This scaling is the expected physics: at a given
value of $\varphi$, lighter sources have steeper gradients
($\varphi' = \varphi^2/\rs$), hence higher kinetic-energy
density.  The \emph{total} pressure deficit integrated over
all space scales as $M$, recovering the correct gravitational
mass (Sec.~\ref{sec:newton_recovery}).

\textit{Verification.}---Substituting the Coulomb profile
into the right-hand side of
Eq.~(\ref{eq:bernoulli_derivation}) using
$(\varphi')^2 = \varphi^4/\rs^2$,
$1/r = -\varphi/\rs$, and
$(\varphi')^3 = \varphi^6/\rs^3$:
\begin{equation}
  \mathrm{RHS}
  = \frac{e^{\varphi}\,\varphi^6}{32\pi G\,\rs^3}
  + \frac{e^{\varphi}\,\varphi^5}{8\pi G\,\rs^3}
  = \frac{e^{\varphi}\,(\varphi^6 + 4\varphi^5)}{32\pi G\,\rs^3}\,.
\end{equation}
The left-hand side, computed from
$d\mathcal{T}^r{}_r/dr
= (d\mathcal{T}^r{}_r/d\varphi)\,\varphi'$
with
$\mathcal{T}^r{}_r = e^{\varphi}\varphi^4/(32\pi G\,\rs^2)$
and $\varphi' = \varphi^2/\rs$, gives the same expression,
confirming the integral.

\textit{Physical interpretation.}  From the Bernoulli
integral~(\ref{eq:bernoulli_integral}), the local pressure
deficit is
\begin{equation}\label{eq:pressure_deficit}
  \Delta P(r) \equiv -P_{\mathrm{static}}(r)
  = \frac{e^{\varphi_0}}{32\pi G}\,|\nabla\varphi_0|^{2}
  = \rho_\varphi(r) \geq 0.
\end{equation}
The pressure deficit equals the local gradient energy
density: wherever the field has a non-zero gradient,
the static pressure is driven below zero.
Gravity thus arises as a Bernoulli-type pressure
deficit in the scalar-metric medium.

\textit{Time-averaged Bernoulli for oscillating fields.}---The
above derivation assumes a static field configuration
($\partial_t\varphi = 0$).  For particle resonances, which
are time-periodic oscillations $\varphi(t,r) = \varphi_0(r) +
\Phi(r)\cos(\omega t)$, the Bernoulli identity applies to the
time-averaged energy densities over one oscillation period
$T = 2\pi/\omega$.  The kinetic term $\frac{1}{2}(\partial_t\varphi)^2$
averages to $\frac{1}{4}\omega^2\Phi^2$, while the spatial gradient
$\frac{1}{2}(\nabla\varphi)^2$ contributes both a static part
from $\varphi_0$ and an oscillatory part from $\Phi$.  To leading
order in the oscillation amplitude, the time-averaged Bernoulli
identity reads
\begin{equation}
  \langle P_{\mathrm{static}}(\varphi) \rangle_T
  + \frac{e^{\varphi_0}}{32\pi G}
    \left[|\nabla\varphi_0|^2 + \frac{1}{2}|\nabla\Phi|^2
      + \frac{1}{2}\omega^2 e^{-2\varphi_0}\Phi^2\right]
  = 0,
\end{equation}
where $\langle \cdot \rangle_T$ denotes the time average.
The key structural point is that the static gravitational source
is the \emph{zero-frequency projection} of the scalar
stress-energy.  Writing the oscillon field as
\begin{equation}
  \varphi(t,r)
  = \varphi_{\mathrm{grav}}(r)
  + \Phi(r)\cos(\omega t)
  + \delta\varphi_{\ge 2}(t,r),
\end{equation}
where $\delta\varphi_{\ge 2}$ collects higher harmonics and slow
modulation, every term linear in the fast oscillation averages to
zero over one period, while the surviving static contribution
begins at quadratic order in the oscillatory amplitudes.
Accordingly,
\begin{equation}
  \langle T^{00}\rangle
  = \langle T^{00}\rangle_{\mathrm{lead}}
  + \delta T^{00}_{(0)},
\end{equation}
where the leading zero-frequency piece is the explicit source
written in Eq.~(\ref{eq:T00_osc_avg}) below, while
$\delta T^{00}_{(0)}$ contains higher-harmonic and nonlinear
backreaction terms.  In the harmonic/metastable regime used in
this paper, $\delta T^{00}_{(0)}$ is subleading: it starts at
higher order in the oscillon amplitudes and acquires only
adiabatic drift corrections of relative size $\order{\Gamma/\omega}$
over one oscillation period.  The static field is therefore
controlled by the zero-frequency projection of the oscillon
stress-energy, not by the oscillatory harmonics themselves.
The oscillatory contributions represent the kinetic and potential
energy of the particle resonance, which drain the static pressure
below zero through the Bernoulli mechanism.

%-------------------------------------------------------------------
\subsection{Microscopic origin: resonances as the kinetic-energy source}
\label{sec:micro_bernoulli}
%-------------------------------------------------------------------

In the absence of matter, $\varphi = 0$ everywhere and
$|\nabla\varphi| = 0$: the medium is ``still'' and the
static pressure vanishes ($P_{\mathrm{static}} = 0$).

When matter is present, each constituent particle is a
resonant excitation (Sec.~\ref{sec:excitations}) that
oscillates the scalar-metric medium in its vicinity.
These oscillations generate nonvanishing gradients
$\nabla\varphi \neq 0$, transferring energy from the
static (pressure) channel to the kinetic (gradient) channel.

For a single particle of mass $m$ at the origin, the
scalar field solution is $\varphi = -r_m/r$ with
$r_m = 2Gm/c^2$.  The gradient energy in a shell
between $r$ and $r + dr$ is
\begin{equation}\label{eq:shell_energy}
  dE_{\mathrm{grad}} = \frac{r_m^{2}}{8G\,r^{2}}\,
  e^{-r_m/r}\,dr
  \approx \frac{r_m^{2}}{8G\,r^{2}}\,dr
  \quad (r \gg r_m).
\end{equation}
For a macroscopic body composed of oscillons centered at
$\mathbf{x}_a$, the far-zone field is fixed not by pointwise
addition of the local pressure deficits, but by the integrated
coarse-grained source that enters the weak-field scalar
equation.  In the non-overlap regime,
\begin{equation}
  \langle T^{00}\rangle_{\mathrm{body}}(\mathbf{x})
  = \sum_{a=1}^{N}
    \langle T^{00}_{a}\rangle(\mathbf{x}-\mathbf{x}_a)
  + \delta T^{00}_{\mathrm{int}}(\mathbf{x}),
\end{equation}
where $\delta T^{00}_{\mathrm{int}}$ collects localized overlap
and binding corrections.  The total gravitational mass is
therefore
\begin{equation}
  M_{\mathrm{body}}
  = \frac{1}{c^2}\int_{\mathbb{R}^3}
    \langle T^{00}\rangle_{\mathrm{body}}\,d^3x
  = \sum_{a=1}^{N} m_a + \frac{E_{\mathrm{int}}}{c^2},
  \qquad
  E_{\mathrm{int}} \equiv \int \delta T^{00}_{\mathrm{int}}\,d^3x .
\end{equation}
Outside the body ($r \gg R_{\mathrm{body}}$), a multipole
expansion of the scalar Gauss-law solution gives
\begin{equation}
  \varphi_{\mathrm{body}}(r)
  = -\frac{2G M_{\mathrm{body}}}{c^2 r}
    \left[1 + \order{\frac{R_{\mathrm{body}}^2}{r^2}}\right],
\end{equation}
after choosing the origin at the center of mass so that the
dipole term vanishes.  Thus the precise many-body closure needed
for gravity is this: microscopic Bernoulli deficits determine a
localized total source, and the exterior field depends only on
its monopole mass.  The interior interference terms renormalize
$M_{\mathrm{body}}$ but do not alter the universal $1/r$
far-zone law.

\textit{Dynamic macroscopic closure.}---The same conclusion survives
when the constituent oscillons move adiabatically.  Let
$\mathbf{x}_a(t)$ denote their centers and let each constituent obey
the single-source control statement
(\ref{eq:modulated_source_decomposition})--(\ref{eq:moving_source_time_derivative}).
Then the coarse-grained body source can be written as
\begin{equation}\label{eq:body_dynamic_source}
  \langle T^{00}\rangle_{\mathrm{body}}(\mathbf{x},t)
  = \sum_{a=1}^{N}
    \langle T^{00}_{a}\rangle\bigl(\mathbf{x}-\mathbf{x}_a(t);\Phi_{c,a}(t)\bigr)
  + \delta T^{00}_{\mathrm{int}}(\mathbf{x},t)
  + \delta T^{00}_{\mathrm{dyn}}(\mathbf{x},t),
\end{equation}
where $\delta T^{00}_{\mathrm{dyn}}$ collects the controlled
non-rigid corrections from the remainders $\eta_a$ and from
nonrelativistic source motion.  Writing
\begin{equation}\label{eq:body_remainder_control}
  \delta T^{00}_{\mathrm{dyn}}
  = \order{v_{\mathrm{body}}^2}
  + \order{E_{\mathrm{lin,body}}^{1/2}},
  \qquad
  v_{\mathrm{body}} \equiv \max_a \frac{|\dot{\mathbf{x}}_a|}{c},
\end{equation}
with $E_{\mathrm{lin,body}}$ the coarse-grained analogue of the
single-source remainder energy, and choosing the center-of-mass frame
\begin{equation}\label{eq:body_center_of_mass}
  \mathbf{R}_{\mathrm{CM}}(t)
  \equiv \frac{1}{M_{\mathrm{body}}(t)}
  \sum_{a=1}^{N} m_a\,\mathbf{x}_a(t),
\end{equation}
the far-zone field obeys
\begin{equation}\label{eq:body_dynamic_far_zone}
  \varphi_{\mathrm{body}}(\mathbf{x},t)
  = -\frac{2G\,M_{\mathrm{body}}(t)}
      {c^2\,|\mathbf{x}-\mathbf{R}_{\mathrm{CM}}(t)|}
    \left[1 + \order{\frac{R_{\mathrm{body}}^2}{r^2}}\right]
    + \order{e^{-\kappa_{\min} d_{\min}}}
    + \order{v_{\mathrm{body}}^2}
    + \order{E_{\mathrm{lin,body}}^{1/2}},
\end{equation}
where $d_{\min}$ is the minimum pairwise separation in the non-overlap
regime.  Thus the entire body inherits the same translated monopole
field as a single oscillon: internal interference only renormalizes
$M_{\mathrm{body}}(t)$, while the exterior field is carried by the
center-of-mass collective coordinate $\mathbf{R}_{\mathrm{CM}}(t)$.

\textit{The waterfall analogy.}---Consider two bodies
of water: a calm lake and the base of a waterfall.
The calm lake (no matter, $\nabla\varphi = 0$) has
high static pressure.  The turbulent pool at the waterfall's
base ($10^{57}$ atomic resonances agitating the medium)
has the same total energy, but much of it has been
converted to kinetic energy of the flow, drastically
reducing the static pressure.  In the theory itself a test
body does not feel a separate hydrodynamic force; it follows
the geodesic of the metric defined by the same scalar profile.
The analogy is therefore interpretive: the pressure deficit
explains why the profile is gravitational, while the geodesic
equation supplies the motion.

%-------------------------------------------------------------------
\subsection{From pressure deficit to Newton's law}
\label{sec:newton_recovery}
%-------------------------------------------------------------------

We now separate two statements that must not be conflated:
the Bernoulli pressure deficit gives the mechanism-level
interpretation of the scalar profile, while test-particle
motion is fixed by geodesics of the bi-conformal metric.

The logic has two distinct steps.  First, the Bernoulli
identity identifies the \emph{physical source} of gravity:
the oscillatory gradient energy of a localized resonance
reduces the static pressure of the scalar-metric medium.
Second, once this pressure deficit is localized, the exterior
field it generates is fixed by the vacuum scalar equation.
Bernoulli therefore explains \emph{why} the scalar profile is
interpreted as gravitational, while the exterior field equation
and boundary normalization determine \emph{how} that field
propagates at large distances.  The two steps share the same
scalar profile, not an additional pressure-gradient equation of
motion for matter.

The static field equation in the weak-field limit
($|\varphi| \ll 1$) is $\nabla^2\varphi = 0$
outside the source, with the boundary condition set by
the integrated source.  For a localized oscillon, matching the
exterior solution $\varphi_{\mathrm{grav}}(r) = -\rs/r$ to the
source integral gives
\begin{equation}\label{eq:gauss_phi}
  4\pi\rs
  = \frac{8\pi G}{c^4}\int_{B_R}\!\langle T^{00}\rangle\,d^3x
  = \frac{8\pi G}{c^2}\,M_{\mathrm{grav}},
  \qquad
  M_{\mathrm{grav}} \equiv \frac{1}{c^2}
  \int_{\mathbb{R}^3}\!\langle T^{00}\rangle\,d^3x .
\end{equation}
The fundamental sourcing relation in the main theory is
$\Box\varphi = -8\pi G\,T/c^{4}$, where $T$ is the \emph{trace}
of the matter stress-energy tensor~\cite{Lotishvili2026}.  For
pressureless matter (dust) one has $T = -\rho c^{2}$ and
$\langle T^{00}\rangle = \rho c^{2}$, so $\langle T^{00}\rangle =
-T$ and Eq.~(\ref{eq:gauss_phi}) is equivalent to the
trace-sourced form in this limit.  For traceless sources
($T = 0$, e.g.~electromagnetic radiation) the trace-sourced
equation gives no direct scalar response: photons contribute to
$\langle T^{00}\rangle$ but not to $T$, so they do not source
$\varphi$ directly, and their gravitational deflection arises
through geodesic motion in the bi-conformal metric determined by
matter-sourced $\varphi$.  The $\langle T^{00}\rangle$
notation used throughout this section should therefore be read
as the effective energy-density source appropriate to the
non-relativistic, trace-dominant regime in which the oscillon
bodies live.
The Bernoulli pressure deficit already fixes the sign of the
field: positive localized energy lowers the static pressure and
therefore produces $\varphi_{\mathrm{grav}} < 0$.  Gauss's
theorem then fixes the magnitude of the $1/r$ coefficient,
yielding
\begin{equation}
  \varphi_{\mathrm{grav}}(r)
  = -\frac{\rs}{r}
  = -\frac{2GM_{\mathrm{grav}}}{c^2 r}.
\end{equation}

The geodesic acceleration experienced by a test
particle in this field is~\cite{Lotishvili2026}
\begin{equation}\label{eq:grav_accel}
  \mathbf{a}_{\mathrm{geo}}
  = -\frac{c^2}{2}\,\nabla\varphi_{\mathrm{grav}}
  = -\frac{GM_{\mathrm{grav}}}{r^{2}}\,\hat{r}.
\end{equation}
The Newtonian potential is
$\Phi_N = -GM_{\mathrm{grav}}/r
= \frac{1}{2}\varphi_{\mathrm{grav}}\,c^2$, and
Newton's law $\mathbf{a} = -\nabla\Phi_N
= -GM_{\mathrm{grav}}\hat{r}/r^2$
is exactly reproduced by the geodesic equation.

The force on the test particle is the far-zone geodesic
acceleration sourced by the same scalar profile that carries
the Bernoulli pressure deficit:
\begin{equation}\label{eq:force_bernoulli}
  \mathbf{F} = -m\,\tfrac{c^{2}}{2}\,\nabla\varphi
  = -m\,\nabla\Phi_{N}.
\end{equation}
\textit{Firewall: minimal coupling and no fifth force.}
Matter is minimally coupled to $g_{\mu\nu}[\varphi]$ alone,
so test particles follow geodesics of the bi-conformal
metric.  Eq.~\eqref{eq:force_bernoulli} is the weak-field
geodesic acceleration in disguise; it is \emph{not} an
independent pressure-gradient force law on matter.  A
literal substitution $\mathbf{F}=-m\nabla(\Delta P/\rho_{\varphi})$
would scale as $r^{-1}$ on the Coulomb branch (since
$\Delta P\propto r^{-4}$, $\rho_{\varphi}\propto r^{-4}$,
so $\nabla(\Delta P)/\rho_{\varphi}\propto r^{-1}$), whereas
the Newtonian acceleration scales as $r^{-2}$; the two are
not equivalent and the literal pressure-gradient force law
is forbidden.  The Bernoulli identity supplies the
\emph{physical interpretation} of why the metric is
gravitational (the same scalar profile carries a static
pressure deficit), but it does not enter the equation of
motion for matter.  The test-particle weak equivalence
principle is preserved exactly because matter couples only
to $g_{\mu\nu}$.

In this interpretive sense the test body moves \emph{toward}
regions of lower static pressure (near the source) along the
geodesic that the same scalar profile defines: the pressure
deficit is the mechanism, the geodesic is the equation of
motion.  Gravity is not a mysterious attraction; it is the
geodesic response of matter to a metric whose scalar profile
also carries a Bernoulli pressure deficit in the medium.

\textit{Two-component structure of the oscillon
field.}---The preceding derivation reveals a fundamental
duality in the field structure of particle-like
resonances.  The full scalar field near an oscillon
decomposes into two distinct components:
\begin{equation}\label{eq:two_component}
  \varphi(r,t) = \underbrace{\varphi_{\mathrm{grav}}(r)}_{\sim\,-M/r}
  + \underbrace{\Phi_0(r)\cos(\Omega t)}_{\sim\,e^{-\kappa r}/r}\,,
\end{equation}
where $\kappa = \sqrt{1 - \Omega^2}$.  This split is fixed,
to leading order in the oscillon harmonic expansion, by
separating the time average from the first oscillatory mode:
\begin{equation}
  \varphi_{\mathrm{grav}}(r) \equiv \langle \varphi(r,t)\rangle_t,
  \qquad
  \Phi_0(r)\cos(\Omega t)
  \equiv \varphi(r,t) - \langle \varphi(r,t)\rangle_t
  + \order{\Phi^2},
\end{equation}
where $\langle\cdot\rangle_t$ denotes averaging over one
period $2\pi/\Omega$.  The static component is therefore the
part of the field that survives coarse-graining over the fast
oscillation, while the localized oscillatory component carries
the bound internal resonance.
The oscillating component $\Phi_0(r)\cos\Omega t$ is the
particle's \emph{body}---the localized resonant mode whose
amplitude decays exponentially at the rate $\kappa$ set
by the bound-state condition $\Omega < 1$.  The static
component $\varphi_{\mathrm{grav}} = -\rs/r$ is the
particle's \emph{gravitational field}---the $1/r$ potential
sourced by the time-averaged energy density of the
oscillation through the Poisson equation
$\nabla^{2}\varphi_{\mathrm{grav}}
= -(8\pi G/c^{4})\langle T^{00}\rangle$.
The mechanism by which the oscillatory mode $\Phi_0\cos\Omega t$
acts as an effective source for the static Poisson equation is
the nonlinear self-coupling of the scalar sector: the bi-conformal
metric introduces $\varphi$-dependent factors in
$\sqrt{-g}\,g^{\mu\nu}\partial_\mu\varphi\partial_\nu\varphi$ that
mix the fast oscillatory and slow static components, and the
time-averaged residue is the effective $\langle T^{00}\rangle$
entering Eq.~(\ref{eq:T00_osc_avg}).  We take that effective
source as given here and verify the resulting $1/r$ far-field
profile numerically below; the full derivation of the effective
source from the scalar action is carried out in the classical
companion paper~\cite{Lotishvili2026}.
These two components obey fundamentally different
equations: the oscillating mode satisfies the massive
Klein--Gordon equation (Sec.~\ref{sec:resonances}) and
decays exponentially; the gravitational potential
satisfies the massless Laplace equation $\nabla^{2}\varphi = 0$
outside the source and decays as $1/r$.  The particle is
thus an exponentially localized excitation dressed by a
long-range gravitational field---a composite object whose
internal structure ($e^{-\kappa r}$) is invisible at
distances $r \gg 1/\kappa$, while its gravitational
influence ($1/r$) extends to infinity.

Combining the time-averaged Bernoulli identity with the
decomposition~(\ref{eq:two_component}), let
$\delta\varphi(r,t) \equiv \Phi_0(r)\cos(\Omega t)$.  The
oscillatory part of the scalar-field energy density measured
by a static observer is then, to leading order in the
oscillon amplitude,
\begin{equation}
  \rho_{\mathrm{osc}}(r,t)
  = \frac{1}{32\pi G}\left[
      e^{-\varphi_{\mathrm{grav}}}(\partial_t\delta\varphi)^2
      + e^{\varphi_{\mathrm{grav}}}|\nabla\delta\varphi|^2
    \right].
\end{equation}
The mixed term
$2\nabla\varphi_{\mathrm{grav}}\!\cdot\!\nabla\delta\varphi
\propto \cos(\Omega t)$ averages to zero over one period, so
the source component entering the weak-field Poisson equation
is
\begin{equation}\label{eq:T00_osc_avg}
  \langle T^{00}\rangle
  = \langle \rho_{\mathrm{osc}}\rangle_t
  = \frac{1}{64\pi G}\left[
      e^{-\varphi_{\mathrm{grav}}}\Omega^2\Phi_0^2
      + e^{\varphi_{\mathrm{grav}}}(\Phi_0')^2
    \right].
\end{equation}
In the Newtonian matching regime $|\varphi_{\mathrm{grav}}| \ll 1$,
this becomes the dimensionless source profile
\begin{equation}\label{eq:rho_osc_numeric}
  \rho(r) = \tfrac{1}{2}\Omega^2\Phi_0^2
  + \tfrac{1}{2}(\Phi_0')^2,
\end{equation}
up to the overall factor $1/(32\pi G)$.  Thus the same
time-averaged Bernoulli deficit
$\langle\Delta P\rangle_t = \langle\rho_{\mathrm{osc}}\rangle_t$
is exactly the localized source whose Poisson reconstruction
produces $\varphi_{\mathrm{grav}}$.

\textit{Numerical verification: the oscillon gravitational
potential.}---We verify this decomposition
by solving the Poisson equation sourced by the density
profile in Eq.~(\ref{eq:rho_osc_numeric}):
\begin{equation}\label{eq:phi_grav_poisson}
  \varphi_{\mathrm{grav}}(r) =
  -\frac{M_{\mathrm{enc}}(r)}{r}
  - \int_{r}^{\infty}4\pi\rho(r')\,r'\,dr',\qquad
  M_{\mathrm{enc}}(r) = 4\pi\!\int_{0}^{r}\!\rho(r')\,r'^{2}\,dr'.
\end{equation}
Consequently the far-zone coefficient is fixed by the same
source integral:
\begin{equation}
  M_{\mathrm{num}}
  \equiv \lim_{r\to\infty}\!\bigl[-r\,\varphi_{\mathrm{grav}}(r)\bigr]
  = 4\pi\!\int_{0}^{\infty}\!\rho(r)\,r^{2}\,dr.
\end{equation}
This is the natural-unit counterpart of
Eq.~(\ref{eq:gauss_phi}): once the localized source profile is
specified, the asymptotic $1/r$ coefficient is fixed.
This is a direct reconstruction of the static field from the
same oscillon profile that defines the localized resonance:
once $\Phi_0(r)$ is fixed, there is no additional freedom in
the asymptotic potential.  The numerical test therefore checks
not only the $1/r$ falloff but also the correct mass
normalization carried by the source.
For the oscillon at $\Phi_0 = 2.35$ ($\Omega = 0.866$,
$M_{\mathrm{num}} = 295.5$ in natural units), the ratio
$\varphi_{\mathrm{grav}}(r)/(-M_{\mathrm{num}}/r)$ equals unity to
$5 \times 10^{-7}$ at $r = 15$ and to machine precision
($< 10^{-13}$) at $r = 30$, confirming the $1/r$ profile
with the correct mass coefficient.

\textit{Two-body interaction: Newton's law from
oscillons.}---For two oscillons separated by distance $d$
(with $d$ much larger than the oscillon radius), the
gravitational interaction energy is
\begin{equation}\label{eq:U_two_body}
  U(d) = \int\!\rho_2(\mathbf{r}_2)\,
  \varphi_{\mathrm{grav},1}(|\mathbf{r}_2 + d\hat{z}|)\,
  d^{3}r_2\,.
\end{equation}
Since $\rho_1$ is spherically symmetric,
Gauss's theorem gives
$\varphi_{\mathrm{grav},1}(r) = -M_1/r$ for
$r > R_{\mathrm{osc}}$, and the angular average of
$1/|\mathbf{r}_2 + d\hat{z}|$ over the second
oscillon's density distribution yields
$U(d) = -M_1 M_2/d$ at leading order in the tail overlap.
This is precisely the regime relevant to the Newtonian
limit: each oscillon acts as a localized source dressed by
its long-range field, while short-distance core overlap is
deliberately excluded from the present claim.
More precisely, spherical symmetry removes all exterior
multipoles of each compact core, so the compact-support part of
the interaction is \emph{exactly} Newtonian rather than merely
asymptotically so.  The only correction comes from the
exponentially small non-compact tail of the oscillon profile.
Writing
\begin{equation}
  \rho_i(r) = \rho_i^{(<R)}(r) + \rho_i^{(>R)}(r), \qquad
  M_{i,\mathrm{tail}}(R)
  \equiv 4\pi\!\int_R^\infty\!\rho_i(r)\,r^2\,dr,
\end{equation}
with $R < d/2$, one obtains
\begin{equation}
  U(d) = -\frac{M_1 M_2}{d}
  + \order{\frac{M_{1,\mathrm{tail}}(R)\,M_2
      + M_1\,M_{2,\mathrm{tail}}(R)}{d}}.
\end{equation}
For a bound oscillon,
$\Phi_0(r) \sim A e^{-\kappa r}/r$, so
$\rho_i(r) \sim e^{-2\kappa_i r}/r^2$ and therefore
\begin{equation}
  M_{i,\mathrm{tail}}(R) = \order{e^{-2\kappa_i R}}.
\end{equation}
Choosing $R = d/2$ gives the explicit non-overlap estimate
\begin{equation}
  U(d) = -\frac{M_1 M_2}{d}
  \left[1 + \order{e^{-\kappa_{\min} d}}\right],
  \qquad
  \kappa_{\min} \equiv \min(\kappa_1,\kappa_2).
\end{equation}
Thus no algebraic finite-size correction survives in the
spherically symmetric core sector; the leading deviation from
Newton's law is exponentially suppressed by tail overlap.
Numerically, we compute $U(d)$ by evaluating the
full angle-averaged integral
for five pairs of oscillons with amplitudes
$\Phi_0 \in \{1.0,\,1.5,\,2.0,\,2.35,\,3.0\}$
(masses $M \in \{127,\,183,\,246,\,296,\,402\}$).
No Newtonian law is inserted at this stage: the input is only
the localized oscillon density profile and the static field
obtained from Eq.~(\ref{eq:phi_grav_poisson}).  The inverse
distance law then re-emerges from the full source integral in
the non-overlap regime.
The results confirm Newton's law at three levels:
\begin{enumerate}
  \item \textit{Potential:} $|U(d)| \propto 1/d^{n}$
    with $n = 1.000011$ (power-law fit for $d \geq 15$;
    deviation $0.001\%$ from unity).
  \item \textit{Force:} $|F| = |dU/dd| \propto 1/d^{n}$
    with $n = 2.000157$ (deviation $0.008\%$ from the
    inverse-square law).
  \item \textit{Mass proportionality:}
    $U(d)/(M_1 M_2/d) = 1.000000$ for all eight tested
    mass pairs (six significant figures), confirming
    $F \propto M_1\,M_2$.
\end{enumerate}
Newton's gravitational law $F = G\,M_1 M_2/d^{2}$
is therefore not an input of the theory but an
\emph{output}: the $1/r^{2}$ \emph{form} of the force follows
from the field equation $\Box\varphi = 0$ applied to localized
energy sources.  The numerical value of $G$ appearing as the
coefficient is the same $G$ that sets the prefactor
$1/(16\pi G)$ in the ISPG action~(\ref{eq:action}); it is
therefore a parameter of the theory, not derived from deeper
principles in this paper.

%-------------------------------------------------------------------
\subsection{Strong-field regime and singularity avoidance}
\label{sec:strong_bernoulli}
%-------------------------------------------------------------------

In the strong-field regime ($r \sim \rs$), the exponential
factors become significant.  The gradient energy density
Eq.~(\ref{eq:rho_schwarz}) peaks at $r = \rs/4$ and
vanishes as $r \to 0$ due to the $e^{-\rs/r}$ suppression.

This has a profound consequence: the pressure deficit
is \emph{self-limiting}.  As the static pressure drops
(increasing $|\varphi|$), the exponential factor
$e^{\varphi}$ in Eq.~(\ref{eq:bernoulli_scalar})
suppresses the gradient energy, slowing the pressure
drop.  In the limit $\varphi \to -\infty$ (approaching
the classical singularity), $e^{\varphi} \to 0$ and the
gradient energy vanishes entirely: the medium reaches a
state of ``maximal depletion'' where no further energy
transfer from static to kinetic channels is possible.

This is the Bernoulli-mechanism explanation for the
singularity resolution established in the classical
theory~\cite{Lotishvili2026}: the gravitational potential
well has an infinite proper depth
($\int_0^\infty e^{-\varphi_0/2}\,dr = \infty$) but
finite coordinate-volume energy integral
($\int \rho_\varphi\,d^3x_{\rm coord} < \infty$; the
proper-volume measure $\sqrt{h}\,d^3x = e^{-3\varphi/2}d^3x$
weighs the strong-field interior differently and is treated
separately),
because the bi-conformal structure automatically regulates
the pressure deficit.  No singularity forms because the
Bernoulli mechanism has a built-in saturation.

%-------------------------------------------------------------------
\subsection{Comparison with general relativity}
\label{sec:gr_comparison}
%-------------------------------------------------------------------

In GR, gravity is described as the curvature of spacetime,
and the geodesic equation $\ddot{x}^\mu + \Gamma^\mu_{\alpha\beta}
\dot{x}^\alpha\dot{x}^\beta = 0$ determines particle
trajectories.  The question ``\emph{why} does mass curve
spacetime?'' is not answered within GR; it is an axiom
of the Einstein field equations.

In ISPG, this question receives a physical answer:
\begin{enumerate}
  \item Mass consists of resonant excitations of the
    scalar-metric medium (Sec.~\ref{sec:excitations}).
  \item These excitations generate gradient (kinetic)
    energy in the medium via their oscillatory tails.
  \item By the Bernoulli identity~(\ref{eq:bernoulli_scalar}),
    this kinetic energy lowers the static pressure.
  \item The same scalar profile that carries this pressure
    deficit defines the bi-conformal metric.
  \item Test particles then follow geodesics of that metric;
    in the weak-field limit this gives Newtonian gravity.
\end{enumerate}

The causal chain
\begin{equation}\label{eq:chain}
  \text{resonance} \xrightarrow{\text{oscillation}}
  |\nabla\varphi|^{2} \xrightarrow{\text{Bernoulli}}
  \Delta P
  \quad\text{and}\quad
  \varphi \xrightarrow{\text{metric}}
  g_{\mu\nu}[\varphi]\xrightarrow{\text{geodesic}}
  \mathbf{a}_{\mathrm{geo}}
\end{equation}
provides the mechanical underpinning of gravity that
Einstein's geometric description, by design, does not
address.  Both descriptions are compatible: the
bi-conformal metric~(\ref{eq:metric}) encodes the
pressure distribution geometrically, and geodesic motion
in this geometry is the equation of motion.  The pressure
deficit supplies the mechanism-level interpretation, not a
separate force law.

%===================================================================
\section{Inertia, hysteresis, and the equivalence principle}
\label{sec:inertia}
%===================================================================

The Bernoulli mechanism (Sec.~\ref{sec:bernoulli_identity})
explains \emph{gravity}---the static pressure deficit
around matter.  We now address the complementary
question: \emph{why does matter resist acceleration?}
In GR, the equivalence of inertial and gravitational mass
is postulated.  In ISPG, it is derived: both arise from
the same resonance--medium coupling.

%-------------------------------------------------------------------
\subsection{Boosted resonance: uniform motion}
\label{sec:uniform_motion}
%-------------------------------------------------------------------

A resonance at rest is described by the oscillon profile
$\Phi(t,\mathbf{x}) = \phi(r)\,e^{-i\omega t}$
(Sec.~\ref{sec:resonances}).  Under a Lorentz boost
with velocity $\mathbf{v}$, the locally measured speed
of light remains $c$ (Eq.~\ref{eq:c_local}), so the
boosted profile is obtained by the standard replacement
$t \to \gamma(t - \mathbf{v}\cdot\mathbf{x}/c^2)$,
$\mathbf{x} \to \mathbf{x} + (\gamma - 1)
(\hat{v}\cdot\mathbf{x})\hat{v} - \gamma\mathbf{v}t$:
\begin{equation}\label{eq:boosted_oscillon}
  \Phi_v(t,\mathbf{x})
  = \phi\!\left(\sqrt{x_\perp^2
    + \gamma^2(x_\parallel - vt)^2}\right)
    \exp\!\left[-i\gamma\omega
    \left(t - \frac{v\,x_\parallel}{c^2}\right)\right],
\end{equation}
where $x_\parallel$ and $x_\perp$ are the components
parallel and perpendicular to $\mathbf{v}$, and
$\gamma = (1 - v^2/c^2)^{-1/2}$.

The spatial profile is Lorentz-contracted by $\gamma$
along the direction of motion.  Two levels of
justification apply:

\begin{enumerate}
  \item \textit{Free particle (self-generated field).}---For
    an isolated oscillon, $\varphi_{0}$ is generated
    self-consistently by the resonance itself, and the
    full coupled system (metric $g_{\mu\nu}$ + scalar
    $\varphi$ + excitation $\Phi$) transforms covariantly
    under Lorentz boosts, because the ISPG action is
    diffeomorphism-invariant.  The boosted oscillon
    \emph{together with its boosted $\varphi_{0}$} is
    therefore an \emph{exact} solution.
  \item \textit{Particle in an external field.}---When
    $\varphi_{0}(\mathbf{x})$ is sourced by other masses,
    the background breaks global Lorentz symmetry.
    Nevertheless, local Lorentz invariance is guaranteed
    by the equivalence principle: in the freely falling
    frame of the oscillon, the background is locally flat
    over scales $\ll L_{\varphi}$, where
    $L_{\varphi} = |\varphi_{0}/\nabla\varphi_{0}|$ is
    the gradient scale of the external field.  Since
    $R_{c} \ll L_{\varphi}$ for all known particles in
    astrophysical fields, the boosted profile is an
    \emph{adiabatic} (WKB) solution with corrections of
    order $(R_{c}/L_{\varphi})^{2} \ll 1$.
\end{enumerate}

In both cases, no net force acts on the resonance during
uniform motion.

\textit{Physical interpretation: motion as pattern
propagation.}---An oscillon is not a rigid body
translating through the medium; it is a resonant
pattern \emph{of} the medium.  ``Motion'' means that
the excitation amplitude is transferred from one set
of field degrees of freedom to the neighbouring set,
exactly as an ocean wave propagates without net
transport of water.  No substance moves from $A$ to
$B$; only the self-sustaining resonance pattern does.

This ontological point has a direct dynamical
consequence.  A resonance moving at constant velocity
through the medium is indistinguishable
from a resonance at rest in a boosted frame.  The
pressure distribution adjusts self-consistently:
the Lorentz-contracted profile generates a
Lorentz-contracted pressure deficit, and the static
Bernoulli balance~(\ref{eq:bernoulli_scalar}) is
preserved in the comoving frame.  No drag, no friction,
no energy dissipation---because the pattern does not
displace the medium, it \emph{is} the medium.  This is
the superfluid property of the scalar-metric medium.

The maximum propagation speed of any pattern is set by
the speed at which the medium transmits information:
$c_{\mathrm{coord}} = c\,e^{\varphi_{0}}$.  No
oscillon can outrun its own field, so
$|\mathbf{v}| < c_{\mathrm{coord}}$ is automatic---the
relativistic speed limit is not imposed externally but
follows from the medium's finite response rate.

%-------------------------------------------------------------------
\subsection{Accelerated resonance: the origin of inertia}
\label{sec:inertia_mechanism}
%-------------------------------------------------------------------

Section~\ref{sec:uniform_motion} established that
oscillon motion is the propagation of a resonant pattern
through the medium, limited by the medium's information
speed~$c$.  Inertia is the direct consequence: when the
pattern is forced to accelerate, the medium cannot
redistribute its pressure instantaneously.

Consider a resonance whose velocity changes:
$\mathbf{v}(t) = \mathbf{v}_0 + \mathbf{a}\,\delta t$
over a short interval $\delta t$.  The oscillon profile
must continuously readjust its Lorentz contraction to
match the instantaneous velocity.  This readjustment
propagates at the local speed of light $c$ from the
center of the oscillon outward.

\textit{Finite readjustment time.}---The oscillon has
a spatial extent of order the Compton wavelength
$R_c = \hbar/(m_0 c)$.  The time required for the
profile to fully readjust is
\begin{equation}\label{eq:t_readjust}
  t_{\mathrm{adj}} \sim \frac{R_c}{c}
  = \frac{\hbar}{m_0 c^2}
  = \frac{1}{\omega_0},
\end{equation}
which is the oscillation period of the resonance.
During acceleration, if the velocity change occurs
on a timescale shorter than $t_{\mathrm{adj}}$, the
profile cannot follow instantaneously.

\textit{Pressure asymmetry.}---The result is a
transient asymmetry in the pressure distribution
around the resonance:
\begin{enumerate}
  \item \textbf{Ahead} (in the direction of
    acceleration): the resonance enters a region of
    undisturbed medium whose pressure has not yet been
    reduced by the Bernoulli mechanism.  The resonance
    encounters its own pressure deficit
    \emph{incompletely formed}, resulting in a net
    restoring pressure.
  \item \textbf{Behind}: the pressure deficit from the
    previous position persists for a time
    $\sim t_{\mathrm{adj}}$ before the medium
    re-equilibrates.  This trailing deficit exerts a
    backward pull.
\end{enumerate}

We now compute the net force rigorously by expanding
the retarded scalar field to first order in the
acceleration.

\textit{Retarded field of an accelerated source.}---Let
the oscillon center follow a trajectory
$\mathbf{r}(t) = \frac{1}{2}\mathbf{a}\,t^{2}$ with
$|\mathbf{a}| \ll c\omega_{0}$ (non-relativistic
acceleration).  The scalar field satisfies the sourced
wave equation $\Box\varphi = -(8\pi G/c^{4})\,T$ with the
trace of the oscillon stress-energy concentrated at
$\mathbf{r}(t)$.  The retarded Green's function solution
is
\begin{equation}\label{eq:retarded_green}
  \varphi(\mathbf{x},t) = \frac{2G}{c^{4}}\int d^{3}x'\,
  \frac{T\bigl(t_{r},\,\mathbf{x}'\bigr)}
       {|\mathbf{x}-\mathbf{x}'|},
  \qquad
  t_{r} = t - \frac{|\mathbf{x}-\mathbf{x}'|}{c}.
\end{equation}
For a compact oscillon ($R_{c} \ll r$) the source is
effectively pointlike:
$T = Mc^{2}\,\delta^{(3)}(\mathbf{x}'-\mathbf{r}(t))$.
Performing the spatial integration in
Eq.~(\ref{eq:retarded_green}) gives
\begin{equation}\label{eq:phi_point}
  \varphi(\mathbf{x},t)
  = -\frac{r_{s}}{|\mathbf{x} - \mathbf{r}(t_{r})|},
\end{equation}
where $t_{r} = t - |\mathbf{x} - \mathbf{r}(t_{r})|/c$
is the retarded time.  For non-relativistic acceleration,
$\mathbf{r}(t_{r}) = \mathbf{r}(t)
- \mathbf{v}(t)\,|\mathbf{x}|/c + \order{a^{2}}$,
so the displacement from the instantaneous position is
\begin{equation}\label{eq:displacement}
  \mathbf{r}(t) - \mathbf{r}(t_{r})
  = \mathbf{v}(t)\,\frac{|\mathbf{x}|}{c}
  = \mathbf{a}\,t\,\frac{|\mathbf{x}|}{c}.
\end{equation}
Expanding $|\mathbf{x} - \mathbf{r}(t_{r})|^{-1}$
around the instantaneous position
$\boldsymbol{\xi} = \mathbf{x} - \mathbf{r}(t)$:
\begin{equation}\label{eq:coulomb_expand}
  \frac{1}{|\boldsymbol{\xi}
    + \mathbf{a}\,t\,|\mathbf{x}|/c|}
  = \frac{1}{|\boldsymbol{\xi}|}
  - \frac{(\mathbf{a}\cdot\hat{\xi})\,t\,|\mathbf{x}|}
         {c\,|\boldsymbol{\xi}|^{2}}
  + \order{a^{2}}.
\end{equation}
The first-order correction to the equilibrium potential
$\varphi_{\mathrm{eq}}(\xi) = -r_{s}/\xi$ is therefore
\begin{equation}\label{eq:phi_ret_expand}
  \varphi(\mathbf{x},t)
  = \varphi_{\mathrm{eq}}(|\boldsymbol{\xi}|)
  + \delta\varphi(\mathbf{x},t)
  + \order{a^{2}},
\end{equation}
with
\begin{equation}\label{eq:delta_phi_accel}
  \delta\varphi(\mathbf{x},t)
  = \frac{(\mathbf{a}\cdot\hat{x})\,|\mathbf{x}|}
         {2c^{2}}\,
    \varphi_{\mathrm{eq}}^{\prime}(|\mathbf{x}|),
\end{equation}
where the factor $1/2$ comes from
$v(t) = at$ evaluated at the retardation-averaged time
$t/2$.
\emph{Notational remark.}---The combination
$|\mathbf{x}|\,\varphi_{\mathrm{eq}}^{\prime}(|\mathbf{x}|) =
r_{s}/|\mathbf{x}|$ on the Coulomb profile should be read as
the dimensionless shorthand for the dipole amplitude of the
retardation-induced displacement; the directly retarded
expansion gives $\delta\varphi(\mathbf{x}) =
-r_{s}(\mathbf{a}\cdot\hat{x})/(2c^{2})$ at leading order in
$a\,|\mathbf{x}|/c^{2}$, and the apparent $|\mathbf{x}|^{-1}$
factor is reabsorbed by the $4\pi r^{2}\,dr$ measure of the
subsequent volume integration.  The final inertial force
law~(\ref{eq:F_ma}) is dimensionally consistent regardless of
this intermediate convention.
This dipolar perturbation is odd under
$\hat{x} \to -\hat{x}$: the field is enhanced
in the direction opposite to $\mathbf{a}$ (trailing
side) and depleted ahead.

\textit{Pressure perturbation.}---The Bernoulli
pressure is
$P = P_{0} - \frac{e^{\varphi}}{32\pi G}
|\nabla\varphi|^{2}$.
Expanding to first order in $\delta\varphi$:
\begin{equation}\label{eq:P_asymmetry}
  \delta P
  = -\frac{e^{\varphi_{\mathrm{eq}}}}{16\pi G}\,
    \nabla\varphi_{\mathrm{eq}} \cdot
    \nabla\!\left(\delta\varphi\right)
  - \frac{e^{\varphi_{\mathrm{eq}}}}{32\pi G}\,
    |\nabla\varphi_{\mathrm{eq}}|^{2}\,
    \delta\varphi.
\end{equation}

\textit{Force integral.}---The net force on the oscillon
is $\mathbf{F} = -\oint d\mathbf{S}\;\delta P$.
Integration by parts converts the volume integral
$\int d^{3}x\,\nabla(\delta P)$ into the surface form;
since $\delta\varphi \propto (\mathbf{a}\cdot\hat{x})$
is purely $\ell = 1$, only the dipolar harmonic survives.
Substituting
Eq.~(\ref{eq:delta_phi_accel}) into
Eq.~(\ref{eq:P_asymmetry}) gives
\begin{equation}\label{eq:delta_P_explicit}
  \delta P
  = -\frac{e^{\varphi_{\mathrm{eq}}}}{16\pi G}\,
    \frac{(\mathbf{a}\cdot\hat{x})}{2c^{2}}
    \Bigl[
      r\,\bigl(\varphi_{\mathrm{eq}}^{\prime}\bigr)^{2}
      + r\,\varphi_{\mathrm{eq}}^{\prime}\,
        \varphi_{\mathrm{eq}}^{\prime\prime}
      + |\nabla\varphi_{\mathrm{eq}}|^{2}\,
        r\,\varphi_{\mathrm{eq}}^{\prime}
    \Bigr],
\end{equation}
where all derivatives are evaluated at $r = |\mathbf{x}|$.
The angular integration projects
$(\mathbf{a}\cdot\hat{x})\,\hat{x}$ onto $\mathbf{a}$:
\begin{equation}\label{eq:F_angular}
  \int d\Omega\;(\mathbf{a}\cdot\hat{x})\,\hat{x}
  = \frac{4\pi}{3}\,\mathbf{a}.
\end{equation}
Combining the radial terms (the first two in
Eq.~(\ref{eq:delta_P_explicit}) add to a total
derivative $\partial_{r}[r(\varphi_{\mathrm{eq}}')^{2}]$
whose boundary terms vanish, leaving the third)
yields
\begin{equation}\label{eq:F_radial}
  \mathbf{F}_{\mathrm{inertia}}
  = -\frac{\mathbf{a}}{c^{2}}
    \int_{0}^{\infty}\!dr\;4\pi r^{2}\,
    \frac{e^{\varphi_{\mathrm{eq}}}}{32\pi G}\,
    \bigl(\varphi_{\mathrm{eq}}^{\prime}\bigr)^{2}
  \equiv
  -\frac{E_{\varphi}}{c^{2}}\,\mathbf{a},
\end{equation}
where
\begin{equation}\label{eq:E_phi_def}
  E_{\varphi}
  \equiv \int d^{3}x\,
    \frac{e^{\varphi_{\mathrm{eq}}}}{32\pi G}\,
    |\nabla\varphi_{\mathrm{eq}}|^{2}
\end{equation}
is the total gradient (kinetic) energy of the
equilibrium oscillon---precisely the quantity
identified in Eq.~(\ref{eq:rho_schwarz}) as the
scalar-field energy density integrated over all space.

\textit{Identification with effective mass.}---By the
Bernoulli identity~(\ref{eq:bernoulli_scalar}),
$E_{\varphi}$ equals the total pressure deficit, which
is the gravitational binding energy.  From
Eq.~(\ref{eq:meff_quantum}), $E_{\varphi} = \meff c^{2}$.
Therefore:
\begin{equation}\label{eq:F_ma}
  \boxed{
  \mathbf{F}_{\mathrm{inertia}}
  = -\meff\,\mathbf{a}
  = -m_{0}\,e^{\varphi/2}\,\mathbf{a}.
  }
\end{equation}

This is Newton's second law, derived from the causal
(retarded) propagation of the scalar field during
acceleration.  The derivation is perturbative in
$a/(c\omega_{0}) \ll 1$; at higher orders, one expects
radiation-reaction corrections of the
Abraham--Lorentz type by analogy with the standard
electromagnetic and gravitational self-force structure,
scaling as $\mathbf{F}_{\mathrm{rad}} \sim
(2e^{\varphi}\,r_{s}/3c^{3})\,\dot{\mathbf{a}}$.
\emph{Status.}---This scaling is quoted by analogy
with the canonical Abraham--Lorentz expression for a
charged particle and is not derived from the scalar
self-field in this paper; a controlled derivation of
the third-derivative term within the ISPG retarded
expansion is left to future work.

For the localized resonance in the weak-field regime considered here,
the inertial mass $m_{\mathrm{inertial}} = m_{0}\,e^{\varphi/2}$
and the gravitational mass assignment
(Eq.~\ref{eq:meff_quantum}) are tied to the same scalar profile and
Bernoulli pressure-deficit bookkeeping.

%-------------------------------------------------------------------
\subsection{The test-particle equivalence principle as an identity}
\label{sec:equivalence}
%-------------------------------------------------------------------

In GR, the equality $m_{\mathrm{inertial}} =
m_{\mathrm{gravitational}}$ is the weak equivalence
principle, elevated to a foundational axiom.  In ISPG, the
test-particle version is first secured by minimal coupling:
matter follows geodesics of the single physical metric
$g_{\mu\nu}[\varphi]$.  Within the localized resonance
bookkeeping used here, the same equality is also reflected in the
unified scalar-field mechanism: under the identifications
established earlier in this section---namely that the
gravitational mass is the Bernoulli pressure-deficit
integral~(\ref{eq:meff_quantum}) and that the inertial
force takes the form~(\ref{eq:F_radial}) with the same
integrand---both masses reduce to the single quantity
\begin{equation}\label{eq:equivalence}
  m_{\mathrm{grav}} = \frac{E_{\mathrm{Bernoulli}}}{c^2}
  = \frac{1}{c^2}\int d^3x\,
    \frac{e^{\varphi}}{32\pi G}\,|\nabla\varphi|^2
  = m_{\mathrm{inertial}}.
\end{equation}
The equality is therefore not an independent postulate
within this approximation: once the gravitational and inertial
identifications above are accepted, both rely on the same
Bernoulli integral over the same equilibrium field configuration.
The actual equation of motion for test matter remains the
geodesic equation of the metric defined by this profile; the
Bernoulli integral supplies the pressure-deficit interpretation
of why that profile is gravitational.  Inertia is the dynamic
manifestation of the same profile bookkeeping, through resistance
to profile readjustment during acceleration.
\emph{Status.}---Derived here under the assumptions of
weak field, slow acceleration $a \ll c\omega_{0}$, and
spherical equilibrium oscillon profile; the
\emph{logical} status is that of a structural identity
within these approximations, not of an exact theorem
of the full nonlinear theory.

\textit{Connection to the MOND transition: a common
dynamical mechanism.}---The inertial mechanism derived
above has a direct galactic-scale counterpart.  In both
cases, \emph{acceleration} creates a pressure deficit
that the medium cannot replenish fast enough.

\begin{itemize}
\item \textbf{Inertia} (local): an oscillon undergoes
  linear acceleration $\mathbf{a}$.  The pressure
  deficit behind the resonance cannot be refilled
  instantaneously because the field propagates at
  finite speed $c$ $\to$ fore--aft asymmetry $\to$
  restoring force $F = -m\,a$.
\item \textbf{MOND} (galactic): the galactic vortex
  subjects the medium to centrifugal acceleration.
  Archion excitation nodes are driven outward, while
  the peripheral nodes cannot flow inward fast enough to
  compensate $\to$ central rarefaction (vortex
  pressure deficit)~\cite{Lotishvili2026MOND} $\to$
  gravitational pull stronger than the Newtonian
  $1/r^{2}$ extrapolation from baryonic mass alone.
\end{itemize}

The structural similarity is qualitative: in both
cases an acceleration field acts on the medium, the
medium's finite response rate prevents equilibrium,
and the resulting pressure deficit contributes either to the
metric/geodesic response or to an effective inertial reaction.  The
analogy should not be read as a quantitative
equivalence---the inertial calculation above is a
microscopic perturbative result around a single
oscillon profile, whereas the MOND derivation
in Ref.~\cite{Lotishvili2026MOND} is a macroscopic
fluid argument around a galactic vortex.  The
qualitative parallel is in the mechanism (causal
delay $\to$ pressure deficit $\to$ metric or inertial response);
the quantitative difference lies in the source of
acceleration (oscillon recoil vs.\ galactic
rotation) and the spatial scale ($R_{c} \sim
\hbar/(mc)$ vs.\ $\ell \sim$ kpc).  The MOND
acceleration scale $a_{0} = cH/(2\pi)$
itself reflects the largest distance over which
the field can sustain a coherent pressure gradient:
$\lambda_{H} = 2\pi c/H$, beyond which the Hubble
damping renders the response
overdamped~\cite{Lotishvili2026MOND}.  This sets
the boundary between the Newtonian regime (where the
medium responds fully) and the MOND regime (where
the centrifugal rarefaction is no longer compensated).

\textit{Contrast with Mach's principle.}---The
inertial mechanism derived above is intrinsically
\emph{local}: the restoring force on an accelerating
oscillon arises from the retarded readjustment of
its own scalar self-field within a region of order
$R_{c} \sim \hbar/(mc)$, not from a coupling to the
rest of the matter content of the universe.  In
Mach's original conception, the inertia of a body is
determined by the gravitational influence of all
distant masses; in the present picture, distant matter
enters only through the asymptotic value
$\varphi_{\infty}$ that fixes the conformal factor
$e^{\varphi/2}$ in $\meff = m_{0}\,e^{\varphi/2}$.
The dynamical content of inertia---the proportionality
between force and acceleration---is set entirely by
the local retardation of the oscillon's own field.
ISPG is therefore Machian only in the weak sense that
$\meff$ inherits an environmental factor through
$\varphi_{\infty}$; it is non-Machian in the strong
sense that the inertial response itself does not
require knowledge of the cosmological matter
distribution.

%-------------------------------------------------------------------
\subsection{Hysteresis: the memory of the medium}
\label{sec:hysteresis_quantum}
%-------------------------------------------------------------------

The readjustment time $t_{\mathrm{adj}} \sim 1/\omega_0$
(Eq.~\ref{eq:t_readjust}) implies that the medium
retains a memory of the resonance's previous state.
When a body accelerates and then stops, the pressure
distribution does not instantaneously snap to the new
equilibrium---it relaxes over a finite timescale.
This is hysteresis: the delayed response of the medium
to changes in the source configuration.

In the classical theory~\cite{Lotishvili2026}, this
memory effect is captured by the hysteresis field $\chi$,
an effective collective field with standalone covariant dynamics in
the quasi-static bound-structure sector.  In that one-way effective
sector its source is treated as an external quasi-static functional
of the background scalar, and the covariant action template is
\begin{equation}\label{eq:S_chi_quantum}
  S_\chi = \int d^4x\,\sqrt{-g}\left[
    -\frac{1}{2}\,\nabla_\mu\chi\,\nabla^\mu\chi
    - \frac{\lambda}{2}\bigl(\chi - \beta\,Y_{\rm qs}\bigr)^2
  \right],
\end{equation}
where
$Y_{\rm qs} \equiv \sqrt{\nabla_\mu\varphi\,\nabla^\mu\varphi}$
is used only in the spacelike-gradient, quasi-static
bound-structure regime, and $\lambda > 0$ is the stiffness
parameter.  The standalone $\chi$
field equation is
\begin{equation}\label{eq:chi_kg_quantum}
  \Box\chi - \lambda\,\chi
  = -\lambda\,\beta\,Y_{\rm qs},
\end{equation}
a Klein--Gordon equation with mass $m_\chi = \sqrt{\lambda}$
driven by the quasi-static gradient invariant $Y_{\rm qs}$.  The signs follow
the $(-,+,+,+)$ signature convention adopted throughout
this paper ($\Box = g^{\mu\nu}\nabla_\mu\nabla_\nu$),
reproducing MAIN~\cite{Lotishvili2026} Eq.~(chi\_KG)
in its effective-template form.  The fully varied backreacting
$(\varphi,\chi)$ principal-symbol analysis is not closed here.

We now provide the quantum-mechanical motivation
for this effective field and for the microscopic tracking scale.

\textit{Origin of $\chi$ from resonance dynamics.}---Each
oscillon generates a pressure deficit that extends to
$r \sim R_c$.  When the oscillon moves, the pressure
deficit must be re-established at the new position and
dissolved at the old one.  The dissolution occurs via
outgoing scalar waves---the same mechanism as the
resonant tails of Sec.~\ref{sec:entropy_mech}---with
a characteristic timescale $t_{\mathrm{adj}}$.

The hysteresis field $\chi$ is the coarse-grained
dynamical variable that tracks the dissolution process.
Its static equilibrium configuration is fixed by the
source term in Eq.~(\ref{eq:chi_kg_quantum}): setting
$\partial_t\chi = 0$ and $\nabla^2\chi \to 0$ at
long wavelengths gives $\chi_{\rm eq} = \beta\,Y_{\rm qs}$, which
coincides with MAIN~\cite{Lotishvili2026}
Eq.~(chieq).  The \emph{dynamical} content of $\chi$
lives in the deviation from this equilibrium,
\begin{equation}\label{eq:chi_interpretation}
  \delta\chi(\mathbf{x}, t)
  \equiv \chi(\mathbf{x}, t) - \chi_{\rm eq}(\mathbf{x}, t)
  = \chi(\mathbf{x}, t) - \beta\,Y_{\rm qs}(\mathbf{x}, t),
\end{equation}
which is the quantity that can relax to zero when the
source motion stops, provided damping, outgoing radiation, or
coarse-graining is included.  Equivalently, when the source has
just moved, $Y_{\rm qs}$ evaluated on the actual (lagging) field
$\varphi$ differs from $Y_{\rm qs}$ evaluated on the instantaneous
equilibrium field $\varphi_{\rm eq}$; this difference is
what drives $\delta\chi \neq 0$ during the transient and
is dissolved on the relaxation timescale discussed below.
The full field $\chi = \chi_{\rm eq} + \delta\chi$ is the
object that appears in the action~(\ref{eq:S_chi_quantum})
and in all subsequent macroscopic formulas, in agreement
with MAIN.

\textit{Microscopic and macroscopic stiffness scales.}---The
microscopic tracking stiffness and the macroscopic collective
relaxation scale are distinct quantities.  From
Sec.~\ref{sec:resonances}, a quasi-bound oscillon has
a complex quasi-normal frequency
$\omega = \omega_R - i\Gamma/2$, where the real part
$\omega_R = m_0 c^2/\hbar$ gives the mass and the
imaginary part $\Gamma$ gives the decay width.  The
decay width represents the rate at which the oscillon
loses energy to outgoing scalar waves---precisely the
process that re-equilibrates the pressure distribution
when the source moves.

For a stable particle (proton, electron), the barrier
penetration is exponentially suppressed (Eq.~\ref{eq:gamma_wkb}):
\begin{equation}\label{eq:gamma_stable}
  \Gamma \sim \omega_R \exp\!\left[-2\int_{r_1}^{r_2}
    \sqrt{V_{\mathrm{eff}} - \omega_R^2 e^{-2\varphi_0}}\,dr\right]
  \ll \omega_R.
\end{equation}
The relaxation timescale is $t_{\mathrm{rel}} \sim
1/\Gamma \gg t_{\mathrm{adj}}$.  However, for the
\emph{coherent} readjustment of the pressure profile
around a \emph{macroscopic} body composed of $N \sim 10^{68}$
oscillons, the relaxation is dominated not by individual
decays but by the collective dephasing of the oscillon
ensemble.

For a \emph{single} particle, the microscopic stiffness
$\lambda_{\mathrm{micro}}$ is set by the inverse Compton
wavelength:
\begin{equation}\label{eq:lambda_micro}
  m_{\chi}^{(\mathrm{micro})}
  = \sqrt{\lambda_{\mathrm{micro}}}
  = \frac{m_{0}\,c}{\hbar}\,e^{\varphi/2}.
\end{equation}
For a proton ($m_{0} \approx 1\;\mathrm{GeV}/c^{2}$),
$m_{\chi}^{(\mathrm{micro})} \sim 5 \times 10^{15}\;\mathrm{m}^{-1}$,
corresponding to
$\lambda_{\mathrm{micro}} \sim 10^{31}\;\mathrm{m}^{-2}$.
This extremely large stiffness means that for individual
particles, $\chi$ tracks the equilibrium $\beta Y_{\rm qs}$
almost instantaneously: the readjustment is completed in
one Compton period $t_{\mathrm{adj}} \sim 10^{-24}$~s.

For a \emph{macroscopic body} (galaxy, cluster) composed
of $N$ oscillons, the collective relaxation is governed
by a different scale.  Each oscillon radiates outgoing
scalar waves at its individual (exponentially suppressed)
decay rate $\Gamma_{\mathrm{single}}$
(Eq.~\ref{eq:gamma_stable}).  Two limiting regimes are
available for how the $N$ single-particle rates combine
into a collective relaxation rate
$\Gamma_{\mathrm{macro}}$:
\begin{itemize}
  \item the fully incoherent regime, where independent
    random phases give a $\sqrt{N}$ central-limit
    enhancement of the radiated amplitude and hence a
    $\sqrt{N}$ enhancement of the rate;
  \item the Dicke super-radiant regime, where phase
    coherence across the ensemble gives a linear-in-$N$
    enhancement.
\end{itemize}
Parametrically we therefore write
\begin{equation}\label{eq:gamma_collective}
  \Gamma_{\mathrm{macro}}
  = \mathcal{C}(N)\,\Gamma_{\mathrm{single}},
\end{equation}
with $\mathcal{C}(N) \in [\sqrt{N},\,N]$ bounded above
by the Dicke limit and below by the fully-incoherent
limit.  The effective macroscopic relaxation timescale
is
\begin{equation}\label{eq:t_macro}
  t_{\mathrm{rel}}^{(\mathrm{macro})}
  = \frac{1}{\Gamma_{\mathrm{macro}}}
  = \frac{1}{\mathcal{C}(N)\,\Gamma_{\mathrm{single}}}.
\end{equation}
\textit{Status of the cosmological matching.}---A
closed-form first-principles determination of
$\mathcal{C}(N)$ in the oscillon medium (whose
wavelength-to-body-size ratio lies between the two
limits above) is a self-contained technical task and
the topic of a separate work.  Consequently the
quantitative matching of
$t_{\mathrm{rel}}^{(\mathrm{macro})}$ to the Hubble
time---which, taking $\Gamma_{\mathrm{single}}
\sim \omega_R e^{-2S_{\mathrm{barrier}}}$ with
$\omega_R \sim m_p c^2/\hbar$ and
$S_{\mathrm{barrier}} \gtrsim 76$ (the lower bound fixed
by the proton lifetime $\tau_p > 10^{34}$~yr, combined
with the WKB expression~\ref{eq:gamma_wkb})
---cannot be claimed as a derived prediction at the
present level of the analysis.  What is derived is the
parametric structure~(\ref{eq:gamma_collective}) and the
two regime bounds; matching the Hubble-time target
requires either an independent calculation of
$\mathcal{C}(N)$ or an independent determination of
$S_{\mathrm{barrier}}$ beyond the proton-stability
lower bound.

The macroscopic Klein--Gordon mass governing
galaxy-scale $\chi$ dynamics is therefore
\begin{equation}\label{eq:lambda_macro}
  \boxed{
  m_{\chi}^{(\mathrm{macro})}
  = \sqrt{\lambda_{\mathrm{macro}}}
  = \frac{\mathcal{C}(N)\,\Gamma_{\mathrm{single}}}{c}
  = \frac{\mathcal{C}(N)\,e^{-2S_{\mathrm{barrier}}}\,m_{0}\,c}
         {\hbar}\,e^{\varphi/2}.
  }
\end{equation}
This mass scale is suppressed relative to the
microscopic value by the factor
$\mathcal{C}(N)\,e^{-2S_{\mathrm{barrier}}} \ll 1$,
placing $m_{\chi}^{(\mathrm{macro})}$ in the intermediate
regime between particle-physics energies and the Hubble
scale---precisely the window relevant for galaxy-halo
phenomenology.  The absolute normalisation within that
window depends on $\mathcal{C}(N)$, whose determination
is deferred as noted above.

\textit{Two regimes of $\chi$.}---The hysteresis
field thus operates in two distinct regimes:
\begin{itemize}
  \item \textit{Microscopic}
    ($\lambda = \lambda_{\mathrm{micro}}$): $\chi$ tracks
    $\beta Y_{\rm qs}$ instantaneously; no observable hysteresis
    at the single-particle level.
  \item \textit{Macroscopic}
    ($\lambda = \lambda_{\mathrm{macro}}$): the
    collective relaxation is slow enough that
    $\chi \neq \beta Y_{\rm qs}$ for timescales comparable to
    galactic dynamical times, producing the halo
    phenomenology modeled in~\cite{Lotishvili2026}.
\end{itemize}

\textit{Universality of $a_0$ under $N$-dependent
$\lambda_{\mathrm{macro}}$.}---The $N$-dependence of
$\lambda_{\mathrm{macro}}$ through $\mathcal{C}(N)$
(Eq.~\ref{eq:lambda_macro}) does not contradict the
universal MOND acceleration scale
$a_0$~\cite{Lotishvili2026}: the covariant static
limit $\chi \to \chi_{\rm eq} = \beta Y_{\rm qs}$ makes the
BTFR-type halo normalisation depend only on the
closure normalization $\beta$ (to be fixed by the
bound-structure coarse-graining and the cosmological scale
$a_0 = cH_0/2\pi$), while
$\lambda$ controls only the \emph{approach} to that
equilibrium.  Thus the system-specific $\lambda_{\rm
macro}$ governs transients and relaxation spectra,
whereas the time-averaged rotation-curve phenomenology tests the
closure normalization.

\textit{Ghost freedom.}---The $\chi$ sector is ghost-free:
the Hamiltonian density
\begin{equation}\label{eq:H_chi_q}
  \mathcal{H}_\chi
  = \frac{c^2}{2}\,\pi_\chi^2
  + \frac{1}{2}\,(\partial_i\chi)^2
  + \frac{\lambda}{2}\,(\chi - \beta Y_{\rm qs})^2
  \geq 0,
\end{equation}
is manifestly non-negative for $\lambda > 0$, ensuring
that no negative-norm states propagate and the theory
is stable at the quantum level.

%-------------------------------------------------------------------
\subsection{Microscopic motivation of the hysteresis action template}
\label{sec:chi_derivation}
%-------------------------------------------------------------------

The action $S_\chi$~(Eq.~\ref{eq:S_chi_quantum}) was
introduced in the classical theory~\cite{Lotishvili2026}
as an effective closure.  We now motivate its microscopic
tracking scale and KG structure from resonance dynamics.  This is a
coarse-graining template, not a completed first-principles derivation
of the full macroscopic hysteresis sector.

\textit{Non-equilibrium decomposition.}---Let
$\varphi_{\mathrm{eq}}(\mathbf{x},t)$ denote the
instantaneous equilibrium field---the solution of
$\Box\varphi = -(8\pi G/c^4)\,T$ with the source at its
current position and all transients decayed.  The actual
field lags behind due to finite propagation speed:
\begin{equation}\label{eq:phi_noneq}
  \varphi(\mathbf{x},t)
  = \varphi_{\mathrm{eq}}(\mathbf{x},t)
  + \delta\varphi(\mathbf{x},t),
\end{equation}
where $\delta\varphi$ is the retardation-induced excess.
This is the same decomposition used in the inertia
derivation (Eq.~\ref{eq:phi_ret_expand}), where
$\delta\varphi$ was shown to be a dipolar perturbation of
order $a\,r/c^2$ (Eq.~\ref{eq:delta_phi_accel}).

\textit{Perturbation dynamics inside the oscillon.}---In
the far field ($r \gg R_c$), $\delta\varphi$ propagates as
a free massless wave ($\Box\delta\varphi = 0$) and
disperses.  Inside the oscillon ($r \lesssim R_c$),
however, the perturbation is confined by the effective
potential $V_{\mathrm{eff}}(r)$ of
Eq.~(\ref{eq:quasi_normal}).  At late times (after
initial transients have radiated away), the confined
perturbation is dominated by the longest-lived
quasi-normal modes.  We write the late-time
approximation as
\begin{equation}\label{eq:qnm_expansion}
  \delta\varphi(\mathbf{x},t)
  \approx \sum_n a_n(t)\,\phi_n(\mathbf{x}),
\end{equation}
where the sum runs over the discrete quasi-normal modes
$\phi_n$ of the confining potential, and each amplitude
$a_n$ satisfies the damped oscillator equation
\begin{equation}\label{eq:a_n_eom}
  \ddot{a}_n + \Gamma_n\,\dot{a}_n
  + \omega_{n}^2\,a_n = f_n(t).
\end{equation}
Here $\omega_n = \omega_{n,R}$ is the real part of the
$n$-th quasi-normal frequency, $\Gamma_n$ is the decay
width (Eq.~\ref{eq:gamma_wkb}), and $f_n(t)$ is the
driving force projected onto mode $n$ from the source
motion.

\textit{Fundamental mode dominance.}---For a slow
perturbation (acceleration timescale $\gg 1/\omega_0$),
the driving force couples predominantly to the fundamental
mode ($n = 0$) with locally measured frequency
$\omega_0 = (m_0 c^2/\hbar)\,e^{\varphi_0/2}$
(Eq.~\ref{eq:lambda_micro}).
Higher overtones ($n \geq 1$) are suppressed by the
ratio of the perturbation timescale to the mode spacing.
The pressure deficit deviation is therefore controlled by
the fundamental amplitude $a_0$:
\begin{equation}\label{eq:chi_a0}
  \chi(\mathbf{x},t)
  \approx \beta\,
    \frac{\nabla_\mu\varphi_{\mathrm{eq}}\,
      \nabla^\mu[\,a_0(t)\,\phi_0(\mathbf{x})\,]}
         {Y_{{\rm qs},\mathrm{eq}}},
\end{equation}
where $Y_{{\rm qs},\mathrm{eq}} =
\sqrt{\nabla_\mu\varphi_{\mathrm{eq}}\,
\nabla^\mu\varphi_{\mathrm{eq}}}$ and the linearization
$Y_{\rm qs} \approx Y_{{\rm qs},\mathrm{eq}} +
\hat{n}_\mu\nabla^\mu\delta\varphi$
(with $\hat{n}_\mu = \nabla_\mu\varphi_{\mathrm{eq}}
/Y_{{\rm qs},\mathrm{eq}}$) has been used.

\textit{Covariant effective equation.}---Equation
(\ref{eq:a_n_eom}) governs the temporal amplitude $a_0$
of a single oscillon.  For a macroscopic body, $\chi$ is
the coarse-grained superposition of the pressure
deviations from all $N$ constituent oscillons,
each located at $\mathbf{r}_a$:
$\chi(\mathbf{x},t)
= \sum_a \chi_a(\mathbf{x} - \mathbf{r}_a, t)$.
In the continuum limit, this sum defines a smooth field.
Because each individual perturbation $\delta\varphi_a$
propagates via the scalar wave equation
($\Box\delta\varphi_a = 0$ outside the source), the
coarse-grained field inherits a spatial d'Alembertian;
the temporal part retains the damping and restoring
force from the confining potential.  The resulting
equation for $\chi$ is
\begin{equation}\label{eq:chi_derived}
  \Box\chi
  - \frac{\Gamma_0}{c^2}\,\partial_t\chi
  - \frac{\omega_0^2}{c^2}\,\chi
  = -\frac{\omega_0^2}{c^2}\,\beta\,Y_{\rm qs}.
\end{equation}
For stable particles (proton, electron), the decay width
is exponentially suppressed:
$\Gamma_0 \sim \omega_0\,e^{-2S_{\mathrm{barrier}}} \ll
\omega_0$ (Eq.~\ref{eq:gamma_stable}).  The damping term
$(\Gamma_0/c^2)\,\partial_t\chi$ is therefore negligible
compared to the mass term
$(\omega_0^2/c^2)\,\chi$, and
Eq.~(\ref{eq:chi_derived}) reduces to
\begin{equation}\label{eq:chi_KG_derived}
  \Box\chi - \lambda\,\chi
  = -\lambda\,\beta\,Y_{\rm qs},
\end{equation}
with
\begin{equation}\label{eq:lambda_derived}
  \lambda
  = \frac{\omega_0^2}{c^2}
  = \frac{m_0^2\,c^2}{\hbar^2}\,e^{\varphi},
\end{equation}
which is precisely the microscopic stiffness
$\lambda_{\mathrm{micro}}$ of
Eq.~(\ref{eq:lambda_micro}).
Equation~(\ref{eq:chi_KG_derived}) is identical to
the phenomenological field
equation~(\ref{eq:chi_kg_quantum}).

\textit{Effective-action reconstruction.}---The motivated
microscopic tracking equation~(\ref{eq:chi_KG_derived}) is reproduced
by the Euler--Lagrange equation of the effective action template
\begin{equation}\label{eq:S_chi_derived}
  S_\chi = \int d^4x\,\sqrt{-g}\left[
    -\frac{1}{2}\,\nabla_\mu\chi\,\nabla^\mu\chi
    - \frac{\lambda}{2}\bigl(\chi - \beta\,Y_{\rm qs}\bigr)^2
  \right],
\end{equation}
which matches $S_\chi$~(Eq.~\ref{eq:S_chi_quantum})
in the quasi-static source sector.  The kinetic term arises from the
wave propagation of the confined perturbation; the microscopic stiffness
$\lambda = \omega_0^2/c^2$ is the square of the inverse
Compton wavelength, reflecting the confining scale of
the oscillon; and the source $\beta Y_{\rm qs}$ is the equilibrium
value toward which the medium relaxes.
This $\chi$ is the dynamical coarse-grained hysteresis field; in the quasi-static galactic limit it is the field-level representative of the local transported-response multiplier written in the companion MOND paper~\cite{Lotishvili2026MOND} for the coherent vortex branch.

\textit{Parameter status.}---The microscopic argument identifies
some, but not all, of the previously phenomenological inputs:
\begin{enumerate}
\item[(i)] The microscopic stiffness is
  $\lambda_{\mathrm{micro}} = \omega_0^2/c^2
  = (m_0 c/\hbar)^2\,e^{\varphi}$, set by the oscillon's
  fundamental frequency and hence by the particle mass.
  This is the origin of the hierarchy between microscopic
  and macroscopic hysteresis: the microscopic relaxation
  is completed in one Compton period
  $t_{\mathrm{adj}} \sim 1/\omega_0
  \sim 10^{-24}\;\mathrm{s}$ (proton), while the
  macroscopic relaxation depends on the collective decay
  of $N$ oscillons (Eq.~\ref{eq:lambda_macro}).
\item[(ii)] The dissipation rate $\Gamma_0$ is
  exponentially suppressed
  (Eq.~\ref{eq:gamma_stable}) for stable particles,
  justifying the neglect of the damping term.  For
  unstable resonances ($\Gamma \sim \omega_0$), the full
  damped equation~(\ref{eq:chi_derived}) applies, and
  the hysteresis field decays on the particle's
  lifetime---as expected physically.
\item[(iii)] The coupling $\beta$ is motivated by the
  Bernoulli relation
  (Sec.~\ref{sec:stress_energy}): the pressure deficit
  $\delta P \propto Y_{{\rm qs},\mathrm{eq}}\,\delta Y_{\rm qs}$,
  so $\beta$ normalizes $\chi$ to the pressure deviation
  per unit gradient.  Its numerical value remains a
  phenomenological closure normalization until the full
  bound-structure coarse-graining is completed.
\end{enumerate}

\textit{Status of the construction.}---The hysteresis
action $S_\chi$ is therefore a microscopically motivated effective
template.  The quasi-normal mode structure, fundamental-mode dominance
for slow perturbations, and exponential suppression of the decay width
for stable particles motivate the microscopic tracking stiffness
$\lambda_{\rm micro}$.  They do not by themselves close the
macroscopic collective relaxation scale, the full backreacting
principal-symbol analysis, or the first-principles value of $\beta$.
Those remain Route~B/coarse-graining tasks to be tested against
galaxy-halo rotation curves~\cite{Lotishvili2026}.

%-------------------------------------------------------------------
\subsection{Summary: the three faces of the medium}
\label{sec:three_faces}
%-------------------------------------------------------------------

The scalar-metric medium exhibits three distinct
mechanical properties, all tied to the same
resonance--medium coupling:

\begin{enumerate}
  \item \textbf{Pressure} ($\nabla\varphi \neq 0$
    around matter): the Bernoulli deficit that produces
    gravity (Sec.~\ref{sec:bernoulli_identity}).
  \item \textbf{Zero drag at constant velocity}:
    Lorentz invariance of the wave equation ensures
    that uniformly moving resonances experience no
    friction (Sec.~\ref{sec:uniform_motion}).
  \item \textbf{Memory} (hysteresis): the finite
    readjustment time of the oscillon profile during
    acceleration produces inertia and, on macroscopic
    scales, the $\chi$-field dynamics responsible for
    galaxy-halo phenomenology
    (Secs.~\ref{sec:hysteresis_quantum}
    and~\ref{sec:chi_derivation}).
\end{enumerate}

These three properties---pressure without friction,
combined with memory---characterize a unique physical
medium: a \emph{superfluid with hysteresis}.  The
equivalence principle is not an axiom but the
inevitable consequence of the fact that gravity and
inertia are the static and dynamic aspects of a
single quantity: the Bernoulli pressure deficit of a
resonant excitation.

%===================================================================
\section{Entropy, self-regulation, and dark energy}\label{sec:entropy}
%===================================================================

The Bernoulli mechanism (Sec.~\ref{sec:bernoulli_identity}) explains how
matter generates a local pressure deficit.  We now turn to the
cosmological consequences: the irreversible spreading of resonant
tails produces a global entropy increase, which is dynamically
regulated by a negative-feedback loop.  This mechanism provides a
\emph{structural channel} for estimating the observed dark-energy
scale, but it is not yet a completed solution of the full
cosmological-constant problem: the cumulative tail-energy budget is
itself small (see Eq.~(\ref{eq:E_cumulative}) below), and the
relevance for $\rho_{\Lambda}$ enters through a
perturbative-branch feedback-regulated equilibrium developed in
Sec.~\ref{sec:cc_problem}.  The same dynamics predicts an evolving
dark-energy equation of state $w(z) \neq -1$ in the effective
fluid description.  The quantitative success of these predictions
hinges on the magnitude of the per-particle tail rate
$\mathcal{P}_{\mathrm{tail}}$; that magnitude depends on
non-perturbative oscillon stability physics whose precise
evaluation is beyond the scope of this paper, as explained below.

%-------------------------------------------------------------------
\subsection{Resonant tails and irreversible entropy production}
\label{sec:entropy_mech}
%-------------------------------------------------------------------

A localized resonant excitation (particle) consists of a concentrated
core and extended oscillatory tails that propagate outward at the
local speed of light (Sec.~\ref{sec:resonances}).  These tails carry
energy away from the core and into the surrounding medium.  Once
emitted, a tail cannot be recalled: the process is irreversible
by the same argument that makes electromagnetic radiation irreversible
(retarded boundary conditions).

The energy flux carried by the tail of a single particle of mass $m$
is determined by the radiative damping of the quasi-bound oscillon.
We compute this explicitly from the outgoing energy flux through a
spherical surface surrounding the oscillon core.

\textit{Energy flux for the scalar field.}---The stress-energy
tensor component $T^{0r}$ gives the radial energy flux.  For
the bi-conformal metric with $\varphi = \varphi_0 + \delta\varphi$,
\begin{equation}\label{eq:T_0r}
  T^{0r} = \frac{1}{16\pi G}\, g^{00}\,g^{rr}\,
    \partial_0(\delta\varphi)\,\partial_r(\delta\varphi)
  = -\frac{e^{2\varphi_0}}{16\pi G}\,
    \dot{\Phi}\,\Phi',
\end{equation}
where $\dot{\Phi} \equiv \partial_t\Phi$ and $\Phi' \equiv \partial_r\Phi$.

\textit{Outgoing wave at infinity.}---Far from the oscillon core
($r \gg R_{\mathrm{core}}$), the field behaves as an outgoing
spherical wave:
\begin{equation}\label{eq:outgoing_wave}
  \Phi(r,t) \sim \frac{A_{\infty}}{r}\,
    \cos(\omega t - kr), \qquad k = \omega/c.
\end{equation}
The radial derivative is $\Phi' \sim -\frac{A_{\infty}}{r^2}\cos(\omega t - kr)
+ \frac{kA_{\infty}}{r}\sin(\omega t - kr)$; at large $r$, the dominant
term is $\Phi' \approx k\Phi$ (in phase quadrature with $\dot{\Phi}$).

\textit{Time-averaged power.}---The time-averaged flux through a
sphere of radius $r$ is
\begin{equation}
  \langle T^{0r} \rangle = \frac{e^{2\varphi_0}}{32\pi G}\,
    \omega k \,\frac{A_{\infty}^2}{r^2}.
\end{equation}
The total radiated power is the flux times the area $4\pi r^2$:
\begin{equation}\label{eq:P_wave}
  \mathcal{P}_{\mathrm{rad}} = 4\pi r^2 \langle T^{0r} \rangle
  = \frac{e^{2\varphi_0}}{8G}\,\omega k \,A_{\infty}^2.
\end{equation}

\textit{Amplitude matching.}---The asymptotic amplitude $A_{\infty}$
is determined by matching to the oscillon core.  For a core of
radius $R$ with amplitude $A_{\mathrm{core}}$, continuity requires
$A_{\infty} \sim A_{\mathrm{core}} R$.  The core amplitude is fixed
by the oscillon mass: the core energy is $E_{\mathrm{core}} \sim
A_{\mathrm{core}}^2 \omega^2 R^3 \sim mc^2$, giving $A_{\mathrm{core}} \sim
\sqrt{m}/(\omega R^{3/2})$.  Substituting:
\begin{equation}
  A_{\infty}^2 \sim A_{\mathrm{core}}^2 R^2
  \sim \frac{m}{\omega^2 R}.
\end{equation}

Combining with Eq.~(\ref{eq:P_wave}), using $\omega = mc^2/\hbar$
and $k = \omega/c$, and tracking the $c$-factors of the kinetic
prefactor in the action carefully (an explicit re-derivation
including the $c^{4}/(16\pi G)$ overall normalisation of the
gradient sector is left as a technical exercise; see the
discussion below), the perturbative scaling reduces to the
dimensionally consistent form
\begin{equation}\label{eq:P_tail_derived}
  \mathcal{P}_{\mathrm{tail}}^{\mathrm{(pert)}} \sim
  \frac{e^{2\varphi_0}}{8G}\,\frac{m^{3}c^{3}}{\hbar^{2}\,R}\;
  \xrightarrow{R\sim\hbar/(mc)}\;
  \frac{G\,m^{4}\,c^{3}}{\hbar^{2}}\,e^{2\varphi_0},
\end{equation}
where in the last step the core radius $R$ is set by the
Compton wavelength $R \sim \hbar/(mc)$ and the leading
$O(1)$ numerical coefficients have been absorbed into the
``$\sim$''.  The boxed asymptotic scaling is therefore
\begin{equation}\label{eq:P_tail}
  \boxed{
  \mathcal{P}_{\mathrm{tail}}^{\mathrm{(pert)}} \sim
  \frac{G\,m^{4}\,c^{3}}{\hbar^{2}}\,e^{2\varphi_0}\,.
  }
\end{equation}
Direct evaluation of this perturbative expression for an
electron at $\varphi_0 \approx 0$ gives
$\mathcal{P}_{\mathrm{tail}}^{\mathrm{(pert)}} \sim 10^{-37}$~W,
and for a nucleon $\sim 10^{-24}$~W.

This derivation confirms the $m^4$ scaling from dimensional and
scaling analysis: the oscillon radiates as an outgoing scalar wave,
with amplitude set by the core mass and frequency.  The exponential
factor $e^{2\varphi_0}$ reflects the gravitational redshift of the
radiation in the local field.

\textit{Perturbative vs.\ non-perturbative estimate.}---The
expression~(\ref{eq:P_tail}) is the perturbative scaling: it
captures only the gross power-counting of the radiating-oscillon
calculation and treats the core as a coherent classical mode of
amplitude fixed by $E_{\mathrm{core}} \sim mc^{2}$.  In the
oscillon literature, however, long-lived quasi-bound scalar
configurations decay much more slowly than the perturbative
$m^{4}$ scaling suggests: their lifetime is controlled by an
exponentially small overlap integral between the bound
profile and the free continuum, of the schematic form
$\mathcal{F}_{\mathrm{nonpert}} \sim e^{-S_{\mathrm{osc}}/\hbar}$
with $S_{\mathrm{osc}}$ the oscillon stability action.  An
estimate of the resulting suppression for ISPG-resonance
parameters is
\begin{equation}\label{eq:P_tail_nonpert}
  \mathcal{P}_{\mathrm{tail}}^{\mathrm{(np)}} \equiv
  \mathcal{F}_{\mathrm{nonpert}}\cdot
  \mathcal{P}_{\mathrm{tail}}^{\mathrm{(pert)}}
  \sim 10^{-107}\,\mathrm{W}\times e^{2\varphi_0}
  \quad\text{(schematic suppressed per-particle benchmark)},
\end{equation}
where the numerical value is not yet a derived
species-specific prediction.  The detailed evaluation of
$\mathcal{F}_{\mathrm{nonpert}}$
from the full nonlinear ISPG oscillon profile---and in particular
the verification that it is consistent both with the per-particle
energy budget here \emph{and} with the cosmological-constant
match of Sec.~\ref{sec:cc_problem}, which uses the perturbative
formula at face value in Planck units---is an open technical
task.  The remainder of this subsection uses the schematic value
$\mathcal{P}_{\mathrm{tail}}^{\mathrm{(np)}}
\sim 10^{-107}$~W\,$\times e^{2\varphi_0}$ only as a
suppressed-branch per-particle energy-budget benchmark; it is not
the normalization used for the perturbative cosmological scale
estimate below.

The total outgoing \emph{power} of the observable
universe into the background scalar field, for
$N_b \sim 10^{80}$ baryons, is
\begin{equation}\label{eq:S_dot}
  \dot{E}_{\mathrm{bg,\,tot}} \sim
  N_b \cdot \mathcal{P}_{\mathrm{tail}}^{\mathrm{(np)}}
  \sim 10^{80} \cdot 10^{-107} \sim 10^{-27}\;\mathrm{W},
\end{equation}
an exceedingly small but strictly positive outgoing power in the
suppressed branch.  The calculation establishes positive outgoing
energy flux under retarded boundary conditions.  Interpreting this
flux as entropy production requires a coarse-grained scalar bath and
an effective temperature; those thermodynamic variables are not
derived here, so we quote energy power rather than a closed
entropy-production law.
Over the age of the universe ($t_0 \sim 4 \times 10^{17}$~s),
the cumulative energy deposited into the background scalar
field is
\begin{equation}\label{eq:E_cumulative}
  E_{\mathrm{bg}} \sim \dot{E}_{\mathrm{bg,\,tot}} \cdot t_0
  \sim 10^{-10}\;\mathrm{J},
\end{equation}
which is negligible compared to the total energy budget.
This confirms that the deposition is a slow, secular
process---consistent with the near-constancy of
fundamental parameters over cosmic history.

\textit{Physical picture.}---Every particle, every moment,
broadcasts a faint resonant ``echo'' into the medium.
These echoes superpose to form a stochastic background
vibration that fills all of space---including regions devoid
of matter (Sec.~\ref{sec:micro_bernoulli}).  This background
vibration is energy distributed in the medium that does not
constitute localized particles (it has no resonant structure)
and does not produce a localized Newtonian attraction: a
spatially uniform $\langle(\nabla\varphi)^{2}\rangle$ has
vanishing gradient of the average, so no local
gravitational pull on a test mass arises from it.  It is,
however, energetically real, and the corresponding uniform
energy density does enter the cosmological Friedmann equation
on horizon scales; this cosmological role is precisely what
is exploited in the cosmological-constant analysis of
Sec.~\ref{sec:cc_problem}.

%-------------------------------------------------------------------
\subsection{The negative-feedback loop}
\label{sec:feedback}
%-------------------------------------------------------------------

The resonant-tail mechanism might seem to predict a
runaway process: more accumulated tail energy $\to$ lower pressure
$\to$ more tail emission.  In fact, the opposite occurs.  The
system is self-regulating through a negative-feedback loop:

\begin{enumerate}
  \item Particles emit resonant tails, increasing the
    background vibration level $\langle\delta\varphi^2\rangle$.
  \item The increased background vibration, through the
    Bernoulli mechanism, lowers the global static pressure
    $P_{\mathrm{static}}$.
  \item Lower pressure reduces the effective mass of each
    particle: $\meff = m_0\,e^{\varphi/2}$, where $\varphi$
    becomes more negative.
  \item A lighter particle has lower resonance amplitude
    and emits \emph{less} energy per unit time:
    $\mathcal{P}_{\mathrm{tail}} \propto m^4$.
  \item The outgoing tail power, and therefore any associated
    coarse-grained entropy-production proxy, decreases.
\end{enumerate}

This can be modeled as a first-order dynamical system.
Let $\varepsilon(t)$ denote the mean background vibration
energy density.  The feedback loop gives
\begin{equation}\label{eq:feedback_ode}
  \dot{\varepsilon}
  = \lambda\,n\,m_0^{4}\,e^{2\varphi(\varepsilon)},
\end{equation}
where $\lambda = Gc^{3}/\hbar^{2}$ is the dimensionally
consistent coupling constant inherited from the boxed
form~(\ref{eq:P_tail}) (subject to the same non-perturbative
factor $\mathcal{F}_{\mathrm{nonpert}}$ noted there),
$n$ is the particle number density, and
$\varphi(\varepsilon) < 0$ is the global scalar-field shift
induced by the accumulated background energy.  At leading
order the Bernoulli relation between accumulated background
energy density and the global pressure shift gives
\begin{equation}\label{eq:phi_of_eps}
  \varphi(\varepsilon) \simeq -\,\varepsilon/\varepsilon_{*},
  \qquad
  \varepsilon_{*} \equiv
  \frac{c^{2}\,e^{-\varphi_{0}}}{16\pi G},
\end{equation}
to leading order around the local equilibrium.  Higher-order
corrections to this functional form modify the late-time
$\varepsilon(t)$ asymptotics quantitatively but not
qualitatively: in every case the right-hand side of
Eq.~(\ref{eq:feedback_ode}) is monotonically decreasing in
$\varepsilon$ and approaches zero asymptotically.  Since
$\varphi(\varepsilon)$ becomes more negative as $\varepsilon$
grows (more background vibration $\to$ lower pressure),
the exponential factor $e^{2\varphi}$ monotonically decreases,
providing the negative feedback.

The fixed point $\dot{\varepsilon} = 0$ is approached only
asymptotically ($\varepsilon \to \infty$, $\varphi \to -\infty$,
$e^{2\varphi} \to 0$): the process never truly stops, but
decelerates without bound.  The solution behaves as
\begin{equation}\label{eq:epsilon_late}
  \varepsilon(t) \sim \varepsilon_0 + A\,\ln(t/t_0)
  \quad (t \gg t_0),
\end{equation}
a logarithmic growth---cosmologically slow.

%-------------------------------------------------------------------
\subsection{Perturbative feedback estimate for the dark-energy scale}
\label{sec:cc_problem}
%-------------------------------------------------------------------

The cosmological constant problem is the $10^{120}$ discrepancy
between the QFT prediction of vacuum energy density
($\rho_{\mathrm{QFT}} \sim M_{\mathrm{Pl}}^4$) and the
observed value
($\rho_{\Lambda} \sim 10^{-120}\,M_{\mathrm{Pl}}^4$)~%
\cite{Weinberg1989,Martin2012}.

In QFT, the vacuum energy is computed by summing zero-point
energies of all field modes up to some cutoff $\Lambda$:
\begin{equation}\label{eq:rho_qft}
  \rho_{\mathrm{QFT}}
  = \int_0^{\Lambda} \frac{4\pi k^{2}}{(2\pi)^{3}}\,
    \frac{\hbar\omega_k}{2}\,dk
  \propto \Lambda^{4}.
\end{equation}
This integral has no built-in regulator; it grows as the
fourth power of the cutoff.

In the ISPG sector considered here, the gravitating dark-energy
component is the \emph{dynamical equilibrium} value of the
accumulated echo-background energy $\varepsilon(t)$---the stochastic
vibration deposited by resonant tails (Sec.~\ref{sec:entropy_mech}),
distinct from the primordial $\nu_0$-mode reservoir that depletes as
matter forms.  This statement does not, by itself, derive a
cancellation or decoupling of the conventional QFT zero-point
contribution in Eq.~(\ref{eq:rho_qft}); that decoupling must be
treated as a renormalization assumption or as part of a future
completion.  Within the echo-background sector, the negative-feedback
loop (Sec.~\ref{sec:feedback}) provides a regulator: the system cannot
accumulate arbitrarily large echo energy because doing so would
suppress the very sources (particle resonances) that generate it.

\textit{Cosmological feedback channel.}---On cosmological
scales, the local feedback of Sec.~\ref{sec:feedback}
is supplemented by a second, dominant channel: the
de~Sitter expansion dilutes the very sources of vacuum
energy.  The particle number density is
$n(t) = n_0/a^{3}(t)$, so the cosmological feedback
ODE becomes
\begin{equation}\label{eq:feedback_cosmo}
  \dot{\varepsilon}
  = \frac{\mathcal{P}_{\mathrm{tail},0}^{\mathrm{(pert)}}\,n_0}
         {a^{3}(t)},
  \qquad
  \mathcal{P}_{\mathrm{tail},0}^{\mathrm{(pert)}} \equiv
  \frac{G\,m_0^{4}\,c^{3}}{\hbar^{2}},
\end{equation}
where $\mathcal{P}_{\mathrm{tail},0}^{\mathrm{(pert)}}$
is the perturbative tail power per particle
(Eq.~\ref{eq:P_tail}).  It has natural-unit dimension $M^2$
and SI dimension W; the dimensionally consistent $c^{3}$ form is
used here, in agreement with Sec.~\ref{sec:entropy_mech}.  As written,
$\dot{\varepsilon}$ is a derivative with respect to the cosmic/coordinate
time $t$, so no additional process-time factor $\mathcal{C}(z)$ is applied
in this source law.  The single-species
approximation is justified because nucleons dominate
the sum
$\sum_i \mathcal{P}_{\mathrm{tail},0,i}^{\mathrm{(pert)}}\,n_{0,i}$
over particle
species: the $m^{4}$ scaling of
$\mathcal{P}_{\mathrm{tail},0,i}^{\mathrm{(pert)}}$
suppresses lighter-species contributions by
$\gtrsim 10^{14}$ for electrons and by $\sim 10^{40}$
for neutrinos at comparable number densities.

When $\varepsilon$ dominates the energy budget, the Friedmann
equation gives $H = \sqrt{8\pi G\varepsilon/(3c^{2})}$
and $a(t) \propto e^{Ht}$.  The sources are then
exponentially diluted:
\begin{equation}
  \dot{\varepsilon}
  = \mathcal{P}_{\mathrm{tail},0}^{\mathrm{(pert)}}\,n_0\,e^{-3Ht}.
\end{equation}
Integrating and taking $t \gg H^{-1}$:
\begin{equation}\label{eq:epsilon_equil}
  \varepsilon \;\longrightarrow\;
  \varepsilon_{\mathrm{eq}}
  = \frac{\mathcal{P}_{\mathrm{tail},0}^{\mathrm{(pert)}}\,n_0}{3H}.
\end{equation}

Self-consistency ($H$ depends on $\varepsilon$) yields a
closed equation for the equilibrium density:
\begin{equation}\label{eq:rho_vac_eq}
  \boxed{
  \rho_{\mathrm{vac}}
  = \varepsilon_{\mathrm{eq}}
  = \left(\frac{
    \mathcal{P}_{\mathrm{tail},0}^{\mathrm{(pert)}}\,n_0}{3}\,
    \sqrt{\frac{3c^{2}}{8\pi G}}\right)^{2/3}.
  }
\end{equation}

\textit{Numerical estimate.}---Taking $m_0$ to be a representative
nucleon/proton mass scale and $n_0$ to be the present baryon number
density (for example $n_b\sim 0.25\,\mathrm{m}^{-3}$), in Planck units
($G = \hbar = c = 1$),
$\mathcal{P}_{\mathrm{tail},0}^{\mathrm{(pert)}} = m_0^{4}$,
$n_0 = n_0^{\mathrm{phys}}\,\ell_{\mathrm{Pl}}^{3}
\sim 10^{-105}$,
$m_0/\Mpl \sim 10^{-19}$, so
$\mathcal{P}_{\mathrm{tail},0}^{\mathrm{(pert)}} n_0
\sim 10^{-76}\times 10^{-105} = 10^{-181}$.
The $2/3$ power gives
\begin{equation}
  \rho_{\mathrm{vac}} \sim (10^{-181})^{2/3}
  \sim 10^{-121}\,\Mpl^{4}
  \sim H_0^{2}\,\Mpl^{2},
\end{equation}
within an order of magnitude of the observed
$\rho_{\Lambda} \sim 10^{-120}\,\Mpl^{4}$.
Equivalently, a direct SI estimate with the same proton-scale
perturbative power and baryon density gives the observed scale only
at the broad order-of-magnitude level, not as a precision fit.

The key point within this sector is that $\rho_{\mathrm{vac}}$ is set
by the self-consistent balance between tail production and de~Sitter
dilution rather than by inserting a UV cutoff into the
echo-background energy.  This sector statement should not be read as
a derivation that ordinary zero-point contributions fail to gravitate.
The suppression
$(m_0/\Mpl)^{8/3}\,(n_0\,\ell_{\mathrm{Pl}}^{3})^{2/3}
\sim 10^{-121}$ arises from two small dimensionless
ratios: the particle mass relative to the Planck mass,
and the cosmological particle density measured in
Planck volumes.

\textit{Status of the no-fine-tuning claim.}---The estimate above is
a perturbative-branch scale estimate, not a completed solution of the
full cosmological-constant problem.  It uses the perturbative scaling
\[
\mathcal{P}_{\mathrm{tail},0}^{\mathrm{(pert)}}
= G m_0^{4} c^{3}/\hbar^{2}
\]
at face value.  Explicitly: the cosmological estimate in this
subsection and the per-particle budget of
Sec.~\ref{sec:entropy_mech} use numerically different values of the
tail power.  Section~\ref{sec:entropy_mech} adopts the
non-perturbative value
$\mathcal{P}_{\mathrm{tail}}^{\mathrm{(np)}} \sim
10^{-107}$~W as a schematic suppressed per-particle benchmark
(at $\varphi_0 \approx 0$), whereas the Planck-unit combination
$\mathcal{P}_{\mathrm{tail},0}^{\mathrm{(pert)}} n_0
\sim 10^{-181}$
entering Eq.~(\ref{eq:rho_vac_eq}) corresponds to the perturbative
$\mathcal{P}_{\mathrm{tail}}^{(\mathrm{pert})}
\sim 10^{-24}$~W---a factor
$\mathcal{F}_{\mathrm{nonpert}}^{-1} \sim 10^{83}$ apart.  If
\[
\mathcal{P}_{\mathrm{tail}}^{\mathrm{(np)}} =
\mathcal{F}_{\mathrm{nonpert}}\,
\mathcal{P}_{\mathrm{tail}}^{\mathrm{(pert)}} ,
\]
then Eq.~(\ref{eq:rho_vac_eq}) rescales as
\[
\rho_{\mathrm{vac}}\rightarrow
\mathcal{F}_{\mathrm{nonpert}}^{2/3}\rho_{\mathrm{vac}}^{\mathrm{(pert)}} .
\]
The structural mechanism---feedback-regulated balance between source
and de~Sitter dilution---does \emph{not} depend on the precise value
of $\mathcal{F}_{\mathrm{nonpert}}$; what does is the quantitative
coincidence with the observed $10^{-120}\,\Mpl^{4}$.  In addition,
the treatment of conventional zero-point energy requires a
renormalization/decoupling assumption not derived in this paper.  A
controlled non-perturbative calculation that simultaneously satisfies
the per-particle energy budget of Sec.~\ref{sec:entropy_mech}, the
perturbative dark-energy scale estimate above, and the zero-point
decoupling condition is an open task.

%-------------------------------------------------------------------
\subsection{Dark-energy equation of state: \texorpdfstring{$w(z) \neq -1$}{w(z) != -1}}
\label{sec:dark_energy}
%-------------------------------------------------------------------

A true cosmological constant has equation-of-state parameter
$w \equiv P/\rho = -1$ exactly, for all time.  In ISPG,
the observable dark-energy density is a dynamical
quantity---the accumulated resonant-tail echo
background $\varepsilon(t)$---that evolves according to
Eq.~(\ref{eq:feedback_ode}).  Its effective equation of
state is therefore time-dependent.

To derive $w(z)$, note that the echo-background sector is not
separately conserved: it receives a continuous source
$Q =
\mathcal{P}_{\mathrm{tail},0}^{\mathrm{(pert)}}n_0/a^3 > 0$
from resonant-tail power
while having intrinsic pressure $p = -\varepsilon$
(the quasi-static phantom condensate satisfies
$\rho + p \approx 0$).  An observer who models the
dark-energy sector as a conserved effective fluid
applies the standard continuity equation
$\dot{\varepsilon} + 3H(1+w)\varepsilon = 0$, obtaining
\begin{equation}\label{eq:w_eff}
  w_{\mathrm{eff}}(z)
  = -1 - \frac{\dot{\varepsilon}(z)}
    {3H(z)\,\varepsilon(z)}\,.
\end{equation}
Since $\dot{\varepsilon} > 0$ at all epochs
(Eq.~\ref{eq:feedback_cosmo}), the echo-background density
\emph{grows}, which in conservation language means
$w < -1$: a growing energy density dilutes less than
$\Lambda$, placing it on the phantom side.  In the
present \emph{effective fluid} description this is not
a kinetic-ghost instability but a \emph{source-driven}
effect: matter continuously feeds the vacuum through
irreversible tail emission, and any observer unaware
of the source term would infer phantom behaviour.  The
underlying scalar action of ISPG separately carries
phantom kinetic sign ($\kappa = -1/2$ in
Eq.~(\ref{eq:action})); its all-orders quantum
stability is established in Sec.~\ref{sec:stability}
via cosmological redshift damping and the higher-derivative
cutoff.  The effective $w<-1$ inferred above is therefore
compatible with the kinetic-stability analysis of the
underlying field, not in addition to it.

The correction $\delta w(z) \equiv |w + 1|
= \dot{\varepsilon}/(3H\varepsilon) > 0$ is small.
At early times ($z \gg 1$), when particles
were more massive ($\meff$ larger due to higher pressure)
and more numerous (pre-recombination), $|\delta w|$ was larger.
At late times ($z \to 0$), the feedback suppression makes
$\delta w \to 0^{+}$: the dark energy asymptotically
approaches a cosmological constant from the phantom side.

The predicted behavior is:
\begin{equation}\label{eq:w_prediction}
  \boxed{
  w(z) = -1 - \delta w_0\,\frac{(1+z)^{3}}{E(z)},
  \quad E(z) \equiv \frac{H(z)}{H_0},
  \quad \delta w_0 \sim 10^{-2}\text{--}10^{-1},
  }
\end{equation}
with $w < -1$ at all epochs (source-driven phantom,
no ghost instability).  The sign follows from the positive
echo-background source term; the magnitude of $\delta w_0$
inherits the tail-power normalization discussed above and
therefore remains conditional until the non-perturbative
normalization is closed.

\textit{Observational status.}---In 2024, the Dark Energy
Spectroscopic Instrument (DESI) reported first-year
baryon acoustic oscillation data preferring time-varying
dark energy over a cosmological constant at the
$\sim 2.6\sigma$ level~\cite{DESI2024}.  The subsequent
Data Release~2~\cite{DESIDR2}, based on $1.4\times 10^7$ tracers,
strengthened the preference to $3.1\sigma$ (BAO$+$CMB),
rising to $4.2\sigma$ with DES-SN5YR supernovae.
The CPL fit
yields $w_0 = -0.45^{+0.34}_{-0.21}$ and
$w_a = -1.79^{+0.48}_{-1.00}$, while the constant-$w$
fit gives $w = -0.99^{+0.15}_{-0.13}$.

The ISPG prediction is $w_{0}^{\mathrm{ISPG}}
= -1 - \delta w_0$ with $\delta w_0 \sim 10^{-2}$--$10^{-1}$
(Eq.~\ref{eq:w_prediction}); for definiteness we use
the central value $\delta w_0 \simeq 0.06$, giving
$w_0 \simeq -1.06$ and $w_a \simeq -0.15$.  This value
lies on the \emph{phantom} side of $-1$ and within the
DESI DR2 \emph{constant-$w$} interval at $1\sigma$.  Under
the CPL parametrisation, however, the DESI central value
$w_0 \simeq -0.45$ sits on the \emph{quintessence} side
at $>4\sigma$ from $-1$; the ISPG prediction is therefore
in tension with the DR2 CPL fit by virtue of an opposite
sign of $w_0 + 1$, even though the sign of $w_a$ is
reproduced (with a magnitude smaller by roughly an
order of magnitude).  Neither parametrisation is yet
definitive: the constant-$w$ and CPL fits make different
assumptions about the time-evolution of $w$, and current
data do not discriminate between an ISPG-type
phantom-side $w$ that approaches $-1$ at $z=0$ and a
CPL-type quintessence-side $w$ with strongly negative
$w_a$.  The Euclid space telescope and the Vera Rubin
Observatory, with projected $\sigma(w_0) \sim 0.02$,
are expected to discriminate between these scenarios
within the current decade.

\paragraph{Cumulative process-time and the process age of the universe.}
The clock-rate postulate $d\tau = e^{\varphi_0/2}\,dt$
(Sec.~\ref{sec:excitations}, Eq.~\ref{eq:c_local}) implies that the total process
duration accumulated by pressure-controlled matter-sector processes since redshift
$z_f$ exceeds the coordinate lookback time whenever past pressure
was higher than present pressure:
\begin{equation}\label{eq:tau_proc_q}
  \tau_{\mathrm{proc}}(z_f)
  = \int_0^{z_f}
    \left[\frac{P_{\mathrm{stat}}(z)}{P_{\mathrm{stat}}(0)}\right]^{1/2}
    \frac{dz}{(1+z)\,H(z)}\,.
\end{equation}
The pressure ratio is fully determined by the feedback ODE
\eqref{eq:feedback_ode} and the dark-energy parametrisation
\eqref{eq:w_prediction}: $P_{\mathrm{stat}}(z)/P_{\mathrm{stat}}(0)
= [e^{2\varphi(z)}/e^{2\varphi(0)}]^{1/2}$, with the ratio supplied
by the self-consistent cosmological solution.  Here $\alpha$
parametrises the cosmological evolution of the per-particle
tail power
$\mathcal{P}_{\mathrm{tail}}(z)
= \mathcal{P}_{\mathrm{tail},0}(1+z)^{\alpha}$ entering
Eq.~(\ref{eq:feedback_cosmo}), in the convention used by the
companion MOND paper~\cite{Lotishvili2026MOND} (their Eq.~70);
$\alpha = 1$ corresponds to the matter-dominated tracking
of the proper-density factor $e^{-\varphi/2}$ at leading order.
For the fiducial parameter set
($z_f=10$, $\alpha=1$, and a high-amplitude closure benchmark
$\delta w_0\simeq 0.24$ used here as a parameter-study input
rather than a central dark-energy prediction
$\delta w_0\simeq 0.06$), the process-time age is
$\tau_{\mathrm{proc}}(10)\simeq 21\;\mathrm{Gyr}$, roughly
$1.6\times$ the coordinate lookback of $\simeq 13.3\;
\mathrm{Gyr}$.  Substituting the central
$\delta w_0 = 0.06$ instead reduces the enhancement to
$\sim 1.1$--$1.2\times$.  The companion MOND
paper~\cite{Lotishvili2026MOND} exploits the upper-end
enhancement to resolve the vortex-formation timeline;
the present derivation shows that the enhancement is a
direct consequence of the feedback mechanism rather than
a separate assumption, and that its quantitative size is
controlled by the chosen $\delta w_0$ benchmark; the
central quantum echo-background range gives the smaller
enhancement quoted above.

%-------------------------------------------------------------------
\subsection{Cosmological fine-structure constant variation}
\label{sec:alpha_var}
%-------------------------------------------------------------------

At the classical level, the bi-conformal symmetry ensures
that the fine-structure constant $\alpha = e^2/(4\pi\epsilon_0
\hbar c)$ is exactly invariant: all factors of $e^{\varphi/2}$
cancel in the dimensionless ratio~\cite{Lotishvili2026}.

At the quantum level, virtual particle loops contribute
corrections to electromagnetic vertices.  Since virtual
particles have effective masses $\meff = m_0\,e^{\varphi/2}$
that depend on the background pressure, these loop corrections
introduce a residual $\varphi$-dependence in the
\emph{effective} fine-structure constant:
\begin{equation}\label{eq:alpha_eff}
  \alpha_{\mathrm{eff}}(\varphi)
  = \alpha_0\left[1
    + \frac{2\alpha_0}{3\pi}\,\varphi
    + \order{\alpha_0^{2}\varphi^{2}}\right],
\end{equation}
where $\alpha_0 = 1/137.036$ is the present-day value
and the leading correction comes from the electron
vacuum-polarization loop (Uehling term) with a
$\varphi$-dependent electron mass.

For the cosmological variation over a lookback time
corresponding to redshift $z$, the scalar field shift
$\Delta\varphi \sim \order{H_0 t_{\mathrm{lookback}}}$
yields
\begin{equation}\label{eq:delta_alpha}
  \boxed{
  \frac{\Delta\alpha}{\alpha}
  \sim \frac{2\alpha_0}{3\pi}\,\Delta\varphi
  \sim 10^{-7}\text{--}10^{-6}
  \quad (z \sim 1\text{--}3).
  }
\end{equation}
Equation~\eqref{eq:delta_alpha} should be read as a
schematic loop estimate until the renormalization prescription,
comparison scale, and sign are fixed.  This caveat is essential
because local dimensionless observables cancel under a uniform
bi-conformal scaling; cosmological comparisons of
$\alpha_{\mathrm{eff}}(\varphi)$ must specify what reference
scale is held fixed across epochs.  Current low-systematics
quasar constraints are already at sub-ppm precision, so a
generic $10^{-7}$--$10^{-6}$ signal is close to existing
sensitivity rather than an observed effect.  The sharp test is
therefore a sign and redshift profile for ESPRESSO and
ELT-ANDES, not only an order-of-magnitude estimate.

%-------------------------------------------------------------------
\subsection{Cosmological expansion as resonant-tail propagation}
\label{sec:tail_expansion}
%-------------------------------------------------------------------

The results of
Secs.~\ref{sec:entropy_mech}--\ref{sec:alpha_var} imply a
natural interpretation of cosmological expansion within the
ISPG framework.  The ingredients are already in place:
\begin{enumerate}
  \item Every particle continuously emits resonant tails
    that propagate outward at the local speed of light
    (Sec.~\ref{sec:entropy_mech}).
  \item The tails are irreversible and fill all of space,
    including regions devoid of localized matter
    (Sec.~\ref{sec:micro_bernoulli}).
  \item The medium is infinite: a medium has no boundary,
    and the tails propagate indefinitely.
  \item In a region where no resonance has ever arrived,
    the scalar field is homogeneous
    ($\nabla\varphi = 0$) and no localized excitation
    exists.  Without a resonant oscillator---an
    atom, a particle, a clock---proper time has no
    operational definition.
\end{enumerate}

Taken together, these statements yield a picture of
cosmological expansion that does not invoke an inflaton
field or an initial singularity:

\textit{The observable universe is the region that resonant
tails have reached since the onset of the first excitations.}
Beyond the outermost tail front, the medium exists in its
quiescent state: homogeneous, devoid of resonances, and
operationally timeless.  The tail front advances at $c$
into the quiescent medium; where it arrives, the scalar
field acquires a nonzero gradient, the first pressure
deficit forms, and operational proper time becomes
definable through localized resonant oscillators.

This provides a physical mechanism for the cosmological
horizon: we observe precisely the region within which
resonant information has had time to propagate.  The
horizon grows monotonically as tails continue to spread,
and the process is irreversible (Sec.~\ref{sec:entropy_mech}).

The quiescent state is defined not by a particular value
of $\varphi$ (which is gauge-dependent under the global
shift $\varphi \to \varphi + C$) but by the
gauge-invariant conditions $\nabla\varphi = 0$ and
the absence of localized resonant excitations.

We now show that this picture is quantitatively
consistent with the standard FLRW expansion history.

\textit{Friedmann equations.}---The ISPG action
$S = (16\pi G)^{-1}\!\int d^4x\sqrt{-g}\,[R +
\frac{1}{2}\,g^{\mu\nu}\partial_\mu\varphi\,
\partial_\nu\varphi] + S_{\mathrm{m}}$
is Einstein gravity minimally coupled to a scalar
field.  On a homogeneous, isotropic background the
scalar field decomposes as $\varphi(t) = \bar\varphi(t)
+ \delta\varphi_{\mathrm{bg}}(t)$, where $\bar\varphi$
is the piece sourced by the mean matter density (mediating
gravity, exactly as in GR at 1PN~\cite{Lotishvili2026})
and $\delta\varphi_{\mathrm{bg}}$ is the accumulated
background vibration from resonant tails
(Sec.~\ref{sec:entropy_mech}).  The Friedmann equation
retains its standard form:
\begin{equation}\label{eq:friedmann_ispg}
  H^{2}(z) = H_0^{2}\!\left[\Omega_m(1+z)^{3}
    + \Omega_r(1+z)^{4}
    + \Omega_{\mathrm{DE}}(z)\right],
\end{equation}
where the dark-energy density $\Omega_{\mathrm{DE}}(z)$
is set by the dynamical echo-background energy
$\varepsilon(t)$ of Eq.~(\ref{eq:epsilon_late}),
rather than by a cosmological constant.

\textit{Particle horizon = tail front.}---In FLRW
cosmology, the comoving particle horizon at cosmic
time $t$ is
\begin{equation}\label{eq:particle_horizon}
  d_H(t) = a(t)\int_0^{t}\frac{c\,dt'}{a(t')}.
\end{equation}
In ISPG, the locally measured speed of light is
exactly $c$ everywhere (Eq.~\ref{eq:c_local}).
The outermost resonant tail, emitted at the onset
of the first excitations and propagating at $c$,
covers precisely the proper distance
$d_H(t)$.  Therefore,
\begin{equation}\label{eq:tail_horizon}
  \boxed{
  d_{\mathrm{tail\;front}}(t) = d_H(t).
  }
\end{equation}
The tail-propagation boundary and the standard
particle horizon are not merely analogous;
they are \emph{mathematically identical}.  The
observable universe is exactly the causal domain
of resonant-tail propagation.

\textit{Modified dark-energy density.}---From
Eq.~(\ref{eq:w_prediction}) and the continuity equation,
the dark-energy density evolves as
\begin{equation}\label{eq:rho_DE_z}
  \frac{\rho_{\mathrm{DE}}(z)}{\rho_{\mathrm{DE}}(0)}
  = \exp\!\left[
    -3\,\delta w_0 \int_0^z
    \frac{(1+z')^{2}}{E(z')}\,dz'
  \right].
\end{equation}
Since the exponent is negative, $\rho_{\mathrm{DE}}$ was
\emph{smaller} in the past---the hallmark of phantom
behaviour.  At low redshift:
\begin{equation}
  \frac{\rho_{\mathrm{DE}}(z)}{\rho_{\mathrm{DE}}(0)}
  \approx 1 - 3\,\delta w_0\,z + \order{z^2}
  \approx 1 - 0.18\,z,
\end{equation}
a percent-level deviation per unit redshift from
$\Lambda$CDM---precisely the sensitivity range of
DESI (year-3 and beyond) and Euclid.

\textit{Hubble law.}---At low redshift ($z \ll 1$),
the Taylor expansion $v = cz \approx H_0 d$ holds
for any dark-energy model.  The ISPG expansion
history therefore reproduces the Hubble law exactly.

\textit{CMB compatibility.}---The angular diameter
distance to last scattering,
\begin{equation}\label{eq:d_A_CMB}
  d_A = \int_0^{z_*}\frac{c\,dz}{H(z)},
\end{equation}
is sensitive to $\Omega_{\mathrm{DE}}(z)$ primarily
at $z \lesssim 2$, where dark energy dominates.
At $z_* \approx 1090$, the dark-energy fraction is
negligible ($\Omega_{\mathrm{DE}}(z_*) \ll
\Omega_m(z_*)$), so the acoustic-peak structure
is preserved.  The residual modification of $d_A$
at low $z$ shifts the inferred $H_0$ by
$\delta H_0/H_0 \sim \delta w_0 \sim \text{few}\;\%$---within
current Planck--SH0ES tension bounds and resolvable
by next-generation CMB experiments (CMB-S4, LiteBIRD).

\textit{Summary.}---The tail-propagation picture
reproduces the standard expansion history to leading
order (Friedmann equations, Hubble law, particle
horizon, CMB peaks), with calculable percent-level
deviations encoded in $w(z) \neq -1$.  We now address
the dynamics of the tail front at the earliest epochs.

%-------------------------------------------------------------------
\subsection{Tail-front dynamics at the earliest epochs}
\label{sec:early_epoch}
%-------------------------------------------------------------------

The preceding subsection identifies the cosmological
horizon with the outermost resonant tail front but
leaves open the dynamics at $t \lesssim t_{\mathrm{Pl}}$.
We now show that the ISPG framework provides a
self-consistent picture of the onset---one that avoids
the initial singularity and addresses the standard
cosmological puzzles without invoking an inflaton field.

\textit{The quiescent state.}---Before any resonance
forms, the medium is in its quiescent configuration:
$\varphi = \mathrm{const}$ everywhere, with
$\nabla\varphi = 0$ and
$P_{\mathrm{static}} = 0$, $\nabla P = 0$.  There are
no localized excitations, no pressure deficits, and---since
proper time is defined by the oscillation of a resonant
mode (Sec.~\ref{sec:excitations})---no operational clock.
The quiescent state is not a moment ``before the Big
Bang'' but a timeless ground configuration of the
infinite medium.

\textit{Quantum nucleation of the first excitations.}---The
Wheeler--DeWitt equation
(Sec.~\ref{sec:canonical_quantization}) describes the
full quantum state of the scalar-metric system,
including configurations where no classical resonance
exists.  The wave functional $\Psi[\varphi, h_{ij}]$
assigns a finite amplitude to configurations with
$\nabla\varphi \neq 0$---that is, to states containing
localized excitations.  The transition from the quiescent
state to a configuration containing one oscillon is a
quantum tunneling event, analogous to Schwinger pair
production in QED or Coleman--De~Luccia bubble
nucleation.

The tunneling amplitude can be estimated semiclassically.
The Euclidean instanton describing oscillon nucleation
is a localized, $O(4)$-symmetric field configuration
$\varphi_{\mathrm{inst}}(\rho)$ (with
$\rho = \sqrt{\tau^2 + |\mathbf{x}|^2}$) that
interpolates between the quiescent state
$\varphi = \mathrm{const}$ at $\rho \to \infty$
and the oscillon core profile at $\rho = 0$.  The
Euclidean action of this instanton is
\begin{equation}\label{eq:S_nucleation}
  S_E
  = \int d^4x_E\,\sqrt{g_E}\left[
    \frac{R_E}{16\pi G}
    + \frac{1}{2}(\nabla_E\varphi)^2
  \right]_{\mathrm{inst}}.
\end{equation}
Dimensional analysis fixes
$S_E \sim (R_c/\ell_{\mathrm{Pl}})^2$, where
$R_c = \hbar/(m_0 c)$ is the oscillon radius and
$\ell_{\mathrm{Pl}} = \sqrt{\hbar G/c^3}$ the Planck
length.  For Planck-mass oscillons
($m_0 \sim \Mpl$, $R_c \sim \ell_{\mathrm{Pl}}$),
$S_E \sim \order{1}$ and the nucleation rate per Planck
four-volume is unsuppressed:
$\Gamma_{\mathrm{nucl}} \sim
\ell_{\mathrm{Pl}}^{-4}\,e^{-S_E}
\sim \ell_{\mathrm{Pl}}^{-4}$.
For lighter oscillons ($m_0 \ll \Mpl$,
$R_c \gg \ell_{\mathrm{Pl}}$), $S_E \gg 1$ and the
rate is exponentially suppressed; the first excitations
are therefore generically Planck-scale.

Since the quiescent medium is spatially infinite,
nucleation occurs independently at every point---not
as a single event but as a spatially uniform process.
The resulting ensemble of Planck-mass oscillons is
statistically homogeneous and isotropic by construction:
every nucleation event involves the same instanton, in
the same quiescent background, with no preferred
direction or location.  This \emph{universal nucleation}
replaces the single-bubble picture and directly
explains the observed large-scale homogeneity of the
CMB.

\textit{Energy density at the onset.}---The first
oscillons form with Planck-scale energy densities:
\begin{equation}\label{eq:rho_onset}
  \rho_{\mathrm{onset}}
  \sim \frac{\Mpl\,c^2}{\ell_{\mathrm{Pl}}^3}
  \sim \frac{\Mpl^4}{c^2}
  \sim 10^{113}\;\mathrm{J/m^3}.
\end{equation}
This is the natural initial condition for the feedback
ODE~(\ref{eq:feedback_cosmo}): $\varepsilon(0) = 0$
in a region of Planck-scale extent, with the first
resonant tails carrying Planck-scale energy.  The
negative-feedback loop (Sec.~\ref{sec:feedback}) together
with the de~Sitter dilution of sources
(Eq.~\ref{eq:feedback_cosmo}) drives, in the
perturbative-branch estimate, the echo-background energy to
the self-consistent equilibrium
$\rho_{\mathrm{vac}} \sim H_0^2 \Mpl^2$ at the
present epoch (Eq.~\ref{eq:rho_vac_eq}).

\textit{Coarse-grained onset closure.}---The onset can be
closed at the homogeneous-background level by a simple
kinetic description.  Let $n_{\mathrm{osc}}(t)$ denote the
number density of newly nucleated Planck-mass oscillons and
$\rho_r(t)$ the radiation bath generated when their
overlapping tails exchange energy.  To leading order,
\begin{equation}\label{eq:onset_rate_eqs}
  \dot n_{\mathrm{osc}} + 3H n_{\mathrm{osc}}
  = \mathcal{J}_{\mathrm{nucl}}
    - \Gamma_{\mathrm{th}}\,n_{\mathrm{osc}},
  \qquad
  \dot\rho_r + 4H\rho_r
  = \Mpl c^2\,\Gamma_{\mathrm{th}}\,n_{\mathrm{osc}},
\end{equation}
where the nucleation source is
\begin{equation}\label{eq:onset_nucleation_source}
  \mathcal{J}_{\mathrm{nucl}}
  \sim \Gamma_{\mathrm{nucl}}
  \sim \ell_{\mathrm{Pl}}^{-4}e^{-S_E}
  \sim \ell_{\mathrm{Pl}}^{-4}
\end{equation}
for the unsuppressed Planck-scale instantons with
$S_E = \order{1}$.  The overlap interaction of neighboring
Planckian cores is geometric, so the scattering cross section
is
\begin{equation}\label{eq:onset_cross_section}
  \sigma_{\mathrm{sc}} = \order{\ell_{\mathrm{Pl}}^{2}}.
\end{equation}
Once the occupancy reaches its natural onset value
$n_{\mathrm{osc}} \sim \ell_{\mathrm{Pl}}^{-3}$, the
thermalization rate is
\begin{equation}\label{eq:onset_gamma_th}
  \Gamma_{\mathrm{th}}
  \sim n_{\mathrm{osc}}\sigma_{\mathrm{sc}}\,c
  \sim \ell_{\mathrm{Pl}}^{-3}\,\ell_{\mathrm{Pl}}^2\,c
  = \frac{c}{\ell_{\mathrm{Pl}}}
  = t_{\mathrm{Pl}}^{-1}.
\end{equation}
Thus both the onset cascade and the conversion of localized
oscillon energy into a thermal radiation bath occur within an
order-one number of Planck times.  At that stage,
\begin{equation}\label{eq:onset_radiation_density}
  \rho_r
  \sim \Mpl c^2\,n_{\mathrm{osc}}
  \sim \frac{\Mpl c^2}{\ell_{\mathrm{Pl}}^3}
  \sim \rho_{\mathrm{onset}},
\end{equation}
so the hot Big-Bang background follows directly from the
nucleation kinetics rather than from an additional reheating
postulate.  From this point onward, the standard Friedmann
evolution (Eq.~\ref{eq:friedmann_ispg}) takes over, with the
radiation density scaling as $(1+z)^4$ and the dark-energy
density governed by the feedback mechanism.

\textit{Coarse-grained homogeneity and flatness.}---Because
the nucleation events are statistically identical and occur
independently in Planck-size cells, the background smooths by
the central-limit law.  On a coarse-graining scale
$L \gg \ell_{\mathrm{Pl}}$, the number of independent onset
cells is
\begin{equation}\label{eq:onset_cell_count}
  N_L \sim n_{\mathrm{osc}}\,L^3
  \sim \left(\frac{L}{\ell_{\mathrm{Pl}}}\right)^3.
\end{equation}
Hence the fractional density fluctuation of the homogeneous
background on that scale is
\begin{equation}\label{eq:onset_density_fluctuation}
  \frac{\delta\rho_L}{\rho_{\mathrm{onset}}}
  \sim N_L^{-1/2}
  \sim \left(\frac{\ell_{\mathrm{Pl}}}{L}\right)^{3/2}.
\end{equation}
The corresponding curvature contribution obeys the same
parametric suppression,
\begin{equation}\label{eq:onset_curvature_suppression}
  \left|\frac{k_{\mathrm{eff}}}{a^2H^2}\right|_L
  = \order{\frac{\delta\rho_L}{\rho_{\mathrm{onset}}}}
  = \order{\left(\frac{\ell_{\mathrm{Pl}}}{L}\right)^{3/2}},
\end{equation}
so for every cosmological smoothing scale
$L \gg \ell_{\mathrm{Pl}}$ the onset background is
homogeneous, isotropic, and spatially flat to extremely high
accuracy.  This is the precise background-level closure used
below: the homogeneous FLRW initial data are analytically
controlled, while the detailed two-point perturbation spectrum
is a separate problem.

\textit{Singularity avoidance.}---In standard GR, the
Penrose--Hawking theorems imply a curvature singularity
at the ``beginning.''  In ISPG, these theorems do not
apply, for two independent reasons:
\begin{enumerate}
\item The phantom kinetic sign produces NEC violation in
  the effective gravitational source (the canonical scalar
  stress-energy tensor itself satisfies the standard energy
  conditions; the negative coupling $\kappa=-1/2$ flips the
  sign of the source-side NEC contraction; cf.\ the
  energy-conditions section of the main theory~\cite{Lotishvili2026}),
  invalidating the focusing theorem
  that underpins the singularity
  theorems~\cite{Lotishvili2026}.  This is the same
  mechanism that resolves the Schwarzschild singularity
  in the classical theory.
\item The Wheeler--DeWitt equation
  (Sec.~\ref{sec:canonical_quantization}) is
  \emph{hyperbolic} in ISPG (due to the phantom sign),
  with the scalar field $\varphi$ playing the role of
  an internal time variable.  The wave functional
  propagates smoothly through the configuration
  $a \to 0$ (zero scale factor) without encountering
  a singular boundary.
\end{enumerate}
The ``Big Bang'' in ISPG is therefore not a singularity
but a smooth quantum transition from the quiescent state
to the excited (expanding) state.

\textit{Standard cosmological puzzles.}---The
tail-front onset addresses the puzzles traditionally
solved by inflation:
\begin{enumerate}
\item \textit{Horizon problem.}---The uniformity of the
  CMB follows from the universality of the nucleation
  mechanism: every point in the quiescent medium
  nucleates through the same instanton, in the same
  background, producing statistically identical initial
  conditions.  On a smoothing scale $L$, the residual
  inhomogeneity is already suppressed by
  Eq.~(\ref{eq:onset_density_fluctuation}).  The tail-front identity
  $d_{\mathrm{tail\;front}} = d_H$
  (Eq.~\ref{eq:tail_horizon}) then ensures that the
  observable universe is the causal domain within
  which these identical seeds have had time to
  thermalize.  Homogeneity at last scattering is
  therefore a consequence of the universality of
  quantum nucleation, not a fine-tuned coincidence.
\item \textit{Flatness problem.}---The quiescent state is
  exactly flat ($k = 0$): a spatially infinite,
  homogeneous medium admits no intrinsic curvature.
  Nucleation is spatially uniform (same instanton
  everywhere), so it cannot generate a preferred
  curvature scale.  Perturbations of the spatial
  geometry are localized within each oscillon's Compton
  radius $R_c \ll d_H$ and average to zero on
  cosmological scales; more precisely,
  Eq.~(\ref{eq:onset_curvature_suppression}) gives
  $|k_{\mathrm{eff}}|/(a^2H^2)
  = \order{(\ell_{\mathrm{Pl}}/L)^{3/2}} \ll 1$ on every
  macroscopic scale.  The resulting
  Friedmann evolution therefore inherits $k = 0$ as an
  \emph{exact} initial condition, not one that must be
  fine-tuned.
\item \textit{Monopole problem.}---ISPG contains a single
  scalar field with an Abelian shift symmetry and no
  spontaneous breaking of a non-Abelian gauge group.
  The theory therefore produces no topological defects
  (monopoles, domain walls, cosmic strings) at any
  epoch.
\item \textit{Primordial perturbations.}---The quantum
  fluctuations of the nucleation epoch provide the seeds
  for large-scale structure.  The spectrum is determined
  by the nonlinear dynamics of the oscillon--tail
  interaction at the Planck epoch; its detailed form
  requires numerical simulation of the coupled
  scalar-metric system.  We note, however, that the
  analytic closure above concerns the homogeneous
  background only.  The feedback mechanism still provides
  a natural tilt: modes
  that enter the horizon later experience a slightly
  different feedback suppression, generically producing
  a nearly scale-invariant but not exactly flat spectrum.
\end{enumerate}

\textit{Summary.}---The ISPG framework provides a
closed homogeneous-background picture of the earliest
epoch: the
quiescent medium undergoes spatially uniform quantum
nucleation of Planck-mass oscillons; the overlapping
tails thermalize to produce the hot Big Bang; the
singularity is avoided by the effective-source NEC violation
(canonical $\tau_{\mu\nu}$ satisfies the energy conditions; the
$\kappa=-1/2$ coupling flips the sign of the source-side
contraction, cf.\ the energy-conditions section of the main
theory~\cite{Lotishvili2026}) and the hyperbolic WDW equation; and the standard cosmological
puzzles are addressed by the universality of the
nucleation mechanism and the geometric inheritance of
$k = 0$.  The quantitative prediction of the primordial
perturbation spectrum requires numerical simulation of
the Planck-epoch dynamics---a task identified for future
lattice work.

%===================================================================
\section{Mass spectrum of three-dimensional resonances}
\label{sec:mass_spectrum}
%===================================================================

The identification of particles as resonant modes of the
scalar-metric medium (Sec.~\ref{sec:excitations}) raises a
quantitative question: \emph{why do the charged leptons---electron,
muon, and tau---have the specific mass ratios they do?}
This section separates three roles: the three-dimensional cavity
supplies structural mode labels, the companion MASS/Mathieu ladder
supplies the quantitative charged-lepton mass ratios, and the
structure-preserving reduction from cavity modes to the Mathieu
index $N$ remains an open calculation.  The cavity analysis then
connects this structure to the empirical Koide formula as a
consistency constraint.

%-------------------------------------------------------------------
\subsection{Three-dimensional resonance eigenvalue problem}
\label{sec:3d_resonance}
%-------------------------------------------------------------------

Unlike string theory, where particles are one-dimensional objects
and the resonance spectrum follows from the harmonics of a string,
ISPG treats particles as three-dimensional resonant structures
embedded in the scalar-metric medium.  The eigenvalue problem
is therefore that of a three-dimensional cavity, not a
one-dimensional string.

Consider a self-consistent resonant mode
(Sec.~\ref{sec:resonances}) with a core of effective radius
$R_{\mathrm{core}}$ and angular frequency $\omega$.  The
spatial profile satisfies the Helmholtz equation
(Eq.~\ref{eq:helmholtz}) which, for the fundamental mode
structure, reduces to
\begin{equation}\label{eq:3d_helmholtz}
  \nabla^{2}\psi + k^{2}\psi = 0,
  \qquad k^{2} = \omega^{2}\,e^{-2\varphi_0},
\end{equation}
inside the resonant core, with outgoing-wave (radiative)
boundary conditions at large $r$.

In spherical coordinates, the angular dependence is given
by spherical harmonics $Y_{lm}(\theta,\phi)$ and the radial
dependence by spherical Bessel functions:
\begin{equation}\label{eq:radial_bessel}
  R_{nl}(r) = j_{l}(k_{nl}\,r),
\end{equation}
where $k_{nl}$ are determined by the boundary conditions
at the effective cavity radius.  For a resonant mode
confined to radius $R$, the eigenfrequencies are
\begin{equation}\label{eq:eigenfreq}
  \omega_{nl} = \frac{\chi_{nl}}{R}\,e^{\varphi_0},
\end{equation}
where $\chi_{nl}$ is the $n$-th zero of the spherical
Bessel function $j_l$ (or a related transcendental equation
for the quasi-bound case).

The first few eigenvalues for $l = 0$ are:
\begin{equation}
  \chi_{10} = \pi, \quad
  \chi_{20} = 2\pi, \quad
  \chi_{30} = 3\pi.
\end{equation}
These would give mass ratios $1:2:3$, which clearly do not
match the lepton masses ($m_e : m_\mu : m_\tau \approx
1 : 207 : 3477$).  The resolution lies in the fact that
the three charged leptons are \emph{not} successive
$s$-wave harmonics of the same cavity.  They correspond
to modes with different angular quantum numbers $(n,l)$
in a three-dimensional resonator whose effective potential
is not a simple hard-wall cavity but the self-consistent
nonlinear profile of Sec.~\ref{sec:resonances}.

%-------------------------------------------------------------------
\subsection{Self-consistent mass determination}
\label{sec:mass_determination}
%-------------------------------------------------------------------

The linearized eigenvalue
problem~(\ref{eq:3d_helmholtz})--(\ref{eq:eigenfreq})
captures the angular structure of the modes but cannot
fix the absolute mass scale, because the confining
potential is created by the nonlinear self-interaction
and must be determined self-consistently.  We now cast
the oscillon equation~(\ref{eq:self_consistent}) as a
dimensionless eigenvalue problem and identify the
parameters that control the spectrum.

\textit{The nonlinear eigenvalue problem.}---The
self-consistent oscillon
equation~(\ref{eq:self_consistent}), derived from
$\Box\varphi = 0$ on the bi-conformal background,
constitutes a nonlinear eigenvalue problem for the
angular frequency $\omega$.  In spherical symmetry,
define $u(r) \equiv r\,\Phi(r)$ and write it as
\begin{equation}\label{eq:mode_eq}
  u'' + \left[\omega^2 e^{-2\varphi_0}
    - \frac{\ell(\ell+1)}{r^2}
    - V_{\mathrm{eff}}(r;\,\Phi)\right]u = 0,
\end{equation}
where the effective potential
\begin{equation}\label{eq:Veff_explicit}
  V_{\mathrm{eff}}(r;\,\Phi) =
    -\frac{\varphi_0'}{r}
    - \tfrac{1}{2}\omega^2 e^{-2\varphi_0}\,\Phi
    + \order{\Phi^2}
\end{equation}
collects the nonlinear backreaction terms from
Eqs.~(\ref{eq:self_consistent})
and~(\ref{eq:U_effective}).  The problem is nonlinear
because $V_{\mathrm{eff}}$ depends on $\Phi$ itself,
and the static background $\varphi_0(r)$ is generated
self-consistently by the time-averaged stress-energy
of the oscillation.

The confining potential arises from the nonlinear coupling
between $\varphi$ and the bi-conformal metric: the metric
$g_{\mu\nu} = \mathrm{diag}(-e^{\varphi},\,
e^{-\varphi}\delta_{ij})$ depends on $\varphi$, so the
oscillating mode modifies its own propagation environment.
Newton's constant $G$ does not appear in the wave equation
$\Box\varphi = 0$; it enters only through the gravitational
self-energy of the oscillon (a negligible correction
discussed below).  The resonant cavity is therefore a
\emph{geometric} property of the nonlinear wave equation,
not a gravitational effect.

The boundary conditions for bound states are:
\begin{equation}\label{eq:bc_mass}
  u(0) = 0,\quad u'(0) = 1;\qquad
  u \to 0 \;\text{as}\; r \to \infty.
\end{equation}
The eigenvalue is $\omega$; the oscillon amplitude
$\Phi_0 \equiv \Phi(0)$ is the primary parameter that
controls the depth of the confining potential.

\textit{Dimensionless formulation.}---Introduce the
Compton radius $R_c \equiv \hbar/(\tilde{m}\,c)$ as the
length scale, where $\tilde{m}$ is a reference mass (to
be identified with the ground-state mass a posteriori),
and define
\begin{equation}\label{eq:dimless_vars}
  x \equiv \frac{r}{R_c},\qquad
  \Omega \equiv \frac{\omega R_c}{c},\qquad
  f(x) \equiv \frac{u(r)}{R_c}.
\end{equation}
The mode equation becomes
\begin{equation}\label{eq:dimless_mode}
  f'' + \left[\Omega^2 e^{-2\varphi_0}
    - \frac{\ell(\ell+1)}{x^2}
    - \hat{V}_{\mathrm{eff}}\right]f = 0,
\end{equation}
with $\hat{V}_{\mathrm{eff}}$ the dimensionless form
of~(\ref{eq:Veff_explicit}).  Within the formal cavity
ansatz, the mass of the $n$-th resonance is taken as
\begin{equation}\label{eq:mass_from_omega}
  m_n = \frac{\hbar\,\Omega_n}{R_c\,c}
  = \tilde{m}\,\Omega_n.
\end{equation}

\textit{Amplitude--frequency relation.}---The
one-dimensional reduction of
Sec.~\ref{sec:resonances} provides the explicit oscillon
solution $\Phi = A\,\mathrm{sech}^2(Bx)$ with
\begin{equation}\label{eq:A_Omega}
  A = \tfrac{3}{2}(1 - \Omega^2),\qquad
  B = \tfrac{1}{2}\sqrt{1 - \Omega^2},
\end{equation}
valid for $\Omega^2 < 1$ (bound-state condition).  This
links the oscillon amplitude to the eigenfrequency: a
deeper potential well (larger $A$, smaller $\Omega$)
supports more tightly bound modes.  The three-dimensional
generalization modifies the numerical coefficients through
the angular barrier $\ell(\ell+1)/x^2$ and the background
gradient $\varphi_0'$, but the qualitative
structure persists: a one-parameter family of oscillons
parametrized by $\Phi_0$ (or equivalently $\Omega$).

\textit{Formal cavity mass-ratio ansatz.}---Within
the self-consistent confining potential, different angular
channels ($\ell$) and radial overtones ($n$) yield
distinct eigenfrequencies $\Omega_{n\ell}$.  Under the same
formal cavity ansatz, the mass ratios are
\begin{equation}\label{eq:mass_ratios}
  \frac{m_{n_1\ell_1}}{m_{n_2\ell_2}}
  = \frac{\Omega_{n_1\ell_1}}{\Omega_{n_2\ell_2}}
\end{equation}
and would be geometric: they depend on the mode numbers
and the cavity profile but are independent of $G$, $\hbar$,
or the absolute mass scale.  This relation is retained as
a diagnostic of the direct-cavity programme; its physical
status for charged leptons is fixed by the numerical audit
below.

\textit{Epistemic status of
Eq.~(\ref{eq:mass_ratios}).}---Direct numerical evaluation
of the self-consistent ISPG oscillon with the canonical
nonlinearity $\Phi\,e^{-\alpha\Phi}$ ($\Phi_0 = 2.35,\,
\alpha = 0.5$) yields exactly four bound cavity modes
$(n,\ell)\in\{(0,0),(1,0),(0,1),(0,2)\}$ with dimensionless
eigenfrequencies $\Omega_{n\ell}\in(0.66,\,1.00)$; the
resulting linear ratios $\Omega_{n\ell}/\Omega_{n'\ell'}$
do \emph{not} reproduce the observed charged-lepton
hierarchy $m_\tau:m_\mu:m_e\approx 3477:207:1$, nor does
any monotone power of $\Omega$ alone match this hierarchy
to better than $\sim 20\%$ (verified in the companion
code \texttt{QUANTUM/Python/cavity\_vs\_mass\_audit.py}
and \texttt{QUANTUM/Python/cavity\_quantum\_oscillon.py}).
The identification
of the three charged leptons with the specific $(n,\ell)$
modes listed above is therefore a \emph{structural}
labelling of particle modes inherited from the cavity
spectrum; the \emph{quantitative} mass ratios are
provided by the Mathieu mass formula
$m_N = m_1\,b_N(q{=}1.853)$ of the companion MASS
paper~\cite{Lotishvili2026Mass}, with
$(N_e,N_\mu,N_\tau) = (5,72,295)$ reproducing the
observed ratios to better than $0.2\%$.  A formal,
structure-preserving reduction from the three-dimensional
cavity spectrum to the one-dimensional Mathieu index $N$
is not established here and is deferred as future work.
The Koide relation (Sec.~\ref{sec:koide}) below functions
as an independent algebraic consistency check on any
candidate identification.

\textit{Absolute mass scale.}---The physical mass
$m_n = \hbar\omega_n/c^2$ is fixed once the effective
cavity radius $R_{\mathrm{core}}$ is known:
\begin{equation}\label{eq:R_core}
  R_{\mathrm{core}}
  \sim \frac{\hbar}{m_n c}
  = \frac{1}{\Omega_n}\,
    \frac{\hbar}{\tilde{m}\,c},
\end{equation}
the Compton wavelength of the resonance, as expected
on dimensional grounds.  The cavity radius is set by
the nonlinear self-consistency condition: the oscillon
amplitude $\Phi_0$ must be large enough to create a
confining potential through the $\Phi^2$ and $\Phi^3$
terms in~(\ref{eq:self_consistent}), but small enough
to remain in the quasi-bound regime ($\Omega < 1$).  For
the $\mathrm{sech}^2$ oscillon of the 1D model, the core
radius scales as $R_{\mathrm{core}} \sim 1/B =
2/\sqrt{1-\Omega^2}$ in Compton units, confirming
$R_{\mathrm{core}} \sim \hbar/(mc)$ to within
order-unity factors.

\textit{Gravitational self-energy correction.}---Beyond
the wave-equation dynamics, the oscillon's mass $m$
generates a gravitational field
$\varphi_{\mathrm{grav}} \sim -Gm/(c^2 r)$ that modifies
the background at the core by a relative amount
\begin{equation}
  \frac{|\varphi_{\mathrm{grav}}|}{1}
  \bigg|_{r \sim R_{\mathrm{core}}}
  \sim \frac{Gm}{c^2 R_{\mathrm{core}}}
  = \frac{1}{2}\,\frac{m^2}{\Mpl^2}
  \equiv \frac{\kappa}{2}.
\end{equation}
For the electron,
$\kappa = (m_e/\Mpl)^2 \approx 1.7 \times 10^{-45}$.
The resulting eigenfrequency correction is
\begin{equation}\label{eq:mass_correction}
  \frac{\delta m_n}{m_n}
  \sim \kappa
  \sim \left(\frac{m_n}{\Mpl}\right)^2,
\end{equation}
a relative shift of order $10^{-45}$ for the
electron---utterly negligible.  Within the formal cavity
ansatz of Eqs.~(\ref{eq:mass_from_omega})--(\ref{eq:mass_ratios}),
the ratios are geometric and the gravitational self-energy
correction is therefore negligible.  The physical
charged-lepton ratios are supplied by the Mathieu mass formula
of the companion MASS paper~\cite{Lotishvili2026Mass}, as
recorded in the epistemic-status remark below
Eq.~(\ref{eq:mass_ratios}).

\textit{Lattice formulation.}---The nonlinear eigenvalue
problem~(\ref{eq:dimless_mode}) admits a natural lattice
discretization.  Replace the continuum by a cubic lattice
with spacing $a$ and sites
$\mathbf{n} \in \mathbb{Z}^3$.  The spatial Laplacian
becomes the discrete operator
\begin{equation}\label{eq:lattice_laplacian}
  (\nabla^2_a\,\varphi)_\mathbf{n}
  = \frac{1}{a^2}\sum_{|\mathbf{e}|=1}
    (\varphi_{\mathbf{n}+\mathbf{e}}
    - \varphi_\mathbf{n}),
\end{equation}
and the nonlinear eigenvalue problem reduces to finding
localized solutions of a finite-dimensional system of
coupled nonlinear equations.  For $N \sim 10^2$ sites
per side (sufficient to resolve the Compton wavelength
with $a \sim R_c/N$), the system is tractable with
standard iterative eigenvalue algorithms.

Separately, the ultralocal kinetic structure of the
canonical theory (Sec.~\ref{sec:physical_interp})
ensures that the \emph{quantum} lattice
formulation---discretizing the Hamiltonian
constraint~(\ref{eq:H_constraint}) rather than the
classical mode equation---inherits the same
nearest-neighbour coupling structure, enabling Monte
Carlo evaluation of the constrained path integral as
an independent non-perturbative check on the mass
spectrum.

%-------------------------------------------------------------------
\subsection{The Koide formula as a resonance constraint}
\label{sec:koide}
%-------------------------------------------------------------------

In 1981, Koide discovered an empirical relation among the
three charged lepton masses~\cite{Koide1983,Foot1994,Rivero2005}:
\begin{equation}\label{eq:koide}
  \boxed{
  Q \equiv \frac{m_e + m_\mu + m_\tau}
    {(\sqrt{m_e} + \sqrt{m_\mu} + \sqrt{m_\tau})^{2}}
  = \frac{2}{3},
  }
\end{equation}
which is satisfied to better than $0.01\%$ with current
mass measurements.  This precision is striking: a random
triple of masses would satisfy $Q = 2/3$ with probability
$\ll 10^{-4}$.  No explanation for this relation exists
within the Standard Model.

In the ISPG resonance framework, the Koide relation is an
\emph{algebraic consistency condition} that any candidate
3D-resonator eigenvalue spectrum must satisfy if it is to
match the observed charged-lepton hierarchy.  To see
this, consider the general eigenvalue problem
\begin{equation}\label{eq:gen_eigenvalue}
  \hat{H}\,\psi_n = \lambda_n\,\psi_n,
\end{equation}
where $\hat{H}$ is the effective Hamiltonian of the
resonant mode in the self-consistent potential
(Sec.~\ref{sec:resonances}).  The eigenvalue gives
the rest energy directly, $\lambda_n = m_n c^{2}$.

The Koide ratio $Q = 2/3$ can be rewritten as a
constraint on the eigenvalues:
\begin{equation}\label{eq:koide_constraint}
  \frac{\sum_{n=1}^{3} \lambda_n}
    {\left(\sum_{n=1}^{3} \lambda_n^{1/2}\right)^{2}}
  = \frac{2}{3}.
\end{equation}
This is equivalent to requiring that the three eigenvalues
lie on a specific curve in $(\lambda_1, \lambda_2, \lambda_3)$
space.  We now prove that this curve is the locus of eigenvalues
of a \textbf{symmetric rank-1 perturbation} of a degenerate
system.

\textit{Proof.}---Consider the unperturbed Hamiltonian $\hat{H}_0$
with threefold degenerate eigenvalue $\lambda_0$:
\begin{equation}
  \hat{H}_0 \psi_n^{(0)} = \lambda_0 \psi_n^{(0)}, \qquad n = 1, 2, 3.
\end{equation}
Add a symmetric rank-1 perturbation $\delta\hat{H} = \alpha |v\rangle\langle v|$,
where $|v\rangle$ is a normalized vector in the degenerate subspace.
In the basis where $|v\rangle = (1, 1, 1)/\sqrt{3}$, the perturbation
matrix is
\begin{equation}
  \delta H = \frac{\alpha}{3}
  \begin{pmatrix}
    1 & 1 & 1 \\
    1 & 1 & 1 \\
    1 & 1 & 1
  \end{pmatrix}.
\end{equation}

The perturbed eigenvalues are found by diagonalizing
$\lambda_0 \mathbb{1} + \delta H$.  The eigenvalues of $\delta H$
are: $\alpha$ (non-degenerate, eigenvector $(1,1,1)$) and $0$
(twofold degenerate, eigenvectors orthogonal to $(1,1,1)$).  Thus
the perturbed eigenvalues are:
\begin{equation}
  \lambda_1 = \lambda_0 + \alpha, \qquad
  \lambda_2 = \lambda_3 = \lambda_0.
\end{equation}

The Koide ratio for these eigenvalues is:
\begin{equation}
  Q = \frac{(\lambda_0 + \alpha) + \lambda_0 + \lambda_0}
    {\bigl(\sqrt{\lambda_0 + \alpha} + \sqrt{\lambda_0} + \sqrt{\lambda_0}\bigr)^2}.
\end{equation}

As shown in Appendix~\ref{app:theorem_koide}, defining
$x \equiv \alpha/\lambda_0$, the Koide ratio is the
monotonically increasing function
\begin{equation}
  Q(x) = \frac{3 + x}
    {\bigl(\sqrt{1+x} + 2\bigr)^2},
\end{equation}
which interpolates from $Q(0) = 1/3$ (exact degeneracy) to
$Q \to 1$ ($x \to \infty$).  Solving $Q(x) = 2/3$ yields
$x = 33 + 24\sqrt{2} \approx 66.9$; the physical value
$Q = 2/3$ is therefore attained at a \emph{specific},
non-perturbative perturbation strength determined by
the nonlinear dynamics.  The corrections to $Q = 2/3$ are
controlled by the secondary degeneracy-lifting parameter
$\epsilon$ (the splitting of $\lambda_2$ and $\lambda_3$),
not by $\alpha$ itself.

This establishes the physical picture:
\begin{enumerate}
  \item In the unperturbed system (free medium, no
    self-interaction), all three modes are degenerate at $\lambda_0$.
  \item The nonlinear self-interaction of the resonant
    mode (metric backreaction) acts as a symmetric rank-1 perturbation
    $\delta\hat{H} \propto |\varphi_0\rangle\langle\varphi_0|$ that
    lifts the degeneracy.
  \item The specific form of the bi-conformal coupling,
    through which the scalar field $\varphi$ couples universally
    to all masses via $\meff = m_0 e^{\varphi/2}$, ensures the
    perturbation is approximately rank-1 in the $(e, \mu, \tau)$
    subspace.  Since neither $\alpha$ nor $\lambda_0$ is a free
    parameter---both are fixed by the self-consistent oscillon
    solution---the theory's self-consistent oscillon problem
    is required to produce the value
    $x \equiv \alpha/\lambda_0$ that satisfies $Q = 2/3$.
    Mathematically this happens if and only if
    $x = 33 + 24\sqrt{2} \approx 66.9$; deriving this
    value of $x$ from the self-consistent oscillon problem
    is the residual numerical task.
\end{enumerate}

Computing $x$ from first principles requires solving the
full nonlinear eigenvalue system of
Sec.~\ref{sec:mass_determination}.  As stated there, the
cavity geometry supplies the structural mode labels and
the nonlinear eigenvalue problem; the quantitative
charged-lepton mass ratios are currently carried by the
MASS/Mathieu ladder of the companion MASS
paper~\cite{Lotishvili2026Mass}, with the formal 3D-to-1D
reduction remaining the open numerical task.  Gravitational
self-energy corrections of order $(m/\Mpl)^2 \sim 10^{-45}$
are negligible.
Whether the self-consistent cavity spectrum by itself
yields $x \approx 66.9$ is a concrete, falsifiable
prediction of ISPG.  At present the numerical cavity
eigenfrequencies $\Omega_{n\ell}$ of
Sec.~\ref{sec:mass_determination} do not by themselves
reproduce $Q=2/3$: direct substitution
$\lambda_n \mapsto \Omega_{n\ell}$ gives $Q\simeq 0.34$,
off by $\sim 50\%$.  What \emph{does} hold numerically
is the Koide identity evaluated on the companion MASS
paper's Mathieu-ladder values~\cite{Lotishvili2026Mass},
$(N_e,N_\mu,N_\tau) = (5,72,295)$:
\begin{equation}\label{eq:koide_Nsquare}
  \frac{N_e^{2}+N_\mu^{2}+N_\tau^{2}}
       {(N_e+N_\mu+N_\tau)^{2}}
  \;=\; \frac{92234}{138384}
  \;=\; 0.66651\ldots ,
\end{equation}
matching $2/3$ to $0.024\%$, in direct correspondence with
the experimental Koide ratio $2/3$ (to $0.01\%$).  The rank-1
perturbation argument therefore plays a \emph{motivational}
role: it shows that the Mathieu spectrum of the MASS paper
is consistent with a rank-1 lifting structure and isolates
$x = 33+24\sqrt{2}\approx 66.9$ as the value at which the
cavity and Mathieu frameworks would coincide.  A
first-principles derivation of this value directly from
the cavity eigenvalues of
Sec.~\ref{sec:mass_determination} is an open calculation
(numerical audit: \texttt{QUANTUM/Python/cavity\_vs\_mass\_audit.py}).
The Koide constraint also serves as a powerful algebraic
consistency check: any candidate resonance spectrum that
does not satisfy $Q = 2/3$ can be excluded, even when the
spectrum is generated from a different effective formulation.

\textit{Beyond rank-1.}---The pure rank-1 perturbation
gives $\lambda_2 = \lambda_3$; the physical mass splitting
$m_\tau/m_\mu \approx 17$ requires a secondary perturbation
$\epsilon$ that lifts this degeneracy.  As shown in
Appendix~\ref{app:theorem_koide}
(Eq.~\ref{eq:Q_lifted}), the Koide ratio receives
corrections of order $\epsilon^{2}/(\lambda_{0}\alpha)$,
not $\alpha/\lambda_{0}$.  For the charged leptons,
$\epsilon/\alpha \sim m_\mu/m_\tau \approx 0.06$, giving
corrections $\sim 4 \times 10^{-3}$---consistent with
$Q = 2/3$ holding to $0.01\%$.  The precise values of
$m_e$, $m_\mu$, $m_\tau$ are determined by the full
nonlinear eigenvalue system
(Sec.~\ref{sec:mass_determination}), whose lattice
formulation reduces the problem to a finite-dimensional
nonlinear system tractable with iterative algorithms.

\textit{Numerical solution of the cavity spectrum.}---We have
performed this computation for the exponential nonlinearity
$F(\Phi) = \Phi\bigl(1 - e^{-\alpha\Phi}\bigr)$ with
$\alpha = 1/2$.  We clarify the epistemic status of this
choice.  The Taylor expansion
$F(\Phi) = \alpha\Phi^{2}
 - \tfrac{1}{2}\alpha^{2}\Phi^{3}
 + \mathcal{O}(\Phi^{4})$
has a leading quadratic coefficient $\alpha$ that matches
the $\tfrac{1}{2}\omega^{2}e^{-2\varphi_{0}}\Phi^{2}$ term
of Eq.~(\ref{eq:self_consistent}) in natural units
($\omega^{2}e^{-2\varphi_{0}} = 1$) precisely when
$\alpha = 1/2$; in this sense the leading-order coefficient
of the oscillon nonlinearity is \emph{fixed by the ISPG
action}, not a free parameter.  The specific
\emph{exponential resummation} of higher-order terms
$\Phi^{3},\,\Phi^{4},\ldots$, however, is a model choice
of ultraviolet completion within the Gleiser-type
nonlinear-Klein--Gordon class identified in
Sec.~\ref{sec:resonances}: many UV completions share the
same quadratic leading term but differ at cubic and higher
order, and ISPG's quadratic-only action does not by itself
single out the exponential form.  The cavity spectrum
reported below is therefore specific to this nonlinearity
class; whether the Koide relation $Q = 2/3$ persists under
alternative UV completions compatible with the same
leading quadratic term $\tfrac{1}{2}\Phi^{2}$ is a
robustness question left for future work.
The background oscillon at amplitude $\Phi_0 = 2.35$
(frequency $\Omega_{\rm bg} = 0.866$) creates a confining
potential that supports exactly four bound states:
\begin{center}
\begin{tabular}{llccc}
\hline
Particle & Mode & $\omega$ & $\kappa$ & $E_{1D}$ \\
\hline
$\tau$     & $(n{=}0,\,\ell{=}0)$ & $0.664$ & $0.748$ & $1.156$ \\
$\mu$      & $(n{=}1,\,\ell{=}0)$ & $0.970$ & $0.241$ & $0.067$ \\
$e$        & $(n{=}0,\,\ell{=}2)$ & $0.999$ & $0.044$ & $4.2 \times 10^{-4}$ \\
grav.      & $(n{=}0,\,\ell{=}1)$ & $0.866$ & $0.500$ & --- \\
\hline
\end{tabular}
\end{center}
\textit{The $\ell = 1$ mode controls the source-following
gravitational response.}---We now prove that the $\ell = 1$
cavity mode is not an independent particle but the oscillon's
translational (Goldstone) mode.  Its gravitational role is to
ensure that the long-range field follows the motion of the
localized oscillon core.

\textit{(i)~Exact identity $\psi_1(r) = \Phi_0'(r)$.}---The
background oscillon $\Phi_0(r)$ satisfies the radial field
equation.  An infinitesimal translation by $\delta$ along the
$z$-axis produces
$\Phi_0(|\mathbf{r} - \delta\hat{z}|)
= \Phi_0(r) - \delta\cos\theta\,\Phi_0'(r) + \order{\delta^2}$.
Since the translated oscillon also satisfies the field equation
(by translational invariance of the Lagrangian), the perturbation
$\delta\varphi = \cos\theta\,\Phi_0'(r)\cos(\Omega t)$ is an
\emph{exact} eigenmode at frequency $\omega = \Omega_{\rm bg}$.
Because $\cos\theta = Y_{10}/\sqrt{4\pi/3}$, this is the
$\ell = 1$ mode with radial eigenfunction
$u_1(r) = r\,\Phi_0'(r)$.  We verify this numerically: the
Pearson correlation between the computed eigenvector and
$r\,\Phi_0'(r)$ is $0.9999999999$, with a maximum point-wise
residual of $1.3 \times 10^{-5}$ and
$|\omega_{\ell=1} - \Omega_{\rm bg}| = 1.7 \times 10^{-6}$.
The identity is exact.

\textit{(ii)~Zero excitation energy (Goldstone theorem).}---The
oscillon is localized and thus breaks the continuous translational
symmetry of the field equation.  By Goldstone's theorem, this
broken symmetry produces a zero-energy mode.  Indeed, translating
the oscillon does not change its total energy:
$E[\Phi_0(r)] = E[\Phi_0(|\mathbf{r} - \delta\hat{z}|)]$ for all
$\delta$.  Therefore the $\ell = 1$ perturbation carries
\emph{zero} excitation energy---it is not a particle.
The formal energy $E_{1D}(\omega_{\ell=1})$ represents the
oscillon's translational kinetic energy $\tfrac{1}{2}Mv^2$,
not a new mass eigenstate.
\textit{(iii)~Collective-coordinate gravitational response.}---Let
$\mathbf{R}(t)$ denote the oscillon center in the adiabatic
regime, where the translational collective coordinate varies on
timescales much longer than $\Omega_{\rm bg}^{-1}$.  The
translated family of configurations is then, to leading order,
\begin{equation}
  \varphi(\mathbf{r},t;\mathbf{R})
  \simeq \varphi_{\mathrm{grav}}(|\mathbf{r}-\mathbf{R}(t)|)
  + \Phi_0(|\mathbf{r}-\mathbf{R}(t)|)\cos(\Omega_{\rm bg} t).
\end{equation}
Expanding to first order in $\mathbf{R}(t)$ gives
\begin{equation}
  \delta\varphi_{\mathbf{R}}
  = -\mathbf{R}(t)\!\cdot\!\nabla
    \Bigl[\varphi_{\mathrm{grav}}(r)
    + \Phi_0(r)\cos(\Omega_{\rm bg} t)\Bigr]
  = -\bigl(\mathbf{R}(t)\!\cdot\!\hat{r}\bigr)
    \Bigl[\varphi_{\mathrm{grav}}'(r)
    + \Phi_0'(r)\cos(\Omega_{\rm bg} t)\Bigr].
\end{equation}
Because $\mathbf{R}\!\cdot\!\hat{r}$ is purely $\ell = 1$, the
oscillatory part of this tangent vector is exactly the
Goldstone profile found above, $\psi_1(r)\propto \Phi_0'(r)$.
Time-averaging over the fast internal oscillation leaves
\begin{equation}
  \langle\delta\varphi_{\mathbf{R}}\rangle_t
  = -\mathbf{R}(t)\!\cdot\!\nabla\varphi_{\mathrm{grav}}(r)
  = -\frac{r_s\,\mathbf{R}(t)\!\cdot\!\hat{r}}{r^2},
\end{equation}
which is precisely the dipole term in the translated
gravitational field
\begin{equation}
  -\frac{r_s}{|\mathbf{r}-\mathbf{R}(t)|}
  = -\frac{r_s}{r}
  - \frac{r_s\,\mathbf{R}(t)\!\cdot\!\hat{r}}{r^2}
  + \order{R^2}.
\end{equation}
Thus the same collective coordinate that translates the
localized oscillon core also translates its long-range
gravitational dressing.  We verify numerically that the dipole
component of the displaced oscillon field equals
$\delta\,\Phi_0'(r)$ to $0.04\%$ accuracy.  The $\ell = 1$ mode
is therefore the channel through which the particle's
gravitational field tracks its spatial position.  It is the
source-following gravitational mode of the oscillon, not an
additional mass eigenstate.

Thus the cavity contains exactly \textbf{three physical
particle modes}---matching the three charged lepton
generations---plus one translational mode that carries the
source-following gravitational response of the oscillon.

\textit{Uniqueness of the mode assignment.}---The four cavity
modes admit four possible triples.  Only the
$(\tau,\,\mu,\,e)$ assignment gives $Q$ near $2/3$:
\begin{center}
\begin{tabular}{lcc}
\hline
Triple & $Q$ & $|Q - 2/3|$ \\
\hline
$(\tau,\,\ell{=}1,\,\mu)$ & $0.414$ & $0.25$ \\
$(\tau,\,\ell{=}1,\,e)$   & $0.510$ & $0.16$ \\
$(\tau,\,\mu,\,e)$        & $0.6667$ & $6.8 \times 10^{-5}$ \\
$(\ell{=}1,\,\mu,\,e)$    & $0.583$ & $0.084$ \\
\hline
\end{tabular}
\end{center}
The assignment is not chosen by hand: it is the unique
combination consistent with $Q \approx 2/3$, and the excluded
$\ell = 1$ mode is independently proven to be the
translational Goldstone mode~(see above).

\textit{The mass formula $E_{1D}$.}---The one-dimensional
oscillon equation $\Phi'' + (\Omega^2 - 1)\Phi + \Phi^2 = 0$
admits the exact solution
$\Phi(x) = \tfrac{3}{2}\kappa^2\,\mathrm{sech}^2(\kappa x/2)$
with $\kappa^2 = 1 - \Omega^2$.  Its kinetic energy is
\begin{equation}\label{eq:E1D_derivation}
  E_{\rm kin} = \int_{-\infty}^{\infty}\!
    \bigl[\tfrac{1}{2}\Omega^2\Phi^2
    + \tfrac{1}{2}(\Phi')^2\bigr]\,dx
  = \tfrac{3}{5}\,\kappa^3\,(4\Omega^2 + 1),
\end{equation}
where the integrals
$\int\!\Phi^2\,dx = 6\kappa^3$ and
$\int\!(\Phi')^2\,dx = \tfrac{6}{5}\kappa^5$ follow
from the standard identities
$\int_{-\infty}^{\infty}\!\mathrm{sech}^4\!u\,du = \tfrac{4}{3}$
and $\int_{-\infty}^{\infty}\!\mathrm{sech}^4\!u\,\tanh^2\!u\,du
= \tfrac{4}{15}$.  The factor $(4\Omega^2 + 1)$ is not
arbitrary: it encodes the exact ratio of temporal oscillation
energy ($3\Omega^2\kappa^3$) to spatial gradient energy
($\tfrac{3}{5}\kappa^5$).  We define
\begin{equation}
  E_{1D}(\Omega) \equiv \kappa^3\,(4\Omega^2 + 1)
  = \tfrac{5}{3}\,E_{\rm kin},
\end{equation}
where the constant $5/3$ cancels in the Koide ratio.
The Koide ratio evaluates to
\begin{equation}
  Q_{\text{E1D-on-cavity}} = 0.66674, \qquad
  |Q - \tfrac{2}{3}| = 6.8 \times 10^{-5},
\end{equation}
reproducing the Koide relation to $0.009\%$ accuracy.
Here $E_{1D}$ denotes the effective one-dimensional mass
prescription evaluated on the cavity data; the test
\texttt{SM/python/test\_e3d\_vs\_e1d\_background.py}
shows that $\langle H\rangle/E_{1D}$ has coefficient of
variation $0.82$ over the oscillon family.  Thus $E_{1D}$
is not a constant multiple of the full 3D oscillon energy,
and the full 3D-energy reduction remains open.

This formula is singled out by the data: among fourteen
candidate mass prescriptions---including
$\kappa^p$ for $p = 0.5,\,1,\,\ldots,\,3$,
$\kappa^p(4\Omega^2 + 1)$ for various $p$,
and the full 3D energy
$E_{3D} = 4\pi\!\int[\tfrac{1}{2}\Omega^2\Phi^2
+ \tfrac{1}{2}(\Phi')^2]\,r^2\,dr$---only
$E_{1D}$ reproduces $Q = 2/3$ to better than $1\%$;
the next-best formula deviates by $2.2 \times 10^{-2}$,
two orders of magnitude worse.  In a Monte Carlo test of
$10^5$ random eigenvalue triples, only $0.001\%$ produce
$|Q - 2/3| < 10^{-3}$ under $E_{1D}$, confirming that
the cavity spectrum is non-generic.

The raw mass ratios $1 : 159 : 2740$ are $\sim\!79\%$ of
the experimental values $1 : 207 : 3477$.  A systematic
scan of $\Phi_0$ reveals a remarkable structural coincidence:
the amplitude at which $Q$ is closest to $2/3$ ($\Phi_0 = 2.350$)
is the \emph{same} amplitude at which the mass-ratio correction
factor equals $4/\pi$ ($\Phi_0 = 2.3501$, with $Q = 0.66670$).
This is not an independent geometric adjustment; it is a
consequence of the Koide condition itself.

In the standard Koide parametrisation
$\sqrt{m_i} = M\bigl(1 + \sqrt{2}\cos(\theta_0 + 2\pi i/3)\bigr)$,
the experimental angle is $\theta_0^{\rm exp} = 0.2223 \approx 2/9$
(to $0.02\%$), while the $E_{1D}$ cavity spectrum gives
$\theta_0^{E_{1D}} = 0.2167$.  The small shift
$\Delta\theta_0 = 0.006$~rad, amplified by the extreme
sensitivity of $m_e$ to $\theta_0$ near the cosine
minimum, produces the correction factor $4/\pi$ for all
ratios involving the electron:
\begin{equation}
  \frac{m_\tau}{m_e}\bigg|_{\rm corr} = 3490
  \quad(\text{expt: } 3477,\ 0.37\%),
  \qquad
  \frac{m_\mu}{m_e}\bigg|_{\rm corr} = 202
  \quad(\text{expt: } 207,\ 2.1\%).
\end{equation}
Equivalently, for the $\ell = 2$ electron mode the energy
exponent shifts from $\kappa^3$ to $\kappa^{3.076}$,
and $\kappa_e^{0.076} = 0.789 \approx \pi/4$ ($0.39\%$).
The $2.1\%$ residual in $m_\mu/m_e$ is the $\tau/\mu$ error
(both $\ell = 0$), reflecting a sub-leading curvature correction
$\mathcal{O}(\kappa)$ from the 1D-to-3D mapping.
A formal derivation of $\theta_0 = 2/9$ from the cavity
spectrum---equivalent to the exact $4/\pi$ correction---remains
an open problem.  The first step toward its resolution is
described next.

\textit{Hartree self-consistency and $\theta_0 = 2/9$.}---We stress
at the outset that the target $\theta_0 = 2/9$ is an
\emph{empirical benchmark} inherited from the data (the experimental
value $\theta_0^{\rm exp} = 0.2223$ lies $0.02\%$ from $2/9$, see
above); the Hartree analysis that follows does \emph{not} derive
$2/9$ from first principles, but asks whether the bare cavity angle
$\theta_0^{E_{1D}} = 0.2167$ can be shifted to the empirical value by
a perturbative mechanism intrinsic to the ISPG oscillon, without
introducing new sectors or couplings. The
cavity modes backreact on the effective potential through the
mean-field (Hartree) correction
\begin{equation}\label{eq:hartree}
  \delta V(r) = -\tfrac{1}{2}\,
    F''[\Phi_{\rm bg}(r)]\,|\psi_\tau(r)|^{2}\,\eta,
\end{equation}
where $\psi_\tau$ is the ground-state eigenfunction (normalized
to unit $L^{2}$ norm), $F'' = d^{2}F_{\rm NL}/d\Phi^{2}$,
and $\eta$ measures the mode occupation strength.
Because the electron eigenfrequency
$\omega_e \approx 0.999$ is extremely close to the continuum
threshold, even a small perturbation $\delta V$ produces a
large relative shift in $m_e$.

Numerical evaluation shows that at $\eta^{*} = 0.074$ the
Koide angle shifts from $\theta_{0} = 0.2167$ (bare) to
$\theta_{0} = 0.2222 \approx 2/9$ with $0.07\%$ accuracy.
The Koide ratio simultaneously remains close to $2/3$: a
joint optimization over $(\Phi_{0}, \eta)$ yields
$Q = 0.6659$ ($0.11\%$ from $2/3$) and
$\theta_{0} = 0.2221$ ($0.07\%$ from $2/9$), reducing the
combined cost function by a factor of $300$ relative to the
bare ($\eta = 0$) spectrum.  The small coupling
$\eta^{*} \approx 0.07$ confirms that the backreaction is
perturbative: the first-order Hartree correction captures the
dominant effect.  Whether the residual $0.11\%$ in $Q$
closes at second order is an open expectation; the
channel-separated self-consistent Hartree--Fock run
described below improves $\theta_0$ but does not improve
$Q$, so the joint $(Q,\theta_0)$ closure beyond first
order remains a numerical task.

The Hartree mechanism thus provides a candidate physical
interpretation of the $4/\pi$ correction as a first-order
self-energy shift of the electron eigenfrequency due to the
ground-state backreaction on the confining potential.  Whether
this is the unique mechanism---as opposed to, for example, a
cavity-geometry correction at higher order---is not
established here.

\textit{Quantum determination of $\eta$.}---In the QFT
framework, the coupling $\eta$ is constrained (rather than
free) by the vacuum zero-point fluctuations of the cavity
modes, under the selected mode-content and overlap-weighting
prescription.  For a mode with angular momentum
$\ell$ and frequency $\omega_n$, the contribution to
the spherically symmetric (monopole) Hartree potential is
\begin{equation}\label{eq:eta_qft}
  \eta_n = \frac{2\ell + 1}{8\pi\,\omega_n}\,,
\end{equation}
where the factor $(2\ell+1)$ accounts for the magnetic
degeneracy.  For the tau mode ($\ell = 0$,
$\omega_\tau = 0.664$), this gives
$\eta_\tau = 0.060$.  This single contribution shifts
$\theta_0$ from $0.217$ to $0.221$---capturing $81\%$ of
the required correction.  Including the muon's
zero-point contribution, weighted by its spatial overlap
with the electron eigenfunction (overlap ratio $0.27$),
yields
\begin{equation}\label{eq:eta_predicted}
  \eta_{\rm pred} = \eta_\tau
    + \frac{\langle\psi_\mu^{2}|\psi_e^{2}\rangle}
           {\langle\psi_\tau^{2}|\psi_e^{2}\rangle}
    \,\eta_\mu
    = 0.060 + 0.27 \times 0.041
    = 0.071\,,
\end{equation}
in $4\%$ agreement with the empirical $\eta^{*} = 0.074$.
The predicted Koide angle is
$\theta_0 = 0.2220 \approx 2/9$ ($0.09\%$).
The iterative Hartree--Fock procedure converges in a
single iteration, confirming that the backreaction is
genuinely perturbative.  The residual $4\%$ discrepancy
between $\eta_{\rm pred}$ and $\eta^{*}$ is attributed
to second-order corrections and the angular coupling
of the $\ell = 2$ electron mode, which enters the
monopole Hartree potential only through its
angle-averaged density. We stress that the electron
mode does not contribute to its own Hartree potential
(standard self-interaction exclusion in mean-field theory)
and enters the $\ell = 0$ channel potential only through the
angle-averaged density of its $\ell = 2$ wavefunction;
had the electron's $(2\ell+1) = 5$ degeneracy entered
\eqref{eq:eta_qft} naively it would give $\eta_e \approx 0.199$
and overshoot the empirical value. A rigorous derivation of
the cross-channel suppression factor from angular-momentum
coupling is an open task and is deferred to future work.

\textit{Channel-separated Hartree--Fock.}---The cavity
modes reside in different angular-momentum channels:
$\tau$ and $\mu$ in $\ell = 0$, the electron in
$\ell = 2$.  A self-consistent Hartree--Fock treatment
must respect this structure.  In the full channel-separated
scheme, each channel receives backreaction only from modes
in other channels (excluding self-interaction):
\begin{align}\label{eq:channel_hf}
  \delta V_{\ell=0}(r) &= -\tfrac{1}{2}\,F''[\Phi_{\rm bg}]\,
    |\psi_{e}|^{2}\,\eta_{e}\,,\\
  \delta V_{\ell=2}(r) &= -\tfrac{1}{2}\,F''[\Phi_{\rm bg}]\,
    \bigl(|\psi_{\tau}|^{2}\,\eta_{\tau}
    + |\psi_{\mu}|^{2}\,\eta_{\mu}\bigr)\,.
\end{align}
However, the large degeneracy factor $(2\ell+1) = 5$
of the electron's $\ell = 2$ mode causes an excessive
shift in the $\ell = 0$ channel, degrading $Q$.

A physically motivated alternative is the
\emph{semiclassical} prescription: the heavy modes
$\tau$ and $\mu$ ($\kappa \gg 0$, deeply bound)
are treated as classical, receiving no quantum
correction, while only the light, nearly unbound
electron receives the backreaction from $\tau + \mu$:
\begin{equation}\label{eq:semiclassical_hf}
  \delta V_{\ell=2}^{\rm (sc)}(r) =
  -\tfrac{1}{2}\,F''[\Phi_{\rm bg}]\,
  \bigl(|\psi_{\tau}|^{2}\,\eta_{\tau}
  + |\psi_{\mu}|^{2}\,\eta_{\mu}\bigr),\qquad
  \delta V_{\ell=0}^{\rm (sc)} = 0\,.
\end{equation}
This is the standard semiclassical hierarchy: heavy
(classical) modes source the potential for light
(quantum) modes, not vice versa.

A scan over the oscillon amplitude $\Phi_0$ with four
prescriptions---bare, semiclassical~(\ref{eq:semiclassical_hf}),
full channel-separated~(\ref{eq:channel_hf}), and
same-channel (old)---shows that the semiclassical scheme
minimizes the combined cost
$(\Delta Q/Q_{\rm target})^{2} + (\Delta\theta_0/\theta_{\rm target})^{2}$
by a factor of $132$ relative to the bare spectrum,
achieving $Q = 0.6662$ ($0.06\%$) and
$\theta_0 = 0.2215$ ($0.3\%$) simultaneously.  The
full channel-separated scheme improves $\theta_0$ but
overshoots $Q$; the semiclassical prescription provides
the optimal balance, consistent with the expectation
that deeply bound modes are well approximated classically.
We note that this selection of the semiclassical scheme is
physically motivated (classical treatment of deeply bound
heavy modes is a standard hierarchy argument) but remains
a discrete choice among the four candidate prescriptions.
A first-principles derivation that uniquely selects the
semiclassical hierarchy from the underlying QFT---for
example, through an explicit $1/\omega_n$ expansion of the
effective action---would eliminate this residual ambiguity
and is an open task.

\textit{Structural origin of the amplitude $\Phi_0$.}---The
electron mode ($\ell = 2$) first becomes bound at the
threshold amplitude $\Phi_0 = 2.29$.  Below this threshold,
the cavity supports only three modes ($n{=}0$ and $n{=}1$
for $\ell = 0$, plus the translational $\ell = 1$), and the
best Koide ratio is $Q = 0.42$---far from $2/3$.  Above the
threshold, $Q$ is closest to $2/3$ at $\Phi_0 = 2.35$, just
$2.6\%$ above the threshold.  The Koide relation thus follows
from the oscillon amplitude sitting close to the electron
binding threshold; the small offset
$\Phi_0 - \Phi_{\rm thr} = 2.6\%$ is a continuous parameter
not fixed by any symmetry argument in the present framework,
and the Hartree coupling $\eta$ is determined partly by QFT
($\eta_{\rm pred} = 0.071$, with $81\%$ from $\tau$ and a small
muon contribution, see~\eqref{eq:eta_predicted}) and partly by
empirical optimization ($\eta^{*} = 0.074$, $4\%$ residual).
The mechanism is therefore best characterized as a physically
motivated self-consistency construction rather than a strict
parameter fit, but it does involve tunable choices that must
be acknowledged. The ratio is extremely sensitive to the
electron eigenfrequency: a $1\%$ shift in $\omega_e$ degrades
$Q$ from $6.8 \times 10^{-5}$ to $8.5 \times 10^{-2}$, which
supports the interpretation of the threshold proximity as the
structural origin of the Koide relation, but does not by itself
determine why Nature selects exactly $2.6\%$ above threshold.

\textit{Parameter count and epistemic status.}---For
transparency we summarize the effective parametric content of
the joint $(Q, \theta_0)$ reproduction: two continuous
parameters ($\Phi_0, \eta$) plus one discrete prescription
choice (semiclassical among four candidates), against two
observables. The $0.07$--$0.11\%$ residuals reported above
reflect the numerical optimization tolerance on a finite
$(\Phi_0, \eta)$ grid rather than an independent theoretical
prediction. A zero-parameter determination of $\Phi_0$ from
first principles, together with a full QFT derivation of
$\eta$ that includes the electron cross-channel coupling and
uniquely selects the semiclassical prescription, remains an
open programme of the theory.

%-------------------------------------------------------------------
\subsection{Particle stability and the evolving spectrum}
\label{sec:stability_spectrum}
%-------------------------------------------------------------------

The resonance eigenvalues $\omega_{nl}$ depend on the
background pressure through Eq.~(\ref{eq:eigenfreq}).
As the background pressure decreases over cosmic time
(Sec.~\ref{sec:entropy}), the resonance spectrum evolves:

\textit{Higher modes become unstable.}---A mode with
eigenfrequency $\omega_{nl}$ is stable if it lies below
the threshold set by the background pressure.  As the
pressure drops, the threshold descends, and modes that were
formerly stable become quasi-bound: their lifetime decreases
from $\infty$ to a finite value, and eventually they decay
into lower modes.

This provides a natural explanation for the evolving
particle content of the universe:
\begin{itemize}
  \item \textbf{Early universe} (high pressure): the
    quark-gluon plasma phase, with free quarks and gluons,
    was stable.  The experimental recreation of this phase
    at RHIC and LHC~\cite{ALICE2010} confirms that
    high-energy resonant modes that are unstable today
    were once stable constituents of matter.
  \item \textbf{Present epoch} (intermediate pressure):
    protons, neutrons, electrons, and photons are stable.
    Heavier particles (muon, tau, W/Z, top quark) are
    unstable and decay on timescales from $10^{-25}$~s
    to $2.2\,\mu$s.
  \item \textbf{Far future} (low pressure): if the
    pressure continues to decrease, eventually even the
    proton may become unstable.  Current experimental
    bounds ($\tau_p > 10^{34}$~yr from
    Super-Kamiokande~\cite{SuperK2020}) are consistent
    with an extremely slow destabilization driven by
    the cosmological pressure decline.
\end{itemize}

The rate of change of the stability threshold is set by
the entropy production rate (Sec.~\ref{sec:entropy_mech}),
which is logarithmically slow.  This ensures that the
proton lifetime, if finite, is vastly longer than the
current age of the universe---consistent with
observations.

%===================================================================
\section{The Higgs mechanism as conformal symmetry breaking}
\label{sec:higgs}
%===================================================================

In the Standard Model, particles acquire mass through the
Higgs mechanism: the electroweak vacuum spontaneously breaks
$SU(2)_L \times U(1)_Y \to U(1)_{\mathrm{EM}}$, and
the vacuum expectation value (VEV) $v = 246$~GeV sets the
mass scale.  At first glance, this appears to be an
independent mechanism from ISPG's pressure-dependent mass.
In fact, the two are deeply connected.

%-------------------------------------------------------------------
\subsection{The conformal Higgs potential}
\label{sec:conformal_higgs}
%-------------------------------------------------------------------

The standard Higgs potential is
\begin{equation}\label{eq:V_higgs}
  V(H) = -\mu^{2}\,|H|^{2} + \lambda\,|H|^{4},
\end{equation}
with the VEV $v = \mu/\sqrt{\lambda}$.  The
$\varphi$-dependence of the VEV is not an extra
postulate: it follows from the ISPG dimensional rule
applied to the parameters of Eq.~(\ref{eq:V_higgs}).

\paragraph{Derivation from dimensional scaling.}
In ISPG, any quantity with mass dimension $d$ carries
a factor $e^{d\,\varphi/2}$ when translated between a
local frame and the asymptotic reference frame (this is
the same rule that gives $\meff = m_0\,e^{\varphi/2}$ for
$d=1$).  Applying it to the two parameters of the Higgs
potential:
\begin{itemize}
  \item $\mu^{2}$ has mass dimension $2$, so
    $\mu^{2}(\varphi) = \mu_0^{2}\,e^{\varphi}$;
  \item $\lambda$ is dimensionless and, by the
    bi-conformal invariance of dimensionless couplings
    (a structural property of the theory, not an
    assumption introduced here), is
    $\varphi$-independent.
\end{itemize}
Therefore
\begin{equation}\label{eq:v_phi}
  \boxed{
  v^{2}(\varphi) \;=\; \frac{\mu^{2}(\varphi)}{\lambda}
  \;=\; v_0^{2}\,e^{\varphi},
  \qquad
  v(\varphi) = v_0\,e^{\varphi/2},
  }
\end{equation}
where $v_0 = 246$~GeV is the VEV measured in the local
frame.  No new coupling and no independent
$\varphi$-dependent parameter has been added: the
scaling in Eq.~(\ref{eq:v_phi}) is forced by the ISPG
dimensional rule together with the bi-conformal
invariance of $\lambda$.

Locally, $v$ is undetectably constant: an observer
using local mass-dimension units measures $v_0 = 246$~GeV
regardless of $\varphi$.  The $\varphi$-dependence in
Eq.~(\ref{eq:v_phi}) is only visible when comparing the
VEV at different positions through a common (asymptotic)
reference frame.

This means the Higgs mechanism is not independent of
gravity: it is the \textbf{electroweak projection} of the
same bi-conformal symmetry that governs the
scalar-metric sector.

\paragraph{Scope of this identification.}
What Eq.~(\ref{eq:v_phi}) establishes is the
$\varphi$-\emph{scaling} of the electroweak VEV: once
the local value $v_0 = 246$~GeV is taken as input, the
functional form $v(\varphi) = v_0\,e^{\varphi/2}$ is
fixed by the ISPG dimensional rule.  What it does
\emph{not} provide is a microscopic derivation of the
numerical value $v_0$ itself from the ISPG condensate
parameter $\Phi_0$.  Fixing $v_0$ in terms of
$\Phi_0$ and the underlying hysteresis sector would
require a full coarse-graining of the electroweak
condensate on top of the Route B microscopic closure,
and is left as a topic for a dedicated paper.  The
present section uses only the \emph{scaling} of $v$,
which is what enters every observable consequence
derived below.

%-------------------------------------------------------------------
\subsection{Bi-conformal link between electroweak and gravitational masses}
\label{sec:ew_grav}
%-------------------------------------------------------------------

What is being linked here is the \emph{mass-scaling}
sector, not the gauge sector.  The electroweak gauge
group, its couplings, and the Yukawa hierarchy are
taken as Standard-Model inputs; ISPG does not derive
them.  What ISPG fixes is that every Standard-Model
mass inherits the same $\varphi$-scaling as a
gravitational mass in Eq.~(\ref{eq:meff}), because both
originate from the same bi-conformal rule applied in
Eq.~(\ref{eq:v_phi}).  The identification
$v(\varphi) = v_0\,e^{\varphi/2}$ has three immediate
consequences:

\textit{(i) Mass generation is universal.}---Every Standard
Model mass (fermions via Yukawa couplings $m_f = y_f v/\sqrt{2}$,
and gauge bosons via $m_W = gv/2$) inherits the
$e^{\varphi/2}$ factor:
\begin{equation}\label{eq:all_masses}
  m_f(\varphi) = \frac{y_f\,v_0}{\sqrt{2}}\,e^{\varphi/2},
  \qquad
  m_W(\varphi) = \frac{g\,v_0}{2}\,e^{\varphi/2},
\end{equation}
precisely reproducing Eq.~(\ref{eq:meff}).  The ISPG
effective mass and the Higgs-generated mass are the
\emph{same} quantity viewed from different perspectives.

\textit{(ii) The Higgs boson mass is position-dependent.}---The
Higgs boson mass $m_H = \sqrt{2\lambda}\,v$ also scales as
$e^{\varphi/2}$, but this is locally undetectable: the ratio
$m_H/m_W = \sqrt{2\lambda}/g$ is dimensionless and therefore
$\varphi$-independent.

\textit{(iii) Electroweak phase transition in the early
universe.}---In the early universe, when the background
pressure was higher ($\varphi$ closer to zero), the VEV
was larger: $v_{\mathrm{early}} > v_{\mathrm{today}}$ (as
measured by an asymptotic observer).  The electroweak
phase transition temperature $T_{\mathrm{EW}} \propto v$
was correspondingly higher.  However, locally measured
temperatures also scale by $e^{\varphi/2}$, so the
dimensionless ratio $T_{\mathrm{EW}}/T_{\mathrm{local}}$
is preserved: the phase transition occurred at the same
locally measured temperature as in the standard scenario.

%-------------------------------------------------------------------
\subsection{Prediction: gravitational correction to the Higgs mass}
\label{sec:higgs_prediction}
%-------------------------------------------------------------------

The coupling between the Higgs field and the scalar-metric
sector introduces a gravitational correction to the Higgs
self-energy.  At one-loop level, the gravitational correction
to the Higgs mass parameter is
\begin{equation}\label{eq:delta_mH}
  \delta m_H^{2}
  \sim \frac{m_H^{2}}{16\pi^{2}}\,
    \frac{m_H^{2}}{\Mpl^{2}}
  \sim 10^{-34}\,m_H^{2},
\end{equation}
i.e.\ $\delta m_H/m_H \sim 10^{-17}$.  This is many
orders of magnitude below the sensitivity of any currently
envisaged collider (LHC: $\sim 10^{-3}$; FCC-ee/CEPC:
$\sim 10^{-4}$--$10^{-5}$), so Eq.~(\ref{eq:delta_mH})
is \emph{not} a near-term experimental discriminant.
Its role here is theoretical: it records the
gravitational piece of the Higgs self-energy that the
ISPG sector unavoidably contributes, and confirms that
this piece is Planck-suppressed and therefore consistent
with the observed Higgs mass without requiring tuning
of new parameters.  Any observable $\varphi$-dependent
signature of the Higgs sector has to come from the
universal scaling in Eq.~(\ref{eq:v_phi}), not from this
loop correction.

%-------------------------------------------------------------------
\subsection{Gauge boson spectrum on the bi-conformal background}
\label{sec:gauge_spectrum}
%-------------------------------------------------------------------

Section~\ref{sec:ew_grav} established that every Standard
Model mass inherits the universal factor $e^{\varphi/2}$.
Once the electroweak gauge algebra and couplings are taken as
the experimentally fixed internal data, the broken-phase gauge
sector on the bi-conformal background is fully determined.
We now derive the curved-space electroweak mass matrix, the
resulting gauge-boson spectrum, and its observational
consequences.

\textit{Broken-phase mass matrix on the bi-conformal
background.}---Start from the curved-space electroweak action
\begin{equation}\label{eq:ew_action_bg}
  S_{\mathrm{EW}}
  = \int d^4x\,\sqrt{-g}\left[
      -\frac{1}{4}W^a_{\mu\nu}W^{a\mu\nu}
      -\frac{1}{4}B_{\mu\nu}B^{\mu\nu}
      + (D_\mu H)^\dagger D^\mu H
      - \lambda\!\left(H^\dagger H
        - \frac{v(\varphi)^2}{2}\right)^2
    \right],
\end{equation}
with
\begin{equation}\label{eq:ew_covariant_derivative}
  D_\mu H
  = \left(
      \partial_\mu
      - \frac{i g}{2}\tau^a W^a_\mu
      - \frac{i g'}{2} B_\mu
    \right)H .
\end{equation}
In unitary gauge,
\begin{equation}\label{eq:higgs_unitary_bg}
  H(x) = \frac{1}{\sqrt{2}}
  \begin{pmatrix}
    0\\
    v(\varphi) + h(x)
  \end{pmatrix},
\end{equation}
so the quadratic gauge contribution from
$(D_\mu H)^\dagger D^\mu H$ is
\begin{equation}\label{eq:ew_mass_matrix_bg}
  \mathcal{L}_{\mathrm{mass}}
  = \sqrt{-g}\,\frac{v(\varphi)^2}{8}\left[
      g^2 W_\mu^1 W^{1\mu}
      + g^2 W_\mu^2 W^{2\mu}
      + \bigl(g W_\mu^3 - g' B_\mu\bigr)^2
    \right].
\end{equation}
Defining the usual broken-phase combinations
\begin{equation}\label{eq:ew_broken_basis}
  W_\mu^\pm \equiv \frac{W_\mu^1 \mp iW_\mu^2}{\sqrt{2}},\qquad
  Z_\mu \equiv \frac{g W_\mu^3 - g' B_\mu}{\sqrt{g^2+g'^2}},\qquad
  A_\mu \equiv \frac{g' W_\mu^3 + g B_\mu}{\sqrt{g^2+g'^2}},
\end{equation}
diagonalizes the mass matrix:
\begin{equation}\label{eq:ew_mass_matrix_diag}
  \mathcal{L}_{\mathrm{mass}}
  = \sqrt{-g}\left[
      m_W(\varphi)^2\,W_\mu^+ W^{-\mu}
      + \frac{1}{2}m_Z(\varphi)^2\,Z_\mu Z^\mu
    \right],
\end{equation}
with
\begin{equation}\label{eq:ew_mass_closure}
  m_W(\varphi) = \frac{g\,v(\varphi)}{2},\qquad
  m_Z(\varphi) = \frac{\sqrt{g^2+g'^2}}{2}\,v(\varphi),\qquad
  m_\gamma = 0.
\end{equation}
Thus, once the internal electroweak algebra is specified,
the entire $\varphi$-dependence of the broken-phase mass matrix
is fixed by the single bi-conformal scaling law
$v(\varphi)=v_0 e^{\varphi/2}$.

\textit{Electroweak mass relations.}---After spontaneous
symmetry breaking
$SU(2)_L \times U(1)_Y \to U(1)_{\mathrm{EM}}$, the
four electroweak gauge bosons have masses
\begin{align}
  \label{eq:mW}
  m_{W^{\pm}}(\varphi)
    &= \frac{g\,v_0}{2}\,e^{\varphi/2}
    \approx 80.4\;\mathrm{GeV}\times e^{\varphi/2},\\[4pt]
  \label{eq:mZ}
  m_Z(\varphi)
    &= \frac{m_W}{\cos\theta_W}
    \approx 91.2\;\mathrm{GeV}\times e^{\varphi/2},\\[4pt]
  \label{eq:mgamma}
  m_\gamma &= 0,
\end{align}
where $g \approx 0.653$ is the $SU(2)_L$ coupling and
$\theta_W$ is the Weinberg angle with
$\sin^2\theta_W \approx 0.231$.  The $\varphi$-dependence
enters solely through $v(\varphi) = v_0\,e^{\varphi/2}$
(Eq.~\ref{eq:v_phi}).

The dimensionless mass ratios are $\varphi$-independent:
\begin{equation}\label{eq:gauge_ratios}
  \frac{m_W}{m_Z} = \cos\theta_W,\qquad
  \frac{m_W}{m_H} = \frac{g}{2\sqrt{2\lambda}},\qquad
  \frac{m_W}{m_f} = \frac{g}{y_f\sqrt{2}}.
\end{equation}
These are identical to the Standard Model relations, with
no corrections from the scalar-metric coupling.  Every
electroweak precision observable measured at LEP, the
Tevatron, and the LHC is therefore reproduced exactly.

\textit{Local Lorentz invariance of the gauge sector.}---On
the bi-conformal background, with the MAIN convention
$ds^{2} = -e^{\varphi}c^{2}dt^{2} + e^{-\varphi}d\vec x^{\,2}$
(so $g^{00}=-e^{-\varphi}$, $g^{ii}=+e^{\varphi}$), the
mass-shell condition $g^{\mu\nu}k_\mu k_\nu + m^2 = 0$
gives the coordinate dispersion relation
\begin{equation}\label{eq:disp_gauge}
  -e^{-\varphi}\,\omega^2 + e^{\varphi}\,k^2 + m^2(\varphi)
  = 0.
\end{equation}
Converting to locally measured quantities
($\omega_{\mathrm{loc}} = \omega\,e^{-\varphi/2}$ from
$d\tau = e^{\varphi/2}dt$,
$k_{\mathrm{loc}} = k\,e^{\varphi/2}$ from
$d\ell = e^{-\varphi/2}dx$,
$m_{\mathrm{loc}} \equiv m(\varphi)$, the locally
measured mass):
\begin{equation}\label{eq:disp_local}
  \omega_{\mathrm{loc}}^2
  = k_{\mathrm{loc}}^2 + m_{\mathrm{loc}}^2.
\end{equation}
This is the standard special-relativistic dispersion
relation.  Although $m(\varphi)$ varies with position,
every local observer measures the \emph{same} numerical
mass value (e.g., $m_W = 80.4\;\mathrm{GeV}$) because
all reference scales---atomic masses, rulers, clocks---carry
the same $e^{\varphi/2}$ factor.  The bi-conformal metric
thus preserves local Lorentz invariance of the gauge sector
exactly.

\textit{Photon mass protection.}---The masslessness of the
photon is protected primarily by the unbroken
$U(1)_{\mathrm{EM}}$ gauge invariance, which forbids a
mass term at any loop order (Ward identity).  In ISPG,
the shift symmetry
$\varphi \to \varphi + \mathrm{const}$ of the
action~(\ref{eq:action}) provides a secondary safeguard:
because $\varphi$ enters the action only through
derivatives and through the metric, the scalar-metric
sector cannot generate a St\"uckelberg-type coupling
$(\partial_\mu\varphi)\,A^\mu$ that would give the photon
an effective mass while preserving diffeomorphism
invariance.  The two symmetries are not independent---the
shift symmetry ensures that the scalar sector does not
spoil the gauge protection---but their combined effect
makes $m_\gamma = 0$ robust against all radiative
corrections from both the Standard Model and the
gravitational sector.

\textit{Gauge couplings as input parameters.}---The
couplings $g$, $g'$, and the strong coupling $g_s$ are
dimensionless and therefore $\varphi$-independent by the
bi-conformal symmetry of dimensionless constants.  In the
present analysis, their values are empirical inputs:
\begin{equation}\label{eq:gauge_couplings}
  g \approx 0.653, \quad
  g' \approx 0.350, \quad
  g_s(m_Z) \approx 1.22.
\end{equation}
The derivation of these couplings from the medium's
internal structure---specifically, why $g/g'$ takes
its observed value---requires solving the nonlinear
resonance equations for the vector perturbation modes
(Sec.~\ref{sec:internal_dof}) and is the subject of
future work.  The ISPG framework ensures that whatever
values the couplings take, they are universal constants
(field-independent), in agreement with precision
electroweak data.

\textit{Status of the gauge completion.}---The derivation
achieved here is therefore precise but two-tiered.  What is
\emph{derived here} is the full broken-phase electroweak mass
matrix on the bi-conformal background, the local Lorentzian
dispersion relations, the field-independence of the gauge
couplings, and the radiative protection of the photon mass.
What is \emph{not yet derived from the medium itself} is the
microscopic origin of the internal gauge algebra
$SU(2)_L \times U(1)_Y$, the multiplicity of its generators,
and the numerical values of $g$ and $g'$.  In that sense the
remaining open step is group-theoretic, not propagational.

\textit{Minimal medium interpretation.}---Within the
resonance framework of Sec.~\ref{sec:internal_dof}, the
scalar-metric medium supplies the local spin-1 carrier through
its vector perturbation sector, while the broken Higgs sector
supplies the standard longitudinal polarization for each broken
generator.  The local polarization content is therefore the
usual one:
\begin{equation}\label{eq:gauge_polarization_count}
  N_{\mathrm{pol}}(W^\pm)=N_{\mathrm{pol}}(Z)=3,\qquad
  N_{\mathrm{pol}}(\gamma)=2.
\end{equation}
Thus the bi-conformal background already closes the kinematics
and mass structure of the electroweak gauge sector once the
internal gauge algebra is specified, while the deeper question
of why the medium realizes precisely that algebra remains
future work.

The massive gauge bosons ($W$, $Z$) are therefore interpreted
as localized spin-1 resonances whose masses are fixed by the
Higgs VEV through the standard Higgs mechanism, now
understood as the electroweak projection of bi-conformal
symmetry breaking (Sec.~\ref{sec:conformal_higgs}).

The massless photon is the unbroken $U(1)_{\mathrm{EM}}$
vector mode.  Its identification with the surviving
massless combination of the electroweak gauge fields is
standard; the ISPG framework adds the prediction that
the locally measured fine-structure constant
$\alpha = e^2/(4\pi\hbar c)$ is position-independent
(since $e = g\sin\theta_W$ and all factors are
dimensionless).  The cosmological variation of $\alpha$
discussed in Sec.~\ref{sec:alpha_var} arises from
vacuum-polarization corrections, not from a running of
the bare coupling.

\textit{The strong sector.}---The eight gluons of quantum
chromodynamics (QCD) mediate the strong interaction between
quarks.  On the bi-conformal background, the QCD Lagrangian
retains its standard form with $\varphi$-independent coupling
$g_s$; all QCD predictions---asymptotic freedom, confinement,
the hadron spectrum---are preserved.  The origin of the
$SU(3)_c$ gauge group from the medium's tensor and nonlinear
scalar excitation sectors, together with the derivation of
$g_s$ from the medium dynamics, is the subject of future work.

%-------------------------------------------------------------------
\subsection{Scope and deferred items of the electroweak treatment}
\label{sec:ew_scope}
%-------------------------------------------------------------------

The present section derives only the \emph{mass-scaling}
sector of the electroweak theory on the bi-conformal
background.  Four adjacent questions that would require
new machinery are explicitly set aside here and flagged
as topics for dedicated work:

\begin{itemize}
  \item \emph{Proton mass and the QCD contribution.}---Most
    of the proton mass comes from the QCD confinement
    scale $\Lambda_{\mathrm{QCD}}$ via dimensional
    transmutation, not from the Higgs VEV.  The
    $\varphi$-scaling of $\Lambda_{\mathrm{QCD}}$ is
    controlled by the running of $g_s$ and by the QCD
    trace anomaly; demonstrating that the Higgs-sector
    scaling $v(\varphi) = v_0\,e^{\varphi/2}$ extends
    consistently to the full hadronic mass spectrum is a
    separate derivation left to future work.
  \item \emph{Hierarchy problem.}---The ISPG framework
    does not address the electroweak hierarchy problem:
    the smallness of $v_0$ relative to $\Mpl$ is taken
    here as an empirical input.  Eq.~(\ref{eq:delta_mH})
    shows that the gravitational correction to the Higgs
    mass is Planck-suppressed and therefore does not
    generate a naturalness issue on its own, but it also
    does not \emph{resolve} the Standard-Model hierarchy
    problem.  A mechanism fixing $v_0$ in terms of the
    condensate parameter $\Phi_0$ is left open.
  \item \emph{Conformal / trace anomaly.}---The
    classical bi-conformal invariance of dimensionless
    couplings is broken at one loop by the standard
    gauge-theory trace anomaly.  Quantifying how the
    resulting anomalous $\varphi$-dependence feeds back
    into running couplings is an open item; it is
    expected to contribute only at loop level and not to
    affect the tree-level scaling used in
    Eq.~(\ref{eq:v_phi}).
  \item \emph{Electroweak phase transition.}---The
    statement in item~(iii) of
    Sec.~\ref{sec:ew_grav} that the phase transition
    occurs at the same \emph{locally measured}
    temperature relies on the classical bi-conformal
    symmetry of the dimensionless ratio
    $T_{\mathrm{EW}}/v$.  Running-coupling and
    finite-temperature anomaly corrections to that
    ratio are likewise deferred.
\end{itemize}

None of these items modifies the results derived above,
which depend only on the $\varphi$-scaling of $v$ and on
the $\varphi$-independence of dimensionless couplings.
They identify the direction in which the electroweak
chapter of ISPG must be extended.

%===================================================================
\section{Entanglement and gravitational decoherence}
\label{sec:entanglement}
%===================================================================

Quantum entanglement---the nonlocal correlation between
spatially separated particles---is one of the most precisely
tested features of quantum mechanics.  In ISPG, entanglement
acquires a concrete physical picture as correlated resonance
structure, and the theory predicts a specific, testable
mechanism for gravitational decoherence.

%-------------------------------------------------------------------
\subsection{Entanglement as correlated resonance structure}
\label{sec:entangle_interp}
%-------------------------------------------------------------------

When two particles are produced from a single quantum event
(e.g., parametric down-conversion or pair production), the
resulting state is a single, spatially extended resonance
pattern with two localized cores.

In the ISPG framework, the two-particle state is described by
a single resonance field configuration
$\delta\varphi_{12}(t,\mathbf{x})$ that cannot be factored
into independent single-particle configurations:
\begin{equation}\label{eq:entangled_state}
  \delta\varphi_{12} \neq \delta\varphi_1 \cdot \delta\varphi_2.
\end{equation}
The cores propagate apart, but the correlations encoded in
the phase structure of the resonance field are preserved as
long as the field evolves unitarily
(i.e., without decoherence).

\textit{Compatibility with Bell's theorem.}---ISPG
respects signal locality: both the background scalar
$\varphi$ and its resonance excitations $\delta\varphi$
obey causal wave equations derived from the
action~\eqref{eq:action}, with characteristic cones
bounded by the speed of light on the
bi-conformal metric, so no perturbation propagates
faster than $c$.  (The massless background $\varphi$
propagates at $c$; massive resonance excitations
propagate with group velocity $\leq c$ set by the
standard relativistic dispersion.)  ISPG therefore
cannot be used to signal superluminally.  The
entangled correlations arise because the two cores were
\emph{never independent}: they share a common resonance
structure established at creation.  The ``nonlocality''
is not in the dynamics but in the initial
conditions---the correlated state was created locally
and subsequently separated.

This position is analogous to the pilot-wave (de~Broglie--Bohm)
interpretation, where nonlocal correlations arise from a
globally defined wave function, and to quantum field theory,
where entanglement is encoded in the Fock-space structure
rather than in local hidden variables.

%-------------------------------------------------------------------
\subsection{Gravitational decoherence rate}
\label{sec:decoherence}
%-------------------------------------------------------------------

The ISPG framework predicts a specific decoherence mechanism
absent in standard quantum mechanics: when entangled particles
are placed at different gravitational potentials
(different values of $\varphi$), the bi-conformal rescaling
introduces a differential phase evolution that degrades the
entangled state.

Consider two entangled particles $A$ and $B$ at positions with
scalar-field values $\varphi_A$ and $\varphi_B$, respectively.
Their local frequencies are (Eq.~\ref{eq:omega_local}):
\begin{equation}
  \omega_A = \omega_0\,e^{-\varphi_A/2}, \qquad
  \omega_B = \omega_0\,e^{-\varphi_B/2},
\end{equation}
where $\omega_0$ is the coordinate frequency at production.

The differential phase accumulated over coordinate time $t$ is
\begin{equation}\label{eq:diff_phase}
  \Delta\Phi(t) = (\omega_A - \omega_B)\,t
  = \omega_0\,t\,
    \left(e^{-\varphi_A/2} - e^{-\varphi_B/2}\right).
\end{equation}
For $|\Delta\varphi| \equiv |\varphi_A - \varphi_B| \ll 1$:
\begin{equation}\label{eq:diff_phase_weak}
  \Delta\Phi(t) \approx \frac{\omega_0\,t}{2}\,\Delta\varphi.
\end{equation}

When the differential phase reaches $\Delta\Phi \sim \pi$,
the entangled state has evolved to an orthogonal (distinguishable)
configuration, and the quantum correlations are lost.  The
decoherence time is therefore
\begin{equation}\label{eq:t_dec}
  t_{\mathrm{dec}} \sim \frac{2\pi}{\omega_0\,|\Delta\varphi|}
  = \frac{h}{m_0 c^{2}\,|\Delta\varphi|},
\end{equation}
and the decoherence rate is
\begin{equation}\label{eq:Gamma_dec}
  \boxed{
  \Gamma_{\mathrm{dec}}
  = \frac{m_0 c^{2}\,|\Delta\varphi|}{h}
  = \frac{m_0 c^{2}}{h}\,\frac{2|\Delta\Phi_N|}{c^{2}}
  = \frac{2m_0\,|\Delta\Phi_N|}{h},
  }
\end{equation}
where $\Delta\Phi_N = -GM\Delta r/r^2$ is the
Newtonian potential difference and the factor of $2$
follows from $\varphi = 2\Phi_N/c^2$.

\textit{Dephasing versus decoherence.}---Eq.~(\ref{eq:Gamma_dec})
describes a differential phase evolution, which in isolation
would be reversible \emph{dephasing}: if the two particles
were brought back to the same potential, the relative phase
could be undone.  In ISPG, however, the process becomes
genuine (irreversible) \emph{decoherence} for two reasons:

\begin{enumerate}
  \item \textit{Resonant-tail emission.}---Each oscillon
    continuously emits resonant tails
    (Sec.~\ref{sec:entropy_mech}) that carry away
    which-path information into the scalar-field vacuum.
    In different $\varphi$-backgrounds, the tail spectra
    differ: the emitted radiation from particle $A$ at
    $\varphi_{A}$ is distinguishable (in principle) from
    that of particle $B$ at $\varphi_{B}$.  When these
    tail modes are traced out, the reduced density matrix
    of the two-particle system acquires off-diagonal
    suppression.
  \item \textit{Stochastic fluctuations.}---The scalar
    field $\varphi$ undergoes vacuum fluctuations
    $\langle\delta\varphi^{2}\rangle \neq 0$.  These
    fluctuations randomly modulate the differential phase
    $\Delta\Phi(t)$, converting the deterministic
    dephasing into a stochastic process.  For Gaussian
    fluctuations the ensemble average of the phase
    factor evaluates to
    $\langle e^{i\Delta\Phi(t)}\rangle
    = \exp\!\left[-\tfrac12\langle\Delta\Phi(t)^{2}\rangle
    \right]$.  In the Markovian limit, in which the
    correlation time $\tau_{c}$ of $\varphi$-fluctuations
    is much shorter than the observation time
    ($\tau_{c}\ll t$), this reduces to the exponential
    form
    $\langle e^{i\Delta\Phi(t)}\rangle
    \simeq e^{-\Gamma_{\mathrm{dec}}\,t}$,
    and Eq.~(\ref{eq:Gamma_dec}) gives the rate.  When
    the Markovian condition is not satisfied, the
    exponential form is replaced by the Gaussian form
    $\exp[-\tfrac12\sigma^{2}t^{2}]$ with the same
    leading coefficient.
\end{enumerate}

The rate~(\ref{eq:Gamma_dec}) thus represents the
\emph{leading-order} decoherence rate in the Markovian
regime; the two mechanisms above ensure irreversibility
while fixing the leading coefficient.  A detailed
derivation of the correlation time $\tau_{c}$ and the
fluctuation spectrum of $\delta\varphi$, which is needed
to confirm that the Markovian limit applies to the
experimental regimes considered below, is deferred to
dedicated work.  At this leading
post-Newtonian order, the rate coincides with the
gravitational dephasing predicted by GR through the
redshift--COW correspondence
$\varphi \simeq 2\Phi_N/c^2$; both depend on the
orientation of the pair relative to the local
gravitational gradient.  A genuine discriminator between
ISPG and GR for this observable therefore lies beyond the
1PN level, in the 2PN response of the bi-conformal factor
(Appendix~4 of the main
theory~\cite{Lotishvili2026}).  The present section
quantifies the leading-order rate with the understanding
that it reproduces the GR baseline at that order.

\textit{Numerical estimates.}---For relativistic particles
such as photons, the rate formula \eqref{eq:Gamma_dec}
admits a formal evaluation but is of limited practical
use.  A photon of energy $E = 1$~eV in a gradient with
$|\Delta\varphi| = 2g\,\Delta h/c^{2} \approx
2.2\times 10^{-15}$ (corresponding to $\Delta h = 10$~m in
Earth's field, $g = 9.8\;\mathrm{m/s^2}$) yields the formal
rate
\begin{equation}
  \Gamma_{\mathrm{dec}}
  \sim \frac{E\,|\Delta\varphi|}{h}
  \approx 0.5\;\mathrm{Hz}.
\end{equation}
However, the photon spends only
$t_{\mathrm{flight}} \sim \Delta h/c \sim 33$~ns in the
gradient region, so the accumulated coherent phase
contribution from the bi-conformal factor is
$\Delta\Phi_{\mathrm{geom}} =
E\,|\Delta\varphi|\,t_{\mathrm{flight}}/\hbar
\sim 10^{-24}$~rad---far below any foreseeable
polarimetric or interferometric sensitivity.  The
decoherence rate \eqref{eq:Gamma_dec} is therefore
physically meaningful only for particles confined in the
gradient region on timescales
$t \gtrsim \Gamma_{\mathrm{dec}}^{-1}$; in practice this
restricts testable cases to nonrelativistic massive
particles (treated below).

For massive particles the rate scales with the rest
energy $m_0 c^2$.  Taking a laboratory-scale
superposition with a height difference
$\Delta h \sim 1$~m in Earth's field gives
$\Delta\Phi_N/c^{2}\sim 10^{-16}$.  For an
optomechanical mass $m\sim 10^{-14}$~kg,
\begin{equation}
  \Gamma_{\mathrm{dec}}
  = \frac{m c^{2}\cdot 2|\Delta\Phi_N|/c^{2}}{h}
  \sim \frac{9\times 10^{2}\;\mathrm{J}\cdot
    2\times 10^{-16}}{6.6\times 10^{-34}\;\mathrm{J\cdot s}}
  \sim 3\times 10^{20}\;\mathrm{Hz},
\end{equation}
i.e.\ decoherence is effectively instantaneous on any
accessible timescale.  For nanospheres
($m\sim 10^{-20}$~kg, $\sim 10^{7}$~amu) at
$\Delta\Phi_N/c^{2}\sim 10^{-16}$ the rate drops to
$\sim 3\times 10^{14}$~Hz---still far beyond any
plausible coherence time, but the scaling
$\Gamma\propto m$ means that mesoscopic masses in the
$10^{-22}$--$10^{-23}$~kg range (proposed in
MAQRO~\cite{Kaltenbaek2012} and STE-QUEST) begin to
enter the kHz--MHz regime where experimental
coherence times are already achievable.  This is the
parameter window in which the rate formula becomes a
falsifiable prediction rather than a consistency
check.

%-------------------------------------------------------------------
\subsection{Comparison with Penrose--Di\'osi models}
\label{sec:penrose_comparison}
%-------------------------------------------------------------------

The prediction~(\ref{eq:Gamma_dec}) is structurally similar
to the gravitational decoherence rates proposed by
Penrose~\cite{Penrose1996} and
Di\'osi~\cite{Diosi1987}:
\begin{equation}
  \Gamma_{\mathrm{P}} \sim \frac{E_G}{\hbar},
  \qquad
  \Gamma_{\mathrm{D}} \sim
    \frac{G m^{2}}{\hbar\,\Delta x},
\end{equation}
where $E_G$ is a gravitational self-energy and $\Delta x$
is a spatial superposition size.  The key differences are:

\begin{enumerate}
  \item \textbf{Origin.}  Penrose and Di\'osi postulate
    gravitational decoherence as a fundamental modification
    of quantum mechanics.  In ISPG, it is a \emph{consequence}
    of the bi-conformal structure: it arises from the
    $\varphi$-dependent rescaling of quantum frequencies,
    without modifying the Schr\"odinger equation.
  \item \textbf{Mechanism.}  In Penrose's model, the
    decoherence comes from the ``ill-definedness'' of time
    in a superposition of different geometries.  In ISPG,
    it comes from the concrete, calculable differential
    phase evolution in an inhomogeneous $\varphi$-field.
  \item \textbf{Testability.}  ISPG predicts a specific
    \emph{directional} dependence: the decoherence rate
    depends on the orientation of the entangled pair
    relative to the gravitational gradient $\nabla\varphi$.
    This anisotropy is absent in the Penrose--Di\'osi models
    and constitutes a unique experimental signature.
\end{enumerate}

%-------------------------------------------------------------------
\subsection{Quantum tunneling in an inhomogeneous field}
\label{sec:tunneling}
%-------------------------------------------------------------------

The resonance picture also modifies quantum tunneling.
A particle's resonant tail extends beyond the classically
forbidden region (the potential barrier).  If the medium
on the far side of the barrier has favorable conditions
(e.g., a local minimum in the effective potential),
the tail ``anchors'' there and the resonance core
transitions through the barrier.

In an inhomogeneous $\varphi$-field, the tunneling
probability acquires a directional dependence.  The
WKB tunneling amplitude through a barrier of height $V_0$
and width $d$ is modified to
\begin{equation}\label{eq:tunnel_aniso}
  T \propto \exp\!\left[
    -\frac{2d}{\hbar}\sqrt{2m_0\,e^{\varphi/2}\,(V_0 - E)}
  \right],
\end{equation}
where $\varphi = \varphi(\mathbf{x})$ varies across the
barrier.  If $\varphi$ decreases in the tunneling direction
(toward a mass), $e^{\varphi/2}$ decreases, the effective
barrier is lowered, and tunneling is \emph{enhanced}.
In the opposite direction, tunneling is suppressed.

\begin{equation}\label{eq:tunnel_prediction}
  \boxed{
  \frac{T_{\downarrow}}{T_{\uparrow}}
  = \exp\!\left[
    \frac{d\,\Delta\varphi}{\hbar}\,
    \sqrt{2m_0(V_0 - E)}
  \right] > 1,
  }
\end{equation}
where $T_{\downarrow}$ and $T_{\uparrow}$ denote tunneling
rates toward and away from the gravitational source,
respectively.

This \textbf{anisotropic tunneling} is a distinctive
ISPG prediction with no analog in standard quantum
mechanics, where tunneling is isotropic in a uniform
barrier.  The relevant quantity in
Eq.~(\ref{eq:tunnel_prediction}) is the change of
$\varphi$ \emph{across the barrier width} $d$, not
across the host system's size.  For nuclear barriers
the barrier width is of order a femtometre, and the
gravitational gradient varies over macroscopic
distances, so the across-barrier shift is
$\Delta\varphi_{\mathrm{barrier}}
\sim (d/L)\,\Delta\varphi_{\mathrm{system}}$ where $L$
is the gradient scale length.  For a white-dwarf
interior with $\Delta\varphi_{\mathrm{system}}\sim
10^{-4}$ across a radius $L\sim 10^{6}$~m, this gives
$\Delta\varphi_{\mathrm{barrier}}
\sim 10^{-25}$ over a fm-scale barrier---negligible
for individual nuclear reactions.  A detectable signal
would require either very long barriers (e.g.\
macroscopic tunneling experiments in trapped cold
atoms) or a cumulative effect integrated over many
events; quantitative feasibility study of these
channels is left for dedicated work.  The conservative
statement retained here is that the anisotropy is
\emph{formally non-zero and directionally signed}; its
amplitude in any concrete experimental context has to
be estimated with the correct barrier-scale
$\Delta\varphi_{\mathrm{barrier}}$, not with the
system-scale value.

%===================================================================
\section{Topological structure: spin, antimatter, and exclusion}
\label{sec:topology}
%===================================================================

The resonance interpretation of particles
(Sec.~\ref{sec:excitations}) requires a mechanism
to generate half-integer spin, distinguish matter from
antimatter, and enforce the Pauli exclusion principle.
In this section we show that all three properties
emerge from the topological structure of vortex
excitations in the scalar-metric medium.

%-------------------------------------------------------------------
\subsection{Spin as vortex winding number}
\label{sec:spin}
%-------------------------------------------------------------------

In a fluid, a vortex carries angular momentum determined by
its circulation $\oint \mathbf{v}\cdot d\mathbf{l} = n\,\kappa$,
where $n$ is the (integer) winding number and $\kappa$ is the
circulation quantum.  The angular momentum is a topological
invariant: it cannot be changed by smooth deformations of the
flow and can only be removed by annihilating the vortex with
an anti-vortex of opposite winding.

In the scalar-metric medium, the vector perturbation
$\delta g_{0i}$ (the frame-dragging or gravitomagnetic sector)
satisfies the linearized field equation obtained by perturbing
the ISPG field equations (derived from the action~\eqref{eq:action}) around the background
$g_{\mu\nu}^{(0)} = \mathrm{diag}(-e^{\varphi_0}, e^{-\varphi_0}, e^{-\varphi_0}, e^{-\varphi_0})$.
For stationary perturbations ($\partial_t = 0$) with only
spatial components $\delta g_{0i}$, the linearized equation is
\begin{equation}\label{eq:vector_pert}
  \nabla^2 (e^{-\varphi_0} \delta g_{0i}) - \partial_i \partial_j (e^{-\varphi_0} \delta g_{0j}) = 0,
\end{equation}
where $\nabla^2$ is the flat-space Laplacian.  This is the
equation for a vector field in a medium with position-dependent
``permeability'' $e^{-\varphi_0}$.  The nonlinear rotational
response of this medium in the galactic regime is developed
in the companion MOND paper~\cite{Lotishvili2026MOND}; here
we focus on the topological content relevant to spin.

\textit{Vortex solution.}---In cylindrical coordinates
$(r, \phi, z)$, we seek axisymmetric solutions with only an
azimuthal component.  Setting $\delta g_{0r} = \delta g_{0z} = 0$
and $\delta g_{0\phi} = A(r) e^{in\phi}$ with integer $n$, the
equation~\eqref{eq:vector_pert} reduces to a Bessel equation
for the radial profile $A(r)$:
\begin{equation}\label{eq:bessel_vortex}
  \frac{1}{r}\frac{d}{dr}\left(r \frac{dA}{dr}\right) - \frac{n^2}{r^2}A = 0.
\end{equation}
The regular solution at $r = 0$ is $A(r) \propto r^{|n|}$; the
solution that decays at infinity requires a cutoff or core
structure (analogous to the oscillon core in Sec.~\ref{sec:resonances}).
The phase factor $e^{in\phi}$ implies that a $2\pi$ traversal
around the vortex accumulates phase $2\pi n$, defining the
\textbf{winding number} $n \in \mathbb{Z}$.

\textit{Angular momentum calculation.}---The angular momentum
of a stationary spacetime perturbation is given by the Komar
integral over a surface at infinity:
\begin{equation}\label{eq:komar_J}
  J = \frac{c^{2}}{16\pi G} \oint_{\infty} \nabla^\mu \xi^\nu \, dS_{\mu\nu},
\end{equation}
where $\xi^\mu = \partial_\phi$ is the azimuthal Killing vector.
The scalar vortex with winding $n$ sources an axially symmetric
frame-dragging field $\delta g_{0\phi} = \kappa/(2\pi r)$ at
large $r$, where $\kappa$ is the circulation quantum.
Evaluating the covariant derivatives on the bi-conformal
background and performing the azimuthal integration, we obtain
\begin{equation}\label{eq:J_vortex_derived}
  J = \frac{c^{2}\,\kappa}{16\pi G}\, e^{\varphi_{0}/2}.
\end{equation}

\textit{Derivation of the circulation quantum from
topological quantization.}---The circulation $\kappa$ is
\emph{not} a free parameter; it is fixed by the
single-valuedness of the scalar perturbation.

\paragraph{What carries the phase.}
The background scalar $\varphi$ of MAIN is real, so a
$U(1)$ winding cannot be imposed on $\varphi$ itself.
The phase structure used here lives on the
oscillon sector: a localized time-oscillating source of
the form $\Phi(t,\mathbf x) = \Phi_0(r)\cos[\Omega t + n\phi]$
introduced in Sec.~\ref{sec:resonances} naturally
combines into the complex envelope
$\Psi \equiv \Phi_0(r)\,e^{i(\Omega t + n\phi)}$, whose
argument is what winds $n$ times around the core.  The
symbol $\delta\varphi$ in Eq.~(\ref{eq:two_vortex}) and
below is used as a shorthand for this complex oscillon
envelope, not for a naive complexification of the real
background field.  With this identification, the
following derivation is the standard
Onsager--Feynman argument for a $U(1)$-phase condensate
applied to the oscillon sector.

Consider a vortex configuration in which the complex
envelope
$\delta\varphi = |\delta\varphi(r)|\,e^{in\phi}$ winds
$n$ times around the core.  Single-valuedness requires the
phase to return to its initial value after a $2\pi$ circuit:
\begin{equation}\label{eq:single_valued}
  \oint d\phi\;\frac{\partial}{\partial\phi}
    (\mathrm{arg}\,\delta\varphi)
  = 2\pi n, \qquad n \in \mathbb{Z}.
\end{equation}
The circulation of the effective velocity field
$\mathbf{v}_{\mathrm{eff}} = (\hbar/m_0)\,\nabla(\mathrm{arg}\,
\delta\varphi)$ is therefore quantized:
\begin{equation}\label{eq:circ_quantized}
  \oint \mathbf{v}_{\mathrm{eff}}\cdot d\mathbf{l}
  = \frac{2\pi n\,\hbar}{m_{0}}.
\end{equation}
This is the Onsager--Feynman quantization condition,
identical to that of superfluid vortices.

Now, the gravitomagnetic sector $\delta g_{0\phi}$ is
sourced by the angular momentum current of the scalar
vortex.  In the bi-conformal metric, the frame-dragging
field is related to the circulation by
$\delta g_{0\phi} = \kappa/(2\pi r)$, where
\begin{equation}\label{eq:kappa_quantized}
  \kappa
  = \frac{4G\,m_{0}}{c^{2}\,e^{\varphi_{0}/2}}
    \cdot \frac{2\pi n\,\hbar}{m_{0}}
  = \frac{8\pi G\,n\,\hbar}{c^{2}\,e^{\varphi_{0}/2}}.
\end{equation}
The factor $4Gm_{0}/(c^{2}e^{\varphi_{0}/2})$ is the
gravitomagnetic coupling of a source of mass $\meff$
in the bi-conformal geometry; the remaining factor is
the quantized circulation from
Eq.~(\ref{eq:circ_quantized}).  Crucially, no assumption
about the value of $J$ has been made---only single-valuedness
of $\delta\varphi$.

\textit{Result.}---Substituting Eq.~(\ref{eq:kappa_quantized})
into Eq.~(\ref{eq:J_vortex_derived}):
\begin{equation}\label{eq:J_vortex}
  \boxed{J = n\,\frac{\hbar}{2}.}
\end{equation}

\textit{Origin of the $1/2$ factor.}---The factor of
$1/2$ is not put in by hand, but it is also not a
specifically bi-conformal effect: it is the
\emph{standard} numerical ratio between the Komar
prefactor $c^{2}/(16\pi G)$ and the gravitomagnetic
coupling of a quantized circulation, namely
$8\pi G/(c^{2})$ per unit quantum of circulation.
Inserting Eq.~(\ref{eq:kappa_quantized}) into
Eq.~(\ref{eq:J_vortex_derived}), the bi-conformal
factors cancel identically:
$e^{\varphi_{0}/2}$ from the Komar evaluation on the
bi-conformal background multiplies $e^{-\varphi_{0}/2}$
inside $\kappa$, leaving the geometry-independent
combinatorial factor
$(c^{2}/16\pi G)\cdot(8\pi G/c^{2}) = 1/2$.
The same calculation performed on a Minkowski
background with the same quantized circulation yields
the same result; a scalar vortex with an integer
winding $n\in\mathbb Z$ and Onsager--Feynman
circulation $2\pi n\hbar/m_{0}$ therefore
\emph{universally} sources an angular momentum
$J = n\hbar/2$, independently of whether the metric is
Minkowski or bi-conformal.  What the bi-conformal
structure contributes is a non-trivial
\emph{$\varphi_{0}$-cancellation}, i.e.\ the result that
$J$ depends only on the topological winding and not on
the local value of $\varphi_{0}$; it does not by itself
explain the $1/2$.  The half-integer spectrum of $J$
therefore reflects the quantized-circulation structure
of the vortex, not a metric-specific halving.

\textit{Fermions.}---Excitations with $n = \pm 1$ carry
angular momentum $J = \hbar/2$, as derived in
Eq.~(\ref{eq:J_vortex}), and are identified with
fermions (electrons, quarks).

\paragraph{Relation to the $4\pi$ spinor rotation.}
The magnitude $J = \hbar/2$ is the defining angular-momentum
eigenvalue of a spin-$1/2$ state, but it does not by
itself capture the $4\pi$ rotation symmetry observed in
neutron interferometry~\cite{Rauch1975}.  The two
statements are logically distinct.  Going once around
the vortex of winding $n = 1$ accumulates a
\emph{vortex-phase} $2\pi n = 2\pi$ (topologically
trivial), not $\pi$; it is not an active spatial rotation
of the state.  The $4\pi$ symmetry is instead captured by
the Berry-phase argument of Sec.~\ref{sec:pauli}: a
\emph{half-exchange} of two $n=1$ vortices accumulates a
phase of $\pi$, which is the hallmark of a spin-$1/2$
state under $2\pi$ spatial rotation via the
spin-statistics connection.  Thus the vortex construction
supplies (i) the correct half-integer angular-momentum
eigenvalue and (ii), through the Berry-exchange
calculation below, the correct $4\pi$ rotation symmetry;
these two ingredients together reproduce the
spin-$1/2$ behaviour of Rauch~\cite{Rauch1975}.

\textit{Bosons.}---Excitations with $n = 0$ (no vortex
structure) are bosons.  Photons, for example, are
transverse tensor perturbations $h_{ij}^{TT}$ with no
rotational topology; their spin-1 character comes from
the two independent transverse polarizations, not from
a vortex winding.

\paragraph{Scope of the vortex picture for spin.}
The vortex construction above rigorously yields the
angular-momentum quantum $J = \hbar/2$ for an
excitation with scalar winding $n=\pm 1$, and is
therefore a derivation of spin for the fermion sector.
For bosonic species the angular-momentum content comes
from a different part of the perturbation spectrum:
photons from the transverse tensor polarizations of
$\delta g_{ij}^{TT}$, and massive gauge bosons from the
vector sector $\delta g_{0i}$ combined with the Higgs
Goldstone mode.  The pair-exchange phases computed in
Sec.~\ref{sec:pauli} for $n = 0, \pm 1, \pm 2$ are
therefore offered as an \emph{illustrative} pattern that
the spin-statistics correspondence is compatible with
this topological labeling; they should not be read as a
derivation that $W^{\pm}$ and $Z$ are literally $n=\pm 2$
scalar vortices.  A unified topological origin covering
fermions and all bosonic species on the same footing
would require a simultaneous treatment of the scalar,
vector, and tensor perturbation sectors and is left for
future work.

%-------------------------------------------------------------------
\subsection{Antimatter as opposite winding}
\label{sec:antimatter}
%-------------------------------------------------------------------

If a particle is a vortex with winding number $n = +1$,
its antiparticle is the vortex with $n = -1$: the same
spatial structure, rotating in the opposite direction.
The total winding number is conserved under smooth
deformations, so particle--antiparticle pair creation
and annihilation correspond to the appearance and
disappearance of $n = +1$, $n = -1$ vortex pairs with
$n_{\mathrm{total}} = 0$.

The vortex picture naturally suggests further
identifications---charge quantization from the
integer-valuedness of $n$, the electromagnetic coupling
from the sign of the winding, and the detailed energy
budget of annihilation---but deriving these connections
from the bi-conformal field equations requires a dedicated
analysis that lies beyond the scope of the present work.

\paragraph{Open structural items (deferred).}
Three further structural items are explicitly set aside
here and flagged as topics for a dedicated continuation
rather than as claims of the present work.  They do not
affect the result $J = n\hbar/2$ derived above but they
circumscribe what the vortex construction does and does
not yet cover:
\begin{itemize}
  \item \emph{Scalar winding $\leftrightarrow$
    gravitomagnetic slot bridge.}---Sec.~\ref{sec:spin}
    uses both the scalar phase $\delta\varphi$ (whose
    winding is quantized by Onsager--Feynman) and the
    gravitomagnetic field $\delta g_{0\phi}$ (which
    carries the Komar angular momentum).  The
    identification between them is made at the level of
    a \emph{source relation}
    ($\delta g_{0\phi}\propto\kappa/r$), not derived
    diagram-by-diagram from the full coupled field
    equations that connect scalar and vector
    perturbations.  A field-equation-level derivation
    of the scalar-phase $\to$ off-diagonal-metric
    channel is left for future work.
  \item \emph{CPT and charge content of the
    antimatter sector.}---The identification
    ``$n\to -n$ is the antiparticle'' correctly captures
    the sign-flip of angular momentum and (schematically)
    of any $U(1)$-type charge tied to the winding.  A
    full CPT argument built on top of this picture, and
    a derivation of the numerical value of the
    electromagnetic charge from the bi-conformal field
    equations, are not attempted here and remain open
    items for a dedicated paper.
  \item \emph{Non-Abelian structure.}---An integer
    winding $n\in\mathbb{Z}$ labels only the abelian
    $U(1)$-type topological charge.  The non-Abelian
    structures of the Standard Model
    ($SU(2)_{L}$, $SU(3)_{c}$) require additional
    internal degrees of freedom beyond the single
    scalar-plus-metric sector considered here and are
    not addressed by this vortex construction.  Whether
    such non-Abelian structure can be realized as
    multi-component vortex configurations on a richer
    medium is a separate question left for future work.
\end{itemize}

\paragraph{Notational convention for $\kappa$.}
Two circulation quanta appear in Sec.~\ref{sec:spin}:
$\kappa_{\mathrm{sf}} \equiv 2\pi n\hbar/m_{0}$
(the Onsager--Feynman superfluid circulation, units
$[\mathrm m^{2}/\mathrm s]$, Eq.~\ref{eq:circ_quantized})
and $\kappa_{\mathrm{gm}}\equiv\oint\delta g_{0\phi}\,d\ell$
(the quantized gravitomagnetic circulation, units
$[\mathrm m]$, Eq.~\ref{eq:kappa_quantized}).  They are
linked by the gravitomagnetic coupling
$\kappa_{\mathrm{gm}} = [4Gm_{0}/(c^{2}e^{\varphi_{0}/2})]\,
\kappa_{\mathrm{sf}}$.  Where the context is
unambiguous, $\kappa$ alone refers to
$\kappa_{\mathrm{gm}}$; elsewhere the subscripted forms
should be used to avoid dimensional confusion.

%-------------------------------------------------------------------
\subsection{Pauli exclusion from vortex topology}
\label{sec:pauli}
%-------------------------------------------------------------------

The Pauli exclusion principle---that two identical fermions
cannot occupy the same quantum state---receives a geometric
explanation in the vortex picture.

Two vortices of the same winding number $n$ cannot occupy
the same spatial location, because the resulting configuration
would have winding number $2n$ at that point---a topologically
distinct object (not two copies of the original particle).
This is the same reason two fluid vortices of the same sign
cannot merge: they orbit each other but maintain their
individual cores.

\paragraph{Status of this hydrodynamic remark.}
The merging/orbiting observation is a
\emph{heuristic} consistency check, not the Pauli
principle.  The quantum Pauli statement is a property of
the many-body wavefunction (antisymmetry under particle
exchange), not of field configurations overlapping in
space.  The rigorous argument below is the Berry-phase
exchange calculation; the hydrodynamic analogy only
motivates why two identical vortices behave as
non-coincident objects.

We now show that the exchange phase is $\pi$ for
$n = \pm 1$ vortices.

\textit{Berry phase calculation.}---Consider two
identical vortices at positions $\mathbf{x}_{1}$ and
$\mathbf{x}_{2}$ in the two-dimensional plane transverse
to the vortex axes.  The two-vortex scalar field
configuration is
\begin{equation}\label{eq:two_vortex}
  \delta\varphi(\mathbf{x})
  = f(|\mathbf{x} - \mathbf{x}_{1}|)\,
    e^{in\,\mathrm{arg}(\mathbf{x}-\mathbf{x}_{1})}
  + f(|\mathbf{x} - \mathbf{x}_{2}|)\,
    e^{in\,\mathrm{arg}(\mathbf{x}-\mathbf{x}_{2})},
\end{equation}
where $f(r)$ is the radial profile.  Exchange of the two
vortices corresponds to adiabatically transporting
vortex~1 along a closed path that encircles vortex~2
and returns to the starting position of vortex~2
(a half-circuit in the relative coordinate
$\mathbf{r} = \mathbf{x}_{1} - \mathbf{x}_{2}$).

During this transport, the phase of the scalar field at
any fixed observation point $\mathbf{x}$ changes by the
angle subtended by the path as seen from $\mathbf{x}$.
The Berry connection for the adiabatic transport of the
$k$-th vortex is
\begin{equation}\label{eq:berry_connection}
  \mathcal{A}_{k}^{i}
  = -i\!\int\! d^{2}x\;\delta\varphi^{*}
    \frac{\partial\,\delta\varphi}
         {\partial x_{k}^{i}}.
\end{equation}
For the half-exchange path
$\mathbf{x}_{1}(\theta) = \frac{R}{2}(\cos\theta,
\sin\theta)$, $\theta: 0 \to \pi$, the Berry phase is
\begin{equation}\label{eq:berry_integral}
  \gamma_{\mathrm{exchange}}
  = \oint \mathcal{A}_{1}^{i}\,dx_{1}^{i}
  = n \cdot \pi.
\end{equation}
The integral evaluates to $n\pi$ because the transported
vortex of winding $n$ accumulates a phase $n\,\Delta\phi$
for each $\Delta\phi$ of angular displacement around
the other vortex, and the half-circuit contributes
$\Delta\phi = \pi$.  This is the Aharonov--Bohm
analog for vortices: each vortex acts as a
topological ``flux tube'' of strength $2\pi n$ for
every other vortex.

For $n = \pm 1$ (fermions):
\begin{equation}\label{eq:exchange_fermion}
  \gamma_{\mathrm{exchange}} = \pm\pi
  \quad\Rightarrow\quad
  e^{i\gamma} = -1,
\end{equation}
and therefore
\begin{equation}\label{eq:antisymmetry}
  \boxed{
  \psi(\mathbf{x}_2, \mathbf{x}_1)
  = e^{i\gamma}\,\psi(\mathbf{x}_1, \mathbf{x}_2)
  = -\psi(\mathbf{x}_1, \mathbf{x}_2).
  }
\end{equation}

For $n = 0$ (bosons, no vortex):
$\gamma_{\mathrm{exchange}} = 0$ and the wavefunction is
symmetric---no exclusion.

For $n = \pm 2$ (even winding):
$\gamma_{\mathrm{exchange}} = \pm 2\pi \equiv 0$,
symmetric.  This is an illustrative consistency check
that the Berry-phase formula $\gamma = n\pi$ respects
Bose statistics for any even winding; it is \emph{not}
a claim that the physical $W^{\pm}, Z$ bosons are
realized as $n = \pm 2$ scalar vortices (cf.\ the scope
paragraph in Sec.~\ref{sec:spin}).

This topological derivation of the spin-statistics
connection complements the standard proof via relativistic
quantum field theory~\cite{Pauli1940,Streater1964} and
provides a physical, geometric picture of \emph{why}
fermions are antisymmetric and bosons are symmetric:
the exchange phase is determined by the winding number,
which is the same quantum number that determines the spin
(Sec.~\ref{sec:spin}).

\textit{Dimensional caveat.}---The Berry-phase calculation
above is carried out in the two-dimensional plane
transverse to the vortex axis.  In $2{+}1$ dimensions,
the braid group permits arbitrary exchange phases
(anyons); in $3{+}1$ dimensions, only $\gamma = 0$ or
$\pi$ survive because exchange paths are contractible
around the vortex line.  The result $\gamma = n\pi$ is
therefore consistent with the $3{+}1$-dimensional
requirement, but the passage from the 2D calculation to
the full 3D vortex-line topology (linking, knotting,
Hopf invariants) requires a dedicated treatment and is
left for future work.

In $3{+}1$ dimensions, the rigorous origin of fermion
antisymmetry is the standard QFT spin--statistics
theorem, which follows from Lorentz invariance,
causality and positivity of the
energy~\cite{Pauli1940,Streater1964}.  The vortex
construction above is therefore \emph{complementary}
to, not a replacement for, that theorem: it supplies a
geometric picture of why the exchange phase takes the
value it does, on top of the symmetry argument that
guarantees the result.

%===================================================================
\section{Gravitational birefringence}\label{sec:birefringence}
%===================================================================

At the classical level, ISPG preserves local Lorentz invariance
exactly: the locally measured speed of light is $c$ for all
polarizations~\cite{Lotishvili2026}, consistent with the
classical light-propagation analysis in Appendix~5 of the
main theory (PPN $\gamma=1$, no birefringence at the
tree level).  At the quantum level, however, vacuum
polarization by virtual particle--antiparticle
pairs introduces a polarization-dependent correction to the
photon propagation speed in the presence of a gravitational
gradient.  The one-loop effect computed in this section
lies beyond the tree-level regime of Appendix~5 and
therefore does not contradict its conclusions.  The
precise comparison to the curvature-coupled GR$+$QED
effect is deferred to \S\ref{sec:biref_distinction}.

%-------------------------------------------------------------------
\subsection{Vacuum polarization in an inhomogeneous field}
\label{sec:vac_pol}
%-------------------------------------------------------------------

In QED, a photon propagating through an external electromagnetic
field interacts with virtual $e^+e^-$ pairs (the Euler--Heisenberg
effective Lagrangian~\cite{EulerHeisenberg1936}).  In ISPG, the
analogous process occurs in a gravitational gradient:
a photon propagating through a region with $\nabla\varphi \neq 0$
interacts with virtual pairs whose effective mass
$\meff = m_0\,e^{\varphi/2}$ varies across the pair's
spatial extent.  This effect---\textbf{gravitational
birefringence}---arises from the dependence of particle
masses on the scalar field and is structurally distinct
from (though numerically comparable to) the one-loop
GR$+$QED effect of
Drummond--Hathrell~\cite{DrummondHathrell1980}; see
\S\ref{sec:biref_distinction} for a detailed comparison.

The one-loop vacuum polarization tensor in the presence of
a slowly varying $\varphi$-field acquires a
$\nabla\varphi$-dependent correction.  For a photon with
four-momentum $k^\mu$ and polarization $\epsilon^\mu$,
the leading correction to the dispersion relation is
\begin{equation}\label{eq:dispersion_biref}
  \omega^{2} = |\mathbf{k}|^{2}\,c^{2}\left[1
    + \frac{\alpha_{\mathrm{EM}}}{45\pi}\,
      \frac{(\nabla\varphi)^{2}}{m_e^{2}c^{2}/\hbar^{2}}\,
      \eta(\hat{\epsilon}, \hat{\nabla}\varphi)\right],
\end{equation}
where $\alpha_{\mathrm{EM}} = 1/137$ is the fine-structure
constant and $\eta$ is a polarization-dependent factor:
\begin{equation}\label{eq:eta_factor}
  \eta =
  \begin{cases}
    7 & \text{if } \hat{\epsilon} \parallel \hat{\nabla}\varphi
        \;\;(\parallel\text{-mode}),\\
    4 & \text{if } \hat{\epsilon} \perp \hat{\nabla}\varphi
        \;\;(\perp\text{-mode}).
  \end{cases}
\end{equation}

\textit{Note on coefficients and status.}---The numerical
values $\eta = 7$ and $\eta = 4$ are the Euler--Heisenberg
coefficients for QED in an external electromagnetic
field~\cite{EulerHeisenberg1936}; they are \emph{adopted here
by analogy}, not derived from ISPG first principles.  The full
ISPG coefficients require an explicit one-loop calculation of
the $\varphi$-$e^+e^-$ triangle diagrams with the
conformal-coupling vertex $m_0\,\partial\varphi\cdot\bar\psi\psi$
specified by the action of the main
theory~\cite{Lotishvili2026}; this derivation is a
self-contained technical task that is the topic of a separate
work.  The present section should be read as a
\emph{parametric} estimate: only the scaling
$(\nabla\varphi)^2/m_e^2$ and the existence of
polarization-dependent propagation are stated here as ISPG
consequences; the precise numerical coefficients remain an
open derivation target.

\textit{Scope of the two-case parametrization.}---Equation~\eqref{eq:eta_factor}
records only the two \emph{extreme} polarization geometries
(electric field exactly parallel or exactly perpendicular to
$\hat{\nabla}\varphi$).  For a generic propagation direction
$\hat k$ and generic linear polarization, the one-loop effective
action produces an angular dependence that interpolates between
the two extremes and also involves the relative orientation
between $\hat k$ and $\hat{\nabla}\varphi$; a general mode carries
both components and experiences a weighted average.  The $(7,4)$
numbers should therefore be read as bounding parallel/perpendicular
values for the leading-order estimate, not as a complete angular
template.

The two polarization modes propagate at slightly
different phase velocities:
\begin{equation}\label{eq:delta_c}
  \frac{\Delta c}{c}
  = \frac{c_{\parallel} - c_{\perp}}{c}
  = \frac{3\,\alpha_{\mathrm{EM}}}{90\pi}\,
    \frac{(\nabla\varphi)^{2}}{m_e^{2}c^{2}/\hbar^{2}}.
\end{equation}

%-------------------------------------------------------------------
\subsection{Order-of-magnitude estimate}
\label{sec:biref_estimate}
%-------------------------------------------------------------------

For a neutron star with $M = 1.4\,M_\odot$ and radius
$R = 10$~km:
\begin{equation}
  |\nabla\varphi|
  \sim \frac{\rs}{R^{2}}
  = \frac{2GM}{c^{2}R^{2}}
  \sim 4 \times 10^{-5}\;\mathrm{m^{-1}},
\end{equation}
and
\begin{equation}
  \frac{(\nabla\varphi)^{2}}
    {(m_e c/\hbar)^{2}}
  \sim \frac{(4 \times 10^{-5})^{2}}
    {(2.6 \times 10^{12})^{2}}
  \sim 2 \times 10^{-34}.
\end{equation}
The birefringence is therefore
\begin{equation}\label{eq:biref_NS}
  \frac{\Delta c}{c}
  \sim \frac{\alpha_{\mathrm{EM}}}{30\pi}
    \cdot 2 \times 10^{-34}
  \sim 10^{-38}.
\end{equation}

For a photon traversing a path of length $L \sim R = 10$~km
near the neutron star surface, the accumulated differential
phase between the two polarizations is
\begin{equation}\label{eq:biref_phase}
  \Delta\Phi_{\mathrm{biref}}
  = \frac{2\pi\nu}{c}\,\frac{\Delta c}{c}\,L
  \sim 2\pi\nu \cdot 10^{-38} \cdot 10^{4}\;\mathrm{m}/c.
\end{equation}
For X-ray photons ($\nu \sim 10^{18}$~Hz):
$\Delta\Phi_{\mathrm{biref}} \sim 10^{-24}$~rad---utterly
negligible with any foreseeable technology.

For the much stronger fields near black holes
($\nabla\varphi \sim 1/\rs$, $L \sim \rs$):
\begin{equation}\label{eq:biref_BH}
  \Delta\Phi_{\mathrm{biref}}
  \sim \frac{\alpha_{\mathrm{EM}}}{30\pi}\,
    \frac{(m_e c/\hbar)^{-2}}{\rs^{2}}\cdot
    \frac{2\pi\nu\,\rs}{c}
  \sim \frac{\alpha_{\mathrm{EM}}\,\nu}
    {30\pi\,m_e^{2}c^{3}/\hbar^{2}\,\rs}.
\end{equation}
For a stellar-mass black hole ($\rs \sim 10$~km) and
gamma-ray photons ($\nu \sim 10^{20}$~Hz):
$\Delta\Phi_{\mathrm{biref}} \sim 10^{-21}$~rad---far
below any conceivable measurement sensitivity.

\textit{Regime caveat.}---For the BH estimate,
$|\nabla\varphi|\,\rs \sim 1$, so the weak-field linearization
underlying Eq.~\eqref{eq:dispersion_biref} ceases to be rigorously
controlled: higher-order terms in $\nabla\varphi/(m_e c/\hbar)^{-1}$
and curvature corrections to the photon self-energy are no longer
formally suppressed.  In this regime the number
$\Delta\Phi_{\mathrm{biref}}\sim 10^{-21}$ should therefore be read
as a parametric extrapolation of the weak-field formula to its
boundary of validity rather than a rigorous prediction; the
non-perturbative strong-field calculation (one-loop self-energy
on a full Schwarzschild-like background with $\nabla\varphi$ of
order $1/\rs$) is deferred.  The quantitative smallness of the
effect---already negligible in the controlled weak-field regime
around neutron stars, Eq.~\eqref{eq:biref_NS}---does not depend
on this caveat.

%-------------------------------------------------------------------
\subsection{Distinction from other birefringence sources}
\label{sec:biref_distinction}
%-------------------------------------------------------------------

Gravitational birefringence in ISPG differs from other
proposed sources:
\begin{itemize}
  \item \textbf{QED vacuum birefringence}
    (Euler--Heisenberg~\cite{EulerHeisenberg1936}):
    requires a strong \emph{electromagnetic} field;
    scales as $B^{2}/B_{\mathrm{cr}}^{2}$.
    ISPG birefringence requires a \emph{gravitational}
    gradient; it scales as $(\nabla\varphi)^{2}$.
    The two effects are orthogonal and additive.
  \item \textbf{Lorentz-violation birefringence}
    (Standard Model Extension~\cite{Colladay1998}):
    produces birefringence that is direction-dependent
    relative to a preferred frame.  ISPG birefringence
    is direction-dependent relative to the local
    gravitational gradient $\nabla\varphi$ and vanishes
    in the absence of gravity.
  \item \textbf{GR\,$+$\,QED prediction:} At the classical
    level, GR predicts zero vacuum birefringence (the
    equivalence principle forbids polarization-dependent
    propagation of classical electromagnetic waves in
    vacuum).  At one-loop in curved spacetime, however,
    Drummond \& Hathrell~\cite{DrummondHathrell1980} showed
    that integrating out the electron produces an effective
    action with curvature-coupled terms of the schematic
    form $R_{\mu\nu\rho\sigma}F^{\mu\nu}F^{\rho\sigma}$,
    yielding polarization-dependent phase velocities that
    scale as $R_{\mu\nu\rho\sigma}/(m_e^{2}c^{2}/\hbar^{2})$.
    The ISPG effect \eqref{eq:dispersion_biref} is
    structurally distinct: it scales as $(\nabla\varphi)^{2}$
    rather than as a curvature component, and does not
    coincide with the Drummond--Hathrell tensorial pattern.
    A detection that correlates with $|\nabla\varphi|$ in
    regions where the Drummond--Hathrell curvature pattern
    is suppressed (or vice versa) would therefore
    discriminate ISPG from the GR$+$QED baseline; a generic
    ``detection of gravitational birefringence'' by itself
    does not falsify GR once one-loop QED corrections are
    included.
\end{itemize}

%===================================================================
\section{Summary of predictions}\label{sec:predictions}
%===================================================================

Table~\ref{tab:predictions} collects the quantitative
predictions of the quantum ISPG framework developed in
this paper.  Each prediction specifies the observable,
the ISPG prediction, the standard (GR/QFT) expectation,
and the experimental channel most sensitive to the effect.
The gravity-core source closure derived in
Sec.~\ref{sec:resonances} is structural rather than observational:
after projection away from the neutral collective coordinates, the
localized source sector has a positive reduced gap and a coercive
energy bound.  For that reason it is part of the logical foundation of
the table below, not a separate row inside it.

\begin{table*}[ht]
\centering
\caption{Falsifiable predictions and prediction-status ledger
  of the quantum ISPG framework.  ``GR/SM'' denotes the
  prediction of general relativity and/or the Standard Model.}
\label{tab:predictions}
\resizebox{\textwidth}{!}{%
\begin{tabular}{llll}
\toprule
\textbf{Observable} &
\textbf{ISPG prediction} &
\textbf{GR/SM} &
\textbf{Experiment} \\
\midrule
Dark-energy EoS $w(z)$ &
$w < -1$, source-driven phantom; DESI DR2 CPL comparison mixed &
$w = -1$ exactly &
DESI, Euclid, Rubin \\
\addlinespace
Fine-structure constant $\Delta\alpha/\alpha$ &
$\Delta\alpha/\alpha \sim 10^{-7}$--$10^{-6}$ at $z \sim 1$--$3$
(future calculation; ESPRESSO ppm now, $10^{-7}$ ELT--ANDES target) &
$0$ &
ESPRESSO, ELT-ANDES \\
\addlinespace
Gravitational decoherence rate &
$\Gamma = 2m|\Delta\Phi_N|/h$ &
$0$ (no mechanism) &
MAQRO, STE-QUEST \\
\addlinespace
Decoherence anisotropy &
Directional ($\propto \nabla\varphi$) &
Absent &
Space quantum links \\
\addlinespace
Anisotropic tunneling &
$T_\downarrow/T_\uparrow > 1$ &
$T_\downarrow = T_\uparrow$ &
White dwarf seismology \\
\addlinespace
Gravitational birefringence &
$\Delta c/c \sim 10^{-38}$ (NS scale), $\propto (\nabla\varphi)^2/m_e^2$
(ISPG channel; coefficient future one-loop calculation) &
Drummond--Hathrell: $\propto R_{\mu\nu\rho\sigma}/m_e^2$ &
Beyond current sensitivity \\
\addlinespace
Koide relation $Q = 2/3$ &
$Q = 0.66673$ (E1D/rank-1 numerical support; $3$D-to-$1$D bridge open)$^{\dagger}$ &
Unexplained coincidence &
$0.009\%$ accuracy \\
\addlinespace
Proton lifetime &
Finite (very long) &
Stable or GUT decay &
Hyper-Kamiokande \\
\addlinespace
Pair-production threshold shift &
$\delta E_{\mathrm{thr}}/E \propto \varphi$ &
No shift &
Black hole X-ray spectra \\
\addlinespace
2PN Shapiro delay (ISPG$-$GR difference) &
$\Delta B = \pi/4$ (closed; standalone absolute coefficients convention-dependent) &
$0$ (by construction) &
SKA pulsar--BH \\
\addlinespace
ISCO orbital frequency &
$f_{\mathrm{ISCO}} = 0.931\,f_{\mathrm{ISCO}}^{\mathrm{GR}}$ &
GR reference &
LIGO--Virgo--KAGRA, ET \\
\addlinespace
Black-hole shadow size &
$b_c = e\,\rs$ ($+4.6\%$ vs.\ GR) &
$b_c = \tfrac{3\sqrt{3}}{2}\,\rs$ &
EHT, ngEHT \\
\addlinespace
Equivalence principle &
Derived: $m_{\mathrm{inert}} \equiv m_{\mathrm{grav}}$ &
Postulated &
MICROSCOPE, STEP \\
\addlinespace
Scalar dipole radiation &
$\dot{E}_{\mathrm{dipole}} \simeq 0$ (vacuum exact, leading-order in matter; matter-filled exact-cancellation open per MAIN/16) &
Theory-dependent &
Pulsar timing \\
\bottomrule
\end{tabular}}%

\smallskip
{\footnotesize ${}^{\dagger}$Requires numerical verification
  that the self-consistent oscillon solution yields
  $x = \alpha/\lambda_0 \approx 66.9$
  (Sec.~\ref{sec:koide}).}
\end{table*}

The predictions span a wide range of energy scales
and experimental techniques.  The most immediately
testable are:
\begin{enumerate}
  \item \textbf{Dark-energy evolution} ($w \neq -1$):
    DESI year-3 and Euclid data, expected by 2027.
  \item \textbf{Fine-structure constant variation}:
    ESPRESSO at the VLT is already operating at ppm
    precision; the $10^{-7}$ regime targeted here is the
    future ELT--ANDES-level test.
  \item \textbf{Gravitational decoherence}:
    proposed space missions (MAQRO) could detect the
    predicted rate within the next decade.
\end{enumerate}

%===================================================================
\section{Remark on the ontology of the medium}\label{sec:ontology}
%===================================================================

The formalism developed in the preceding sections treats space
as a dynamical, pressure-carrying medium whose properties are
encoded in a single scalar field~$\varphi$.  This ontology
naturally raises a question: if space is a medium, what defines
the boundaries of its existence?  The remarks below do not
constitute a formal derivation; they record a physical intuition
that is consistent with the formalism and may guide its future
development.

If space is a medium in which pressure varies---decreasing in
gravitational wells, oscillating in particles, propagating as
waves---then this variability logically requires two asymptotic
bounds, not one: a maximum and a minimum.  If the substance of
space is water, it must be poured into a vessel; the vessel is
the region where the substance is absent.  One pole---\emph{substance}---represents
the extreme state in which pressure, mass--energy capacity, the
operational clock rate of resonant structures, and spatial scale
are asymptotically maximal.  The other pole---\emph{void}---represents
the null limit: zero pressure, zero mass--energy content, no
operational clock rate, and zero size.  The void is not an
independent entity; it is the logical complement of
substance---the necessary consequence of something existing.
Physical reality is a bounded oscillation between these two
poles; a vibration that can never exceed either limit.

This motivates the ontological name \emph{Archion}.  An Archion
is the primary quantum of existence: the minimal active unit of
the scalar-metric medium whose measurable life is precisely this
bounded oscillation between the substance pole and the void pole.
The older word ``oscillon'' names only the dynamical profile---a
localized, time-periodic resonance.  \emph{Archion} names the
ontological role: the first active unit by which existence becomes
measurable.

This picture becomes intuitively transparent through a single
analogy: coloured soluble crystals in water.  The water is
space---the background medium.  The crystal is an oscillon---a
localized, concentrated resonant form of matter in the Archion
field.  As the crystal dissolves it releases colour into the
surrounding water, just as an oscillon emits pressure-carrying
tails into the surrounding space.  The
spreading of colour through the water is the spreading of
the scalar profile through space; through the bi-conformal
metric this profile carries gravitational influence.  While
dissolution is underway the water is
mottled: intensely coloured near the crystals, nearly
transparent far from them.  This is the cosmic web---galaxy
clusters where matter is concentrated and the
$\varphi$-gradient is deep; filaments where colour migrates
from crystal to crystal; and voids---vast stretches of water
between crystals that the colour has barely reached.  Gravity,
in this picture, is not an additional colour-gradient force
on matter; it is the geodesic response to the metric whose
scalar profile is visualised by the colour distribution.  No
separate ``dark pigment'' needs to be added, just as no extra
dye is required to explain the
diffusion of colour in water: the dynamics of the medium suffice.

The analogy reveals one further essential property:
self-regulation.  A crystal dissolves from its surface: the
more colour it releases, the smaller it becomes; a smaller
crystal has less surface area in contact with the water and
therefore dissolves more slowly.  The process decelerates
itself.  An oscillon behaves identically: it vibrates in pressure
and emits pressure tails into the
surrounding medium.  These tails reduce the background
pressure around the oscillon.  The reduction forces the
oscillon to shrink in effective size and mass
($\meff = m_0\,e^{\varphi/2}$).  A smaller, lighter oscillon
interacts less with the medium---it emits fewer tails---and
consumes less pressure.  This self-regulating cycle---emission,
contraction, stabilisation---operates on every scale by the
same principle: from oscillons to stars, from galaxy clusters
to the universe as a whole.

When all crystals have fully dissolved, the water will be
uniformly tinted---an intermediate tone, neither the intensity
of the crystal nor the transparency of pure water.  This is
maximum entropy: the vanishing of gradients, the cessation of
process.  In the theory this is reflected by the relation
$a_0 = cH/(2\pi)$: as long as $H>0$, structure persists.
Yet the difference between maximum and minimum continuously
decreases while never reaching zero.  Complete equalisation
would imply the disappearance of both poles---a logical
contradiction, since the arena in which equalisation occurs
would itself cease to exist.  The process is asymptotic: the
universe perpetually tends toward equilibrium but never
attains it.

We do not elevate this picture to a postulate.  We record it
as a heuristic that is consistent with every result derived
in this work and that may serve as a guide for its future
development.

%===================================================================
\section{Conclusion}\label{sec:conclusion}
%===================================================================

We have developed the quantum extension of bi-conformal
scalar-tensor gravity (ISPG), building on the classical
theory established in~\cite{Lotishvili2026}.  The central
thesis---that elementary particles are localized resonant
excitations of the scalar-metric medium---has been shown
to produce a coherent, falsifiable theoretical framework.
The central closed achievement of the present paper is the
quantum-mechanical derivation of the gravitational mechanism itself:
the oscillon source, the Bernoulli pressure deficit, the localized
zero-frequency source, the exterior $1/r$ field, Newtonian recovery,
the source-following collective response, and the coercive separation
of neutral and non-neutral source deformations.  The broader sectors
below show how far the same medium picture extends into cosmology,
particle structure, and quantum phenomena.  They are part of the same
research program, but they do not all have the same logical status in
the present work: where a sector still requires a dedicated numerical
or structural closure, that is stated explicitly and does not weaken
the core gravity derivation.  Within that scope, the principal
results and extensions are:

\begin{enumerate}
  \item \textbf{Particles as spacetime resonances}
    (Sec.~\ref{sec:excitations}).  The scalar field
    perturbation equation on the bi-conformal background
    admits quasi-localized oscillatory solutions whose
    energy reproduces the effective mass formula
    $\meff = m_0\,e^{\varphi/2}$.  Local undetectability
    extends to all quantum observables.

  \item \textbf{Bernoulli mechanism for gravity}
    (Sec.~\ref{sec:bernoulli_identity}).  The kinetic energy of
    particle resonances drains the static pressure of the
    medium via an exact scalar-field analog of Bernoulli's
    equation.  This provides the first mechanical explanation
    for \emph{why} matter curves spacetime.  After projection
    onto the reduced source space, a positive spectral gap and
    coercive source-energy estimate keep non-neutral
    deformations perturbatively controlled, so the exterior
    field follows the same collective coordinate that moves the
    oscillon core.

  \item \textbf{Inertia and the equivalence principle}
    (Sec.~\ref{sec:inertia}).  The retarded potential of
    an accelerated oscillon creates a dipolar pressure
    asymmetry whose integral yields
    $\mathbf{F} = -\meff\,\mathbf{a}$---Newton's second
    law derived, not postulated for the oscillon model.  The
    identity
    $m_{\mathrm{inertial}} \equiv m_{\mathrm{gravitational}}$
    is traced to the same scalar profile and coupling that set
    the Bernoulli energy bookkeeping; exact non-perturbative
    compact-body sensitivity is a separate question.  The
    hysteresis action $S_\chi$ is an effective template
    motivated by the quasi-normal mode structure of the
    oscillon (Sec.~\ref{sec:chi_derivation}): the microscopic
    tracking stiffness $\lambda_{\rm micro} = \omega_0^2/c^2$
    is set by the fundamental confining frequency, and the
    Klein--Gordon structure emerges in the limit
    $\Gamma \ll \omega_0$ appropriate for stable particles.

  \item \textbf{Self-regulated echo-background energy}
    (Sec.~\ref{sec:entropy}).  The irreversible spreading
    of resonant tails gives a positive outgoing energy flux,
    but the resulting
    pressure drop weakens the sources through a negative
    feedback loop.  On the perturbative tail-power branch,
    the equilibrium echo-background (dark-energy) density is
    estimated at
    $\rho_\Lambda \sim H_0^2 \Mpl^2$; the zero-point
    decoupling and non-perturbative normalization required for
    a full cosmological-constant solution remain open.

  \item \textbf{Evolving dark energy}
    (Sec.~\ref{sec:dark_energy}).  The dynamical
    nature of the echo background predicts
    $w(z) = -1 - \delta w_0\,(1+z)^3/E(z)$ with $w < -1$ at all
    epochs (source-driven phantom, no ghost instability),
    with a magnitude conditional on the tail-power
    normalization.  DESI DR2 gives a mixed CPL comparison:
    the sign of $w_a$ is aligned, while the sign of
    $w_0+1$ is opposite, so this remains a near-term test
    rather than a confirmation.

  \item \textbf{Earliest-epoch dynamics without inflation}
    (Sec.~\ref{sec:early_epoch}).  Planck-mass oscillons
    nucleate uniformly from the quiescent medium via
    $O(4)$-symmetric instantons with $S_E \sim \order{1}$;
    their overlapping tails thermalize to produce the
    hot Big Bang.  The initial singularity is avoided by
    the effective-source NEC violation (canonical $\tau_{\mu\nu}$
    satisfies the energy conditions; the $\kappa=-1/2$ coupling
    flips the sign of the source-side contraction, cf.\ the
    energy-conditions section of the main theory~\cite{Lotishvili2026})
    and the hyperbolic WDW equation.
    The horizon and flatness problems are resolved by
    the universality of the nucleation mechanism and
    the geometric inheritance of $k = 0$, without
    invoking an inflaton field.  The homogeneous
    background closes analytically, while the detailed
    primordial perturbation spectrum remains a separate
    numerical problem.

  \item \textbf{Mass hierarchy from 3D resonances}
    (Sec.~\ref{sec:mass_spectrum}).  The charged-lepton
    mass hierarchy is organized by the MASS/Mathieu ladder
    and the E1D/rank-1 cavity structure.  In $N$-space the
    Koide identity holds numerically to $0.024\%$, while the
    rank-1 perturbation argument selects
    $x = 33+24\sqrt{2} \approx 66.9$ as the bridge value.
    A first-principles derivation of this bridge value from
    the self-consistent $3$D oscillon eigenvalue system
    (Sec.~\ref{sec:mass_determination}) is the residual
    numerical task, with gravitational self-energy corrections
    of order $(m/\Mpl)^2 \sim 10^{-45}$.

  \item \textbf{Higgs--gravity unification and electroweak
    spectrum} (Sec.~\ref{sec:higgs}).  The Higgs VEV inherits
    the bi-conformal factor
    $v(\varphi) = v_0\,e^{\varphi/2}$, identifying
    electroweak symmetry breaking as the electroweak
    projection of conformal symmetry breaking.  Given the
    experimentally fixed electroweak gauge algebra, the
    broken-phase gauge-boson mass matrix
    (Sec.~\ref{sec:gauge_spectrum}) reproduces all Standard
    Model mass ratios exactly, with photon masslessness
    protected by $U(1)_{\mathrm{EM}}$ and reinforced by
    the shift symmetry.

  \item \textbf{Scalar sensitivity (structural argument)}
    (Sec.~\ref{sec:scalar_sensitivity}).  The shift-symmetry
    Ward identity supports $s = 1/2$ at all loop orders for
    vacuum compact objects and at leading order for
    matter-filled bodies ($s_{\mathrm{excess}} = 0$ in those
    controlled limits; matter-coupled exact-cancellation
    remains an open task per MAIN/16), arguing for
    suppressed scalar dipole radiation and a perturbatively
    controlled strong-equivalence statement in the scalar
    sector, with the
    exact matter-filled Nordtvedt statement left open.

  \item \textbf{Gravitational decoherence}
    (Sec.~\ref{sec:entanglement}).  Entangled particles
    in inhomogeneous $\varphi$-fields undergo differential
    phase evolution, producing a calculable decoherence
    rate $\Gamma = 2m|\Delta\Phi_N|/h$ with a unique
    directional signature.

  \item \textbf{Topological particle classification}
    (Sec.~\ref{sec:topology}).  Spin, antimatter, and
    the Pauli exclusion principle all emerge from the
    winding-number structure of vortex excitations.

  \item \textbf{Gravitational birefringence}
    (Sec.~\ref{sec:birefringence}).  Vacuum polarization
    in a gravitational gradient identifies a
    polarization-dependent ISPG channel with leading scaling
    $\Delta c/c \propto(\nabla\varphi)^2/m_e^2$ and neutron-star
    scale $\Delta c/c \sim 10^{-38}$---far below current
    measurement sensitivity.  This channel is distinct from
    the GR$+$QED Drummond--Hathrell baseline at the
    operator/scaling level; its numerical coefficient is the
    dedicated one-loop calculation target.
\end{enumerate}

The framework achieves these results without introducing
extra dimensions, supersymmetry, or new fundamental
fields.  It requires only the scalar-metric medium
already present in the classical theory, extended to
the quantum domain through standard techniques (canonical
quantization, perturbation theory, one-loop corrections).
The quantum effective action
(Sec.~\ref{sec:effective_action}) is radiatively stable
due to the shift symmetry $\varphi \to \varphi + \mathrm{const}$,
one-loop controlled with corrections suppressed by
$(\mathrm{curvature}/\Mpl^2)^2$ in the weak-field
regime, and
controlled at all loop orders by the EFT expansion
parameter $(E/\Mpl)^2 \sim 10^{-76}$.

The self-consistent mass determination problem has been
formulated as a dimensionless nonlinear eigenvalue system
(Sec.~\ref{sec:mass_determination}).  The confining potential
is a geometric property of the nonlinear wave equation,
independent of $G$.  The cavity geometry supplies the
structural mode labels and the nonlinear eigenvalue problem;
the quantitative charged-lepton mass ratios are currently
carried by the MASS/Mathieu ladder of the companion MASS
paper~\cite{Lotishvili2026Mass}, with the formal 3D-to-1D
reduction remaining the open numerical task.  Gravitational
self-energy corrections of order $(m/\Mpl)^2 \sim 10^{-45}$
are negligible.  The lattice
formulation reduces the classical eigenvalue problem to a
finite-dimensional nonlinear system, while the quantum
lattice (exploiting the ultralocal kinetic structure) enables
Monte Carlo evaluation of the constrained path integral.

The direction for future work includes: (i)~lattice simulation of
the coupled scalar-metric system (to verify oscillon
stability and decay properties, and to compute the primordial
perturbation spectrum from the Planck-epoch nucleation dynamics of
Sec.~\ref{sec:early_epoch}); (ii)~numerical solution of the
self-consistent oscillon eigenvalue problem to determine
whether the perturbation ratio $x = \alpha/\lambda_0$ attains
the value $33 + 24\sqrt{2} \approx 66.9$ required by the Koide
relation---initial numerical results (Sec.~\ref{sec:koide})
reproduce $Q = 2/3$ to $10^{-4}$ and the
$\tau/e$ mass ratio to $0.37\%$, while the first-order
Hartree correction, under the selected
semiclassical/first-order prescription, recovers the Koide
angle $\theta_0 = 2/9$ to $0.07\%$, with the joint
$(Q,\theta_0)$ closure beyond first order remaining open;
(iii)~detailed confrontation with
forthcoming data from DESI, Euclid, ESPRESSO, and
space-based quantum experiments.

The electroweak gauge boson spectrum ($W^{\pm}$, $Z$,
$\gamma$) has been derived on the bi-conformal background
(Sec.~\ref{sec:gauge_spectrum}) once the internal
electroweak gauge algebra is specified, confirming that all
Standard Model mass ratios and dispersion relations
are preserved locally, with photon masslessness protected
by $U(1)_{\mathrm{EM}}$ and reinforced by the ISPG shift
symmetry.  The derivation of the gauge couplings
themselves ($g$, $g'$, $g_s$) and the origin of the
$SU(3)_c$ colour group from the medium's internal
excitation structure remain open problems for future
work.

The scalar sensitivity $s = 1/2$ and the vanishing of the
excess sensitivity $s_{\mathrm{excess}} = 0$ have been
argued to persist at all loop orders for vacuum compact
objects (Sec.~\ref{sec:scalar_sensitivity}), via the
shift-symmetry Ward identity of the quantum effective action;
exact non-perturbative matter-filled-body cancellation remains
an open task per MAIN/16. This supports suppressed scalar
dipole radiation and a perturbatively controlled
strong equivalence statement in the scalar sector; the exact
matter-filled Nordtvedt statement belongs to the same open
sensitivity programme.

The hysteresis action
$S_\chi$~(Eq.~\ref{eq:S_chi_quantum}) is treated as a
microscopically motivated effective template
(Sec.~\ref{sec:chi_derivation}).  The microscopic tracking
stiffness $\lambda_{\rm micro} = \omega_0^2/c^2$ emerges
from the confining frequency of the oscillon's fundamental
quasi-normal mode, and the Klein--Gordon structure follows
from the damped oscillator equation in the limit
$\Gamma \ll \omega_0$ appropriate for stable particles.
The macroscopic collective relaxation scale and the closure
normalization $\beta$ remain Route~B/coarse-graining tasks,
with rotation curves serving as the calibration and test
channel~\cite{Lotishvili2026}.

One item requires attention in future extensions:
\begin{enumerate}
  \item \textit{Independent numerical verification.}---The
    structural argument for quantum stability of the phantom
    sector within the pure-gravity constrained sector has been
    presented at all orders in the loop expansion
    (Sec.~\ref{sec:stability}), resting on three pillars:
    kinematic ghost freedom from the single-DOF constraint,
    structurally anomaly-free Abelian constraint algebra (full
    bracket-closure/anomaly proof deferred), and
    non-perturbative boundedness of the Euclidean action on the
    constrained background (matter-coupled extension via the
    EFT completion of Appendix~11).
    Independent numerical verification of these results---for
    example via lattice simulation of the constrained path
    integral---would provide a valuable cross-check.
\end{enumerate}

The picture that emerges is one of radical economy:
a single medium, vibrating at different frequencies
and in different topological configurations, generates
all the diversity of matter, the gravitational and
inertial phenomena, and the arrow of time itself.  Gravity is
not a separate interaction added to quantum mechanics;
it is the medium's response to its own internal
vibrations.

%===================================================================
\appendix
%===================================================================

\section{Technical Lemmas}
\label{app:lemmas}

This appendix collects mathematical proofs and technical
results that support the main text.

%-------------------------------------------------------------------
\subsection{Lemma: Local Lorentz invariance of wave propagation}
\label{app:lemma_lorentz}
%-------------------------------------------------------------------

\textbf{Lemma.}  For the bi-conformal metric $ds^2 = -e^{\varphi}c^2dt^2 +
e^{-\varphi}(dx^2+dy^2+dz^2)$, the locally measured speed of scalar
field perturbations is exactly $c$, independent of position.

\textit{Proof.}  The scalar field equation $\Box\varphi = 0$ on the
static background $\varphi = \varphi_0(\mathbf{x})$ becomes (Eq.~\ref{eq:wave_rearranged}):
\begin{equation}
  e^{-\varphi_0}\ddot{\delta\varphi} = e^{\varphi_0}\nabla^2\delta\varphi.
\end{equation}
The coordinate speed is $v_{\mathrm{coord}} = e^{\varphi_0}c$.  The
locally measured speed uses proper time $d\tau = e^{\varphi_0/2}dt$
and proper length $d\ell = e^{-\varphi_0/2}dr$:
\begin{equation}
  v_{\mathrm{local}} = \frac{d\ell}{d\tau}
  = \frac{e^{-\varphi_0/2}dr}{e^{\varphi_0/2}dt}
  = e^{-\varphi_0} \cdot e^{\varphi_0}c = c.
\end{equation}
Thus $v_{\mathrm{local}} = c$ for all $\varphi_0$. \hfill $\square$

%-------------------------------------------------------------------
\subsection{Lemma: Positivity of scalar-field energy density}
\label{app:lemma_positive}
%-------------------------------------------------------------------

\textbf{Lemma.}  For the ISPG phantom scalar field with static
configuration $\varphi_0(\mathbf{x})$, the energy density measured
by a static observer is positive-definite:
\begin{equation}
  \rho_\varphi = \frac{e^{\varphi_0}}{32\pi G}|\nabla\varphi_0|^2 \geq 0.
\end{equation}

\textit{Proof.}  The stress-energy tensor for the scalar field is
(Eq.~\ref{eq:tau_munu}):
\begin{equation}
  \tau_{\mu\nu} = \partial_\mu\varphi\,\partial_\nu\varphi
    - \frac{1}{2}g_{\mu\nu}(\partial\varphi)^2.
\end{equation}
For a static field, the kinetic invariant is purely spatial
(Eq.~\ref{eq:X_static}):
\begin{equation}
  (\partial\varphi)^2 = g^{ij}\partial_i\varphi_0\,\partial_j\varphi_0
  = e^{\varphi_0}|\nabla\varphi_0|^2 \geq 0.
\end{equation}
The static observer has $u^\mu = (e^{-\varphi_0/2}/c, 0, 0, 0)$
(equivalently $u^\mu = (e^{-\varphi_0/2}, 0, 0, 0)$ in
natural units $c = 1$), giving
(Eq.~\ref{eq:rho_phi}):
\begin{equation}
  \rho_\varphi = \frac{1}{16\pi G}\tau_{\mu\nu}u^\mu u^\nu
  = \frac{e^{\varphi_0}}{32\pi G}|\nabla\varphi_0|^2.
\end{equation}
Since $e^{\varphi_0} > 0$ and $|\nabla\varphi_0|^2 \geq 0$, we have
$\rho_\varphi \geq 0$. \hfill $\square$

%-------------------------------------------------------------------
\subsection{Lemma: WKB tunneling for quasi-bound states}
\label{app:lemma_wkb}
%-------------------------------------------------------------------

\textbf{Lemma.}  For the quasi-normal mode equation (Eq.~\ref{eq:quasi_normal})
with effective potential $V_{\mathrm{eff}}(r)$, the decay width is
approximately:
\begin{equation}
  \Gamma \sim \omega_R \exp\left[-2\int_{r_1}^{r_2}
    \sqrt{V_{\mathrm{eff}}(r) - \omega_R^2}\,dr\right].
\end{equation}

\textit{Proof.}  In the WKB approximation, the radial equation
$\phi'' + k^2(r)\phi = 0$ with $k^2(r) = \omega^2 - V_{\mathrm{eff}}(r)$
has turning points where $k^2(r) = 0$.  In the classically forbidden
region ($r_1 < r < r_2$, where $V_{\mathrm{eff}} > \omega^2$), the
solution decays exponentially: $\phi \sim \exp[-\int\sqrt{V-
\omega^2}\,dr]$.  The outgoing wave at infinity has amplitude
reduced by the barrier penetration factor $\exp[-\int_{r_1}^{r_2}
\sqrt{V-\omega^2}\,dr]$.  The decay rate (imaginary part of $\omega$)
is proportional to the squared amplitude, giving the result. \hfill $\square$

%-------------------------------------------------------------------
\subsection{Theorem: Rank-1 perturbation and Koide ratio}
\label{app:theorem_koide}
%-------------------------------------------------------------------

\textbf{Theorem (exact).}  Let $\hat{H} =
\lambda_{0}\mathbb{1} + \alpha|v\rangle\langle v|$ be a
symmetric rank-1 perturbation of a threefold degenerate
Hamiltonian, with $|v\rangle$ normalized.  Then
for any $\alpha > -\lambda_{0}$ (the only
constraint being $\lambda_{1} = \lambda_{0}+\alpha > 0$),
the Koide ratio is the one-parameter function
\begin{equation}
  Q(x) = \frac{3+x}{(\sqrt{1+x}+2)^2}, \qquad
  x \equiv \alpha/\lambda_0,
\end{equation}
which increases monotonically from $Q(0) = 1/3$ to
$Q(\infty) = 1$.  If the remaining twofold degeneracy is
lifted by a fractional amount $\delta$, then
$Q = Q(x) + \mathcal{O}(\delta^2)$.

\textit{Proof.}  The eigenvalues are
$\lambda_{1} = \lambda_{0}+\alpha$ and
$\lambda_{2} = \lambda_{3} = \lambda_{0}$.
Define $x \equiv \alpha/\lambda_{0}$ (which can be
arbitrarily large).  The Koide ratio is
\begin{equation}\label{eq:Q_exact}
  Q(x) = \frac{3 + x}
    {\bigl(\sqrt{1+x} + 2\bigr)^{2}}.
\end{equation}
We verify the two extreme limits:
\begin{itemize}
  \item $x \to 0$ (degenerate):
    $Q = 3/9 = 1/3$. (Note: $Q = 1/3$ for exact
    degeneracy, since all three $\sqrt{\lambda}$ are
    equal.)
  \item $x \to \infty$ (one mass dominates):
    $Q \to x/x = 1$.
\end{itemize}
The function $Q(x)$ interpolates monotonically from
$1/3$ to $1$.  It equals $2/3$ at a specific value of $x$.
Setting $Q = 2/3$:
\begin{equation}
  \frac{3+x}{(\sqrt{1+x}+2)^{2}} = \frac{2}{3}.
\end{equation}
Let $s = \sqrt{1+x}$.  Then $x = s^{2}-1$ and the
equation becomes
\begin{equation}
  \frac{s^{2}+2}{(s+2)^{2}} = \frac{2}{3}
  \quad\Rightarrow\quad
  3s^{2}+6 = 2s^{2}+8s+8
  \quad\Rightarrow\quad
  s^{2}-8s-2 = 0,
\end{equation}
with positive root $s = 4+3\sqrt{2}$, giving
$x = (4+3\sqrt{2})^{2}-1 = 33+24\sqrt{2} \approx 66.9$.
The corresponding mass ratios are
$\sqrt{\lambda_{1}/\lambda_{2}} = s = 4+3\sqrt{2}
\approx 8.24$.  For the actual leptons,
$\sqrt{m_{\tau}/m_{e}} \approx 59$ and
$\sqrt{m_{\mu}/m_{e}} \approx 14.4$, which means the
physical masses are \emph{not} a pure rank-1 perturbation
but require lifting the $\lambda_{2} = \lambda_{3}$
degeneracy.

\textit{Degeneracy-lifting regime.}---Now
consider a rank-1 perturbation plus a small
degeneracy-lifting term:
$\hat{H} = \lambda_{0}\mathbb{1}
+ \alpha|v\rangle\langle v|
+ \epsilon\,\mathrm{diag}(0, +1, -1)$
with $\epsilon \ll \alpha$.  The eigenvalues become
$\lambda_{1} = \lambda_{0}+\alpha$,
$\lambda_{2} = \lambda_{0}+\epsilon$,
$\lambda_{3} = \lambda_{0}-\epsilon$.  A direct
calculation gives
\begin{equation}\label{eq:Q_lifted}
  Q = \frac{2}{3}
  + \mathcal{O}\!\left(\frac{\epsilon^{2}}
    {\lambda_{0}\,\alpha}\right).
\end{equation}
The correction is quadratic in $\epsilon$ and suppressed by
$\epsilon^{2}/(\lambda_{0}\alpha)$, not by $\alpha/\lambda_{0}$.
This establishes the perturbative robustness of the rank-1
Koide constraint in the regime $\epsilon \ll \alpha$.
For the physical charged leptons, $m_{e} \ll m_{\mu}$ implies
$\epsilon \approx \lambda_{0}$, placing them outside the
perturbative regime of this bound.  The agreement of the
measured Koide value with $Q=2/3$ to $0.01\%$ is therefore
the empirical target that the rank-1/E1D consistency argument
is intended to explain, not an independent confirmation of
perturbative robustness in the physical hierarchy.
\hfill $\square$

%===================================================================
% BIBLIOGRAPHY
%===================================================================

\begin{thebibliography}{99}

\bibitem{Lotishvili2026}
  G.~Lotishvili,
  ``Bi-conformal scalar-tensor gravity: exact PPN equivalence with
  general relativity, emergent MOND phenomenology, and falsifiable
  predictions,''
  Zenodo (2026).

\bibitem{Lotishvili2026MOND}
  G.~Lotishvili,
  ``Deriving MOND from Bi-Conformal Gravity: Vortex Closure,
  the Interpolating Function, and the Dark-Energy Connection,''
  Zenodo (2026).

\bibitem{Lotishvili2026Mass}
  G.~Lotishvili,
  ``Mathieu-Harmonic Mass Spectrum of Elementary Particles
  from Induced Scalar-Pressure Gravity,''
  Zenodo (2026).

\bibitem{Penrose1965}
  R.~Penrose,
  ``Gravitational collapse and space-time singularities,''
  Phys.\ Rev.\ Lett.\ \textbf{14}, 57 (1965).

\bibitem{Hawking1970}
  S.~W.~Hawking and R.~Penrose,
  ``The singularities of gravitational collapse and cosmology,''
  Proc.\ R.\ Soc.\ Lond.\ A \textbf{314}, 529 (1970).

\bibitem{Weinberg1989}
  S.~Weinberg,
  ``The cosmological constant problem,''
  Rev.\ Mod.\ Phys.\ \textbf{61}, 1 (1989).

\bibitem{Martin2012}
  J.~Martin,
  ``Everything you always wanted to know about the cosmological
  constant problem (but were afraid to ask),''
  C.\ R.\ Phys.\ \textbf{13}, 566 (2012).

\bibitem{Penrose1996}
  R.~Penrose,
  ``On gravity's role in quantum state reduction,''
  Gen.\ Relativ.\ Gravit.\ \textbf{28}, 581 (1996).

\bibitem{Diosi1987}
  L.~Di\'osi,
  ``A universal master equation for the gravitational violation
  of quantum mechanics,''
  Phys.\ Lett.\ A \textbf{120}, 377 (1987).

\bibitem{Polchinski1998}
  J.~Polchinski,
  \textit{String Theory} (Cambridge University Press, 1998).

\bibitem{Rovelli2004}
  C.~Rovelli,
  \textit{Quantum Gravity} (Cambridge University Press, 2004).

\bibitem{Reuter2012}
  M.~Reuter and F.~Saueressig,
  ``Quantum Einstein gravity,''
  New J.\ Phys.\ \textbf{14}, 055022 (2012).

\bibitem{Thomson1867}
  W.~Thomson (Lord Kelvin),
  ``On vortex atoms,''
  Proc.\ R.\ Soc.\ Edinburgh \textbf{6}, 94 (1867).

\bibitem{Volovik2003}
  G.~E.~Volovik,
  \textit{The Universe in a Helium Droplet}
  (Oxford University Press, 2003).

\bibitem{Zloshchastiev2011}
  K.~G.~Zloshchastiev,
  ``Spontaneous symmetry breaking and mass generation as built-in
  phenomena in logarithmic nonlinear quantum theory,''
  Acta Phys.\ Polon.\ B \textbf{42}, 261 (2011).

\bibitem{Verlinde2011}
  E.~Verlinde,
  ``On the origin of gravity and the laws of Newton,''
  JHEP \textbf{04}, 029 (2011).

\bibitem{Jacobson1995}
  T.~Jacobson,
  ``Thermodynamics of spacetime: the Einstein equation of state,''
  Phys.\ Rev.\ Lett.\ \textbf{75}, 1260 (1995).

\bibitem{Bogolyubsky1976}
  I.~L.~Bogolyubsky and V.~G.~Makhankov,
  ``Lifetime of pulsating solitons in some classical models,''
  JETP Lett.\ \textbf{24}, 12 (1976).

\bibitem{Gleiser1994}
  M.~Gleiser,
  ``Pseudostable bubbles,''
  Phys.\ Rev.\ D \textbf{49}, 2978 (1994).

\bibitem{Copeland1995}
  E.~J.~Copeland, M.~Gleiser, and H.-R.~M\"uller,
  ``Oscillons: resonant configurations during bubble collapse,''
  Phys.\ Rev.\ D \textbf{52}, 1920 (1995).

\bibitem{ALICE2010}
  ALICE Collaboration,
  ``Elliptic flow of charged particles in Pb--Pb collisions
  at $\sqrt{s_{NN}} = 2.76$~TeV,''
  Phys.\ Rev.\ Lett.\ \textbf{105}, 252302 (2010).

\bibitem{Bardeen1980}
  J.~M.~Bardeen,
  ``Gauge-invariant cosmological perturbations,''
  Phys.\ Rev.\ D \textbf{22}, 1882 (1980).

\bibitem{KodamaSasaki1984}
  H.~Kodama and M.~Sasaki,
  ``Cosmological perturbation theory,''
  Prog.\ Theor.\ Phys.\ Suppl.\ \textbf{78}, 1 (1984).

\bibitem{DESI2024}
  DESI Collaboration,
  ``DESI 2024 VI: Cosmological constraints from the measurements
  of baryon acoustic oscillations,''
  Phys.\ Rev.\ D \textbf{110} (2024) 023507.

\bibitem{DESIDR2}
  DESI Collaboration, A.~Abdul-Karim \textit{et al.},
  ``DESI DR2 Results II: Measurements of Baryon Acoustic
  Oscillations and Cosmological Constraints,''
  Phys.\ Rev.\ D \textbf{112} (2025) 083515.

\bibitem{Webb2011}
  J.~K.~Webb \textit{et al.},
  ``Indications of a spatial variation of the fine structure
  constant,''
  Phys.\ Rev.\ Lett.\ \textbf{107}, 191101 (2011).

\bibitem{Kotus2017}
  S.~M.~Kotus, M.~T.~Murphy, and R.~F.~Carswell,
  ``A sub-ppm measurement of the deuterium abundance at
  $z_{\mathrm{abs}}=3.572$,''
  Mon.\ Not.\ R.\ Astron.\ Soc.\ \textbf{464}, 3679 (2017).

\bibitem{Koide1983}
  Y.~Koide,
  ``New view of quark and lepton mass hierarchy,''
  Phys.\ Rev.\ D \textbf{28}, 252 (1983).

\bibitem{Foot1994}
  R.~Foot,
  ``A note on Koide's lepton mass relation,''
  arXiv:hep-ph/9402242 (1994).

\bibitem{Rivero2005}
  A.~Rivero and A.~Gsponer,
  ``The Koide formula and geometry,''
  arXiv:hep-ph/0505220 (2005).

\bibitem{SuperK2020}
  Super-Kamiokande Collaboration,
  ``Search for proton decay via $p \to e^{+}\pi^{0}$ and
  $p \to \mu^{+}\pi^{0}$ with an enlarged fiducial volume,''
  Phys.\ Rev.\ D \textbf{102}, 112011 (2020).

\bibitem{Kaltenbaek2012}
  R.~Kaltenbaek \textit{et al.},
  ``Macroscopic quantum resonators (MAQRO): testing quantum
  and gravitational physics with massive mechanical resonators,''
  Exp.\ Astron.\ \textbf{34}, 123 (2012).

\bibitem{Rauch1975}
  H.~Rauch \textit{et al.},
  ``Verification of coherent spinor rotation of fermions,''
  Phys.\ Lett.\ A \textbf{54}, 425 (1975).

\bibitem{Pauli1940}
  W.~Pauli,
  ``The connection between spin and statistics,''
  Phys.\ Rev.\ \textbf{58}, 716 (1940).

\bibitem{Streater1964}
  R.~F.~Streater and A.~S.~Wightman,
  \textit{PCT, Spin and Statistics, and All That}
  (W.~A.~Benjamin, 1964).

\bibitem{EulerHeisenberg1936}
  W.~Heisenberg and H.~Euler,
  ``Consequences of Dirac's theory of positrons,''
  Z.~Phys.\ \textbf{98}, 714 (1936).

\bibitem{DrummondHathrell1980}
  I.~T.~Drummond and S.~J.~Hathrell,
  ``QED vacuum polarization in a background gravitational
  field and its effect on the velocity of photons,''
  Phys.\ Rev.\ D \textbf{22}, 343 (1980).

\bibitem{Colladay1998}
  D.~Colladay and V.~A.~Kosteleck\'y,
  ``Lorentz-violating extension of the standard model,''
  Phys.\ Rev.\ D \textbf{58}, 116002 (1998).

\bibitem{Arminjon2012}
  M.~Arminjon,
  ``Gravity as Archimedes' thrust and a bifurcation in that theory,''
  Found.\ Phys.\ \textbf{42}, 1341--1363 (2012).

\bibitem{ADM1962}
  R.~Arnowitt, S.~Deser, and C.~W.~Misner,
  ``The dynamics of general relativity,''
  Gen.\ Relativ.\ Gravit.\ \textbf{40}, 1997 (2008);
  originally in \textit{Gravitation: An Introduction to
  Current Research}, ed.~L.~Witten (Wiley, 1962), Ch.~7.

\bibitem{DeWitt1967}
  B.~S.~DeWitt,
  ``Quantum theory of gravity.\ I.\ The canonical theory,''
  Phys.\ Rev.\ \textbf{160}, 1113 (1967).

\bibitem{Kuchar1981}
  K.~V.~Kucha\v{r},
  ``Canonical quantization of gravity,''
  in \textit{Relativity, Astrophysics and Cosmology},
  ed.~W.~Israel (Reidel, Dordrecht, 1973), pp.~237--288;
  ``Canonical methods of quantization,''
  in \textit{Quantum Gravity~2}, eds.~C.~J.~Isham,
  R.~Penrose, and D.~W.~Sciama
  (Oxford Univ.\ Press, 1981), pp.~329--376.

\bibitem{Halliwell1991}
  J.~J.~Halliwell,
  ``Introductory lectures on quantum cosmology,''
  in \textit{Quantum Cosmology and Baby Universes},
  eds.~S.~Coleman, J.~B.~Hartle, T.~Piran, and
  S.~Weinberg (World Scientific, 1991), pp.~159--243;
  arXiv:0909.2566 [gr-qc].

\bibitem{Gotay1986}
  M.~J.~Gotay,
  ``Constraints, reduction, and quantization,''
  J.~Math.\ Phys.\ \textbf{27}, 2051 (1986).

\bibitem{Carroll2003}
  S.~M.~Carroll, M.~Hoffman, and M.~Trodden,
  ``Can the dark energy equation-of-state parameter $w$
  be less than $-1$?''
  Phys.\ Rev.\ D \textbf{68}, 023509 (2003).

\bibitem{Cline2004}
  J.~M.~Cline, S.~Jeon, and G.~D.~Moore,
  ``The phantom menaced: constraints on low-energy
  effective ghosts,''
  Phys.\ Rev.\ D \textbf{70}, 043543 (2004).

\bibitem{Afshordi2007}
  N.~Afshordi, D.~J.~H.~Chung, and G.~Geshnizjani,
  ``Cuscuton: a causal field theory with an infinite
  speed of sound,''
  Phys.\ Rev.\ D \textbf{75}, 083513 (2007).

\bibitem{Chamseddine2013}
  A.~H.~Chamseddine and V.~Mukhanov,
  ``Mimetic dark matter,''
  J.\ High Energy Phys.\ \textbf{2013}(11), 135 (2013).

\bibitem{Cheung2008}
  C.~Cheung, P.~Creminelli, A.~L.~Fitzpatrick,
  J.~Kaplan, and L.~Senatore,
  ``The effective field theory of inflation,''
  J.\ High Energy Phys.\ \textbf{2008}(03), 014 (2008);
  arXiv:0709.0293 [hep-th].

\bibitem{tHooftVeltman1974}
  G.~'t~Hooft and M.~J.~G.~Veltman,
  ``One-loop divergencies in the theory of gravitation,''
  Ann.\ Inst.\ Henri Poincar\'e A \textbf{20}, 69 (1974).

\bibitem{GoroffSagnotti1986}
  M.~H.~Goroff and A.~Sagnotti,
  ``The ultraviolet behavior of Einstein gravity,''
  Nucl.\ Phys.\ B \textbf{266}, 709 (1986).

\bibitem{Nordtvedt1968}
  K.~Nordtvedt,
  ``Equivalence principle for massive bodies.
  II.~Theory,''
  Phys.\ Rev.\ \textbf{169}, 1017 (1968).

\end{thebibliography}

\end{document}
